% Generated mechanically from the frozen categories.tex target.
% Do not edit this generated file; edit the bound locale source instead.


% OK, start here.
%

\setcounter{chapter}{3}
\chapter{范畴}



\phantomsection
\label{categories-section-phantom}


\section{引言}
\label{categories-section-introduction}

\noindent
范畴最早在 \cite{GenEqui} 中引入。
范畴所成的范畴（它是一个真类）是一个 $2$-范畴。
类似地，栈所成的范畴也构成一个 $2$-范畴。
如果你已经熟悉范畴，但尚不熟悉 $2$-范畴，
则应把第 \ref{categories-section-formal-cat-cat} 节作为后文形式定义的导论来读。

\section{定义}
\label{categories-section-definition-categories}

\noindent
我们回顾这些定义，部分目的是固定记号。

\begin{definition}
\label{categories-definition-category}
一个{\it 范畴} $\mathcal{C}$ 由下列数据组成：
\begin{enumerate}
\item 对象集 $\Ob(\mathcal{C})$。
\item 对每一对 $x, y \in \Ob(\mathcal{C})$，给定态射集
$\Mor_\mathcal{C}(x, y)$。
\item 对每个三元组 $x, y, z\in \Ob(\mathcal{C})$，给定复合映射
$ \Mor_\mathcal{C}(y, z) \times \Mor_\mathcal{C}(x, y)
\to \Mor_\mathcal{C}(x, z) $，记作 $(\phi, \psi) \mapsto
\phi \circ \psi$。
\end{enumerate}
这些数据须满足下列规则：
\begin{enumerate}
\item 对每个 $x\in \Ob(\mathcal{C})$，存在态射
$\text{id}_x\in \Mor_\mathcal{C}(x, x)$，使得只要相应复合有定义，
就有 $\text{id}_x \circ \phi = \phi$ 和
$\psi \circ \text{id}_x = \psi $。
\item 复合满足结合律；也就是说，只要相应复合有定义，就有
$(\phi \circ \psi) \circ \chi =
\phi \circ ( \psi \circ \chi)$。
\end{enumerate}
\end{definition}

\noindent
通常要求所有态射集 $\Mor_\mathcal{C}(x, y)$ 两两不交。
这样，态射 $\phi : x \to y$ 就有唯一的{\it 源} $x$
和唯一的{\it 靶} $y$。严格说来，这一要求并非必需；
但若不作此要求，在表述上述条件 (2) 时就必须格外谨慎。
这一约定很方便，我们往往假定它成立。
在这种情形下，若 $\phi$ 的源等于 $\psi$ 的靶，
则称 $\phi$ 与 $\psi$ {\it 可复合}，此时 $\phi \circ \psi$ 有定义。
一个等价的定义是把范畴定义为五元组
$(\text{Ob}, \text{Arrows}, s, t, \circ)$：
它由对象集、态射（箭头）集、源、靶和复合组成，并满足一长列公理。
我们偶尔会采用这一观点。

\begin{remark}
\label{categories-remark-big-categories}
大范畴。在有些文献中，范畴可以有一个真类的对象。
本书也允许这种情形，但仅限于下面列出的情况
（该表将随内容推进而更新）。
特别地，当我们说“设 $\mathcal{C}$ 为范畴”时，
默认 $\Ob(\mathcal{C})$ 是集合。
\begin{enumerate}
\item 集合的范畴 $\textit{Sets}$。
\item 阿贝尔群的范畴 $\textit{Ab}$。
\item 群的范畴 $\textit{Groups}$。
\item 给定群 $G$，带左 $G$-作用的集合所成范畴
$G\textit{-Sets}$。
\item 给定环 $R$，$R$-模的范畴 $\text{Mod}_R$。
\item 给定域 $k$，$k$ 上向量空间的范畴。
\item 环的范畴。
\item 除幂环的范畴，见《除幂代数》第
\ref{dpa-section-divided-power-rings} 节。
\item 概形的范畴。
\item 拓扑空间的范畴 $\textit{Top}$。
\item 给定拓扑空间 $X$，$X$ 上集合预层的范畴
$\textit{PSh}(X)$。
\item 给定拓扑空间 $X$，$X$ 上集合层的范畴
$\Sh(X)$。
\item 给定拓扑空间 $X$，$X$ 上阿贝尔群预层的范畴
$\textit{PAb}(X)$。
\item 给定拓扑空间 $X$，$X$ 上阿贝尔群层的范畴
$\textit{Ab}(X)$。
\item 给定小范畴 $\mathcal{C}$，从 $\mathcal{C}$ 到
$\textit{Sets}$ 的函子所成范畴。
\item 给定范畴 $\mathcal{C}$，$\mathcal{C}$ 上集合预层的范畴。
\item 给定位点 $\mathcal{C}$，$\mathcal{C}$ 上集合层的范畴。
\end{enumerate}
在这里逐项列出它们，原因之一是尽量避免处理类似
“大”范畴的“集合体”这样的东西；那相当于处理所有类的集合体，
而在我们看来，这无疑是元数学对象。
\end{remark}

\begin{remark}
\label{categories-remark-unique-identity}
由定义立即可知，范畴 $\mathcal{A}$ 的对象 $x$ 的任意两个恒等态射都相同。
因此，我们可以并且将会谈论 $x$ 的{\it 这个}恒等态射
$\text{id}_x$。
\end{remark}

\begin{definition}
\label{categories-definition-isomorphism}
若范畴 $\mathcal{C}$ 中的态射 $\phi : x \to y$ 满足：
存在态射 $\psi : y \to x$，使得
$\phi \circ \psi = \text{id}_y$ 且
$\psi \circ \phi = \text{id}_x$，
则称该态射为一个{\it 同构}。
\end{definition}

\noindent
同构 $\phi$ 有时也称为{\it 可逆}态射；定义中的态射 $\psi$
称为其{\it 逆}，记作 $\phi^{-1}$。只要存在，逆就是唯一的。
注意，给定范畴 $\mathcal{A}$ 的对象 $x$，
$\Mor_\mathcal{A}(x, x)$ 中的可逆元素集
$\text{Aut}_\mathcal{A}(x)$ 在复合下构成一个群。
这个群称为 $x$ 在 $\mathcal{A}$ 中的{\it 自同构}群。

\begin{definition}
\label{categories-definition-groupoid}
若一个范畴的每个态射都是同构，则称它为{\it 群胚}。
\end{definition}

\begin{example}
\label{categories-example-group-groupoid}
群 $G$ 给出一个只有单个对象 $x$ 的群胚，
其态射满足 $\Mor(x, x) = G$，复合规则由 $G$ 中的群运算给出。
每个只有单个对象的群胚都具有这种形式。
\end{example}

\begin{example}
\label{categories-example-set-groupoid}
集合 $C$ 给出如下定义的群胚 $\mathcal{C}$：
取 $\Ob(\mathcal{C}) := C$ 为对象集；对态射，
若 $x\neq y$，则令 $\Mor(x, y)$ 为空集，
若 $x = y$，则令它等于 $\{\text{id}_x\}$。
\end{example}

\begin{definition}
\label{categories-definition-functor}
两个范畴 $\mathcal{A}, \mathcal{B}$ 之间的{\it 函子}
$F : \mathcal{A} \to \mathcal{B}$ 由下列数据给出：
\begin{enumerate}
\item 映射 $F : \Ob(\mathcal{A}) \to \Ob(\mathcal{B})$。
\item 对每个 $x, y \in \Ob(\mathcal{A})$，给定映射
$F : \Mor_\mathcal{A}(x, y) \to \Mor_\mathcal{B}(F(x), F(y))$，
记作 $\phi \mapsto F(\phi)$。
\end{enumerate}
这些数据须按如下方式与复合及恒等态射相容：
对 $\mathcal{A}$ 中任意可复合态射对 $(\phi, \psi)$，有
$F(\phi \circ \psi) = F(\phi) \circ F(\psi)$；
并且 $F(\text{id}_x) = \text{id}_{F(x)}$。
\end{definition}

\noindent
注意，每个范畴 $\mathcal{A}$ 都有{\it 恒等}函子
$\text{id}_\mathcal{A}$。
此外，给定函子 $G : \mathcal{B} \to \mathcal{C}$ 和
$F : \mathcal{A} \to \mathcal{B}$，有一个按显然方式定义的
{\it 复合}函子 $G \circ F : \mathcal{A} \to \mathcal{C}$。

\begin{definition}
\label{categories-definition-faithful}
设 $F : \mathcal{A} \to \mathcal{B}$ 为函子。
\begin{enumerate}
\item 若对任意对象 $x, y \in \Ob(\mathcal{A})$，映射
$$
F : \Mor_\mathcal{A}(x, y) \to \Mor_\mathcal{B}(F(x), F(y))
$$
都是单射，则称 $F$ {\it 忠实}。
\item 若这些映射全都是双射，则称 $F$ {\it 全忠实}。
\item
若对任意对象 $y \in \Ob(\mathcal{B})$，都存在对象
$x \in \Ob(\mathcal{A})$，使得 $F(x)$ 在 $\mathcal{B}$ 中
同构于 $y$，则称函子 $F$ {\it 本质满射}。
\end{enumerate}
\end{definition}

\begin{definition}
\label{categories-definition-subcategory}
范畴 $\mathcal{B}$ 的一个{\it 子范畴}是范畴 $\mathcal{A}$，
其对象和箭头分别组成 $\mathcal{B}$ 的对象和箭头的子集，
并且 $\mathcal{A}$ 中的源、靶和复合与 $\mathcal{B}$ 中的相应运算一致，
而 $\mathcal{A}$ 中任一对象的恒等态射也与它在 $\mathcal{B}$ 中的恒等态射相同。
若对所有 $x, y \in \Ob(\mathcal{A})$ 都有
$\Mor_\mathcal{A}(x, y) = \Mor_\mathcal{B}(x, y)$，
则称 $\mathcal{A}$ 为 $\mathcal{B}$ 的{\it 满子范畴}。
若它是满子范畴，并且给定 $x \in \Ob(\mathcal{A})$ 时，
$\mathcal{B}$ 中每个与 $x$ 同构的对象也属于 $\mathcal{A}$，
则称 $\mathcal{A}$ 为 $\mathcal{B}$ 的{\it 严格满}子范畴。
\end{definition}

\noindent
若 $\mathcal{A} \subset \mathcal{B}$ 是子范畴，
则恒等映射给出从 $\mathcal{A}$ 到 $\mathcal{B}$ 的函子。
此外，子范畴 $\mathcal{A} \subset \mathcal{B}$ 是满的，
当且仅当包含函子是全忠实的。
注意，给定范畴 $\mathcal{B}$，$\mathcal{B}$ 的满子范畴所成集合
与 $\Ob(\mathcal{B})$ 的子集所成集合相同。

\begin{remark}
\label{categories-remark-functor-into-sets}
设 $\mathcal{A}$ 为范畴。
从 $\mathcal{A}$ 到 $\textit{Sets}$ 的函子 $F$ 是数学对象
（也就是说，它是集合，而不是类或集合论公式；见
《集合论》第 \ref{sets-section-sets-everything} 节），
尽管集合范畴是“大”的。
事实上，$F$ 在对象上的像 $F(\Ob(\mathcal{A}))$ 是集合；
于是，可以把 $F$ 看作 $\mathcal{A}$ 与集合范畴的一个满子范畴之间的函子，
其中该满子范畴的对象正是 $F(\Ob(\mathcal{A}))$ 的元素。
\end{remark}

\begin{example}
\label{categories-example-group-homomorphism-functor}
群同态 $p : G\to H$ 在例 \ref{categories-example-group-groupoid} 所述的相应群胚之间
给出一个函子。它忠实当且仅当 $p$ 为单射；
它全忠实当且仅当该同态为同构。
\end{example}

\begin{example}
\label{categories-example-category-over-X}
给定范畴 $\mathcal{C}$ 和对象 $X\in \Ob(\mathcal{C})$，
如下定义{\it $X$ 上对象的范畴}，记作 $\mathcal{C}/X$。
$\mathcal{C}/X$ 的对象是某个 $Y\in \Ob(\mathcal{C})$ 所给出的态射
$Y\to X$。对象 $Y\to X$ 与 $Y'\to X$ 之间的态射，
是在 $\mathcal{C}$ 中使显然图表交换的态射 $Y\to Y'$。
注意，有一个函子 $p_X : \mathcal{C}/X\to \mathcal{C}$，
它只是忘掉所给的态射。此外，给定 $\mathcal{C}$ 中的态射
$f : X'\to X$，与 $f$ 复合诱导函子
$F : \mathcal{C}/X' \to \mathcal{C}/X$，并且
$p_X\circ F = p_{X'}$。
\end{example}

\begin{example}
\label{categories-example-category-under-X}
给定范畴 $\mathcal{C}$ 和对象 $X\in \Ob(\mathcal{C})$，
如下定义{\it $X$ 下对象的范畴}，记作 $X/\mathcal{C}$。
$X/\mathcal{C}$ 的对象是某个 $Y\in \Ob(\mathcal{C})$ 所给出的态射
$X\to Y$。对象 $X\to Y$ 与 $X\to Y'$ 之间的态射，
是在 $\mathcal{C}$ 中使显然图表交换的态射 $Y\to Y'$。
注意，有一个函子 $p_X : X/\mathcal{C}\to \mathcal{C}$，
它只是忘掉所给的态射。此外，给定 $\mathcal{C}$ 中的态射
$f : X'\to X$，与 $f$ 复合诱导函子
$F : X/\mathcal{C} \to X'/\mathcal{C}$，并且
$p_{X'}\circ F = p_X$。
\end{example}




\begin{definition}
\label{categories-definition-transformation-functors}
设 $F, G : \mathcal{A} \to \mathcal{B}$ 为函子。
一个{\it 自然变换}，或一个{\it 函子态射}
$t : F \to G$，是一个族 $\{t_x\}_{x\in \Ob(\mathcal{A})}$，满足：
\begin{enumerate}
\item $t_x : F(x) \to G(x)$ 是范畴 $\mathcal{B}$ 中的态射；
\item 对 $\mathcal{A}$ 中每个态射 $\phi : x \to y$，下图交换：
$$
\xymatrix{
F(x) \ar[r]^{t_x} \ar[d]_{F(\phi)} & G(x) \ar[d]^{G(\phi)} \\
F(y) \ar[r]^{t_y} & G(y) }
$$
\end{enumerate}
\end{definition}

\noindent
有时我们用图表
$$
\xymatrix{
\mathcal{A}
\rtwocell^F_G{t}
&
\mathcal{B}
}
$$
表示 $t$ 是从 $F$ 到 $G$ 的态射。

\medskip\noindent
注意，每个函子 $F$ 都有{\it 恒等}变换
$\text{id}_F : F \to F$。
此外，给定函子态射 $t : F \to G$ 与函子态射 $s : E \to F$，
它们的{\it 复合} $t \circ s$ 由规则
$$
(t \circ s)_x = t_x \circ s_x : E(x) \to G(x)
$$
定义，其中 $x \in \Ob(\mathcal{A})$。
容易验证，这确实是从 $E$ 到 $G$ 的函子态射。
这样，给定范畴 $\mathcal{A}$ 与 $\mathcal{B}$，
就得到一个新范畴，即 $\mathcal{A}$ 与 $\mathcal{B}$ 之间的函子所成范畴。

\begin{remark}
\label{categories-remark-functors-sets-sets}
这是当 $\mathcal{A}$ 为“大”范畴时，同样结论不成立的一种情形。
例如，考虑函子 $\textit{Sets} \to \textit{Sets}$。
按照目前的定义，这样的函子是类而不是集合。
换言之，它由一个集合论公式（其中某些变量等于指定集合）给出！
试图把集合论的所有可能公式都纳入数学对象的定义，并不是好主意。
例如，考虑概形范畴上的层时，也会出现同样的问题。
我们稍后将回到这一点。
\end{remark}

\begin{definition}
\label{categories-definition-equivalence-categories}
一个{\it 范畴等价}
$F : \mathcal{A} \to \mathcal{B}$ 是满足下列条件的函子：
存在函子 $G : \mathcal{B} \to \mathcal{A}$，使复合
$F \circ G$ 与 $G \circ F$ 分别同构于恒等函子
$\text{id}_\mathcal{B}$ 与 $\text{id}_\mathcal{A}$。
在这种情形下，称 $G$ 为 $F$ 的一个{\it 拟逆}。
\end{definition}

\begin{lemma}
\label{categories-lemma-construct-quasi-inverse}
设 $F : \mathcal{A} \to \mathcal{B}$ 为全忠实函子。
假设对每个 $X \in \Ob(\mathcal{B})$，都给定
$\mathcal{A}$ 的对象 $j(X)$ 以及同构 $i_X : X \to F(j(X))$。
则存在唯一函子 $j : \mathcal{B} \to \mathcal{A}$，
$j$ 在对象上延拓上述规则，并且这些同构 $i_X$ 定义函子同构
$\text{id}_\mathcal{B} \to F \circ j$。
此外，$j$ 与 $F$ 是互为拟逆的范畴等价。
\end{lemma}

\begin{proof}
构造 $j : \mathcal{B} \to \mathcal{A}$ 分两步进行。
首先，定义映射 $j : \Ob(\mathcal{B}) \to \Ob(\mathcal{A})$，
它把 $X \in \mathcal{B}$ 对应到 $j(X)$。
其次，若 $X,Y \in \Ob(\mathcal{B})$ 且 $\phi : X \to Y$，
则考虑 $\phi' := i_Y \circ\phi\circ i_X^{-1}$。
由于 $F$ 全忠实，存在唯一的 $\varphi$ 满足 $F(\varphi) = \phi'$。
定义 $j(\phi) = \varphi$。我们略去 $j$ 是函子的验证。
由构造，图表
$$
\xymatrix{
X \ar[r]_-{i_X} \ar[d]_{\phi}
&
F(j(X)) \ar[d]^{F\circ j(\phi)}
\\
Y \ar[r]^-{i_Y}
&
F(j(Y))
}
$$
交换。因此，由于每个 $i_X$ 都是同构，
$\{i_X\}_X$ 是函子同构 $\text{id}_\mathcal{B} \to F\circ j$。
最后，还须证明 $j \circ F$ 同构于 $\text{id}_\mathcal{A}$。
然而，由于 $F$ 全忠实，要证明这一点，只需在后面与 $F$ 复合后证明；
也就是说，只需证明 $F \circ j \circ F$ 同构于
$F \circ \text{id}_\mathcal{A}$（略去一个小细节）。
由 $F \circ j \cong \text{id}_\mathcal{B}$，这显然成立。
\end{proof}

\begin{lemma}
\label{categories-lemma-equivalence-categories}
一个函子是范畴等价，当且仅当它既全忠实又本质满射。
\end{lemma}

\begin{proof}
设 $F : \mathcal{A} \to \mathcal{B}$ 本质满射且全忠实。
按照约定，所有范畴都是小的；又因 $F$ 本质满射，
利用选择公理，可以对每个 $X \in \Ob(\mathcal{B})$ 选取
$\mathcal{A}$ 的对象 $j(X)$ 以及同构 $i_X : X \to F(j(X))$。
于是，利用 $F$ 全忠实，应用引理 \ref{categories-lemma-construct-quasi-inverse}。
\end{proof}

\begin{definition}
\label{categories-definition-product-category}
设 $\mathcal{A}$、$\mathcal{B}$ 为范畴。
定义{\it 积范畴} $\mathcal{A} \times \mathcal{B}$：
其对象集为
$\Ob(\mathcal{A} \times \mathcal{B}) =
\Ob(\mathcal{A}) \times \Ob(\mathcal{B})$，并且
$$
\Mor_{\mathcal{A} \times \mathcal{B}}((x, y), (x', y'))
:=
\Mor_\mathcal{A}(x, x')\times
\Mor_\mathcal{B}(y, y').
$$
复合按分量定义。
\end{definition}

\section{对偶范畴与 Yoneda 引理}
\label{categories-section-opposite}

\begin{definition}
\label{categories-definition-opposite}
给定范畴 $\mathcal{C}$，其{\it 对偶范畴}
$\mathcal{C}^{opp}$ 与 $\mathcal{C}$ 有相同的对象，
但所有态射的方向都反转。
\end{definition}

\noindent
换言之，
$\Mor_{\mathcal{C}^{opp}}(x, y) = \Mor_\mathcal{C}(y, x)$。
$\mathcal{C}^{opp}$ 中的复合与 $\mathcal{C}$ 中相同，只是方向相反：
若 $\phi : y \to z$ 和 $\psi : x \to y$ 是
$\mathcal{C}^{opp}$ 中的态射，亦即 $\mathcal{C}$ 中的箭头
$z \to y$ 和 $y \to x$，则 $\phi \circ^{opp} \psi$ 是
$\mathcal{C}^{opp}$ 中的态射 $x \to z$，它对应于
$\mathcal{C}$ 中的复合 $z \to y \to x$。

\begin{definition}
\label{categories-definition-contravariant}
设 $\mathcal{C}$、$\mathcal{S}$ 为范畴。
从 $\mathcal{C}$ 到 $\mathcal{S}$ 的{\it 反变}函子 $F$
是函子 $\mathcal{C}^{opp}\to \mathcal{S}$。
\end{definition}

\noindent
具体地，反变函子 $F$ 由映射
$F : \Ob(\mathcal{C}) \to \Ob(\mathcal{S})$ 以及下列数据给出：
对 $\mathcal{C}$ 中每个态射 $\psi : x \to y$，
给定态射 $F(\psi) : F(y) \to F(x)$。
这些数据须满足：再给定态射 $\phi : y \to z$ 时，
作为态射 $F(z) \to F(x)$，有
$F(\phi \circ \psi) = F(\psi) \circ F(\phi)$。
（注意次序反转了。）

\begin{definition}
\label{categories-definition-presheaf}
设 $\mathcal{C}$ 为范畴。
\begin{enumerate}
\item $\mathcal{C}$ 上的一个{\it 集合预层}，简称{\it 预层}，
是从 $\mathcal{C}$ 到 $\textit{Sets}$ 的反变函子 $F$。
\item 预层范畴记作 $\textit{PSh}(\mathcal{C})$。
\end{enumerate}
\end{definition}

\noindent
当然，预层范畴是一个真类。

\begin{example}
\label{categories-example-hom-functor}
点函子。
对任意 $U\in \Ob(\mathcal{C})$，有反变函子
$$
\begin{matrix}
h_U & : & \mathcal{C}
&
\longrightarrow
&
\textit{Sets} \\
& &
X
&
\longmapsto
&
\Mor_\mathcal{C}(X, U)
\end{matrix}
$$
它把对象 $X$ 送到集合 $\Mor_\mathcal{C}(X, U)$。
换言之，$h_U$ 是一个预层。
给定态射 $f : X\to Y$，相应映射
$h_U(f) : \Mor_\mathcal{C}(Y, U)\to \Mor_\mathcal{C}(X, U)$
把 $\phi$ 送到 $\phi\circ f$。
我们总把这个预层记作 $h_U : \mathcal{C}^{opp} \to \textit{Sets}$。
它称为与 $U$ 相伴的{\it 可表预层}。
如果 $\mathcal{C}$ 是概形范畴，这个函子有时称为 $U$ 的
\emph{点函子}。
\end{example}

\noindent
注意，给定 $\mathcal{C}$ 中的态射 $\phi : U \to V$，
与态射 $U \to V$ 复合给出相应的函子自然变换
$h(\phi) : h_U \to h_V$。
这把 $\mathcal{C}$ 中态射的复合变成函子变换的复合。
换言之，得到函子
$$
h :
\mathcal{C}
\longrightarrow
\textit{PSh}(\mathcal{C})
$$
注意，靶是一个“大”范畴，见注 \ref{categories-remark-big-categories}。
另一方面，$h$ 是一个真正的数学对象（即一个集合），
参见注 \ref{categories-remark-functor-into-sets}。

\begin{lemma}[Yoneda 引理]
\label{categories-lemma-yoneda}
\begin{reference}
它以某种形式出现于 \cite{Yoneda-homology}；Grothendieck 在
\cite{Gr-II} 中使用了其推广形式。
\end{reference}
设 $U, V \in \Ob(\mathcal{C})$。
给定任意函子态射 $s : h_U \to h_V$，
存在唯一态射 $\phi : U \to V$，使得 $h(\phi) = s$。
换言之，函子 $h$ 是全忠实的。
更一般地，给定任意反变函子 $F$ 以及 $\mathcal{C}$ 的任意对象
$U$，有自然双射
$$
\Mor_{\textit{PSh}(\mathcal{C})}(h_U, F) \longrightarrow F(U),
\quad
s \longmapsto s_U(\text{id}_U).
$$
\end{lemma}

\begin{proof}
对第一个断言，只需取
$\phi = s_U(\text{id}_U) \in \Mor_\mathcal{C}(U, V)$。
对第二个断言，给定 $\xi \in F(U)$，如下定义 $s$：
$s_V : h_U(V) \to F(V)$ 把
$h_U(V) = \Mor_\mathcal{C}(V, U)$ 的元素 $f : V \to U$
送到 $F(f)(\xi)$。
\end{proof}

\begin{definition}
\label{categories-definition-representable-functor}
若反变函子 $F : \mathcal{C}\to \textit{Sets}$
同构于 $\mathcal{C}$ 的某个对象 $U$ 的点函子 $h_U$，
则称该函子{\it 可表}。
\end{definition}

\noindent
设 $\mathcal{C}$ 为范畴，并设
$F : \mathcal{C}^{opp} \to \textit{Sets}$ 为可表函子。
选取 $\mathcal{C}$ 的对象 $U$ 以及同构 $s : h_U \to F$。
Yoneda 引理保证序对 $(U, s)$ 在唯一同构意义下是唯一的。
对象 $U$ 称为{\it 表示} $F$ 的对象。
由 Yoneda 引理，变换 $s$ 对应于唯一元素 $\xi \in F(U)$。
该元素称为{\it 泛对象}。它具有下列性质：
对 $V \in \Ob(\mathcal{C})$，映射
$$
\Mor_\mathcal{C}(V, U) \longrightarrow F(V),\quad
(f : V \to U) \longmapsto F(f)(\xi)
$$
是双射。因此，$\xi$ 在如下意义下是泛的：
$F(V)$ 的每个元素都等于 $\xi$ 在 $F(f)$ 下的像，
其中 $f : V \to U$ 是 $\mathcal{C}$ 中唯一的态射。






\section{对象对的乘积}
\label{categories-section-products-pairs}

\begin{definition}
\label{categories-definition-products}
设 $x, y\in \Ob(\mathcal{C})$。
$x$ 与 $y$ 的一个{\it 乘积}是对象
$x \times y \in \Ob(\mathcal{C})$，连同态射
$p\in \Mor_{\mathcal C}(x \times y, x)$ 和
$q\in\Mor_{\mathcal C}(x \times y, y)$，
它们满足下列泛性质：对任意 $w\in \Ob(\mathcal{C})$
以及态射 $\alpha \in \Mor_{\mathcal C}(w, x)$ 和
$\beta \in \Mor_\mathcal{C}(w, y)$，
存在唯一的 $\gamma\in \Mor_{\mathcal C}(w, x \times y)$，
使图表
$$
\xymatrix{
w \ar[rrrd]^\beta \ar@{-->}[rrd]_\gamma \ar[rrdd]_\alpha & & \\
& & x \times y \ar[d]_p \ar[r]_q & y \\
& & x &
}
$$
交换。
\end{definition}

\noindent
若乘积存在，则它在唯一同构意义下是唯一的。
这由 Yoneda 引理得到，因为定义要求 $x \times y$ 是
$\mathcal{C}$ 的对象，并且函子性地依赖于 $w$ 满足
$$
h_{x \times y}(w) = h_x(w) \times h_y(w)
$$
换言之，乘积 $x \times y$ 是表示函子
$w \mapsto h_x(w) \times h_y(w)$ 的对象。

\begin{definition}
\label{categories-definition-has-products-of-pairs}
若对任意 $x, y \in \Ob(\mathcal{C})$，乘积 $x \times y$
都存在，则称范畴 $\mathcal{C}$ {\it 具有对象对的乘积}。
\end{definition}

\noindent
我们采用这一术语，是为了把该概念与“具有乘积”或“具有有限乘积”
区别开来；后二者通常另有所指（特别是，它们总蕴含终对象存在）。





\section{对象对的余积}
\label{categories-section-coproducts-pairs}

\begin{definition}
\label{categories-definition-coproducts}
设 $x, y \in \Ob(\mathcal{C})$。
$x$ 与 $y$ 的一个{\it 余积}，或{\it 并合和}，
是对象 $x \amalg y \in \Ob(\mathcal{C})$，连同态射
$i \in \Mor_{\mathcal C}(x, x \amalg y)$ 和
$j \in \Mor_{\mathcal C}(y, x \amalg y)$，
它们满足下列泛性质：对任意 $w \in \Ob(\mathcal{C})$
以及态射 $\alpha \in \Mor_{\mathcal C}(x, w)$ 和
$\beta \in \Mor_\mathcal{C}(y, w)$，
存在唯一的 $\gamma \in \Mor_{\mathcal C}(x \amalg y, w)$，
使图表
$$
\xymatrix{
& y \ar[d]^j \ar[rrdd]^\beta \\
x \ar[r]^i \ar[rrrd]_\alpha & x \amalg y \ar@{-->}[rrd]^\gamma \\
& & & w
}
$$
交换。
\end{definition}

\noindent
若余积存在，则它在唯一同构意义下是唯一的。
这由 Yoneda 引理（应用于对偶范畴）得到，因为定义要求
$x \amalg y$ 是 $\mathcal{C}$ 的对象，并且函子性地依赖于 $w$ 满足
$$
\Mor_\mathcal{C}(x \amalg y, w) =
\Mor_\mathcal{C}(x, w) \times \Mor_\mathcal{C}(y, w)
$$

\begin{definition}
\label{categories-definition-has-coproducts-of-pairs}
若对任意 $x, y \in \Ob(\mathcal{C})$，余积 $x \amalg y$
都存在，则称范畴 $\mathcal{C}$ {\it 具有对象对的余积}。
\end{definition}

\noindent
我们采用这一术语，是为了把该概念与“具有余积”或“具有有限余积”
区别开来；后二者通常另有所指（特别是，它们总蕴含
$\mathcal{C}$ 中存在始对象）。





\section{纤维积}
\label{categories-section-fibre-products}

\begin{definition}
\label{categories-definition-fibre-products}
设 $x, y, z\in \Ob(\mathcal{C})$，
$f\in \Mor_\mathcal{C}(x, y)$，并且
$g\in \Mor_{\mathcal C}(z, y)$。
$f$ 与 $g$ 的一个{\it 纤维积}是对象
$x \times_y z\in \Ob(\mathcal{C})$，连同态射
$p \in \Mor_{\mathcal C}(x \times_y z, x)$ 和
$q \in \Mor_{\mathcal C}(x \times_y z, z)$，使图表
$$
\xymatrix{
x \times_y z \ar[r]_q \ar[d]_p & z \ar[d]^g \\
x \ar[r]^f & y
}
$$
交换，并满足下列泛性质：对任意 $w\in \Ob(\mathcal{C})$
以及满足 $f \circ \alpha = g \circ \beta$ 的态射
$\alpha \in \Mor_{\mathcal C}(w, x)$ 和
$\beta \in \Mor_\mathcal{C}(w, z)$，
存在唯一的 $\gamma \in \Mor_{\mathcal C}(w, x \times_y z)$，
使图表
$$
\xymatrix{
w \ar[rrrd]^\beta \ar@{-->}[rrd]_\gamma \ar[rrdd]_\alpha & & \\
& & x \times_y z \ar[d]^p \ar[r]_q & z \ar[d]^g \\
& & x \ar[r]^f & y
}
$$
交换。
\end{definition}

\noindent
若纤维积存在，则它在唯一同构意义下是唯一的。
这由 Yoneda 引理得到，因为定义要求 $x \times_y z$ 是
$\mathcal{C}$ 的对象，并且函子性地依赖于 $w$ 满足
$$
h_{x \times_y z}(w) = h_x(w) \times_{h_y(w)} h_z(w)
$$
换言之，纤维积 $x \times_y z$ 是表示函子
$w \mapsto h_x(w) \times_{h_y(w)} h_z(w)$ 的对象。

\begin{definition}
\label{categories-definition-cartesian}
若范畴中的交换图表
$$
\xymatrix{
w \ar[r] \ar[d] &
z \ar[d] \\
x \ar[r] &
y
}
$$
满足：$w$ 以及态射 $w \to x$ 和 $w \to z$
构成态射 $x \to y$ 与 $z \to y$ 的纤维积，
则称该图表是{\it 笛卡尔的}。
\end{definition}

\begin{definition}
\label{categories-definition-has-fibre-products}
若对任意 $f\in \Mor_{\mathcal C}(x, y)$ 和
$g\in \Mor_{\mathcal C}(z, y)$，纤维积都存在，
则称范畴 $\mathcal{C}$ {\it 具有纤维积}。
\end{definition}

\begin{definition}
\label{categories-definition-representable-morphism}
若范畴 $\mathcal{C}$ 中的态射 $f : x \to y$ 满足：
对 $\mathcal{C}$ 中每个态射 $z \to y$，
纤维积 $x \times_y z$ 都存在，则称该态射{\it 可表}。
\end{definition}

\begin{lemma}
\label{categories-lemma-composition-representable}
设 $\mathcal{C}$ 为范畴。
设 $f : x \to y$ 和 $g : y \to z$ 可表。
则 $g \circ f : x \to z$ 可表。
\end{lemma}

\begin{proof}
设 $t \in \Ob(\mathcal C)$ 且
$ \varphi \in \Mor_{\mathcal C}(t,z) $。
由于 $g$ 与 $f$ 可表，得到交换图表
$$
\xymatrix{
y \times_z t \ar[r]_q \ar[d]_p &
t \ar[d]^{\varphi} \\
y \ar[r]^{g} & z
}
\quad\quad
\xymatrix{
x \times_y (y\! \times_z\! t) \ar[r]_{q'} \ar[d]_{p'} &
y \times_z t \ar[d]^p \\
x \ar[r]^f & y
}
$$
它们具有定义 \ref{categories-definition-fibre-products} 中的泛性质。
我们声称，令 $x \times_z t = x \times_y (y \times_z t)$，
并取态射 $q \circ q' : x \times_z t \to t$ 和
$p' : x \times_z t \to x$，就得到一个纤维积。
首先，由上述图表的交换性可知
$\varphi \circ q \circ q' = g \circ f \circ p'$。
为验证泛性质，设 $w \in \Ob(\mathcal C)$，
并设态射 $\alpha : w \to x$ 和 $\beta : w \to t$ 满足
$\varphi \circ \beta = g \circ f \circ \alpha$。
由纤维积的定义，存在唯一态射 $\delta$ 与 $\gamma$，使图表
$$
\xymatrix{
w \ar[rrrd]^\beta \ar@{-->}[rrd]_\delta \ar[rrdd]_{f\circ\alpha} & & \\
& & y \times_z t \ar[d]_p \ar[r]_q & t \ar[d]^{\varphi} \\
& & y \ar[r]^{g} & z
}
$$
以及
$$
\xymatrix{
w \ar[rrrd]^\delta \ar@{-->}[rrd]_\gamma \ar[rrdd]_{\alpha} & & \\
& & x \times_y (y\!\times_z\! t) \ar[d]_{p'} \ar[r]_{q'} &
y \times_z t \ar[d]^{p} \\
& & x \ar[r]^{f} & y
}
$$
交换。于是，$\gamma$ 使图表
$$
\xymatrix{
w \ar[rrrd]^\beta \ar@{-->}[rrd]_\gamma \ar[rrdd]_{\alpha} & & \\
& & x \times_z t \ar[d]_{p'} \ar[r]_{q\circ q'} & t \ar[d]^{\varphi} \\
& & x \ar[r]^{g\circ f} & z
}
$$
交换。为证明它的唯一性，设 $\gamma'$ 满足
$ q\circ q'\circ\gamma' = \beta $ 和
$ p'\circ \gamma' = \alpha $。
由于 $\gamma$ 唯一，只需证明
$ q'\circ\gamma' = \delta $ 和
$ p'\circ\gamma' = \alpha $ 即可。第二个等式正是我们的假设。
对第一个等式，还须使用上述态射的唯一性。
注意，$\delta$ 是满足 $ q\circ\delta = \beta $ 和
$ p\circ\delta = f\circ\alpha $ 的唯一态射。
我们已经假设 $ q\circ (q'\circ\gamma') = \beta $。
此外，由纤维积的定义，知道 $ f\circ p' = p\circ q' $。
因此：
$$
p\circ (q'\circ\gamma') =
(p\circ q')\circ\gamma' =
(f\circ p')\circ\gamma' =
f\circ (p'\circ\gamma') = f\circ\alpha.
$$
所以 $ q'\circ\gamma' = \delta $，证毕。
\end{proof}

\begin{lemma}
\label{categories-lemma-base-change-representable}
设 $\mathcal{C}$ 为范畴。
设 $f : x \to y$ 可表，并设 $y' \to y$ 是
$\mathcal{C}$ 中的态射。
则态射 $x' := x \times_y y' \to y'$ 也可表。
\end{lemma}

\begin{proof}
设 $z \to y'$ 为态射。纤维积 $x' \times_{y'} z$
应当表示函子
\begin{eqnarray*}
w & \mapsto & h_{x'}(w)\times_{h_{y'}(w)} h_z(w) \\
& = & (h_x(w) \times_{h_y(w)} h_{y'}(w)) \times_{h_{y'}(w)} h_z(w) \\
& = & h_x(w) \times_{h_y(w)} h_z(w)
\end{eqnarray*}
而后者由假设可表。
\end{proof}

\section{纤维积的例}
\label{categories-section-example-fibre-products}

\noindent
本节列举并描述一些纤维积的例子。

\medskip\noindent
首先是一个十分平凡的例子：集合范畴具有纤维积，
因而每个态射都可表。具体地，若 $f : X \to Y$ 和
$g : Z \to Y$ 是集合映射，则把 $X \times_Y Z$ 定义为
$X \times Z$ 中所有满足 $f(x) = g(z)$ 的序对 $(x, z)$ 所成的子集。
态射 $p : X \times_Y Z \to X$ 与 $q : X \times_Y Z \to Z$
是投影映射 $(x, z) \mapsto x$ 和 $(x, z) \mapsto z$。
最后，若态射 $\alpha : W \to X$ 与 $\beta : W \to Z$
满足 $f \circ \alpha = g \circ \beta$，则映射
$W \to X \times Z$，$w\mapsto (\alpha(w), \beta(w))$
显然取值于 $X \times_Y Z$，正合所需。

\medskip\noindent
许多范畴的对象是在集合上赋予某类代数结构；在这些范畴中，
底层集合的纤维积也给出该范畴中的纤维积。
例如，假设上述 $X$、$Y$ 与 $Z$ 是群，而 $f$、$g$ 是群同态。
那么，集合论纤维积 $X \times_Y Z$ 继承群结构；
只需用公式 $(x, z) \cdot (x', z') = (xx', zz')$
定义两个序对的乘积。下列范畴都可使用类似论证。
\begin{enumerate}
\item 群的范畴 $\textit{Groups}$。
\item 给定群 $G$ 时，带左 $G$-作用的集合所成范畴
$G\textit{-Sets}$。
\item 环的范畴。
\item 给定环 $R$ 时，$R$-模的范畴。
\end{enumerate}




\section{纤维积与可表性}
\label{categories-section-representable-map-presheaves}

\noindent
本节具体研究从一个范畴到集合范畴的反变函子所成范畴中的纤维积。
以后讨论位点、预层与层时，将用更一般的内容取代这里的讨论。
这类函子之间的“可表态射”概念颇有意义。

\begin{lemma}
\label{categories-lemma-fibre-product-presheaves}
设 $\mathcal{C}$ 为范畴。
设 $F, G, H : \mathcal{C}^{opp} \to \textit{Sets}$ 为函子，
并设 $a : F \to G$ 与 $b : H \to G$ 为函子变换。
则预层范畴 $\textit{PSh}(\mathcal{C})$ 中的纤维积
$F \times_{a, G, b} H$ 存在，而且对 $\mathcal{C}$ 的任意对象
$X$，它由公式
$$
(F \times_{a, G, b} H)(X) =
F(X) \times_{a_X, G(X), b_X} H(X)
$$
给出。
\end{lemma}

\begin{proof}
略。
\end{proof}

\noindent
作为特例，设有态射 $a : F \to G$、对象
$U \in \Ob(\mathcal{C})$ 以及元素 $\xi \in G(U)$。
根据 Yoneda 引理 \ref{categories-lemma-yoneda}，这给出变换
$\xi : h_U \to G$。此时的纤维积由规则
$$
(h_U \times_{\xi, G, a} F)(X) =
\{ (f, \xi') \mid f : X \to U, \ \xi' \in F(X), \ G(f)(\xi) = a_X(\xi')\}
$$
描述。若 $F$、$G$ 也可表，则只要纤维积存在，这就是表示它的函子；
见第 \ref{categories-section-fibre-products} 节。
与定义 \ref{categories-definition-representable-morphism} 的类比促使我们
定义可表变换的概念。

\begin{definition}
\label{categories-definition-representable-map-presheaves}
设 $\mathcal{C}$ 为范畴，并设
$F, G : \mathcal{C}^{opp} \to \textit{Sets}$ 为函子。
若对每个 $U \in \Ob(\mathcal{C})$ 及任意 $\xi \in G(U)$，
函子 $h_U \times_{\xi, G, a} F$ 都可表，
则称态射 $a : F \to G$ {\it 可表}，
或称{\it $F$ 在 $G$ 上相对可表}。
\end{definition}

\begin{lemma}
\label{categories-lemma-representable-over-representable}
设 $\mathcal{C}$ 为范畴。
设 $a : F \to G$ 是从 $\mathcal{C}$ 到 $\textit{Sets}$
的反变函子之间的态射。若 $a$ 可表，并且 $G$ 是可表函子，
则 $F$ 可表。
\end{lemma}

\begin{proof}
略。
\end{proof}

\begin{lemma}
\label{categories-lemma-representable-diagonal}
设 $\mathcal{C}$ 为范畴，并设
$F : \mathcal{C}^{opp} \to \textit{Sets}$ 为函子。
假设 $\mathcal{C}$ 具有对象对的乘积和纤维积。
下列条件等价：
\begin{enumerate}
\item 对角态射 $\Delta : F \to F \times F$ 可表；
\item 对 $\mathcal{C}$ 中每个 $U$ 及任意 $\xi \in F(U)$，
映射 $\xi : h_U \to F$ 可表；
\item 对 $\mathcal{C}$ 中每一对 $U, V$ 及任意
$\xi \in F(U)$、$\xi' \in F(V)$，纤维积
$h_U \times_{\xi, F, \xi'} h_V$ 可表。
\end{enumerate}
\end{lemma}

\begin{proof}
我们继续利用 Yoneda 引理，把 $F(U)$ 与从 $h_U \to F$
的函子变换等同起来。

\medskip\noindent
(2) 与 (3) 的等价性。设 $U, \xi, V, \xi'$ 如 (3) 中所述。
(2) 与 (3) 都恰好断言 $h_U \times_{\xi, F, \xi'} h_V$
可表；唯一差别是 (3) 关于 $U$ 与 $V$ 对称，而 (2) 并不对称。

\medskip\noindent
假设条件 (1) 成立。设 $U, \xi, V, \xi'$ 如 (3) 中所述。
注意，$h_U \times h_V = h_{U \times V}$ 可表。
用 $\eta : h_{U \times V} \to F \times F$ 表示对应于乘积
$\xi \times \xi' : h_U \times h_V \to F \times F$ 的映射。
由假设，纤维积
$F \times_{\Delta, F \times F, \eta} h_{U \times V}$ 可表。
这意味着存在 $W \in \Ob(\mathcal{C})$、态射
$W \to U$、$W \to V$ 以及 $h_W \to F$，使得
$$
\xymatrix{
h_W \ar[d] \ar[r] & h_U \times h_V \ar[d]^{\xi \times \xi'} \\
F \ar[r] & F \times F
}
$$
是笛卡尔图表。使用引理 \ref{categories-lemma-fibre-product-presheaves}
中纤维积的显式描述，读者可见这蕴含
$h_W = h_U \times_{\xi, F, \xi'} h_V$，正合所需。

\medskip\noindent
假设等价的条件 (2) 和 (3) 成立。
设 $U$ 为 $\mathcal{C}$ 的对象，并设
$(\xi, \xi') \in (F \times F)(U)$。
由 (3)，纤维积 $h_U \times_{\xi, F, \xi'} h_U$ 可表。
选取对象 $W$ 以及同构
$h_W \to h_U \times_{\xi, F, \xi'} h_U$。
两个投影 $\text{pr}_i : h_U \times_{\xi, F, \xi'} h_U \to h_U$
由 Yoneda 对应于态射 $p_i : W \to U$。
考虑 $W' = W \times_{(p_1, p_2), U \times U} U$。
因为
$$
h_{W'} = h_W \times_{h_U \times h_U} h_U
= (h_U \times_{\xi, F, \xi'} h_U) \times_{h_U \times h_U} h_U
= F \times_{F \times F} h_U.
$$
形式地可知 $W'$ 表示 $F \times_{\Delta, F \times F} h_U$。
因此 $\Delta$ 可表，证明完成。
\end{proof}









\section{推出}
\label{categories-section-pushouts}

\noindent
与纤维积对偶的概念是推出。

\begin{definition}
\label{categories-definition-pushouts}
设 $x, y, z\in \Ob(\mathcal{C})$，
$f\in \Mor_\mathcal{C}(y, x)$ 且
$g\in \Mor_{\mathcal C}(y, z)$。
$f$ 与 $g$ 的一个 {\it 推出}是对象
$x\amalg_y z\in \Ob(\mathcal{C})$，连同态射
$p\in \Mor_{\mathcal C}(x, x\amalg_y z)$ 和
$q\in\Mor_{\mathcal C}(z, x\amalg_y z)$，使图表
$$
\xymatrix{
y \ar[r]_g \ar[d]_f & z \ar[d]^q \\
x \ar[r]^p & x\amalg_y z
}
$$
交换，并满足下列泛性质：对任意 $w\in \Ob(\mathcal{C})$
及满足 $\alpha \circ f = \beta \circ g$ 的态射
$\alpha \in \Mor_{\mathcal C}(x, w)$ 与
$\beta \in \Mor_\mathcal{C}(z, w)$，存在唯一的
$\gamma\in \Mor_{\mathcal C}(x\amalg_y z, w)$，使图表
$$
\xymatrix{
y \ar[r]_g \ar[d]_f & z \ar[d]^q \ar[rrdd]^\beta & & \\
x \ar[r]^p \ar[rrrd]^\alpha & x \amalg_y z \ar@{-->}[rrd]^\gamma & & \\
& & & w
}
$$
交换。
\end{definition}

\noindent
若三元组 $(x\amalg_y z, p, q)$ 存在，可直接且容易地证明它
在唯一同构意义下是唯一的。另一种做法是把推出看作对偶范畴中的
纤维积，于是第 \ref{categories-section-fibre-products} 节的讨论立即给出这种唯一性。

\begin{definition}
\label{categories-definition-cocartesian}
若范畴中的交换图表
$$
\xymatrix{
y \ar[r] \ar[d] & z \ar[d] \\
x \ar[r] & w
}
$$
满足：$w$ 以及态射 $x \to w$、$z \to w$ 构成态射
$y \to x$ 与 $y \to z$ 的一个推出，则称该图表{\it 余笛卡尔}。
\end{definition}


\section{等化子}
\label{categories-section-equalizers}

\begin{definition}
\label{categories-definition-equalizers}
设 $X$、$Y$ 为范畴 $\mathcal{C}$ 的对象，
$a, b : X \to Y$ 为态射。若 $a \circ e = b \circ e$，且
$(Z, e)$ 满足下列泛性质，则称态射 $e : Z \to X$
为态射对 $(a, b)$ 的一个{\it 等化子}：对 $\mathcal{C}$ 中每个满足
$a \circ t = b \circ t$ 的态射 $t : W \to X$，存在唯一态射
$s : W \to Z$，使得 $t = e \circ s$。
\end{definition}

\noindent
与上面的纤维积情形一样，等化子一旦存在，便在唯一同构意义下唯一。
把此定义推广到多于 $2$ 个态射的情形也是直接的。

\section{余等化子}
\label{categories-section-coequalizers}

\begin{definition}
\label{categories-definition-coequalizers}
设 $X$、$Y$ 为范畴 $\mathcal{C}$ 的对象，
$a, b : X \to Y$ 为态射。若 $c \circ a = c \circ b$，且
$(Z, c)$ 满足下列泛性质，则称态射 $c : Y \to Z$
为态射对 $(a, b)$ 的一个{\it 余等化子}：对 $\mathcal{C}$ 中每个满足
$t \circ a = t \circ b$ 的态射 $t : Y \to W$，存在唯一态射
$s : Z \to W$，使得 $t = s \circ c$。
\end{definition}

\noindent
与上面的推出情形一样，余等化子一旦存在，便在唯一同构意义下唯一；
考虑对偶范畴，即可由等化子的唯一性推出这一点。
把此定义推广到多于 $2$ 个态射的情形也是直接的。

\section{始对象与终对象}
\label{categories-section-initial-final}

\begin{definition}
\label{categories-definition-initial-final}
设 $\mathcal{C}$ 为范畴。
\begin{enumerate}
\item 若对 $\mathcal{C}$ 的每个对象 $y$，恰有一个态射 $x \to y$，
则称 $\mathcal{C}$ 的对象 $x$ 为{\it 始对象}。
\item 若对 $\mathcal{C}$ 的每个对象 $y$，恰有一个态射 $y \to x$，
则称 $\mathcal{C}$ 的对象 $x$ 为{\it 终对象}。
\end{enumerate}
\end{definition}

\noindent
在集合范畴中，空集 $\emptyset$ 是始对象，而且事实上是唯一的始对象。
此外，任意{\it 单点集}，即只有一个元素的集合，都是终对象
（所以终对象并不唯一）。




\section{单态射与满态射}
\label{categories-section-mono-epi}

\begin{definition}
\label{categories-definition-mono-epi}
设 $\mathcal{C}$ 为范畴，$f : X \to Y$ 为 $\mathcal{C}$ 的态射。
\begin{enumerate}
\item 若对每个对象 $W$ 及每一对满足
$f \circ a = f \circ b$ 的态射 $a, b : W \to X$，均有
$a = b$，则称 $f$ 为{\it 单态射}。
\item 若对每个对象 $W$ 及每一对满足
$a \circ f = b \circ f$ 的态射 $a, b : Y \to W$，均有
$a = b$，则称 $f$ 为{\it 满态射}。
\end{enumerate}
\end{definition}

\begin{example}
\label{categories-example-mono-epi-sets}
在集合范畴中，单态射对应于单射映射，满态射对应于满射映射。
\end{example}

\begin{lemma}
\label{categories-lemma-characterize-mono-epi}
设 $\mathcal{C}$ 为范畴，$f : X \to Y$ 为 $\mathcal{C}$ 的态射。则
\begin{enumerate}
\item $f$ 是单态射，当且仅当 $X$ 是纤维积 $X \times_Y X$；
\item $f$ 是满态射，当且仅当 $Y$ 是推出 $Y \amalg_X Y$。
\end{enumerate}
\end{lemma}

\begin{proof}
假设 $f$ 为单态射。设 $W$ 为 $\mathcal C$ 的对象，并设
$\alpha, \beta \in \Mor_\mathcal C(W,X)$ 满足 $f\circ \alpha = f\circ
\beta$。由于 $f$ 是单态射，故 $\alpha = \beta$。此外，有交换图表
$$
\xymatrix{
X \ar[r]^{\text{id}_X} \ar[d]_{\text{id}_X} & X \ar[d]^{f} \\
X \ar[r]^{f} & Y
}
$$
它以 $\gamma := \alpha = \beta$ 满足泛性质。
因此 $X$ 的确是纤维积 $X\times_Y X$。

\medskip\noindent
假设 $X \times_Y X \cong X $。图表
$$
\xymatrix{
X \ar[r]^{\text{id}_X} \ar[d]_{\text{id}_X} & X \ar[d]^{f} \\
X \ar[r]^{f} & Y }
$$
交换；若 $W \in \Ob(\mathcal C)$ 且 $\alpha, \beta : W \to X$
满足 $f \circ \alpha = f \circ \beta$，则存在唯一的
$\gamma$ 满足
$$
\gamma = \text{id}_X\circ\gamma = \alpha = \beta
$$
从而证明 $\alpha = \beta$。

\medskip\noindent
第二点的证明完全相同，只需改用推出 $Y\amalg_X Y = Y$。
\end{proof}





\section{极限与余极限}
\label{categories-section-limits}

\noindent
设 $\mathcal{C}$ 为范畴。$\mathcal{C}$ 中的一个{\it 图表}，
就是一个函子 $M : \mathcal{I} \to \mathcal{C}$。
称 $\mathcal{I}$ 为{\it 指标范畴}，也称 $M$ 为
$\mathcal{I}$-{图表}。我们用记号 $M_i$ 表示
$\mathcal{I}$ 的对象 $i$ 的像。因此，对 $\mathcal{I}$ 中的态射
$\phi : i \to i'$，有 $M(\phi) : M_i \to M_{i'}$。

\begin{definition}
\label{categories-definition-limit}
范畴 $\mathcal{C}$ 中 $\mathcal{I}$-图表 $M$ 的一个{\it 极限}，
由 $\mathcal{C}$ 中的对象 $\lim_\mathcal{I} M$ 连同态射
$p_i : \lim_\mathcal{I} M \to M_i$ 给出，并满足：
\begin{enumerate}
\item 对 $\mathcal{I}$ 中的态射 $\phi : i \to i'$，有
$p_{i'} = M(\phi) \circ p_i$；
\item 对 $\mathcal{C}$ 的任意对象 $W$ 以及由
$i \in \Ob(\mathcal{I})$ 编号的任意态射族 $q_i : W \to M_i$，
若对 $\mathcal{I}$ 中所有 $\phi : i \to i'$ 都有
$q_{i'} = M(\phi) \circ q_i$，则存在唯一态射
$q : W \to \lim_\mathcal{I} M$，使得对 $\mathcal{I}$ 的每个对象
$i$ 均有 $q_i = p_i \circ q$。
\end{enumerate}
\end{definition}

\noindent
若极限 $(\lim_\mathcal{I} M, (p_i)_{i\in \Ob(\mathcal{I})})$ 存在，
则由定义中的唯一性要求，它在唯一同构意义下唯一。
对象对的乘积、纤维积与等化子都是极限的例子。
空图表的极限是 $\mathcal{C}$ 的终对象。
在集合范畴中，所有极限都存在。对偶概念是余极限。

\begin{definition}
\label{categories-definition-colimit}
范畴 $\mathcal{C}$ 中 $\mathcal{I}$-图表 $M$ 的一个{\it 余极限}，
由 $\mathcal{C}$ 中的对象 $\colim_\mathcal{I} M$ 连同态射
$s_i : M_i \to \colim_\mathcal{I} M$ 给出，并满足：
\begin{enumerate}
\item 对 $\mathcal{I}$ 中的态射 $\phi : i \to i'$，有
$s_i = s_{i'} \circ M(\phi)$；
\item 对 $\mathcal{C}$ 的任意对象 $W$ 以及由
$i \in \Ob(\mathcal{I})$ 编号的任意态射族 $t_i : M_i \to W$，
若对 $\mathcal{I}$ 中所有 $\phi : i \to i'$ 都有
$t_i = t_{i'} \circ M(\phi)$，则存在唯一态射
$t : \colim_\mathcal{I} M \to W$，使得对 $\mathcal{I}$ 的每个对象
$i$ 均有 $t_i = t \circ s_i$。
\end{enumerate}
\end{definition}

\noindent
若余极限 $(\colim_\mathcal{I} M, (s_i)_{i\in \Ob(\mathcal{I})})$ 存在，
则由定义中的唯一性要求，它在唯一同构意义下唯一。
对象对的余积、推出与余等化子都是余极限的例子。
空图表的余极限是 $\mathcal{C}$ 的始对象。
在集合范畴中，所有余极限都存在。

\begin{remark}
\label{categories-remark-diagram-small}
（余）极限的指标范畴绝不允许具有真类那么多的对象。
在本项目中，这意味着它不能是注 \ref{categories-remark-big-categories}
所列出的范畴之一。
\end{remark}

\begin{remark}
\label{categories-remark-limit-colim}
我们常用 $\lim_i M_i$、$\colim_i M_i$、
$\lim_{i\in \mathcal{I}} M_i$ 或 $\colim_{i\in \mathcal{I}} M_i$
代替以 $\mathcal{I}$ 为下标的写法。使用这种记号，并利用下文
第 \ref{categories-section-limit-sets} 节对集合的极限与余极限的描述，
可以作如下陈述。设 $M : \mathcal{I} \to \mathcal{C}$ 为图表。
\begin{enumerate}
\item 若对象 $\lim_i M_i$ 存在，则它满足
$$
\Mor_\mathcal{C}(W, \lim_i M_i)
=
\lim_i \Mor_\mathcal{C}(W, M_i)
$$
其中右侧的极限在集合范畴中取。
\item 若对象 $\colim_i M_i$ 存在，则它满足
$$
\Mor_\mathcal{C}(\colim_i M_i, W)
=
\lim_{i \in \mathcal{I}^{opp}} \Mor_\mathcal{C}(M_i, W)
$$
其中右侧取对偶范畴上的极限，并以集合范畴为值域。
\end{enumerate}
由 Yoneda 引理（及其对偶），这些公式分别完全确定极限与余极限。
\end{remark}

\begin{remark}
\label{categories-remark-cones-and-cocones}
设 $M : \mathcal{I} \to \mathcal{C}$ 为图表。
在此情形下，$M$ 的一个{\it 锥}由对象 $W$ 以及态射族
$q_i : W \to M_i$、$i \in \Ob(\mathcal{I})$ 给出，
并要求对 $\mathcal{I}$ 的所有态射 $\phi : i \to i'$，图表
$$
\xymatrix{
& W \ar[dl]_{q_i} \ar[dr]^{q_{i'}} \\
M_i \ar[rr]^{M(\phi)} & & M_{i'}
}
$$
交换。所有锥构成一个范畴，其态射概念是显然的。
显然，若 $M$ 的极限存在，它就是锥范畴中的终对象。
对偶地，$M$ 的一个{\it 余锥}由对象 $W$ 与态射族
$t_i : M_i \to W$ 给出，并要求对 $\mathcal{I}$ 中的所有态射
$\phi : i \to i'$，图表
$$
\xymatrix{
M_i \ar[rr]^{M(\phi)} \ar[dr]_{t_i} & & M_{i'} \ar[dl]^{t_{i'}} \\
& W
}
$$
交换。所有余锥构成一个范畴，其态射概念也是显然的。
与上面类似，$M$ 的余极限存在，当且仅当余锥范畴有始对象。
\end{remark}

\noindent
作为极限与余极限概念的应用，我们定义乘积与余积。

\begin{definition}
\label{categories-definition-product}
设 $I$ 为集合，并且对每个 $i \in I$ 给定范畴 $\mathcal{C}$
的对象 $M_i$。{\it 乘积} $\prod_{i\in I} M_i$ 按定义是
$\lim_\mathcal{I} M$（若其存在），其中 $\mathcal{I}$ 是以
$I$ 的元素为对象、且只以恒等态射为态射的范畴。
\end{definition}

\noindent
一个重要特例是 $I = \emptyset$；此时乘积是该范畴的终对象。
态射 $p_i : \prod M_i \to M_i$ 称为{\it 投影态射}。

\begin{definition}
\label{categories-definition-coproduct}
设 $I$ 为集合，并且对每个 $i \in I$ 给定范畴 $\mathcal{C}$
的对象 $M_i$。{\it 余积} $\coprod_{i\in I} M_i$ 按定义是
$\colim_\mathcal{I} M$（若其存在），其中 $\mathcal{I}$ 是以
$I$ 的元素为对象、且只以恒等态射为态射的范畴。
\end{definition}

\noindent
一个重要特例是 $I = \emptyset$；此时余积是该范畴的始对象。
注意，余积配有态射 $M_i \to \coprod M_i$。
这些态射有时称为{\it 余投影}。

\begin{lemma}
\label{categories-lemma-functorial-colimit}
设 $M : \mathcal{I} \to \mathcal{C}$ 与
$N : \mathcal{J} \to \mathcal{C}$ 为余极限均存在的图表。
设 $H : \mathcal{I} \to \mathcal{J}$ 为函子，并设
$t : M \to N \circ H$ 为函子变换。则存在唯一态射
$$
\theta :
\colim_\mathcal{I} M
\longrightarrow
\colim_\mathcal{J} N
$$
使所有图表
$$
\xymatrix{
M_i \ar[d]_{t_i} \ar[r]
&
\colim_\mathcal{I} M \ar[d]^{\theta}
\\
N_{H(i)} \ar[r]
&
\colim_\mathcal{J} N
}
$$
交换。
\end{lemma}

\begin{proof}
略。
\end{proof}

\begin{lemma}
\label{categories-lemma-functorial-limit}
设 $M : \mathcal{I} \to \mathcal{C}$ 与
$N : \mathcal{J} \to \mathcal{C}$ 为极限均存在的图表。
设 $H : \mathcal{I} \to \mathcal{J}$ 为函子，并设
$t : N \circ H \to M$ 为函子变换。则存在唯一态射
$$
\theta :
\lim_\mathcal{J} N
\longrightarrow
\lim_\mathcal{I} M
$$
使所有图表
$$
\xymatrix{
\lim_\mathcal{J} N \ar[d]^{\theta} \ar[r]
&
N_{H(i)} \ar[d]_{t_i}
\\
\lim_\mathcal{I} M \ar[r]
&
M_i
}
$$
交换。
\end{lemma}

\begin{proof}
略。
\end{proof}


\begin{lemma}
\label{categories-lemma-colimits-commute}
设 $\mathcal{I}$、$\mathcal{J}$ 为指标范畴，并设
$M : \mathcal{I} \times \mathcal{J} \to \mathcal{C}$ 为函子。
假设对所有 $i$，$M_{i, \infty} = \colim_j M_{i,j}$ 都存在。
则所得函子 $M_{-, \infty} : \mathcal{I} \to \mathcal{C}$
有余极限，当且仅当 $M$ 有余极限；此时二者的余极限相同。
特别地，只要所示各余极限都存在，就有
$$
\colim_i \colim_j M_{i, j}
=
\colim_{i, j} M_{i, j}
=
\colim_j \colim_i M_{i, j}
$$
极限的情形类似。
\end{lemma}

\begin{proof}
略。
\end{proof}

\begin{lemma}
\label{categories-lemma-limits-products-equalizers}
设 $M : \mathcal{I} \to \mathcal{C}$ 为图表。
记 $I = \Ob(\mathcal{I})$，$A = \text{Arrows}(\mathcal{I})$，
并以 $s, t : A \to I$ 表示源映射与靶映射。
假设 $\prod_{i \in I} M_i$ 与 $\prod_{a \in A} M_{t(a)}$ 存在，
并假设
$$
\xymatrix{
\prod_{i \in I} M_i
\ar@<1ex>[r]^\phi \ar@<-1ex>[r]_\psi
&
\prod_{a \in A} M_{t(a)}
}
$$
的等化子存在，其中各态射由其分量按如下方式确定：
$p_a \circ \psi = M(a) \circ p_{s(a)}$，且
$p_a \circ \phi = p_{t(a)}$。则这个等化子就是该图表的极限。
\end{lemma}

\begin{proof}
略。
\end{proof}


\begin{lemma}
\label{categories-lemma-colimits-coproducts-coequalizers}
\begin{slogan}
若所有余积与余等化子都存在，则所有余极限都存在。
\end{slogan}
设 $M : \mathcal{I} \to \mathcal{C}$ 为图表。
记 $I = \Ob(\mathcal{I})$，$A = \text{Arrows}(\mathcal{I})$，
并以 $s, t : A \to I$ 表示源映射与靶映射。
假设 $\coprod_{i \in I} M_i$ 与 $\coprod_{a \in A} M_{s(a)}$ 存在，
并假设
$$
\xymatrix{
\coprod_{a \in A} M_{s(a)}
\ar@<1ex>[r]^\phi \ar@<-1ex>[r]_\psi
&
\coprod_{i \in I} M_i
}
$$
的余等化子存在，其中各态射由其分量按如下方式确定：
分量 $M_{s(a)}$ 通过态射 $M(a)$，经 $\psi$ 映到分量
$M_{t(a)}$。分量 $M_{s(a)}$ 经 $\phi$，由恒等态射映到分量
$M_{s(a)}$。则这个余等化子就是该图表的余极限。
\end{lemma}

\begin{proof}
略。
\end{proof}






\section{集合范畴中的极限与余极限}
\label{categories-section-limit-sets}

\noindent
在 $\textit{Sets}$ 中，极限与余极限不仅存在，而且很容易描述。
具体地，设 $M : \mathcal{I} \to \textit{Sets}$、$i \mapsto M_i$
为集合图表，并记 $I = \Ob(\mathcal{I})$。其极限描述为
$$
\lim_\mathcal{I} M
=
\{
(m_i)_{i\in I} \in \prod\nolimits_{i\in I} M_i
\mid
\forall \phi : i \to i' \text{ 于 }\mathcal{I},
M(\phi)(m_i) = m_{i'}
\}.
$$
因此，可把极限的一个元素看作所有集合 $M_i$ 中元素的一个相容系。

\medskip\noindent
另一方面，余极限为
$$
\colim_\mathcal{I} M
=
(\coprod\nolimits_{i\in I} M_i)/\sim
$$
其中等价关系 $\sim$ 由下列关系生成：若 $m_i \in M_i$、
$m_{i'} \in M_{i'}$，且对某个 $\phi : i \to i'$ 有
$M(\phi)(m_i) = m_{i'}$，则令 $m_i \sim m_{i'}$。
换言之，$m_i \in M_i$ 与 $m_{i'} \in M_{i'}$ 等价，
当且仅当 $\mathcal{I}$ 中存在一条态射链
$$
\xymatrix{
&
i_1 \ar[ld] \ar[rd] & &
i_3 \ar[ld] & &
i_{2n-1} \ar[rd] & \\
i = i_0 & &
i_2 & &
\ldots & &
i_{2n} = i'
}
$$
以及元素 $m_{i_j} \in M_{i_j}$，使它们在上述 $\mathcal{I}$
中态射所诱导的映射 $M_{i_{2k-1}} \to M_{i_{2k-2}}$ 与
$M_{i_{2k-1}} \to M_{i_{2k}}$ 下彼此对应。

\medskip\noindent
这种对象用起来并不方便。但若图表是滤过的，其描述便简单得多。
我们将在第 \ref{categories-section-directed-colimits} 节说明这一点。



\section{连通极限}
\label{categories-section-connected-limits}

\noindent
若一个（余）极限的指标范畴是连通的，则称该（余）极限是连通的。

\begin{definition}
\label{categories-definition-category-connected}
若由
$x \sim y \Leftrightarrow \Mor_\mathcal{I}(x, y) \not = \emptyset$
生成的等价关系恰有一个等价类，则称范畴 $\mathcal{I}$
是{\it 连通的}。
\end{definition}

\noindent
这里遵循拓扑学定义 \ref{topology-definition-connected-components}
的约定：连通空间非空。下述结果在某种宽泛意义上刻画连通极限。

\begin{lemma}
\label{categories-lemma-connected-limit-over-X}
设 $\mathcal{C}$ 为范畴，$X$ 为 $\mathcal{C}$ 的对象，并设
$M : \mathcal{I} \to \mathcal{C}/X$ 为 $X$ 上对象范畴中的图表。
若指标范畴 $\mathcal{I}$ 连通，且 $M$ 在 $\mathcal{C}/X$
中的极限存在，则复合
$\mathcal{I} \to \mathcal{C}/X \to \mathcal{C}$
的极限也存在，且二者相同。
\end{lemma}

\begin{proof}
设 $L \to X$ 为表示 $\mathcal{C}/X$ 中极限的对象。考虑函子
$$
W \longmapsto \lim_i \Mor_\mathcal{C}(W, M_i).
$$
设 $(\varphi_i)$ 为右侧集合的元素。由于每个 $M_i$ 都配有态射
$s_i : M_i \to X$，得到态射
$f_i = s_i \circ \varphi_i : W \to X$。
但因为 $\mathcal{I}$ 连通，所有 $f_i$ 都相等；又因
$\mathcal{I}$ 非空，至少有一个 $f_i$。因此这个公共值
$W \to X$ 赋予 $W$ 以 $\mathcal{C}/X$ 中对象的结构，而
$(\varphi_i)$ 给出 $\lim_i \Mor_{\mathcal{C}/X}(W, M_i)$
的一个元素。由此得到唯一态射 $\phi : W \to L$，使每个
$\varphi_i$ 都是 $\phi$ 与 $L \to M_i$ 的复合，正合所需。
\end{proof}

\begin{lemma}
\label{categories-lemma-connected-colimit-under-X}
设 $\mathcal{C}$ 为范畴，$X$ 为 $\mathcal{C}$ 的对象，并设
$M : \mathcal{I} \to X/\mathcal{C}$ 为 $X$ 下对象范畴中的图表。
若指标范畴 $\mathcal{I}$ 连通，且 $M$ 在 $X/\mathcal{C}$
中的余极限存在，则复合
$\mathcal{I} \to X/\mathcal{C} \to \mathcal{C}$
的余极限也存在，且二者相同。
\end{lemma}

\begin{proof}
略。提示：本引理与引理 \ref{categories-lemma-connected-limit-over-X} 对偶。
\end{proof}




\section{共尾范畴与始范畴}
\label{categories-section-cofinal}

\noindent
文献中有时在下述定义里用“终”一词代替“共尾”。

\begin{definition}
\label{categories-definition-cofinal}
设 $H : \mathcal{I} \to \mathcal{J}$ 为范畴之间的函子。
若下列条件成立，则称{\it $\mathcal{I}$ 在 $\mathcal{J}$ 中共尾}，
或称 $H$ {\it 共尾}：
\begin{enumerate}
\item 对所有 $y \in \Ob(\mathcal{J})$，存在
$x \in \Ob(\mathcal{I})$ 及态射 $y \to H(x)$；
\item 给定 $y \in \Ob(\mathcal{J})$、
$x, x' \in \Ob(\mathcal{I})$ 及态射 $y \to H(x)$、
$y \to H(x')$，存在 $\mathcal{I}$ 中的态射序列
$$
x = x_0 \leftarrow x_1 \rightarrow x_2 \leftarrow x_3 \rightarrow \ldots
\rightarrow x_{2n} = x'
$$
以及 $\mathcal{J}$ 中的态射 $y \to H(x_i)$，其中
$y \to H(x_0)$ 与 $y \to H(x_{2n})$ 是给定态射，并且对
$k = 0, \ldots, n - 1$，图表
$$
\xymatrix{
& y \ar[ld] \ar[d] \ar[rd] \\
H(x_{2k}) & H(x_{2k + 1}) \ar[l] \ar[r] & H(x_{2k + 2})
}
$$
交换。
\end{enumerate}
换言之，固定 $\mathcal{J}$ 的对象 $y$，考虑所有序对
$(x, b)$ 构成的集合 $S$，其中 $x$ 是 $\mathcal{I}$ 的对象，
$b : y \to H(x)$ 是态射。在 $S$ 上考虑由下列关系生成的
等价关系：若存在态射 $a : x \to x'$ 满足
$b' = H(a) \circ b$，则令 $(x, b) \sim (x', b')$。
此时 $S$ 应恰好由一个等价类组成。
\end{definition}

\begin{lemma}
\label{categories-lemma-cofinal}
设 $H : \mathcal{I} \to \mathcal{J}$ 为范畴的函子，并假设
$\mathcal{I}$ 在 $\mathcal{J}$ 中共尾。则对每个图表
$M : \mathcal{J} \to \mathcal{C}$，只要任一边存在，就有典范同构
$$
\colim_\mathcal{I} M \circ H
=
\colim_\mathcal{J} M
$$
\end{lemma}

\begin{proof}
略。
\end{proof}

\begin{definition}
\label{categories-definition-initial}
设 $H : \mathcal{I} \to \mathcal{J}$ 为范畴之间的函子。
若下列条件成立，则称{\it $\mathcal{I}$ 在 $\mathcal{J}$ 中是始的}，
或称 $H$ {\it 是始的}：
\begin{enumerate}
\item 对所有 $y \in \Ob(\mathcal{J})$，存在
$x \in \Ob(\mathcal{I})$ 及态射 $H(x) \to y$；
\item 对任意 $y \in \Ob(\mathcal{J})$、
$x , x' \in \Ob(\mathcal{I})$ 以及 $\mathcal{J}$ 中的态射
$H(x) \to y$、$H(x') \to y$，存在 $\mathcal{I}$ 中的态射序列
$$
x = x_0 \leftarrow x_1 \rightarrow x_2 \leftarrow x_3 \rightarrow \ldots
\rightarrow x_{2n} = x'
$$
以及 $\mathcal{J}$ 中的态射 $H(x_i) \to y$，使得对
$k = 0, \ldots, n - 1$，图表
$$
\xymatrix{
H(x_{2k}) \ar[rd] &
H(x_{2k + 1}) \ar[l] \ar[r] \ar[d] &
H(x_{2k + 2}) \ar[ld] \\
& y
}
$$
交换。
\end{enumerate}
\end{definition}

\noindent
这正是“共尾”函子的对偶概念。

\begin{lemma}
\label{categories-lemma-initial}
设 $H : \mathcal{I} \to \mathcal{J}$ 为范畴的函子，并假设
$\mathcal{I}$ 在 $\mathcal{J}$ 中是始的。则对每个图表
$M : \mathcal{J} \to \mathcal{C}$，只要任一边存在，就有典范同构
$$
\lim_\mathcal{I} M \circ H = \lim_\mathcal{J} M
$$
\end{lemma}

\begin{proof}
略。
\end{proof}

\begin{lemma}
\label{categories-lemma-colimit-constant-connected-fibers}
设 $F : \mathcal{I} \to \mathcal{I}'$ 为函子。假设
\begin{enumerate}
\item $\mathcal{I}$ 在 $\mathcal{I}'$ 上的所有纤维范畴
（见定义 \ref{categories-definition-fibre-category}）均连通；
\item 对 $\mathcal{I}'$ 中每个态射 $\alpha' : x' \to y'$，
都存在 $\mathcal{I}$ 中的态射 $\alpha : x \to y$，使得
$F(\alpha) = \alpha'$。
\end{enumerate}
则对每个图表 $M : \mathcal{I}' \to \mathcal{C}$，余极限
$\colim_\mathcal{I} M \circ F$ 存在，当且仅当
$\colim_{\mathcal{I}'} M$ 存在；若存在，则二者相同。
\end{lemma}

\begin{proof}
可以先证明 $\mathcal{I}$ 在 $\mathcal{I}'$ 中共尾，
再应用引理 \ref{categories-lemma-cofinal}。也可如下直接证明。
只需证明，对 $\mathcal{C}$ 的任意对象 $T$，有
$$
\lim_{\mathcal{I}^{opp}} \Mor_\mathcal{C}(M_{F(i)}, T)
=
\lim_{(\mathcal{I}')^{opp}} \Mor_\mathcal{C}(M_{i'}, T)
$$
若 $(g_{i'})_{i' \in \Ob(\mathcal{I}')}$ 是右侧的元素，
令 $f_i = g_{F(i)}$，便得到左侧的元素
$(f_i)_{i \in \Ob(\mathcal{I})}$。反过来，设
$(f_i)_{i \in \Ob(\mathcal{I})}$ 为左侧的元素。
注意，在每个（连通）纤维范畴 $\mathcal{I}_{i'}$ 上，
函子 $M \circ F$ 都是取常值 $M_{i'}$ 的常值函子。
因此，所有满足 $F(i) = i'$ 的 $i \in \Ob(\mathcal{I})$
所对应的态射 $f_i$ 都相同，并确定一个良定义态射
$g_{i'} : M_{i'} \to T$。由假设 (2)，族
$(g_{i'})_{i' \in \Ob(\mathcal{I}')}$ 给出右侧的一个元素。
\end{proof}

\begin{lemma}
\label{categories-lemma-product-with-connected}
设 $\mathcal{I}$ 与 $\mathcal{J}$ 为范畴，并以
$p : \mathcal{I} \times \mathcal{J} \to \mathcal{J}$ 表示投影。
若 $\mathcal{I}$ 连通，则对图表 $M : \mathcal{J} \to \mathcal{C}$，
余极限 $\colim_\mathcal{J} M$ 存在，当且仅当
$\colim_{\mathcal{I} \times \mathcal{J}} M \circ p$ 存在；
若存在，则二者相等。
\end{lemma}

\begin{proof}
这是引理 \ref{categories-lemma-colimit-constant-connected-fibers} 的特例。
\end{proof}






\section{有限极限与余极限}
\label{categories-section-finite-limits}

\noindent
若一个（余）极限的指标范畴有限，即该指标范畴只有有限多个对象
和有限多个态射，则称它为{\it 有限}（余）极限。
若指标范畴非空，则称（余）极限{\it 非空}；
若指标范畴连通，则称（余）极限{\it 连通}，
见定义 \ref{categories-definition-category-connected}。
事实证明，有限指标范畴已经“足够多”。

\begin{lemma}
\label{categories-lemma-finite-diagram-category}
设 $\mathcal{I}$ 为满足下列条件的范畴：
\begin{enumerate}
\item $\Ob(\mathcal{I})$ 有限；
\item 存在有限多个态射
$f_1, \ldots, f_m \in \text{Arrows}(\mathcal{I})$，
使 $\mathcal{I}$ 的每个态射都能写成复合
$f_{j_1} \circ f_{j_2} \circ \ldots \circ f_{j_k}$。
\end{enumerate}
则存在函子 $F : \mathcal{J} \to \mathcal{I}$，使得
\begin{enumerate}
\item[(a)] $\mathcal{J}$ 是有限范畴；
\item[(b)] 对任意图表 $M : \mathcal{I} \to \mathcal{C}$，
$M$ 在 $\mathcal{I}$ 上的（余）极限存在，当且仅当
$M \circ F$ 在 $\mathcal{J}$ 上的（余）极限存在；在这种情形下，
这两个（余）极限典范同构。
\end{enumerate}
此外，$\mathcal{J}$ 连通（相应地，非空），当且仅当
$\mathcal{I}$ 连通（相应地，非空）。
\end{lemma}

\begin{proof}
设 $\Ob(\mathcal{I}) = \{x_1, \ldots, x_n\}$。
以 $s, t : \{1, \ldots, m\} \to \{1, \ldots, n\}$
表示满足 $f_j : x_{s(j)} \to x_{t(j)}$ 的函数。
令 $\Ob(\mathcal{J}) = \{y_1, \ldots, y_n, z_1, \ldots, z_n\}$。
除恒等态射外，再引入态射
$g_j : y_{s(j)} \to z_{t(j)}$、$j = 1, \ldots, m$，
以及态射 $h_i : y_i \to z_i$、$i = 1, \ldots, n$。
由于 $\mathcal{J}$ 中所有非恒等态射都从某个 $y$ 指向某个 $z$，
没有需要定义的复合，也没有需要检验的结合律。
令 $F(y_i) = F(z_i) = x_i$、$F(g_j) = f_j$ 且
$F(h_i) = \text{id}_{x_i}$。显然 $F$ 是函子，
$\mathcal{J}$ 是有限的，并且 $\mathcal{J}$ 连通（相应地，非空）
当且仅当 $\mathcal{I}$ 如此。

\medskip\noindent
设 $M : \mathcal{I} \to \mathcal{C}$ 为图表。
考虑 $\mathcal{C}$ 的对象 $W$ 以及定义 \ref{categories-definition-limit}
所述的态射 $q_i : W \to M(x_i)$。
把它们取作 $q_i : W \to M(F(y_i)) = M(F(z_i)) = M(x_i)$，
便得到定义 \ref{categories-definition-limit} 中图表 $M \circ F$ 所需的映射族。
反过来，设给定定义 \ref{categories-definition-limit} 中图表 $M \circ F$
所需的映射 $qy_i : W \to M(F(y_i))$ 与
$qz_i : W \to M(F(z_i))$。由于
$$
M(F(h_i)) = \text{id} : M(F(y_i)) = M(x_i) \longrightarrow M(x_i) = M(F(z_i))
$$
可知对所有 $i$ 都有 $qy_i = qz_i$。令 $q_i$ 等于这个公共值。
$q_{s(j)} = qy_{s(j)}$ 与 $q_{t(j)} = qz_{t(j)}$
对态射 $M(f_j)$ 的相容性，保证族 $q_i$ 对 $\mathcal{I}$
中的所有态射相容，因为根据假设，每个这样的态射都是态射 $f_j$
的复合。因此得到典范双射
$$
\lim_{B \in \Ob(\mathcal{J})} \Mor_\mathcal{C}(W, M(F(B)))
=
\lim_{A \in \Ob(\mathcal{I})} \Mor_\mathcal{C}(W, M(A))
$$
这蕴含引理关于极限的陈述。余极限的陈述可用同样方式证明
（证明略）。
\end{proof}

\begin{lemma}
\label{categories-lemma-fibre-products-equalizers-exist}
设 $\mathcal{C}$ 为范畴。下列条件等价：
\begin{enumerate}
\item $\mathcal{C}$ 中存在连通有限极限。
\item $\mathcal{C}$ 中存在等化子与纤维积。
\end{enumerate}
\end{lemma}

\begin{proof}
由于等化子与纤维积都是有限连通极限，故 (1) 蕴含 (2)。
反之，设 $\mathcal{I}$ 为有限连通指标范畴，并设
$F : \mathcal{J} \to \mathcal{I}$ 为引理
\ref{categories-lemma-finite-diagram-category} 的证明中构造的指标范畴函子。
于是可用 $\mathcal{J}$ 代替 $\mathcal{I}$。由此可假设
$\Ob(\mathcal{I}) = \{x_1, \ldots, x_n\} \amalg \{y_1, \ldots, y_m\}$，
其中 $n, m \geq 1$，并且 $\mathcal{I}$ 中所有非恒等态射都是
某些 $i$、$j$ 所对应的态射 $f : x_i \to y_j$。

\medskip\noindent
假设 $n > 1$。由于 $\mathcal{I}$ 连通，存在指标
$i_1, i_2$、$j_0$ 以及态射 $a : x_{i_1} \to y_{j_0}$、
$b : x_{i_2} \to y_{j_0}$。考虑范畴
$$
\mathcal{I}' =
\{x\} \amalg \{x_1, \ldots, \hat x_{i_1}, \ldots, \hat x_{i_2}, \ldots x_n\}
\amalg \{y_1, \ldots, y_m\}
$$
其中
$$
\Mor_{\mathcal{I}'}(x, y_j) = \Mor_\mathcal{I}(x_{i_1}, y_j)
\amalg \Mor_\mathcal{I}(x_{i_2}, y_j)
$$
而所有其他态射集都与 $\mathcal{I}$ 中相同。
对任意函子 $M : \mathcal{I} \to \mathcal{C}$，可构造函子
$M' : \mathcal{I}' \to \mathcal{C}$，令
$$
M'(x) = M(x_{i_1}) \times_{M(a), M(y_{j_0}), M(b)} M(x_{i_2})
$$
并且对与某个 $f : x_{i_1} \to y_j$ 对应的态射
$f' : x \to y_j$，令 $M'(f) = M(f) \circ \text{pr}_1$。
于是函子 $M$ 有极限，当且仅当函子 $M'$ 有极限（证明略）。
故由归纳法归约到 $n = 1$ 的情形。

\medskip\noindent
若 $n = 1$，则任意 $M : \mathcal{I} \to \mathcal{C}$ 的极限
是各对映射 $x_1 \to y_j$ 的逐次等化子，因而由假设存在。
\end{proof}

\begin{lemma}
\label{categories-lemma-almost-finite-limits-exist}
设 $\mathcal{C}$ 为范畴。下列条件等价：
\begin{enumerate}
\item $\mathcal{C}$ 中存在非空有限极限。
\item $\mathcal{C}$ 中存在对象对的乘积与等化子。
\item $\mathcal{C}$ 中存在对象对的乘积与纤维积。
\end{enumerate}
\end{lemma}

\begin{proof}
由于对象对的乘积、纤维积与等化子都是指标范畴非空的极限，
故 (1) 同时蕴含 (2) 与 (3)。假设 (2)。则有限非空乘积与
等化子存在。由引理 \ref{categories-lemma-limits-products-equalizers}，
有限非空极限存在，即 (1) 成立。假设 (3)。若
$a, b : A \to B$ 为 $\mathcal{C}$ 的态射，则 $a, b$ 的等化子为
$$
(A \times_{a, B, b} A)\times_{(\text{pr}_1, \text{pr}_2), A \times A, \Delta} A.
$$
因此 (3) 蕴含 (2)，引理得证。
\end{proof}

\begin{lemma}
\label{categories-lemma-finite-limits-exist}
设 $\mathcal{C}$ 为范畴。下列条件等价：
\begin{enumerate}
\item $\mathcal{C}$ 中存在有限极限。
\item 存在有限乘积与等化子。
\item 该范畴有终对象，且纤维积存在。
\end{enumerate}
\end{lemma}

\begin{proof}
由于有限乘积、纤维积、等化子与终对象都是有限指标范畴上的极限，
故 (1) 同时蕴含 (2) 与 (3)。由上面的引理
\ref{categories-lemma-limits-products-equalizers}，可知 (2) 蕴含 (1)。
假设 (3)。注意，乘积 $A \times B$ 是终对象上的纤维积。
若 $a, b : A \to B$ 为 $\mathcal{C}$ 的态射，则 $a, b$ 的等化子为
$$
(A \times_{a, B, b} A)\times_{(\text{pr}_1, \text{pr}_2), A \times A, \Delta} A.
$$
因此 (3) 蕴含 (2)，引理得证。
\end{proof}

\begin{lemma}
\label{categories-lemma-push-outs-coequalizers-exist}
设 $\mathcal{C}$ 为范畴。下列条件等价：
\begin{enumerate}
\item $\mathcal{C}$ 中存在连通有限余极限。
\item $\mathcal{C}$ 中存在余等化子与推出。
\end{enumerate}
\end{lemma}

\begin{proof}
略。提示：这与引理 \ref{categories-lemma-fibre-products-equalizers-exist} 对偶。
\end{proof}

\begin{lemma}
\label{categories-lemma-almost-finite-colimits-exist}
设 $\mathcal{C}$ 为范畴。下列条件等价：
\begin{enumerate}
\item $\mathcal{C}$ 中存在非空有限余极限。
\item $\mathcal{C}$ 中存在对象对的余积与余等化子。
\item $\mathcal{C}$ 中存在对象对的余积与推出。
\end{enumerate}
\end{lemma}

\begin{proof}
略。提示：这与引理 \ref{categories-lemma-almost-finite-limits-exist} 对偶。
\end{proof}

\begin{lemma}
\label{categories-lemma-colimits-exist}
设 $\mathcal{C}$ 为范畴。下列条件等价：
\begin{enumerate}
\item $\mathcal{C}$ 中存在有限余极限。
\item $\mathcal{C}$ 中存在有限余积与余等化子。
\item 该范畴有始对象，且推出存在。
\end{enumerate}
\end{lemma}

\begin{proof}
略。提示：这与引理 \ref{categories-lemma-finite-limits-exist} 对偶。
\end{proof}





\section{滤过余极限}
\label{categories-section-directed-colimits}

\noindent
在滤过图表上取余极限时，余极限较容易计算或描述。定义如下。

\begin{definition}
\label{categories-definition-directed}
若下列条件成立，则称图表 $M : \mathcal{I} \to \mathcal{C}$
是{\it 有向的}，或称{\it 滤过的}：
\begin{enumerate}
\item 范畴 $\mathcal{I}$ 至少有一个对象；
\item 对 $\mathcal{I}$ 的每一对对象 $x, y$，存在对象 $z$
及态射 $x \to z$、$y \to z$；
\item 对 $\mathcal{I}$ 的每一对对象 $x, y$ 以及 $\mathcal{I}$ 中每一对态射
$a, b : x \to y$，存在 $\mathcal{I}$ 中的态射 $c : y \to z$，
使得 $M(c \circ a) = M(c \circ b)$ 作为 $\mathcal{C}$ 中的态射相等。
\end{enumerate}
若 $\text{id} : \mathcal{I} \to \mathcal{I}$ 是滤过的，
则称指标范畴 $\mathcal{I}$ {\it 有向}或{\it 滤过}；
换言之，删去上面第 (3) 项中的 $M$。
\end{definition}

\noindent
可以看到，指标范畴滤过的任意图表都是滤过图表，而滤过余极限通常
正是以这种方式出现。事实上，若 $M : \mathcal{I} \to \mathcal{C}$
是滤过图表，则可把 $M$ 分解为
$\mathcal{I} \to \mathcal{I}' \to \mathcal{C}$，其中
$\mathcal{I}'$ 是滤过指标范畴\footnote{具体地，令
$\mathcal{I}'$ 与 $\mathcal{I}$ 有相同对象，但把
$\Mor_{\mathcal{I}'}(x, y)$ 定义为 $\Mor_\mathcal{I}(x, y)$
关于如下等价关系的商：若 $M(a) = M(b)$，则等同
$a, b : x \to y$。}，并且 $\colim_\mathcal{I} M$ 存在，
当且仅当 $\colim_{\mathcal{I}'} M'$ 存在；此时两个余极限典范同构。

\medskip\noindent
假设 $M : \mathcal{I} \to \textit{Sets}$ 为滤过图表。
此时，公式
$$
\colim_\mathcal{I} M
=
(\coprod\nolimits_{i\in I} M_i)/\sim
$$
中的等价关系可以简单描述为
$$
m_i \sim m_{i'}
\Leftrightarrow
\exists i'', \phi : i \to i'', \phi': i' \to i'',
M(\phi)(m_i) = M(\phi')(m_{i'}).
$$
换言之，两个元素在余极限中相等，当且仅当它们“最终变得相等”。

\begin{lemma}
\label{categories-lemma-directed-commutes}
设 $\mathcal{I}$ 与 $\mathcal{J}$ 为指标范畴。
假设 $\mathcal{I}$ 滤过，$\mathcal{J}$ 有限，并设
$M : \mathcal{I} \times \mathcal{J} \to \textit{Sets}$、
$(i, j) \mapsto M_{i, j}$ 为集合图表所成的图表。则
$$
\colim_i \lim_j M_{i, j}
=
\lim_j \colim_i M_{i, j}.
$$
特别地，$\mathcal{I}$ 上的余极限与集合的有限乘积、纤维积
及等化子可交换。
\end{lemma}

\begin{proof}
略。事实上，证明下述命题是一道有趣的练习：一个范畴滤过，
当且仅当该范畴上的余极限与有限极限（取值于集合范畴）可交换。
\end{proof}

\noindent
下面给出 $\mathcal{J}$ 无限时上述引理的反例。
令 $\mathcal{I}$ 的对象为 $\mathbf{N} = \{1, 2, 3, \ldots\}$，
且当 $i \leq i'$ 时恰有一个态射 $i \to i'$。
令 $\mathcal{J}$ 为离散范畴
$\mathbf{N} = \{1, 2, 3, \ldots\}$，即其态射只有恒等态射。
令 $M_{i, j} = \{1, 2, \ldots, i\}$；当 $i \leq i'$ 时，
态射 $M_{i, j} \to M_{i', j}$ 是显然的包含映射。
此时 $\colim_i M_{i, j} = \mathbf{N}$，因而
$$
\lim_j \colim_i M_{i, j}
=
\prod\nolimits_j \mathbf{N}
=
\mathbf{N}^\mathbf{N}
$$
另一方面，$\lim_j M_{i, j} = \prod\nolimits_j M_{i, j}$，故
$$
\colim_i \lim_j M_{i, j}
=
\bigcup\nolimits_i \{1, 2, \ldots, i\}^{\mathbf{N}}
$$
它比另一个极限更小。

\begin{lemma}
\label{categories-lemma-cofinal-in-filtered}
\begin{reference}
\cite[Proposition 3.2.4]{KS}
\end{reference}
设 $\mathcal{I}$ 为范畴，$\mathcal{J}$ 为其满子范畴。
假设 $\mathcal{I}$ 滤过，并且对 $\mathcal{I}$ 的任意对象 $i$，
都存在到 $\mathcal{J}$ 的某个对象 $j$ 的态射 $i \to j$。
则 $\mathcal{J}$ 滤过，且在 $\mathcal{I}$ 中共尾。
\end{lemma}

\begin{proof}
略。这是对有关概念的一道愉快练习。
\end{proof}

\noindent
事实证明，有时需要更精细地控制可以对指标范畴施加的条件。
因此，我们补充若干关于可能要求的性质的引理。

\begin{lemma}
\label{categories-lemma-preserve-products}
设 $\mathcal{I}$ 为指标范畴，即一个范畴。假设对
$\mathcal{I}$ 的每一对对象 $x, y$，都存在对象 $z$ 及态射
$x \to z$、$y \to z$。则
\begin{enumerate}
\item 若 $M$ 与 $N$ 是 $\mathcal{I}$ 上的集合图表，则
$\colim (M_i \times N_i) \to \colim M_i \times \colim N_i$ 为满射；
\item 一般而言，$\mathcal{I}$ 上集合图表的余极限不与有限非空乘积交换。
\end{enumerate}
\end{lemma}

\begin{proof}
证明 (1)。设 $(\overline{m}, \overline{n})$ 是
$\colim M_i \times \colim N_i$ 的元素。则对某些
$x, y \in \Ob(\mathcal{I})$，可找到 $m \in M_x$ 与 $n \in N_y$，
使 $m$ 映到 $\overline{m}$，而 $n$ 映到 $\overline{n}$。
见第 \ref{categories-section-limit-sets} 节。在 $\mathcal{I}$ 中选取
$a : x \to z$ 与 $b : y \to z$。于是
$(M(a)(m), N(b)(n))$ 是 $(M \times N)_z$ 的元素，
它在 $\colim (M_i \times N_i)$ 中的像映到
$(\overline{m}, \overline{n})$，正合所需。

\medskip\noindent
证明 (2)。设 $G$ 为非平凡群，并令 $\mathcal{I}$ 为只有一个对象、
自同态幺半群为 $G$ 的范畴。则 $\mathcal{I}$ 显然满足引理所述条件。
现在令 $G$ 以平移作用于自身，并把 $G$-集 $G$ 看作取值于集合的
$\mathcal{I}$-图表。则
$$
\colim_\mathcal{I} G \times \colim_\mathcal{I} G \cong G/G \times G/G
$$
不同构于
$$
\colim_\mathcal{I} (G \times G) \cong (G \times G)/G
$$
这个例子说明，即使只关心有限乘积，也不能直接删去引理
\ref{categories-lemma-directed-commutes} 中的附加条件。
\end{proof}

\begin{lemma}
\label{categories-lemma-colimits-abelian-as-sets}
设 $\mathcal{I}$ 为指标范畴，即一个范畴。假设对
$\mathcal{I}$ 的每一对对象 $x, y$，都存在对象 $z$ 及态射
$x \to z$、$y \to z$。设
$M : \mathcal{I} \to \textit{Ab}$ 为 $\mathcal{I}$ 上的阿贝尔群图表。
则 $M$ 在集合范畴中的余极限满射到 $M$ 在阿贝尔群范畴中的余极限。
\end{lemma}

\begin{proof}
回忆集合范畴中的余极限是无交并 $\coprod M_i$ 关于某关系的商，
见第 \ref{categories-section-limit-sets} 节。类似地，阿贝尔群范畴中的余极限
是直和 $\bigoplus M_i$ 的商。引理的假设意味着：给定
$i, j \in \Ob(\mathcal{I})$、$m \in M_i$ 与 $n \in M_j$，
可以找到对象 $k$ 以及态射 $a : i \to k$、$b : j \to k$。
因此，$m + n$ 在余极限中由 $M_k$ 的元素
$M(a)(m) + M(b)(n)$ 表示。故 $\coprod M_i$ 满射到该余极限。
\end{proof}

\begin{lemma}
\label{categories-lemma-split-into-connected}
设 $\mathcal{I}$ 为指标范畴，即一个范畴。假设对
$\mathcal{I}$ 中的每个实线图表
$$
\xymatrix{
x \ar[d] \ar[r] & y \ar@{..>}[d] \\
z \ar@{..>}[r] & w
}
$$
都存在对象 $w$ 和虚线箭头，使图表交换。
则 $\mathcal{I}$ 或为空，或为若干具有同一性质的连通范畴的
非空无交并。
\end{lemma}

\begin{proof}
若 $\mathcal{I}$ 是空范畴，则引理成立。否则，在
$\mathcal{I}$ 的对象上定义关系：若存在 $z$ 以及态射
$x \to z$、$y \to z$，则称 $x \sim y$。
由引理的假设，这是等价关系。因此 $\Ob(\mathcal{I})$
是若干等价类的无交并。令 $\mathcal{I}_j$ 为这些等价类所对应的
满子范畴。于是 $\mathcal{I} = \coprod \mathcal{I}_j$，
且每个 $\mathcal{I}_j$ 都非空，正合所需。
\end{proof}

\begin{lemma}
\label{categories-lemma-preserve-injective-maps}
设 $\mathcal{I}$ 为指标范畴，即一个范畴。假设对
$\mathcal{I}$ 中的每个实线图表
$$
\xymatrix{
x \ar[d] \ar[r] & y \ar@{..>}[d] \\
z \ar@{..>}[r] & w
}
$$
都存在对象 $w$ 和虚线箭头，使图表交换。则
\begin{enumerate}
\item $\mathcal{I}$ 上集合图表的单射态射 $M \to N$，
诱导集合的单射映射 $\colim M_i \to \colim N_i$；
\item 对阿贝尔群图表及其余极限，一般并非如此。
\end{enumerate}
\end{lemma}

\begin{proof}
若 $\mathcal{I}$ 是空范畴，则引理成立。因此可假设
$\mathcal{I}$ 非空。此时可写成
$\mathcal{I} = \coprod \mathcal{I}_j$，其中每个
$\mathcal{I}_j$ 都非空且满足同一性质，见引理
\ref{categories-lemma-split-into-connected}。由于
$\colim_\mathcal{I} M = \coprod_j \colim_{\mathcal{I}_j} M|_{\mathcal{I}_j}$，
这把 (1) 的证明归约到连通情形。

\medskip\noindent
假设 $\mathcal{I}$ 连通且 $M \to N$ 是单射，
即所有映射 $M_i \to N_i$ 都是单射。把 $M_i$ 与
$M_i \to N_i$ 的像等同，也就是把 $M_i$ 看作 $N_i$ 的子集。
下文直接使用第 \ref{categories-section-limit-sets} 节给出的余极限描述。
设 $s, s' \in \colim M_i$ 映到 $\colim N_i$ 的同一元素。
设 $s$ 来自 $M_i$ 的元素 $m$，$s'$ 来自 $M_{i'}$ 的元素 $m'$。
则可找到 $\mathcal{I}$ 的对象序列
$i = i_0, i_1, \ldots, i_n = i'$ 以及态射
$$
\xymatrix{
&
i_1 \ar[ld] \ar[rd] & &
i_3 \ar[ld] & &
i_{2n-1} \ar[rd] & \\
i = i_0 & &
i_2 & &
\ldots & &
i_{2n} = i'
}
$$
和元素 $n_{i_j} \in N_{i_j}$；在上述 $\mathcal{I}$ 中态射
所诱导的映射 $N_{i_{2k-1}} \to N_{i_{2k-2}}$ 与
$N_{i_{2k-1}} \to N_{i_{2k}}$ 下，这些元素彼此对应，
并且 $n_{i_0} = m$、$n_{i_{2n}} = m'$。
我们对 $n$ 归纳证明这蕴含 $s = s'$。初始情形 $n = 0$ 显然。
假设 $n \geq 1$。利用关于 $\mathcal{I}$ 的假设，得到交换图表
$$
\xymatrix{
& i_1 \ar[ld] \ar[rd] \\
i_0 \ar[rd] & & i_2 \ar[ld] \\
& w
}
$$
由于 $m$ 与 $n_{i_2}$ 都是元素 $n_{i_1}$ 的像，可知二者映到
$N_w$ 的同一元素。特别地，该元素是某个 $m'' \in M_w$，
它在 $\colim M_i$ 中给出的元素与 $s$ 相同。于是得到链
$$
\xymatrix{
&
i_3 \ar[ld] \ar[rd] & &
i_5 \ar[ld] & &
i_{2n-1} \ar[rd] & \\
w & &
i_4 & &
\ldots & &
i_{2n} = i'
}
$$
以及 $j \geq 3$ 时的元素 $n_{i_j}$；这条链比原链短。
这证明归纳步骤，从而完成 (1) 的证明。

\medskip\noindent
设 $G$ 为群，并令 $\mathcal{I}$ 为只有一个对象、自同态幺半群为
$G$ 的范畴。则 $\mathcal{I}$ 满足引理所述条件，因为给定
$g_1, g_2 \in G$，总可找到 $h_1, h_2 \in G$，使
$h_1 g_1 = h_2 g_2$。$\textit{Ab}$ 中 $\mathcal{I}$ 上的图表 $M$
等同于带 $G$-作用的阿贝尔群 $M$，而
$\colim_\mathcal{I} M$ 是 $M$ 的余不变量 $M_G$。
令 $G$ 为 $2$ 阶群，它平凡作用于
$M = \mathbf{Z}/2\mathbf{Z}$；把它映入
$N = \mathbf{Z}/2\mathbf{Z} \times \mathbf{Z}/2\mathbf{Z}$
的第一个直和项，并令 $G$ 的非平凡元素按
$(x, y) \mapsto (x + y, y)$ 作用。则
$M_G \to N_G$ 是零映射。
\end{proof}

\begin{lemma}
\label{categories-lemma-split-into-directed}
设 $\mathcal{I}$ 为指标范畴，即一个范畴。假设
\begin{enumerate}
\item 对 $\mathcal{I}$ 中每一对态射 $a : w \to x$ 与
$b : w \to y$，都存在对象 $z$ 及态射 $c : x \to z$、
$d : y \to z$，使得 $c \circ a = d \circ b$；
\item 对每一对态射 $a, b : x \to y$，都存在态射
$c : y \to z$，使得 $c \circ a = c \circ b$。
\end{enumerate}
则 $\mathcal{I}$ 是若干两两无交的滤过指标范畴
$\mathcal{I}_j$ 的并（可以为空）。
\end{lemma}

\begin{proof}
若 $\mathcal{I}$ 是空范畴，则引理成立。否则，在
$\mathcal{I}$ 的对象上定义关系：若存在 $z$ 以及态射
$x \to z$、$y \to z$，则称 $x \sim y$。
由引理的第一个假设，这是等价关系。因此
$\Ob(\mathcal{I})$ 是若干等价类的无交并。
令 $\mathcal{I}_j$ 为这些等价类所对应的满子范畴。
余下结论由定义立即得到。
\end{proof}

\begin{lemma}
\label{categories-lemma-almost-directed-commutes-equalizers}
设 $\mathcal{I}$ 为满足上面引理
\ref{categories-lemma-split-into-directed} 各项假设的指标范畴。
则 $\mathcal{I}$ 上的余极限与集合中的纤维积和等化子交换
（更一般地，与有限连通极限交换）。
\end{lemma}

\begin{proof}
由引理 \ref{categories-lemma-split-into-directed}，可写
$\mathcal{I} = \coprod \mathcal{I}_j$，其中每个
$\mathcal{I}_j$ 都滤过。由引理 \ref{categories-lemma-directed-commutes}，
$\mathcal{I}_j$ 上的余极限与等化子和纤维积交换。
因此，只需证明集合范畴中的等化子和纤维积与余积交换
（包括空余积）。换言之，给定集合 $J$、集合
$A_j, B_j, C_j$ 以及 $j \in J$ 时的集合映射
$A_j \to B_j$、$C_j \to B_j$，需证明
$$
(\coprod\nolimits_{j \in J} A_j)
\times_{(\coprod\nolimits_{j \in J} B_j)}
(\coprod\nolimits_{j \in J} C_j)
=
\coprod\nolimits_{j \in J} A_j \times_{B_j} C_j
$$
并且给定 $a_j, a'_j : A_j \to B_j$，还需证明
$$
\text{Equalizer}(
\coprod\nolimits_{j \in J} a_j,
\coprod\nolimits_{j \in J} a'_j)
=
\coprod\nolimits_{j \in J}
\text{Equalizer}(a_j, a'_j)
$$
即使 $J = \emptyset$，这些结论也成立。细节略。
\end{proof}



\section{余滤过极限}
\label{categories-section-codirected-limits}

\noindent
在余滤过图表上取极限时，极限较容易计算或描述。定义如下。

\begin{definition}
\label{categories-definition-codirected}
若下列条件成立，则称图表 $M : \mathcal{I} \to \mathcal{C}$
是{\it 反向的}或{\it 余滤过的}：
\begin{enumerate}
\item 范畴 $\mathcal{I}$ 至少有一个对象；
\item 对 $\mathcal{I}$ 的每一对对象 $x, y$，存在对象 $z$
及态射 $z \to x$、$z \to y$；
\item 对 $\mathcal{I}$ 的每一对对象 $x, y$ 以及
$\mathcal{I}$ 中的每一对态射 $a, b : x \to y$，
存在 $\mathcal{I}$ 中的态射 $c : w \to x$，使得
$M(a \circ c) = M(b \circ c)$ 作为 $\mathcal{C}$ 中的态射相等。
\end{enumerate}
若 $\text{id} : \mathcal{I} \to \mathcal{I}$ 是余滤过的，
则称指标范畴 $\mathcal{I}$ {\it 反向}或{\it 余滤过}；
换言之，删去上面第 (3) 项中的 $M$。
\end{definition}

\noindent
可以看到，指标范畴余滤过的任意图表都是余滤过图表，
而这种情形通常正是以这种方式出现。

\medskip\noindent
为说明集合的余滤过极限为何比一般极限“容易”，我们提及如下事实：
有限非空集合的余滤过图表具有非空极限
（引理 \ref{categories-lemma-nonempty-limit}）。
这一结论对有限非空集合的一般极限并不成立。










\section{预序集上的极限与余极限}
\label{categories-section-posets-limits}

\noindent
图表的一种特殊情形是预序集上的系。

\begin{definition}
\label{categories-definition-directed-set}
设 $I$ 为集合，$\leq$ 为 $I$ 上的二元关系。
\begin{enumerate}
\item 若 $\leq$ 具有传递性（若 $i \leq j$ 且 $j \leq k$，
则 $i \leq k$）和自反性（对所有 $i \in I$ 都有 $i \leq i$），
则称该关系为{\it 预序}。
\item 赋有预序的集合称为{\it 预序集}。
\item 若预序集 $(I, \leq)$ 满足 $I$ 非空，且
$\forall i, j \in I$，存在 $k \in I$ 使
$i \leq k, j \leq k$，则称其为{\it 有向集}。
\item 若 $\leq$ 是反对称的预序（若 $i \leq j$ 且
$j \leq i$，则 $i = j$），则称该关系为{\it 偏序}。
\item 赋有偏序的集合称为{\it 偏序集}。
\item 序关系为偏序的有向集称为{\it 有向偏序集}。
\end{enumerate}
\end{definition}

\noindent
谈论预序集时，通常从记号中略去 $\leq$；也就是说，
说预序集 $I$，而不说预序集 $(I, \leq)$。
给定预序集 $I$，符号 $\geq$ 由下述规则定义：对所有
$i, j \in I$，$i \geq j \Leftrightarrow j \leq i$。
“偏序集”有时简称为“poset”。

\medskip\noindent
给定预序集 $I$，可以构造一个范畴：对象是 $I$ 的元素；
若 $i \leq i'$，则恰有一个态射 $i \to i'$，否则没有这样的态射。
反过来，给定任意两个对象之间至多有一个箭头的范畴 $\mathcal{C}$，
集合 $\Ob(\mathcal{C})$ 上由规则
$x \leq y \Leftrightarrow \Mor_\mathcal{C}(x, y) \not = \emptyset$
定义了一个预序。

\begin{definition}
\label{categories-definition-system-over-poset}
设 $(I, \leq)$ 为预序集，$\mathcal{C}$ 为范畴。
\begin{enumerate}
\item $\mathcal{C}$ 中 $I$ 上的一个{\it 系}，有时也称
$\mathcal{C}$ 中 $I$ 上的{\it 归纳系}，由 $\mathcal{C}$ 的对象
$M_i$ 以及每当 $i \leq i'$ 时给定的态射
$f_{ii'} : M_i \to M_{i'}$ 构成，并满足 $f_{ii}
= \text{id}$，且每当 $i \leq i' \leq i''$ 时有
$f_{ii''} = f_{i'i''} \circ f_{i i'}$。
\item $\mathcal{C}$ 中 $I$ 上的一个{\it 逆系}，
有时也称 $\mathcal{C}$ 中 $I$ 上的{\it 射影系}，
由 $\mathcal{C}$ 的对象 $M_i$ 以及每当 $i' \leq i$ 时给定的态射
$f_{ii'} : M_i \to M_{i'}$ 构成，并满足 $f_{ii}
= \text{id}$，且每当 $i'' \leq i' \leq i$ 时有
$f_{ii''} = f_{i'i''} \circ f_{i i'}$。（注意不等式的方向反转。）
\end{enumerate}
我们用“$(M_i, f_{ii'})$ 是 $I$ 上的一个（逆）系”表示上述情形。
映射 $f_{ii'}$ 有时称为{\it 过渡映射}。
\end{definition}

\noindent
换言之，$I$ 上的一个系就是图表
$M : \mathcal{I} \to \mathcal{C}$，其中 $\mathcal{I}$ 是上面
与 $I$ 相联系的范畴：其对象是 $I$ 的元素，而且
$\mathcal{I}$ 中存在唯一箭头 $i \to i'$，当且仅当 $i \leq i'$。
逆系则是图表 $M : \mathcal{I}^{opp} \to \mathcal{C}$。
从这个观点看，可以对 $I$ 上的任意（逆）系取（余）极限。
不过，习惯上{\it 只对 $I$ 上的系取余极限}，
并且{\it 只对 $I$ 上的逆系取极限}。更精确地说：
给定 $I$ 上的系 $(M_i, f_{ii'})$，系 $(M_i, f_{ii'})$ 的余极限定义为
$$
\colim_{i \in I} M_i = \colim_\mathcal{I} M,
$$
即相应图表的余极限。给定 $I$ 上的逆系 $(M_i, f_{ii'})$，
逆系 $(M_i, f_{ii'})$ 的极限定义为
$$
\lim_{i \in I} M_i = \lim_{\mathcal{I}^{opp}} M,
$$
即相应图表的极限。

\begin{remark}
\label{categories-remark-preorder-versus-partial-order}
设 $I$ 为预序集。可以由 $I$ 构造一个典范偏序集
$\overline{I}$ 以及保序映射 $\pi : I \to \overline{I}$。
具体地，在 $I$ 上按规则
$$
i \sim j \Leftrightarrow (i \leq j\text{ 且 }j \leq i).
$$
定义等价关系 $\sim$。令 $\overline{I} = I/\sim$，并令
$\pi : I \to \overline{I}$ 为商映射。最后，$\overline{I}$
上存在唯一偏序，使得
$\pi(i) \leq \pi(j) \Leftrightarrow i \leq j$。
注意，若 $I$ 是有向集，则 $\overline{I}$ 是有向偏序集。
给定 $\overline{I}$ 上的（逆）系 $N$，令
$M_i = N_{\pi(i)}$，便得到 $I$ 上的（逆）系 $M$。
这一构造定义了 $I$ 与 $\overline{I}$ 上逆系范畴之间的函子。
事实上，这是一个范畴等价。原因是，若 $i \sim j$，则对
$I$ 上的任意系 $M$，映射 $M_i \to M_j$ 与 $M_j \to M_i$
互为逆同构。更精确地，选取 $\pi$ 的截面
$s : \overline{I} \to I$，上述函子的一个拟逆把 $M$ 送到
满足 $N_{\overline{i}} = M_{s(\overline{i})}$ 的 $N$。
最后，这种对应与系的余极限相容：若 $M$ 与 $N$ 如上相关，
并且 $\colim_{\overline{I}} N$ 或 $\colim_I M$ 任一存在，
则另一个也存在，且
$\colim_{\overline{I}} N = \colim_I M$。
对逆系及逆系的极限，有类似结论。
\end{remark}

\noindent
注 \ref{categories-remark-preorder-versus-partial-order} 的要点是：
计算一个系的余极限或一个逆系的极限时，总可假设该预序为偏序。

\begin{definition}
\label{categories-definition-directed-system}
设 $I$ 为预序集。若 $I$ 是有向集
（定义 \ref{categories-definition-directed-set}），即 $I$ 非空，
且对所有 $i_1, i_2 \in I$，存在 $i\in I$ 使
$i_1 \leq i$ 且 $i_2 \leq i$，则称系（相应地，逆系）
$(M_i, f_{ii'})$ 为{\it 有向系}（相应地，{\it 有向逆系}）。
\end{definition}

\noindent
在这种情形下，余极限有时（很不幸地）称为“直极限”。
我们不采用后一术语。事实证明，在下述意义上，
滤过范畴上的图表并不比有向系更一般。

\begin{lemma}
\label{categories-lemma-directed-category-system}
设 $\mathcal{I}$ 为滤过指标范畴。存在有向集 $I$ 以及
$\mathcal{I}$ 中 $I$ 上的系 $(x_i, \varphi_{ii'})$，
并具有下列性质：
\begin{enumerate}
\item 对任意范畴 $\mathcal{C}$ 以及任意取值于 $\mathcal{C}$ 的图表
$M : \mathcal{I} \to \mathcal{C}$，用
$(M(x_i), M(\varphi_{ii'}))$ 表示相应的 $I$ 上的系。
若 $\colim_{i \in I} M(x_i)$ 存在，则
$\colim_\mathcal{I} M$ 也存在，而且引理
\ref{categories-lemma-functorial-colimit} 中的变换
$$
\theta :
\colim_{i \in I} M(x_i)
\longrightarrow
\colim_\mathcal{I} M
$$
是同构。
\item 对任意范畴 $\mathcal{C}$ 以及 $\mathcal{C}$ 中任意图表
$M : \mathcal{I}^{opp} \to \mathcal{C}$，用
$(M(x_i), M(\varphi_{ii'}))$ 表示相应的 $I$ 上的逆系。
若 $\lim_{i \in I} M(x_i)$ 存在，则
$\lim_{\mathcal{I}^{opp}} M$ 也存在，而且引理
\ref{categories-lemma-functorial-limit} 中的变换
$$
\theta :
\lim_{\mathcal{I}^{opp}} M
\longrightarrow
\lim_{i \in I} M(x_i)
$$
是同构。
\end{enumerate}
\end{lemma}

\begin{proof}
如定义 \ref{categories-definition-system-over-poset} 后的正文所述，
可以把预序集看作范畴，把系看作函子。在整个证明中，
我们将自由地在这两个观点之间切换。为证明第一个陈述，
构造一个对应于有向集的范畴 $\mathcal{I}_0$\footnote{事实上，
我们的构造会得到一个有向偏序集。}以及共尾函子
$M_0 : \mathcal{I}_0 \to \mathcal{I}$。于是由引理
\ref{categories-lemma-cofinal}，图表 $M : \mathcal{I} \to \mathcal{C}$
的余极限与图表 $M \circ M_0 : \mathcal{I}_0 \to \mathcal{C}$
的余极限相同，从而得到所需陈述。第二个陈述与第一个对偶；
把 $\mathcal{C}$ 中的极限解释为 $\mathcal{C}^{opp}$ 中的余极限，
即可证明。细节略。

\medskip\noindent
若范畴 $\mathcal{F}$ 中存在有限箭头集 $F$，使
$\mathcal{F}$ 中每个箭头都能由 $F$ 中的箭头复合而得，
则称 $\mathcal{F}$ {\em 有限生成}。特别地，这蕴含
$\mathcal{F}$ 只有有限多个对象。我们先把证明归约到
$\mathcal{I}$ 具有下述性质的情形：
$\mathcal{I}$ 的每个有限生成子范畴都能扩张为具有唯一终对象的
有限生成子范畴。

\medskip\noindent
以 $\omega$ 表示有限序数所成的有向集，并把它看作滤过范畴。
容易验证，积范畴 $\mathcal{I}\times \omega$ 也滤过，且投影
$\Pi : \mathcal{I} \times \omega \to \mathcal{I}$ 共尾。

\medskip\noindent
现在设 $\mathcal{F}$ 为 $\mathcal{I}\times \omega$
的任意有限生成子范畴。利用滤过范畴的公理，并对
$\mathcal{F}$ 的有限生成元集作简单归纳，可为图表
$\mathcal{F} \to \mathcal{I} \times \omega$ 构造
$\mathcal{I} \times \omega$ 中的一个余锥
$(\{f_i\}, i_\infty)$。也就是说，对 $\mathcal{F}$ 中的每个对象
$i$，有态射 $f_i : i \to i_\infty$，并且对
$\mathcal{F}$ 中每个箭头 $f : i \to i'$，有
$f_i = f_{i'} \circ f$。还可选取 $i_\infty$，使得不存在从
$i_\infty$ 到 $\mathcal{F}$ 中对象的箭头。这是可能的，
因为总可把箭头 $f_i$ 后复合一个箭头；该箭头在
$\mathcal{I}$-分量上为恒等，在 $\omega$-分量上严格递增。
现在令 $\mathcal{F}^+$ 为如下范畴：它包含 $\mathcal{F}$
的所有对象与箭头，并加入对象 $i_\infty$、恒等箭头
$\text{id}_{i_\infty}$ 及箭头 $f_i$。由于在 $\mathcal{F}^+$
中不存在从 $i_\infty$ 到 $\mathcal{F}$ 任一对象的箭头，
$\mathcal{F}^+$ 的箭头集对复合封闭，所以
$\mathcal{F}^+$ 的确是范畴。按构造，它是 $\mathcal{I}$
的有限生成子范畴，并以 $i_\infty$ 为唯一终对象。
又由引理 \ref{categories-lemma-cofinal}，任意图表
$M : \mathcal{I} \to \mathcal{C}$ 的余极限都与
$M\circ\Pi$ 的余极限相同；这就给出了所需归约。

\medskip\noindent
$\mathcal{I}$ 中所有具有唯一终对象的有限生成子范畴所成的集合，
按包含关系自然排序。令 $\mathcal{I}_0$ 为与该集合对应的范畴。
还有函子 $M_0 : \mathcal{I}_0 \to \mathcal{I}$：它把
$\mathcal{I}_0$ 中的箭头 $\mathcal{F} \subset \mathcal{F'}$
送到从 $\mathcal{F}$ 的终对象到 $\mathcal{F}'$ 的终对象的唯一映射。
给定 $\mathcal{I}$ 的任意两个有限生成子范畴，由这两个范畴
生成的范畴仍有限生成。由关于 $\mathcal{I}$ 的假设，它还包含于
$\mathcal{I}$ 的某个具有唯一终对象的有限生成子范畴中。
这说明 $\mathcal{I}_0$ 有向。

\medskip\noindent
最后验证 $M_0$ 共尾。$\mathcal{I}$ 的每个对象，都是只含该对象
及其恒等箭头的子范畴中的终对象，故函子 $M_0$ 在对象上满射。
特别地，定义 \ref{categories-definition-cofinal} 的条件 (1) 成立。
给定 $\mathcal{I}$ 的对象 $i$、$\mathcal{I}_0$ 的对象
$\mathcal{F}_1, \mathcal{F}_2$，以及 $\mathcal{I}$ 中的映射
$\varphi_1 : i \to M_0(\mathcal{F}_1)$ 与
$\varphi_2 : i \to M_0(\mathcal{F}_2)$，可取
$\mathcal{F}_{12}$ 为具有唯一终对象的有限生成范畴，
其中包含 $\mathcal{F}_1$、$\mathcal{F}_2$ 及态射
$\varphi_1, \varphi_2$。所得图表
$$
\xymatrix{
& M_0(\mathcal{F}_{12}) & \\
M_0(\mathcal{F}_{1}) \ar[ru] & & M_0(\mathcal{F}_{2}) \ar[lu] \\
& i \ar[lu] \ar[ru]
}
$$
交换，因为它位于范畴 $\mathcal{F}_{12}$ 中，而
$M_0(\mathcal{F}_{12})$ 是该范畴的终对象。
因此条件 (2) 也成立，证明完成。
\end{proof}

\begin{remark}
\label{categories-remark-trick-needed}
注意，有限有向集 $(I, \geq)$ 总有最大对象 $i_\infty$。
因此，这类集合上的任意系 $(M_i, f_{ii'})$ 的余极限都是平凡的，
其含义是该余极限等于 $M_{i_\infty}$。相比之下，
由有限滤过范畴编号的余极限不一定平凡。例如，令
$\mathcal{I}$ 为只有一个对象 $i$ 和一个满足
$e = e \circ e$ 的非平凡态射 $e$ 的范畴。
图表 $M : \mathcal{I} \to Sets$ 的余极限是幂等映射
$M(e)$ 的像。这说明，在引理
\ref{categories-lemma-directed-category-system} 的证明中，采用类似于传到
$\mathcal{I}\times \omega$ 的技巧是必要的。
\end{remark}

\begin{lemma}
\label{categories-lemma-nonempty-limit}
若 $S : \mathcal{I} \to \textit{Sets}$ 是集合的余滤过图表，
且所有 $S_i$ 都有限非空，则 $\lim_i S_i$ 非空。
换言之，有限非空集合的有向逆系的极限非空。
\end{lemma}

\begin{proof}
由引理 \ref{categories-lemma-directed-category-system}，两个陈述等价。
设 $I$ 为有向集，并设 $(S_i)_{i \in I}$ 为 $I$ 上有限非空集合的逆系。
称一个{\it 子系} $T$ 是非空子集族
$T = (T_i)_{i \in I}$、$T_i \subset S_i$，并满足：
对所有 $i' \geq i$，过渡映射 $S_{i'} \to S_i$ 都把
$T_{i'}$ 映入 $T_i$。以 $\mathcal{T}$ 表示子系所成的集合。
我们按包含关系排列 $\mathcal{T}$。设 $T_\alpha$、$\alpha \in A$
是 $\mathcal{T}$ 中元素的一个全序族，并写
$T_\alpha = (T_{\alpha, i})_{i \in I}$。
令 $T_i = \bigcap_{\alpha \in A} T_{\alpha, i}$，便得到下界
$T = (T_i)_{i \in I}$；显然 $T_i$ 是 $S_i$ 的有限非空子集，
因为所有 $T_{\alpha, i}$ 都非空且 $T_\alpha$ 构成全序族。
因此可应用 Zorn 引理，得知 $\mathcal{T}$ 有极小元。

\medskip\noindent
下面分析极小元 $T \in \mathcal{T}$ 的形状。
首先注意，所有映射 $T_{i'} \to T_i$ 都是满射。
事实上，由于 $I$ 是有向集且 $T_i$ 有限，交
$T'_i = \bigcap_{i' \geq i} \Im(T_{i'} \to T_i)$ 非空。
因此 $T' = (T'_i)$ 是包含于 $T$ 的子系；由极小性，
$T' = T$。最后断言，对每个 $i$，该集合都是单点集。
事实上，若 $x \in T_i$，则可在 $i' \geq i$ 时定义
$T'_{i'} = (T_{i'} \to T_i)^{-1}(\{x\})$，并在
$j \not \geq i$ 时令 $T'_j = T_j$。这又是一个子系，
因为上面已看到子系 $T$ 的过渡映射都是满射。
由极小性，$T = T'$，这的确蕴含 $T_i$ 是单点集。
这对每个 $i \in I$ 都成立，所以对某些 $x_i \in S_i$，
有 $T_i = \{x_i\}$，并且对每个 $i' \geq i$，
在映射 $S_{i'} \to S_i$ 下有 $x_{i'} \mapsto x_i$。
换言之，$(x_i) \in \lim S_i$，引理得证。
\end{proof}






\section{本质常值系}
\label{categories-section-essentially-constant}

\noindent
设 $M : \mathcal{I} \to \mathcal{C}$ 为范畴 $\mathcal{C}$ 中的图表，
并假设指标范畴 $\mathcal{I}$ 滤过。此时有三个逐次增强的概念，
它们选出 $\mathcal{C}$ 的一个对象 $X$。第一个概念仅仅是
$$
X = \colim_{i \in \mathcal{I}} M_i.
$$
此时 $X$ 配有余投影 $M_i \to X$。更强的条件是：要求 $X$
为余极限，并存在 $i \in \mathcal{I}$ 及态射 $X \to M_i$，
使复合 $X \to M_i \to X$ 为 $\text{id}_X$。
更强的条件如下。

\begin{definition}
\label{categories-definition-essentially-constant-diagram}
设 $M : \mathcal{I} \to \mathcal{C}$ 为范畴 $\mathcal{C}$ 中的图表。
\begin{enumerate}
\item 假设指标范畴 $\mathcal{I}$ 滤过，并设
$(X, \{M_i \to X\}_i)$ 为 $M$ 的余锥，见注
\ref{categories-remark-cones-and-cocones}。若存在 $i \in \mathcal{I}$
及态射 $X \to M_i$，使得
\begin{enumerate}
\item $X \to M_i \to X$ 是 $\text{id}_X$；
\item 对所有 $j$，都存在 $k$ 以及态射 $i \to k$、$j \to k$，
使态射 $M_j \to M_k$ 等于复合
$M_j \to X \to M_i \to M_k$，
\end{enumerate}
则称 $M$ 是以 $X$ 为{\it 值}的{\it 本质常值}图表。
\item 假设指标范畴 $\mathcal{I}$ 余滤过，并设
$(X, \{X \to M_i\}_i)$ 为 $M$ 的锥，见注
\ref{categories-remark-cones-and-cocones}。若存在 $i \in \mathcal{I}$
及态射 $M_i \to X$，使得
\begin{enumerate}
\item $X \to M_i \to X$ 是 $\text{id}_X$；
\item 对所有 $j$，都存在 $k$ 以及态射 $k \to i$、$k \to j$，
使态射 $M_k \to M_j$ 等于复合
$M_k \to M_i \to X \to M_j$，
\end{enumerate}
则称 $M$ 是以 $X$ 为{\it 值}的{\it 本质常值}图表。
\end{enumerate}
使用本定义时，请牢记引理
\ref{categories-lemma-essentially-constant-is-limit-colimit}。
\end{definition}

\noindent
上下文会表明所指的是两个版本中的哪一个。若可能混淆，
我们称第一个版本中的 $M$ 作为 {\it ind-对象}本质常值，
称第二个版本中的它作为 {\it pro-对象}本质常值。
注 \ref{categories-remark-ind-category} 与 \ref{categories-remark-pro-category}
会进一步说明这一术语。事实上，我们常使用“本质常值系”这一术语，
而严格说来，它只对有向集上的系作了定义。

\begin{definition}
\label{categories-definition-essentially-constant-system}
设 $\mathcal{C}$ 为范畴。若把有向系 $(M_i, f_{ii'})$ 中的 $M$
视为函子 $I \to \mathcal{C}$ 时，它定义一个本质常值图表，
则称其为{\it 本质常值系}。若把有向逆系 $(M_i, f_{ii'})$ 中的 $M$
视为函子 $I^{opp} \to \mathcal{C}$ 时，它定义一个
本质常值逆图表，则称其为{\it 本质常值逆系}。
\end{definition}

\noindent
若 $(M_i, f_{ii'})$ 是本质常值系，且态射 $f_{ii'}$ 是单态射，
则对充分大的 $i \leq i'$，态射 $f_{ii'}$ 都是同构。
另一方面，考虑系
$$
\mathbf{Z}^2 \to \mathbf{Z}^2 \to \mathbf{Z}^2 \to \ldots
$$
其中映射由 $(a, b) \mapsto (a + b, 0)$ 给出。
该系以 $\mathbf{Z}$ 为值本质常值，但每个过渡映射都有核。

\medskip\noindent
下面给出一个非本质常值系的例子。令
$M = \bigoplus_{n \geq 0} \mathbf{Z}$，并令 $S : M \to M$
为移位算子 $(a_0, a_1, \ldots) \mapsto (a_1, a_2, \ldots)$。
此时，以 $S$ 为过渡映射的系 $M \to M \to M \to \ldots$
具有余极限 $0$，且复合 $0 \to M \to 0$ 是恒等态射，
但该系并非本质常值。

\medskip\noindent
下一个引理是一项合理性检验。

\begin{lemma}
\label{categories-lemma-essentially-constant-is-limit-colimit}
设 $M : \mathcal{I} \to \mathcal{C}$ 为图表。
若 $\mathcal{I}$ 滤过，且 $M$ 作为 ind-对象本质常值，
则 $X = \colim M_i$ 存在，并且 $M$ 以 $X$ 为值本质常值。
若 $\mathcal{I}$ 余滤过，且 $M$ 作为 pro-对象本质常值，
则 $X = \lim M_i$ 存在，并且 $M$ 以 $X$ 为值本质常值。
\end{lemma}

\begin{proof}
略。这是对各项定义的一道好练习。
\end{proof}

\begin{remark}
\label{categories-remark-ind-category}
设 $\mathcal{C}$ 为范畴。存在由 $\mathcal{C}$ 的
{\it ind-对象}构成的大范畴 $\text{Ind-}\mathcal{C}$。
具体地，若 $F : \mathcal{I} \to \mathcal{C}$ 与
$G : \mathcal{J} \to \mathcal{C}$ 是 $\mathcal{C}$ 中的滤过图表，
则可定义
$$
\Mor_{\text{Ind-}\mathcal{C}}(F, G) =
\lim_i \colim_j \Mor_\mathcal{C}(F(i), G(j)).
$$
存在典范函子 $\mathcal{C} \to \text{Ind-}\mathcal{C}$，
它把 $X$ 映到取常值 $X$ 的{\it 常值系}。
这是一个全忠实嵌入。用这种语言可知，图表 $F$ 本质常值，
当且仅当 $F$ 同构于一个常值系。若以后需要这些内容，
我们会在此将它表述成一个引理并加以证明。
\end{remark}

\begin{remark}
\label{categories-remark-pro-category}
设 $\mathcal{C}$ 为范畴。存在由 $\mathcal{C}$ 的
{\it pro-对象}构成的大范畴 $\text{Pro-}\mathcal{C}$。
具体地，若 $F : \mathcal{I} \to \mathcal{C}$ 与
$G : \mathcal{J} \to \mathcal{C}$ 是 $\mathcal{C}$ 中的余滤过图表，
则可定义
$$
\Mor_{\text{Pro-}\mathcal{C}}(F, G) =
\lim_j \colim_i \Mor_\mathcal{C}(F(i), G(j)).
$$
存在典范函子 $\mathcal{C} \to \text{Pro-}\mathcal{C}$，
它把 $X$ 映到取常值 $X$ 的{\it 常值系}。
这是一个全忠实嵌入。用这种语言可知，图表 $F$ 本质常值，
当且仅当 $F$ 同构于一个常值系。若以后需要这些内容，
我们会在此将它表述成一个引理并加以证明。
\end{remark}

\begin{example}
\label{categories-example-pro-morphism-inverse-systems}
设 $\mathcal{C}$ 为范畴，并设 $(X_n)$ 与 $(Y_n)$ 为
$\mathcal{C}$ 中在通常序关系下以 $\mathbf{N}$ 为指标的逆系。图示为
$$
\ldots \to X_3 \to X_2 \to X_1
\quad\text{且}\quad
\ldots \to Y_3 \to Y_2 \to Y_1
$$
设 $a : (X_n) \to (Y_n)$ 为 $\mathcal{C}$ 的 pro-对象的态射。
$a$ 包含什么数据？对每个 $n \in \mathbf{N}$，应存在一个
$m(n)$ 及态射 $a_n : X_{m(n)} \to Y_n$。这些态射应在下述意义下相容：
对所有 $n' \geq n$，存在
$m(n', n) \geq m(n'), m(n)$，使图表
$$
\xymatrix{
X_{m(n, n')} \ar[rr] \ar[d] & & X_{m(n)} \ar[d]^{a_n} \\
X_{m(n')} \ar[r]^{a_{n'}} & Y_{n'} \ar[r] & Y_n
}
$$
交换。把 $m(n)$ 替换为
$\max_{k, l \leq n}\{m(n, k), m(k, l)\}$ 后，可得到
$\ldots \geq m(3) \geq m(2) \geq m(1)$ 以及交换图表
$$
\xymatrix{
\ldots \ar[r] &
X_{m(3)} \ar[d]^{a_3} \ar[r] &
X_{m(2)} \ar[d]^{a_2} \ar[r] &
X_{m(1)} \ar[d]^{a_1} \\
\ldots \ar[r] &
Y_3 \ar[r] &
Y_2 \ar[r] &
Y_1
}
$$
给定满足 $m' \geq m$ 的递增映射
$m' : \mathbf{N} \to \mathbf{N}$，并令
$a'_i : X_{m'(i)} \to X_{m(i)} \to Y_i$，则序对
$(m', a')$ 定义同一个 pro-系态射。反过来，给定如上的两个序对
$(m_1, a_1)$ 与 $(m_1, a_2)$，它们定义同一个 pro-对象态射，
当且仅当存在 $m' \geq m_1, m_2$，使得 $a'_1 = a'_2$。
\end{example}

\begin{remark}
\label{categories-remark-pro-category-copresheaves}
设 $\mathcal{C}$ 为范畴，并设 $F : \mathcal{I} \to \mathcal{C}$、
$G : \mathcal{J} \to \mathcal{C}$ 为 $\mathcal{C}$ 中的余滤过图表。
考虑由下式定义的函子 $A, B : \mathcal{C} \to \textit{Sets}$：
$$
A(X) = \colim_i \Mor_\mathcal{C}(F(i), X)
\quad\text{且}\quad
B(X) = \colim_j \Mor_\mathcal{C}(G(j), X)
$$
我们断言，从 $F$ 到 $G$ 的 pro-系态射等同于函子变换
$t : B \to A$。具体地，给定 $t$，可把 $t$ 作用于
$B(G(j))$ 中 $\text{id}_{G(j)}$ 的类，得到元素的相容系
$\xi_j \in A(G(j)) = \colim_i \Mor_\mathcal{C}(F(i), G(j))$；
这恰是注 \ref{categories-remark-pro-category} 中
$\text{Pro-}\mathcal{C}$ 的态射定义。给定一个从 $F$ 到 $G$
的 pro-对象态射，如何构造变换 $B \to A$，这里从略。
\end{remark}

\begin{lemma}
\label{categories-lemma-image-essentially-constant}
设 $\mathcal{C}$ 为范畴，并设 $M : \mathcal{I} \to \mathcal{C}$
为指标范畴 $\mathcal{I}$ 滤过（相应地，余滤过）的图表。
设 $F : \mathcal{C} \to \mathcal{D}$ 为函子。
若 $M$ 作为 ind-对象（相应地，pro-对象）本质常值，
则 $F \circ M : \mathcal{I} \to \mathcal{D}$ 亦然。
\end{lemma}

\begin{proof}
若 $X$ 是 $M$ 的一个值，则由定义立即可知，$F(X)$ 是
$F \circ M$ 的一个值。
\end{proof}

\begin{lemma}
\label{categories-lemma-characterize-essentially-constant-ind}
设 $\mathcal{C}$ 为范畴，并设 $M : \mathcal{I} \to \mathcal{C}$
为指标范畴 $\mathcal{I}$ 滤过的图表。下列条件等价：
\begin{enumerate}
\item $M$ 是本质常值 ind-对象；
\item 存在余锥 $(X, \{M_i \to X\}_i)$，使得对
$\mathcal{C}$ 中任意 $W$，映射
$\colim_i \Mor_\mathcal{C}(W, M_i) \to \Mor_\mathcal{C}(W, X)$
为双射；
\item $X = \colim_i M_i$ 存在，并且对 $\mathcal{C}$ 中任意 $W$，
映射 $\colim_i \Mor_\mathcal{C}(W, M_i) \to \Mor_\mathcal{C}(W, X)$
为双射；
\item 存在 $\mathcal{I}$ 中的 $i$ 以及态射 $X \to M_i$，
使得对 $\mathcal{C}$ 中任意 $W$，映射 $\Mor_\mathcal{C}(W, X) \to
\colim_{j \in \mathcal{I}} \Mor_\mathcal{C}(W, M_j)$
为双射。
\end{enumerate}
在情形 (2)、(3)、(4) 中，本质常值系的值为 $X$。
\end{lemma}

\begin{proof}
显然 (3) 蕴含 (2)。假设 (2)。于是
$\text{id}_X \in \Mor_\mathcal{C}(X, X)$ 来自某个 $i$ 所对应的态射
$X \to M_i$，即 $X \to M_i \to X$ 是恒等态射。
于是对所有 $W$，两个映射
$$
\Mor_\mathcal{C}(W, X)
\longrightarrow
\colim_i \Mor_\mathcal{C}(W, M_i)
\longrightarrow
\Mor_\mathcal{C}(W, X)
$$
都是双射；其中第一个映射由上面找到的态射 $X \to M_i$ 诱导，
而二者的复合是恒等映射。这意味着复合
$$
\colim_i \Mor_\mathcal{C}(W, M_i)
\longrightarrow
\Mor_\mathcal{C}(W, X)
\longrightarrow
\colim_i \Mor_\mathcal{C}(W, M_i)
$$
也为恒等映射。令 $W = M_j$，并从余极限中的
$\text{id}_{M_j}$ 出发，可知对某个充分大的 $k$，
$M_j \to X \to M_i \to M_k$ 等于 $M_j \to M_k$。
这证明 (1) 成立。

\medskip\noindent
假设 (4)。设 $k$ 为 $\mathcal{I}$ 的对象。令 $W = M_k$，
可推出存在唯一态射 $M_k \to X$，并且存在 $j$ 及
$\mathcal{I}$ 中的态射 $k \to j$、$i \to j$，使
$M_k \to X \to M_i \to M_j$ 等于 $M_k \to M_j$。
唯一性保证得到余锥 $(X, \{M_k \to X\})$。
由此可见 (4) 蕴含 (2)；略去若干细节。

\medskip\noindent
我们略去 (1) 蕴含其他条件的证明。
\end{proof}

\begin{lemma}
\label{categories-lemma-characterize-essentially-constant-pro}
设 $\mathcal{C}$ 为范畴，并设 $M : \mathcal{I} \to \mathcal{C}$
为指标范畴 $\mathcal{I}$ 余滤过的图表。下列条件等价：
\begin{enumerate}
\item $M$ 是本质常值 pro-对象；
\item 存在锥 $(X, \{X \to M_i\})$，使得对
$\mathcal{C}$ 中任意 $W$，映射
$\colim_{i \in \mathcal{I}^{opp}} \Mor_\mathcal{C}(M_i, W) \to
\Mor_\mathcal{C}(X, W)$ 为双射；
\item $X = \lim_i M_i$ 存在，并且对 $\mathcal{C}$ 中任意 $W$，
映射
$\colim_{i \in \mathcal{I}^{opp}} \Mor_\mathcal{C}(M_i, W) \to
\Mor_\mathcal{C}(X, W)$ 为双射；
\item 存在 $\mathcal{I}$ 中的 $i$ 以及态射 $M_i \to X$，
使得对 $\mathcal{C}$ 中任意 $W$，映射 $\Mor_\mathcal{C}(X, W) \to
\colim_{j \in \mathcal{I}^{opp}} \Mor_\mathcal{C}(M_j, W)$ 为双射。
\end{enumerate}
在情形 (2)、(3)、(4) 中，本质常值系的值为 $X$。
\end{lemma}

\begin{proof}
本引理与引理 \ref{categories-lemma-characterize-essentially-constant-ind} 对偶。
\end{proof}

\begin{lemma}
\label{categories-lemma-cofinal-essentially-constant}
设 $\mathcal{C}$ 为范畴，$H : \mathcal{I} \to \mathcal{J}$
为滤过指标范畴之间的函子。若 $H$ 共尾，则任意图表
$M : \mathcal{J} \to \mathcal{C}$ 本质常值，当且仅当
$M \circ H$ 本质常值。
\end{lemma}

\begin{proof}
这由引理 \ref{categories-lemma-characterize-essentially-constant-ind}
与引理 \ref{categories-lemma-cofinal} 形式地推出。
\end{proof}

\begin{lemma}
\label{categories-lemma-essentially-constant-over-product}
设 $\mathcal{I}$ 与 $\mathcal{J}$ 为滤过范畴，并以
$p : \mathcal{I} \times \mathcal{J} \to \mathcal{J}$ 表示投影。
则 $\mathcal{I} \times \mathcal{J}$ 滤过，而且图表
$M : \mathcal{J} \to \mathcal{C}$ 本质常值，当且仅当
$M \circ p : \mathcal{I} \times \mathcal{J} \to \mathcal{C}$
本质常值。
\end{lemma}

\begin{proof}
我们略去 $\mathcal{I} \times \mathcal{J}$ 滤过的验证。
由于 $p$ 共尾（验证略），所述等价性由引理
\ref{categories-lemma-cofinal-essentially-constant} 得到。
\end{proof}

\begin{lemma}
\label{categories-lemma-initial-essentially-constant}
设 $\mathcal{C}$ 为范畴，$H : \mathcal{I} \to \mathcal{J}$
为余滤过指标范畴之间的函子。若 $H$ 是始的，则任意图表
$M : \mathcal{J} \to \mathcal{C}$ 本质常值，当且仅当
$M \circ H$ 本质常值。
\end{lemma}

\begin{proof}
这由引理 \ref{categories-lemma-characterize-essentially-constant-pro}、
\ref{categories-lemma-initial}、\ref{categories-lemma-cofinal} 以及下述事实形式地推出：
若 $\mathcal{I}$ 在 $\mathcal{J}$ 中是始的，则
$\mathcal{I}^{opp}$ 在 $\mathcal{J}^{opp}$ 中共尾。
\end{proof}





\section{正合函子}
\label{categories-section-exact-functor}

\noindent
本节定义正合函子。

\begin{definition}
\label{categories-definition-exact}
设 $F : \mathcal{A} \to \mathcal{B}$ 为函子。
\begin{enumerate}
\item 假设 $\mathcal{A}$ 中所有有限极限都存在。
若 $F$ 与所有有限极限交换，则称其是{\it 左正合}的。
\item 假设 $\mathcal{A}$ 中所有有限余极限都存在。
若 $F$ 与所有有限余极限交换，则称其是{\it 右正合}的。
\item 若 $F$ 同时左正合且右正合，则称其是{\it 正合}的。
\end{enumerate}
\end{definition}

\begin{lemma}
\label{categories-lemma-characterize-left-exact}
设 $F : \mathcal{A} \to \mathcal{B}$ 为函子。假设 $\mathcal{A}$ 中
所有有限极限都存在，见引理 \ref{categories-lemma-finite-limits-exist}。
下列条件等价：
\begin{enumerate}
\item $F$ 左正合；
\item $F$ 与有限积和等化子交换；
\item $F$ 把 $\mathcal{A}$ 的终对象变为 $\mathcal{B}$ 的终对象，
且与纤维积交换。
\end{enumerate}
\end{lemma}

\begin{proof}
引理 \ref{categories-lemma-limits-products-equalizers} 表明 (2) 蕴含 (1)。
假设 (3) 成立。终对象上的纤维积就是积。若
$a, b : A \to B$ 是 $\mathcal{A}$ 的态射，则 $a, b$ 的等化子为
$$
(A \times_{a, B, b} A)\times_{(\text{pr}_1, \text{pr}_2), A \times A, \Delta} A.
$$
故 (3) 蕴含 (2)。最后，(1) 蕴含 (3)，因为空极限是终对象，
而纤维积是极限。
\end{proof}

\begin{lemma}
\label{categories-lemma-characterize-right-exact}
设 $F : \mathcal{A} \to \mathcal{B}$ 为函子。假设 $\mathcal{A}$ 中
所有有限余极限都存在，见引理 \ref{categories-lemma-colimits-exist}。
下列条件等价：
\begin{enumerate}
\item $F$ 右正合；
\item $F$ 与有限余积和余等化子交换；
\item $F$ 把 $\mathcal{A}$ 的始对象变为 $\mathcal{B}$ 的始对象，
且与推出交换。
\end{enumerate}
\end{lemma}

\begin{proof}
这是引理 \ref{categories-lemma-characterize-left-exact} 的对偶。
\end{proof}




\section{伴随函子}
\label{categories-section-adjoint}

\begin{definition}
\label{categories-definition-adjoint}
设 $\mathcal{C}$、$\mathcal{D}$ 为范畴，并设
$u : \mathcal{C} \to \mathcal{D}$ 与
$v : \mathcal{D} \to \mathcal{C}$ 为函子。若存在双射
$$
\Mor_\mathcal{D}(u(X), Y)
\longrightarrow
\Mor_\mathcal{C}(X, v(Y))
$$
且它关于 $X \in \Ob(\mathcal{C})$ 与
$Y \in \Ob(\mathcal{D})$ 具有函子性，则称 $u$ 是 $v$ 的
{\it 左伴随}，或称 $v$ 是 $u$ 的{\it 右伴随}。
\end{definition}

\noindent
换言之，这表示给定了从 $\Mor_\mathcal{D}(u(-), -)$ 到
$\Mor_\mathcal{C}(-, v(-))$ 的函子同构；这些函子定义在
$\mathcal{C}^{opp} \times \mathcal{D} \to \textit{Sets}$ 上。
对 $\mathcal{C}$ 的任意对象 $X$，由 $\text{id}_{u(X)}$ 对应得到态射
$X \to v(u(X))$。类似地，对 $\mathcal{D}$ 的任意对象 $Y$，
由 $\text{id}_{v(Y)}$ 对应得到态射 $u(v(Y)) \to Y$。
这些映射称为{\it 伴随映射}。伴随映射关于 $X$ 与 $Y$ 具有函子性，
因而得到函子态射
$$
\eta : \text{id}_\mathcal{C} \to v \circ u\quad (\text{单位})
\quad\text{且}\quad
\epsilon : u \circ v \to \text{id}_\mathcal{D}\quad (\text{余单位}).
$$
此外，若 $\alpha : u(X) \to Y$ 与 $\beta : X \to v(Y)$ 为态射，
则下列条件等价：
\begin{enumerate}
\item $\alpha$ 与 $\beta$ 通过定义中的双射相互对应；
\item $\beta$ 是复合 $X \to v(u(X)) \xrightarrow{v(\alpha)} v(Y)$；
\item $\alpha$ 是复合 $u(X) \xrightarrow{u(\beta)} u(v(Y)) \to Y$。
\end{enumerate}
这样便可用伴随映射重新表述伴随函子的概念。

\begin{lemma}
\label{categories-lemma-adjoint-exists}
设 $u : \mathcal{C} \to \mathcal{D}$ 为范畴之间的函子。
若对每个 $y \in \Ob(\mathcal{D})$，函子
$x \mapsto \Mor_\mathcal{D}(u(x), y)$ 都可表，则 $u$ 有右伴随。
\end{lemma}

\begin{proof}
对每个 $y$，选择对象 $v(y)$ 以及函子同构
$\Mor_\mathcal{C}(-, v(y)) \to \Mor_\mathcal{D}(u(-), y)$。
由 Yoneda 引理（引理 \ref{categories-lemma-yoneda}），对任意态射
$g : y \to y'$，函子变换
$$
\Mor_\mathcal{C}(-, v(y)) \to \Mor_\mathcal{D}(u(-), y) \to
\Mor_\mathcal{D}(u(-), y') \to \Mor_\mathcal{C}(-, v(y'))
$$
对应于唯一态射 $v(g) : v(y) \to v(y')$。我们略去 $v$ 为函子且
为 $u$ 的右伴随的验证。
\end{proof}

\begin{lemma}
\label{categories-lemma-left-adjoint-composed-fully-faithful}
\begin{reference}
Bhargav Bhatt，私人通信。
\end{reference}
设 $u$ 是定义 \ref{categories-definition-adjoint} 中 $v$ 的左伴随。
\begin{enumerate}
\item 若 $v \circ u$ 全忠实，则 $u$ 全忠实。
\item 若 $u \circ v$ 全忠实，则 $v$ 全忠实。
\end{enumerate}
\end{lemma}

\begin{proof}
证明 (2)。假设 $u \circ v$ 全忠实。设 $X$、$Y$ 是
$\mathcal{D}$ 中的对象。则自然的复合映射
$$
\Mor(X,Y) \to \Mor(v(X),v(Y)) \to \Mor(u(v(X)), u(v(Y)))
$$
为双射，所以 $v$ 至少忠实。为证明全忠实性，只需证明上面的
第二个映射为单射。但 $u$ 与 $v$ 之间的伴随说明
$$
\Mor(v(X), v(Y)) \to \Mor(u(v(X)), u(v(Y))) \to \Mor(u(v(X)), Y)
$$
为双射；其中第一个映射是自然映射，第二个映射来自伴随的余单位
$u(v(Y)) \to Y$。因此
$\Mor(v(X), v(Y)) \to \Mor(u(v(X)), u(v(Y)))$
也为单射，正如所需。(1) 的证明与此对偶。
\end{proof}

\begin{lemma}
\label{categories-lemma-adjoint-fully-faithful}
设 $u$ 是定义 \ref{categories-definition-adjoint} 中 $v$ 的左伴随。则
\begin{enumerate}
\item $u$ 全忠实 $\Leftrightarrow$ $\text{id} \cong v \circ u$
$\Leftrightarrow$ $\eta : \text{id} \to v \circ u$ 是同构；
\item $v$ 全忠实 $\Leftrightarrow$
$u \circ v \cong \text{id}$ $\Leftrightarrow$
$\epsilon : u \circ v \to \text{id}$ 是同构。
\end{enumerate}
\end{lemma}

\begin{proof}
证明 (1)。假设 $u$ 全忠实。我们证明
$\eta_X : X \to v(u(X))$ 是同构。任取态射
$X' \to v(u(X))$。由伴随性，它对应于态射 $u(X') \to u(X)$；
由 $u$ 的全忠实性，它又对应于唯一态射 $X' \to X$。
因此，与 $\eta_X$ 后复合给出双射
$\Mor(X', X) \to \Mor(X', v(u(X)))$，故 $\eta_X$ 为同构。
若存在函子同构 $\text{id} \cong v \circ u$，则 $v \circ u$ 全忠实。
由引理 \ref{categories-lemma-left-adjoint-composed-fully-faithful}，$u$ 全忠实；
由上述论证，这蕴含 $\eta$ 为同构。因此 $3$ 个条件等价
（这些条件也等价于 $v \circ u$ 全忠实）。

\medskip\noindent
(2) 与 (1) 对偶。
\end{proof}

\begin{lemma}
\label{categories-lemma-adjoint-exact}
设 $u$ 是定义 \ref{categories-definition-adjoint} 中 $v$ 的左伴随。
\begin{enumerate}
\item 假设 $M : \mathcal{I} \to \mathcal{C}$ 为图表，且
$\colim_\mathcal{I} M$ 在 $\mathcal{C}$ 中存在。则
$u(\colim_\mathcal{I} M) = \colim_\mathcal{I} u \circ M$。
换言之，$u$ 与（可表的）余极限交换。
\item 假设 $M : \mathcal{I} \to \mathcal{D}$ 为图表，且
$\lim_\mathcal{I} M$ 在 $\mathcal{D}$ 中存在。则
$v(\lim_\mathcal{I} M) = \lim_\mathcal{I} v \circ M$。
换言之，$v$ 与可表极限交换。
\end{enumerate}
\end{lemma}

\begin{proof}
从余极限到一个对象的态射，等价于从该极限各组成对象到该对象的
相容态射系，见注 \ref{categories-remark-limit-colim}。因此
$$
\begin{matrix}
\Mor_\mathcal{D}(u(\colim_{i \in \mathcal{I}} M_i), Y) &
= & \Mor_\mathcal{C}(\colim_{i \in \mathcal{I}} M_i, v(Y)) \\
& = &
\lim_{i \in \mathcal{I}^{opp}}
\Mor_\mathcal{C}(M_i, v(Y)) \\
& = &
\lim_{i \in \mathcal{I}^{opp}}
\Mor_\mathcal{D}(u(M_i), Y)
\end{matrix}
$$
证明了 $u(\colim_{i \in \mathcal{I}} M_i)$ 就是所求余极限。
另一断言可类似证明。
\end{proof}

\begin{lemma}
\label{categories-lemma-exact-adjoint}
设 $u$ 是定义 \ref{categories-definition-adjoint} 中 $v$ 的左伴随。
\begin{enumerate}
\item 若 $\mathcal{C}$ 有有限余极限，则 $u$ 右正合。
\item 若 $\mathcal{D}$ 有有限极限，则 $v$ 左正合。
\end{enumerate}
\end{lemma}

\begin{proof}
这由定义与引理 \ref{categories-lemma-adjoint-exact} 立即得到。
\end{proof}

\begin{lemma}
\label{categories-lemma-unit-counit-relations}
设 $u : \mathcal{C} \to \mathcal{D}$ 是函子
$v : \mathcal{D} \to \mathcal{C}$ 的左伴随。设
$\eta_X : X \to v(u(X))$ 为单位，
$\epsilon_Y : u(v(Y)) \to Y$ 为余单位。则
$$
u(X) \xrightarrow{u(\eta_X)} u(v(u(X))
\xrightarrow{\epsilon_{u(X)}} u(X)
\quad\text{且}\quad
v(Y) \xrightarrow{\eta_{v(Y)}} v(u(v(Y))) \xrightarrow{v(\epsilon_Y)}
v(Y)
$$
都是恒等态射。
\end{lemma}

\begin{proof}
略。
\end{proof}

\begin{lemma}
\label{categories-lemma-transformation-between-functors-and-adjoints}
设 $u_1, u_2 : \mathcal{C} \to \mathcal{D}$ 为具有右伴随
$v_1, v_2 : \mathcal{D} \to \mathcal{C}$ 的函子。
设 $\beta : u_2 \to u_1$ 为函子变换，并设
$\beta^\vee : v_1 \to v_2$ 为相应的伴随函子变换。则
$$
\xymatrix{
u_2 \circ v_1 \ar[r]_\beta \ar[d]_{\beta^\vee} &
u_1 \circ v_1 \ar[d] \\
u_2 \circ v_2 \ar[r] & \text{id}
}
$$
可交换；其中未标记的箭头是余单位变换。
\end{lemma}

\begin{proof}
这是因为 $\beta^\vee_D : v_1D \to v_2D$ 是唯一满足下述条件的态射：
由它诱导的映射 $\Mor(C, v_1D) \to \Mor(C, v_2D)$，正是由
$\beta_C : u_2C \to u_1C$ 诱导的映射
$\Mor(u_1C, D) \to \Mor(u_2C, D)$。也就是说，由
$\beta_{v_1 D}$ 诱导的映射
$$
\Mor(u_1 v_1 D, D') \to \Mor(u_2 v_1 D, D')
$$
与由 $\beta^\vee_{D'}$ 诱导的映射
$$
\Mor(v_1 D, v_1 D') \to \Mor(v_1 D, v_2 D')
$$
相同。取 $D' = D$，可知余单位 $u_1 v_1 D \to D$ 与
$\beta_{v_1D}$ 前复合后，在伴随下对应于 $\beta^\vee_D$。
这恰好说明该图表在 $D$ 上取值后可交换。
\end{proof}

\begin{lemma}
\label{categories-lemma-compose-counits}
设 $\mathcal{A}$、$\mathcal{B}$、$\mathcal{C}$ 为范畴。设
$v : \mathcal{A} \to \mathcal{B}$ 与
$v' : \mathcal{B} \to \mathcal{C}$ 为函子，其左伴随分别为
$u$ 与 $u'$。则
\begin{enumerate}
\item 函子 $v'' = v' \circ v$ 的左伴随为
$u'' = u \circ u'$。
\item 给定 $\mathcal{A}$ 中的 $X$，有
\begin{equation}
\label{categories-equation-compose-counits}
\epsilon_X^v \circ u(\epsilon^{v'}_{v(X)}) = \epsilon^{v''}_X :
u''(v''(X)) \to X
\end{equation}
其中 $\epsilon$ 是各伴随的余单位。
\end{enumerate}
\end{lemma}

\begin{proof}
我们展开 (2) 中的公式，因为这也会立即证明 (1)。首先，伴随对
$(u, v)$ 与 $(u', v')$ 的余单位分别为映射
$\epsilon_X^v : u(v(X)) \to X$ 与
$\epsilon_Y^{v'} : u'(v'(Y)) \to Y$，见定义
\ref{categories-definition-adjoint} 后的讨论。以 (1) 中的 $u''$ 与 $v''$ 展开一切，
$$
u''(v''(X)) = u(u'(v'(v(X)))) \xrightarrow{u(\epsilon_{v(X)}^{v'})}
u(v(X)) \xrightarrow{\epsilon_X^v} X
$$
便得到 (\ref{categories-equation-compose-counits}) 左边的映射。
暂将其记为 $\epsilon_X^{v''}$。要说明它是伴随对
$(u'', v'')$ 的余单位，须证明：给定 $\mathcal{C}$ 中的 $Z$，
把态射 $\beta : Z \to v''(X)$ 映为
$\alpha = \epsilon_X^{v''} \circ u''(\beta) : u''(Z) \to X$
的规则在态射集合上为双射。这确实成立，因为该规则是下述两个规则的复合：
先把 $\beta$ 映为 $\epsilon_{v(X)}^{v'} \circ u'(\beta)$，
这由对 $(u', v')$ 的假设是双射；再把它映为
$\epsilon_X^v \circ u(\epsilon_{v(X)}^{v'} \circ u'(\beta))$，
这由对 $(u, v)$ 的假设也是双射。
\end{proof}






\section{一个可表性判据}
\label{categories-section-representable}

\noindent
下列引理常用于证明大范畴中泛对象的存在性；请参见注
\ref{categories-remark-how-to-use-it} 中的讨论。

\begin{lemma}
\label{categories-lemma-a-version-of-brown}
设 $\mathcal{C}$ 为有极限的大\footnote{见注
\ref{categories-remark-big-categories}。}范畴，并设
$F : \mathcal{C} \to \textit{Sets}$ 为函子。假设：
\begin{enumerate}
\item $F$ 与极限交换；
\item 存在 $\mathcal{C}$ 的一族对象 $\{x_i\}_{i \in I}$，且对每个
$i \in I$ 有元素 $f_i \in F(x_i)$，使得对
$y \in \Ob(\mathcal{C})$ 与 $g \in F(y)$，都存在 $i$ 及态射
$\varphi : x_i \to y$，满足 $F(\varphi)(f_i) = g$。
\end{enumerate}
则 $F$ 可表；也就是说，存在 $\mathcal{C}$ 的对象 $x$，使得
$$
F(y) = \Mor_\mathcal{C}(x, y)
$$
关于 $y$ 具有函子性。
\end{lemma}

\begin{proof}
令 $\mathcal{I}$ 为以下范畴：其对象是偶对 $(x_i, f_i)$，
而态射 $(x_i, f_i) \to (x_{i'}, f_{i'})$ 是
$\mathcal{C}$ 中满足 $F(\varphi)(f_i) = f_{i'}$ 的映射
$\varphi : x_i \to x_{i'}$。令
$$
x = \lim_{(x_i, f_i) \in \mathcal{I}} x_i
$$
（这并不是我们要找的 $x$，见下文）。由假设，该极限存在。
由于 $F$ 与极限交换，有
$$
F(x) = \lim_{(x_i, f_i) \in \mathcal{I}} F(x_i).
$$
因此存在泛元素 $f \in F(x)$，在对投影映射 $x \to x_i$
应用 $F$ 后，它映到 $f_i \in F(x_i)$。利用 $f$，得到函子变换
$$
\xi : \Mor_\mathcal{C}(x, - ) \longrightarrow F(-)
$$
（见第 \ref{categories-section-opposite} 节）。任取 $\mathcal{C}$ 的对象 $y$
及 $g \in F(y)$。由假设可选取 $x_i \to y$，使 $f_i$ 映到 $g$。
于是对映射
$$
x \longrightarrow x_i \longrightarrow y
$$
应用 $F$（第一个映射是定义 $x$ 的极限投影），会把 $f$ 映到 $g$。
故变换 $\xi$ 为满射。

\medskip\noindent
为了找出表示 $F$ 的对象，令 $e : x' \to x$ 为所有满足
$F(\varphi)(f) = f$ 的自映射 $\varphi : x \to x$ 的等化子。
由于 $F$ 与极限交换，它也与等化子交换，因而存在
$f' \in F(x')$ 映到 $F(x)$ 中的 $f$。由于 $\xi$ 为满射且
$f'$ 映到 $f$，可知 $\xi' : \Mor_\mathcal{C}(x', -) \to F(-)$
也为满射。最后，假设 $a, b : x' \to y$ 是两个满足
$F(a)(f') = F(b)(f')$ 的映射。我们须证明 $a = b$。
考虑等化子 $e' : x'' \to x'$。同样可找到映到 $f'$ 的
$f'' \in F(x'')$。选取映射 $\psi : x \to x''$，使
$F(\psi)(f) = f''$。于是
$e \circ e' \circ \psi : x \to x$ 是满足
$F(e \circ e' \circ \psi)(f) = f$ 的态射。因此
$e \circ e' \circ \psi \circ e = e$。由于 $e$ 是单态射，
这蕴含 $e'$ 是满态射，从而如所需有 $a = b$。
\end{proof}

\begin{remark}
\label{categories-remark-how-to-use-it}
上面的引理常用于构造某物上的自由某物，例如集合上的自由阿贝尔群、
集合上的自由群等等。以集合 $E$ 上的自由群为例，其想法是考虑函子
$$
F : \textit{Groups} \to \textit{Sets},\quad
G \longmapsto \text{Map}(E, G)
$$
该函子与极限交换。作为对象族，可以取由群 $G_i$ 及集合映射
$E \to G_i$ 组成的一族 $E \to G_i$，其中这些群的基数至多为
$\max(\aleph_0, |E|)$，并使这种结构的每个同构类至少出现一次。
事实上，若 $E \to G$ 是从 $E$ 到群 $G$ 的映射，则由其像生成的
子群 $G'$ 的基数至多为 $\max(\aleph_0, |E|)$。该引理说明函子可表，
故存在群 $F_E$，使
$\Mor_{\textit{Groups}}(F_E, G) = \text{Map}(E, G)$。
特别地，$F_E$ 的恒等态射对应于一个映射 $E \to F_E$，并且可证明
$F_E$ 由其像生成而不施加任何关系。

\medskip\noindent
另一种典型应用是：一旦已知极限存在，就可用该引理构造余极限。
我们以拓扑空间范畴为例；由拓扑，引理 \ref{topology-lemma-limits}，
该范畴有极限。具体地，假设
$\mathcal{I} \to \textit{Top}$，$i \mapsto X_i$ 是函子。
则可考虑
$$
F : \textit{Top} \longrightarrow \textit{Sets},\quad
Y \longmapsto \lim_\mathcal{I} \Mor_{\textit{Top}}(X_i, Y)
$$
该函子与极限交换。此外，给定任意拓扑空间 $Y$ 及
$F(Y)$ 的元素 $(\varphi_i : X_i \to Y)$，存在基数至多为
$|\coprod X_i|$ 的子空间 $Y' \subset Y$，使态射 $\varphi_i$
都映入 $Y'$。事实上，可在各 $\varphi_i$ 的像之并上取诱导拓扑。
所以显然满足引理的假设，从而得到表示函子 $F$ 的拓扑空间 $X$；
这恰好意味着 $X$ 是图表 $i \mapsto X_i$ 的余极限。
\end{remark}

\begin{theorem}[伴随函子定理]
\label{categories-theorem-adjoint-functor}
设 $G : \mathcal{C} \to \mathcal{D}$ 为大范畴之间的函子。
假设 $\mathcal{C}$ 有极限，$G$ 与极限交换，并且对
$\mathcal{D}$ 的每个对象 $y$，存在一组偶对
$(x_i, f_i)_{i \in I}$，其中 $x_i \in \Ob(\mathcal{C})$、
$f_i \in \Mor_\mathcal{D}(y, G(x_i))$；它满足：对任意偶对
$(x, f)$，其中 $x \in \Ob(\mathcal{C})$、
$f \in \Mor_\mathcal{D}(y, G(x))$，都存在 $i$ 及态射
$h : x_i \to x$，使 $f = G(h) \circ f_i$。则 $G$ 有左伴随 $F$。
\end{theorem}

\begin{proof}
这些假设蕴含：对 $\mathcal{D}$ 的每个对象 $y$，函子
$x \mapsto \Mor_\mathcal{D}(y, G(x))$ 都满足引理
\ref{categories-lemma-a-version-of-brown} 的假设。因此它由某个对象表示，
将该对象记为 $F(y)$。应用 Yoneda 引理（引理 \ref{categories-lemma-yoneda}），
规则 $y \mapsto F(y)$ 变成一个函子；由构造，它是 $G$ 的伴随。
细节略去。
\end{proof}





\section{范畴紧致对象}
\label{categories-section-compact}

\noindent
下面略述范畴中的“小”对象。

\begin{definition}
\label{categories-definition-compact-object}
设 $\mathcal{C}$ 为大\footnote{见注 \ref{categories-remark-big-categories}。}范畴。
若对每个使 $\colim_i M_i$ 存在的滤过图表
$M : \mathcal{I} \to \mathcal{C}$，都有
$$
\Mor_\mathcal{C}(X, \colim_i M_i) =
\colim_i \Mor_\mathcal{C}(X, M_i)
$$
则称 $\mathcal{C}$ 的对象 $X$ 是{\it 范畴紧致}的。
\end{definition}

\noindent
通常只在 $\mathcal{C}$ 有所有滤过余极限的假设下作此定义。

\begin{lemma}
\label{categories-lemma-extend-functor-by-colim}
设 $\mathcal{C}$ 与 $\mathcal{D}$ 为具有滤过余极限的大范畴。
设 $\mathcal{C}' \subset \mathcal{C}$ 为由 $\mathcal{C}$ 的范畴紧致对象
组成的小满子范畴，并且 $\mathcal{C}$ 的每个对象都是
$\mathcal{C}'$ 中对象的滤过余极限。则每个函子
$F' : \mathcal{C}' \to \mathcal{D}$ 都有唯一扩张
$F : \mathcal{C} \to \mathcal{D}$，它与滤过余极限交换。
\end{lemma}

\begin{proof}
对 $\mathcal{C}$ 的每个对象 $X$，可将 $X$ 写成滤过余极限
$X = \colim X_i$，其中 $X_i \in \Ob(\mathcal{C}')$。令
$$
F(X) = \colim F'(X_i)
$$
于 $\mathcal{D}$ 中。下文将证明此构造不依赖于 $X$ 的余极限表示。

\medskip\noindent
设给定 $\mathcal{C}$ 的态射 $\alpha : X \to Y$，且
$X = \colim_{i \in I} X_i$ 与 $Y = \colim_{j \in J} Y_i$
写成 $\mathcal{C}'$ 中对象的滤过余极限。对每个 $i \in I$，
由于 $X_i$ 是 $\mathcal{C}$ 的范畴紧致对象，可找到 $j \in J$
以及交换图表
$$
\xymatrix{
X_i \ar[r] \ar[d] & X \ar[d]^\alpha \\
Y_j \ar[r] & Y
}
$$
于是得到态射 $F'(X_i) \to F'(Y_j) \to F(Y)$，其中第二个态射是到
$F(Y) = \colim F'(Y_j)$ 的余投影。箭头
$\beta_i : F'(X_i) \to F(Y)$ 不依赖于 $j$ 的选择。
对 $i \leq i'$，复合
$$
F'(X_i) \to F'(X_{i'}) \xrightarrow{\beta_{i'}} F(Y)
$$
等于 $\beta_i$。因此由余极限的泛性质得到良定义的箭头
$$
F(\alpha) : F(X) = \colim F(X_i) \to F(Y)
$$
若 $\alpha' : Y \to Z$ 是 $\mathcal{C}$ 的另一个态射，且
$Z = \colim Z_k$ 也写成 $\mathcal{C}'$ 中对象的滤过余极限，
则可作为一项愉快的练习证明：诱导态射
$F(\alpha) : F(X) \to F(Y)$ 与 $F(\alpha') : F(Y) \to F(Z)$
复合为态射 $F(\alpha' \circ \alpha)$。细节略去。

\medskip\noindent
特别地，若给定两个表示
$X = \colim X_i$ 与 $X = \colim X'_{i'}$，它们都是
$\mathcal{C}'$ 中系统的滤过余极限，则得到互逆箭头
$\colim F'(X_i) \to \colim F'(X'_{i'})$ 与
$\colim F'(X'_{i'}) \to \colim F'(X_i)$。换言之，值 $F(X)$
不依赖于将 $X$ 表示成 $\mathcal{C}'$ 中对象的滤过余极限的选择。
结合上一段所讨论的 $F$ 的函子性，可知 $F$ 是函子。
此外，若 $X \in \Ob(\mathcal{C}')$，显然 $F(X) = F'(X)$。

\medskip\noindent
只要证明 $F$ 与滤过余极限交换，引理中的唯一性断言便很清楚
（否则该断言没有意义）。为此，假设
$X = \colim_{\lambda \in \Lambda} X_\lambda$
是 $\mathcal{C}$ 中的滤过余极限。由于 $F$ 是函子，当然得到映射
$$
\colim_\lambda F(X_\lambda) \longrightarrow F(X)
$$
另一方面，把 $X = \colim X_i$ 写成 $\mathcal{C}'$ 中对象的
滤过余极限。如上，对每个 $i \in I$，可选择
$\lambda \in \Lambda$ 及交换图表
$$
\xymatrix{
X_i \ar[rr] \ar[rd] & & X_\lambda \ar[ld] \\
& X
}
$$
同样，这给出与过渡态射相容的良定义态射
$F'(X_i) \to \colim_\lambda F(X_\lambda)$，因而给出态射
$$
F(X) = \colim_i F'(X_i) \longrightarrow \colim_\lambda F(X_\lambda)
$$
此态射与上面的态射互逆（细节略），并如所需证明
$F(X) = \colim_\lambda F(X_\lambda)$。
\end{proof}







\section{范畴中的局部化}
\label{categories-section-localization}

\noindent
本节的基本思想是：给定范畴 $\mathcal{C}$ 和一组箭头 $S$，
构造函子 $F : \mathcal{C} \to S^{-1}\mathcal{C}$，使 $S$ 的所有元素
在 $S^{-1}\mathcal{C}$ 中均可逆，并且 $F$ 在所有具有此性质的函子中
是泛的。本节的参考文献为 \cite[Chapter I, Section 2]{GZ}
与 \cite[Chapter II, Section 2]{Verdier}。

\begin{definition}
\label{categories-definition-multiplicative-system}
设 $\mathcal{C}$ 为范畴。若 $\mathcal{C}$ 的一组箭头 $S$ 具有下列性质，
则称它为{\it 左乘法系}：
\begin{enumerate}
\item[LMS1] $\mathcal{C}$ 每个对象的恒等态射都属于 $S$，且
$S$ 中两个可复合元素的复合仍属于 $S$。
\item[LMS2] 每个实线图表
$$
\xymatrix{
X \ar[d]_t \ar[r]_g & Y \ar@{..>}[d]^s \\
Z \ar@{..>}[r]^f & W
}
$$
若满足 $t \in S$，都可补成虚线交换方形，并使 $s \in S$。
\item[LMS3] 对每一对态射 $f, g : X \to Y$ 以及靶为 $X$ 的
$t \in S$，若 $f \circ t = g \circ t$，则存在源为 $Y$ 的
$s \in S$，使 $s \circ f = s \circ g$。
\end{enumerate}
若 $\mathcal{C}$ 的一组箭头 $S$ 具有下列性质，则称它为
{\it 右乘法系}：
\begin{enumerate}
\item[RMS1] $\mathcal{C}$ 每个对象的恒等态射都属于 $S$，且
$S$ 中两个可复合元素的复合仍属于 $S$。
\item[RMS2] 每个实线图表
$$
\xymatrix{
X \ar@{..>}[d]_t \ar@{..>}[r]_g & Y \ar[d]^s \\
Z \ar[r]^f & W
}
$$
若满足 $s \in S$，都可补成虚线交换方形，并使 $t \in S$。
\item[RMS3] 对每一对态射 $f, g : X \to Y$ 以及源为 $Y$ 的
$s \in S$，若 $s \circ f = s \circ g$，则存在靶为 $X$ 的
$t \in S$，使 $f \circ t = g \circ t$。
\end{enumerate}
若 $\mathcal{C}$ 的一组箭头 $S$ 同时是左乘法系和右乘法系，
则称它为{\it 乘法系}。换言之，这表示 MS1、MS2、MS3 成立，其中
MS1 $=$ LMS1 $+$ RMS1，MS2 $=$ LMS2 $+$ RMS2，且
MS3 $=$ LMS3 $+$ RMS3。（当然，LMS1 $=$ RMS1 $=$ MS1。）
\end{definition}

\noindent
这些条件可用于如下构造范畴 $S^{-1}\mathcal{C}$。

\medskip\noindent
{\bf 左分式演算。}
设 $\mathcal{C}$ 为范畴，$S$ 为左乘法系。如下定义一个新范畴
$S^{-1}\mathcal{C}$（我们将在引理 \ref{categories-lemma-left-localization}
的证明中验证此定义确实成立）：
\begin{enumerate}
\item 令 $\Ob(S^{-1}\mathcal{C}) = \Ob(\mathcal{C})$。
\item $S^{-1}\mathcal{C}$ 中的态射 $X \to Y$ 由偶对
$(f : X \to Y', s : Y \to Y')$ 给出，其中 $s \in S$，
并模去等价关系。（该等价关系定义如下。可以把偶对 $(f, s)$
的等价类看作 $s^{-1}f : X \to Y$。）
\item 称两个偶对 $(f_1 : X \to Y_1, s_1 : Y \to Y_1)$ 与
$(f_2 : X \to Y_2, s_2 : Y \to Y_2)$ 等价，如果存在第三个偶对
$(f_3 : X \to Y_3, s_3 : Y \to Y_3)$ 以及 $\mathcal{C}$ 的态射
$u : Y_1 \to Y_3$ 与 $v : Y_2 \to Y_3$，使它们嵌入交换图表
$$
\xymatrix{
 & Y_1 \ar[d]^u & \\
X \ar[ru]^{f_1} \ar[r]^{f_3} \ar[rd]_{f_2} &
Y_3 &
Y \ar[lu]_{s_1} \ar[l]_{s_3} \ar[ld]^{s_2} \\
& Y_2 \ar[u]_v &
}
$$
\item 偶对 $(f : X \to Y', s : Y \to Y')$ 与
$(g : Y \to Z', t : Z \to Z')$ 的等价类之复合，定义为偶对
$(h \circ f : X \to Z'', u \circ t : Z \to Z'')$ 的等价类；
其中选择 $h$ 与 $u \in S$ 使其嵌入交换图表
$$
\xymatrix{
Y \ar[d]_s \ar[r]_g & Z' \ar[d]^u \\
Y' \ar[r]^h & Z''
}
$$
这样的选择由假设存在。
\item $S^{-1} \mathcal{C}$ 中的恒等态射 $X \to X$ 是偶对
$(\text{id} : X \to X,
\text{id} : X \to X)$ 的等价类。
\end{enumerate}

\begin{lemma}
\label{categories-lemma-left-localization}
设 $\mathcal{C}$ 为范畴，$S$ 为左乘法系。
\begin{enumerate}
\item 上面在偶对上定义的关系是等价关系。
\item 上面给出的复合规则在等价类上良定义。
\item 复合满足结合律（且恒等态射满足恒等公理），
所以 $S^{-1}\mathcal{C}$ 是范畴。
\end{enumerate}
\end{lemma}

\begin{proof}
证明 (1)。称两个偶对
$p_1 = (f_1 : X \to Y_1, s_1 : Y \to Y_1)$ 与
$p_2 = (f_2 : X \to Y_2, s_2 : Y \to Y_2)$ 初等等价，
如果存在 $\mathcal{C}$ 的态射 $a : Y_1 \to Y_2$，使
$a \circ f_1 = f_2$ 且 $a \circ s_1 = s_2$。图示为
$$
\xymatrix{
X \ar@{=}[d] \ar[r]_{f_1} & Y_1 \ar[d]^a & Y \ar[l]^{s_1} \ar@{=}[d] \\
X \ar[r]^{f_2} & Y_2 & Y \ar[l]_{s_2}
}
$$
把此性质记为 $p_1Ep_2$。注意 $pEp$，且
$aEb, bEc \Rightarrow aEc$。（尽管如此命名，$E$ 并不是等价关系。）
断言 (1) 声称关系
$p \sim p' \Leftrightarrow \exists q: pEq \wedge p'Eq$
是等价关系（这里 $q$ 应为满足与 $p$ 和 $p'$ 相同条件的偶对）。
使用上面 $E$ 的性质作简单的形式论证可知，只须证明
$p_3Ep_1, p_3Ep_2 \Rightarrow p_1 \sim p_2$。
因此，假设给定交换图表
$$
\xymatrix{
 & Y_1 & \\
X \ar[ru]^{f_1} \ar[r]^{f_3} \ar[rd]_{f_2} &
Y_3 \ar[u]_{a_{31}} \ar[d]^{a_{32}} &
Y \ar[lu]_{s_1} \ar[l]_{s_3} \ar[ld]^{s_2} \\
& Y_2 &
}
$$
其中 $s_i \in S$。首先应用 LMS2 得到交换图表
$$
\xymatrix{
Y \ar[d]_{s_1} \ar[r]_{s_2} & Y_2 \ar@{..>}[d]^{a_{24}} \\
Y_1 \ar@{..>}[r]^{a_{14}} & Y_4
}
$$
且 $a_{24} \in S$。于是
$$
a_{14} \circ a_{31} \circ s_3 =
a_{14} \circ s_1 =
a_{24} \circ s_2 =
a_{24} \circ a_{32} \circ s_3.
$$
故由 LMS3，存在态射 $s_{44} : Y_4 \to Y'_4$，使
$s_{44} \in S$ 且
$s_{44} \circ a_{14} \circ a_{31}
= s_{44} \circ a_{24} \circ a_{32}$。
因此，把 $Y_4$、$a_{14}$、$a_{24}$ 分别替换为 $Y'_4$、
$s_{44} \circ a_{14}$、$s_{44} \circ a_{24}$ 后，可以假设
$a_{14} \circ a_{31} = a_{24} \circ a_{32}$
（并且仍有 $a_{24} \in S$ 以及
$a_{14} \circ s_1 = a_{24} \circ s_2$）。令
$$
f_4 =
a_{14} \circ f_1 =
a_{14} \circ a_{31} \circ f_3 =
a_{24} \circ a_{32} \circ f_3 =
a_{24} \circ f_2
$$
以及 $s_4 = a_{14} \circ s_1 = a_{24} \circ s_2$。则图表
$$
\xymatrix{
X \ar@{=}[d] \ar[r]_{f_1} & Y_1 \ar[d]^{a_{14}} & Y \ar[l]^{s_1} \ar@{=}[d] \\
X \ar[r]^{f_4} & Y_4 & Y \ar[l]_{s_4}
}
$$
可交换，并有 $s_4 \in S$（由 LMS1）。所以 $p_1 E p_4$，
其中 $p_4 = (f_4, s_4)$。类似地，$p_2 E p_4$。合并二者得到
$p_1 \sim p_2$。

\medskip\noindent
证明 (2)。设 $p = (f : X \to Y', s : Y \to Y')$ 与
$q = (g : Y \to Z', t : Z \to Z')$ 是上面复合定义中的偶对。
为了复合，选择交换图表
$$
\xymatrix{
Y \ar[d]_s \ar[r]_g & Z' \ar[d]^{u_2} \\
Y' \ar[r]^{h_2} & Z_2
}
$$
其中 $u_2 \in S$。先证明偶对
$r_2 = (h_2 \circ f : X \to Z_2, u_2 \circ t : Z \to Z_2)$
的等价类不依赖于 $(Z_2, h_2, u_2)$ 的选择。事实上，假设
$(Z_3, h_3, u_3)$ 是另一选择，相应复合为
$r_3 = (h_3 \circ f : X \to Z_3, u_3 \circ t : Z \to Z_3)$。
由 LMS2，可选择图表
$$
\xymatrix{
Z' \ar[d]_{u_2} \ar[r]_{u_3} & Z_3 \ar[d]^{u_{34}} \\
Z_2 \ar[r]^{h_{24}} & Z_4
}
$$
且 $u_{34} \in S$。有 $h_2 \circ s = u_2 \circ g$，并且类似地
$h_3 \circ s = u_3 \circ g$。现在
$$
u_{34} \circ h_3 \circ s
= u_{34} \circ u_3 \circ g
= h_{24} \circ u_2 \circ g
= h_{24} \circ h_2 \circ s.
$$
所以 LMS3 表明存在 $Z'_4$ 以及 $s_{44} : Z_4 \to Z'_4$，使
$s_{44} \circ u_{34} \circ h_3 = s_{44} \circ h_{24} \circ h_2$。
把 $Z_4$、$h_{24}$、$u_{34}$ 分别替换为 $Z'_4$、
$s_{44} \circ h_{24}$、$s_{44} \circ u_{34}$，可假设
$u_{34} \circ h_3 = h_{24} \circ h_2$。同时，关系
$u_{34} \circ u_3 = h_{24} \circ u_2$ 与 $u_{34} \in S$ 仍成立。
现可令 $h_4 = u_{34} \circ h_3 = h_{24} \circ h_2$ 且
$u_4 = u_{34} \circ u_3 = h_{24} \circ u_2$。于是得到交换图表
$$
\xymatrix{
X \ar@{=}[d] \ar[r]_{h_2\circ f} &
Z_2 \ar[d]^{h_{24}} &
Z \ar[l]^{u_2 \circ t} \ar@{=}[d] \\
X \ar@{=}[d] \ar[r]^{h_4\circ f} &
Z_4 &
Z \ar@{=}[d] \ar[l]_{u_4 \circ t} \\
X \ar[r]^{h_3 \circ f} &
Z_3 \ar[u]^{u_{34}} &
Z \ar[l]_{u_3 \circ t}
}
$$
因此得到偶对
$r_4 =
(h_4 \circ f : X \to Z_4, u_4 \circ t : Z \to Z_4)$；
上图表明 $r_2Er_4$ 且 $r_3Er_4$，从而如所需有
$r_2 \sim r_3$。所以现在可以把 $p \circ q$ 定义为以上方式得到的
所有可能偶对 $r$ 的等价类。

\medskip\noindent
为完成 (2) 的证明，须证明：给定偶对 $p_1, p_2, q$，若
$p_1Ep_2$，则每当复合有意义时，
$p_1 \circ q = p_2 \circ q$ 且 $q \circ p_1 = q \circ p_2$。
为此，写
$p_1 = (f_1 : X \to Y_1, s_1 : Y \to Y_1)$、
$p_2 = (f_2 : X \to Y_2, s_2 : Y \to Y_2)$，并令
$a : Y_1 \to Y_2$ 为 $\mathcal{C}$ 的态射，使
$f_2 = a \circ f_1$ 且 $s_2 = a \circ s_1$。
先假设 $q = (g : Y \to Z', t : Z \to Z')$。
此时选择左侧所示的交换图表
$$
\vcenter{
\xymatrix{
Y \ar[d]_{s_2} \ar[r]^g & Z' \ar[d]^u \\
Y_2 \ar[r]^h & Z''
}
}
\quad
\Rightarrow
\quad
\vcenter{
\xymatrix{
Y \ar[d]_{s_1} \ar[r]^g & Z' \ar[d]^u \\
Y_1 \ar[r]^{h \circ a} & Z''
}
}
$$
（其中 $u \in S$）；它蕴含右侧图表也可交换。利用这些图表可见，
两个复合 $q \circ p_1$ 与 $q \circ p_2$ 都是
$(h \circ a \circ f_1 : X \to Z'', u \circ t : Z \to Z'')$
的等价类。因此 $q \circ p_1 = q \circ p_2$。
另一种情形的证明略去；此时须证明
$p_1 \circ q = p_2 \circ q$。（它与已证明的情形类似。）

\medskip\noindent
证明 (3)。须证明复合的结合律。考虑实线图表
$$
\xymatrix{
& & & Z \ar[d] \\
& & Y \ar[d] \ar[r] & Z' \ar@{..>}[d] \\
& X \ar[d] \ar[r] & Y' \ar@{..>}[d] \ar@{..>}[r] & Z'' \ar@{..>}[d] \\
W \ar[r] & X' \ar@{..>}[r] & Y'' \ar@{..>}[r] & Z'''
}
$$
（其中竖直箭头都属于 $S$），它给出三个可复合的偶对。
利用 LMS2，可选择虚线箭头使各方形可交换，并使竖直箭头属于 $S$。
于是显然，三个偶对的复合是偶对
$(W \to Z''', Z \to Z''')$ 的等价类；该偶对由复合底行的横箭头
与右列的竖箭头得到。

\medskip\noindent
恒等公理的验证留给读者。
\end{proof}

\begin{remark}
\label{categories-remark-motivation-localization}
构造 $S^{-1} \mathcal{C}$ 的动机，是通过人为地为 $S$ 中的态射创造逆，
“迫使”它们可逆（代价是某些已有态射可能彼此等同）。这类似于
交换环关于乘法子集的局部化，更一般地，也类似于非交换环关于右分母集
的局部化（见 \cite[Section 10A]{Lam}）。这不只是一种类比：
$S^{-1} \mathcal{C}$ 的构造（更确切地说，它对加性范畴
$\mathcal{C}$ 的版本）实际上推广了后一种局部化。具体地，
一个非交换环可视为只有一个对象的预加性范畴（态射即环的元素）；
该环的一个乘法子集遂成为满足 LMS1（即 RMS1）的一组态射 $S$。
于是，这个范畴与这个子集 $S$ 的条件 RMS2、RMS3，分别转化为
右分母集的两个条件（“右可置换”与“右可逆”）；LMS 与左分母集的
情形类似。而 $S^{-1} \mathcal{C}$（赋予适当定义的加性结构）
就是对应于环之局部化的单对象范畴。
\end{remark}

\begin{definition}
\label{categories-definition-left-localization-as-fraction}
设 $\mathcal{C}$ 为范畴，$S$ 为 $\mathcal{C}$ 的态射构成的左乘法系。
给定 $\mathcal{C}$ 中任意态射 $f : X \to Y'$ 及 $S$ 中任意态射
$s : Y \to Y'$，用 {\it $s^{-1} f$} 表示偶对
$(f : X \to Y', s : Y \to Y')$ 的等价类。这是
$S^{-1} \mathcal{C}$ 中从 $X$ 到 $Y$ 的态射。
\end{definition}

\noindent
这个记号富有启发性，而且它所暗示的性质确实成立：给定
$\mathcal{C}$ 中任意态射 $f : X \to Y'$ 以及 $S$ 中任意两个态射
$s : Y \to Y'$ 与 $t : Y' \to Y''$，有
$\left(t \circ s\right)^{-1} \left(t \circ f\right) = s^{-1} f$。
此外，对 $\mathcal{C}$ 中任意 $f : X \to Y'$ 与
$g : Y' \to Z'$，以及所有 $S$ 中的 $s : Z \to Z'$，有
$s^{-1} \left(g \circ f\right) = \left(s^{-1} g\right) \circ
\left(\text{id}_{Y'}^{-1} f\right)$。最后，对 $\mathcal{C}$ 中任意
$f : X \to Y'$、所有 $S$ 中的 $s : Y \to Y'$ 及
$S$ 中的 $t : Z \to Y$，有
$\left(s \circ t\right)^{-1} f
= \left(t^{-1} \text{id}_Y\right)
\circ \left(s^{-1} f\right)$。这些都直接由定义可知。
我们可以“把具有相同靶的任意有限族态射写成有公分母的分式”。

\begin{lemma}
\label{categories-lemma-morphisms-left-localization}
设 $\mathcal{C}$ 为范畴，$S$ 为 $\mathcal{C}$ 的态射构成的左乘法系。
给定 $S^{-1}\mathcal{C}$ 中任意有限族态射 $g_i : X_i \to Y$
（以 $i$ 为指标），可以找到 $S$ 的元素 $s : Y \to Y'$ 以及
$\mathcal{C}$ 的一族态射 $f_i : X_i \to Y'$，使每个 $g_i$
都是偶对 $(f_i : X_i \to Y', s : Y \to Y')$ 的等价类。
\end{lemma}

\begin{proof}
对每个 $i$，选择 $g_i$ 的一个代表
$(X_i \to Y_i, s_i : Y \to Y_i)$。若能找到 $S$ 中的态射
$s : Y \to Y'$，使得对每个 $i$ 都存在态射
$a_i : Y_i \to Y'$ 满足 $a_i \circ s_i = s$，引理即得证。
若有两个指标 $i = 1, 2$，则可补全方形
$$
\xymatrix{
Y \ar[d]_{s_1} \ar[r]_{s_2} & Y_2 \ar[d]^{t_2} \\
Y_1 \ar[r]^{a_1} & Y'
}
$$
并使 $t_2 \in S$；定义 \ref{categories-definition-multiplicative-system}
保证这可以做到。此时 $s = t_2 \circ s_2 \in S$ 满足要求。
若有 $n > 2$ 个态射，就用上述技巧将问题归约到
$n - 1$ 个态射的情形，再由归纳法得证。
\end{proof}

\noindent
具有相同分母的态射之相等性有一个简单刻画。

\begin{lemma}
\label{categories-lemma-equality-morphisms-left-localization}
设 $\mathcal{C}$ 为范畴，$S$ 为 $\mathcal{C}$ 的态射构成的左乘法系。
设 $A, B : X \to Y$ 是 $S^{-1}\mathcal{C}$ 的态射，分别为
$(f : X \to Y', s : Y \to Y')$ 与
$(g : X \to Y', s : Y \to Y')$ 的等价类。下列条件等价：
\begin{enumerate}
\item $A = B$；
\item 存在 $S$ 中的态射 $t : Y' \to Y''$，使
$t \circ f = t \circ g$；
\item 存在态射 $a : Y' \to Y''$，使
$a \circ f = a \circ g$ 且 $a \circ s \in S$。
\end{enumerate}
\end{lemma}

\begin{proof}
我们将使用 $S^{-1}\mathcal{C}$ 是范畴这一事实
（引理 \ref{categories-lemma-left-localization}），并使用定义
\ref{categories-definition-left-localization-as-fraction} 的记号以及该定义后的讨论，
来等同 $S^{-1}\mathcal{C}$ 中的某些态射。因此写
$A = s^{-1}f$ 与 $B = s^{-1}g$。

\medskip\noindent
若 $A = B$，则
$(\text{id}_{Y'}^{-1}s) \circ A = (\text{id}_{Y'}^{-1}s) \circ B$。
有 $(\text{id}_{Y'}^{-1}s) \circ A = \text{id}_{Y'}^{-1}f$ 且
$(\text{id}_{Y'}^{-1}s) \circ B = \text{id}_{Y'}^{-1}g$。
$\text{id}_{Y'}^{-1}f$ 与 $\text{id}_{Y'}^{-1}g$ 相等，按定义意味着
存在交换图表
$$
\xymatrix{
 & Y' \ar[d]^u & \\
X \ar[ru]^f \ar[r]^h \ar[rd]_g &
Z &
Y' \ar[lu]_{\text{id}_{Y'}} \ar[l]_t \ar[ld]^{\text{id}_{Y'}} \\
& Y' \ar[u]_v &
}
$$
其中 $t \in S$。特别地，$u = v = t \in S$ 且
$t \circ f = t\circ g$。故 (1) 蕴含 (2)。

\medskip\noindent
(2) $\Rightarrow$ (3) 是立即的。假设 $a$ 如 (3) 中所述，记
$s' = a \circ s \in S$。则 $\text{id}_{Y''}^{-1}s'$ 是范畴
$S^{-1}\mathcal{C}$ 中的同构（逆为 $(s')^{-1}\text{id}_{Y''}$）。
所以为验证 $A = B$，只须验证
$\text{id}_{Y''}^{-1}s' \circ A = \text{id}_{Y''}^{-1}s' \circ B$。
利用定义 \ref{categories-definition-left-localization-as-fraction} 后文所讨论的规则，
计算得到
$\text{id}_{Y''}^{-1}s' \circ A =
\text{id}_{Y''}^{-1}(a \circ s) \circ s^{-1}f =
\text{id}_{Y''}^{-1}(a \circ f) =
\text{id}_{Y''}^{-1}(a \circ g) =
\text{id}_{Y''}^{-1}(a \circ s) \circ s^{-1}g =
\text{id}_{Y''}^{-1}s' \circ B$，故 (1) 成立。
\end{proof}

\begin{remark}
\label{categories-remark-left-localization-morphisms-colimit}
设 $\mathcal{C}$ 为范畴，$S$ 为左乘法系。给定 $\mathcal{C}$ 的对象
$Y$，以 $Y/S$ 表示如下范畴：其对象是满足 $s \in S$ 的
$s : Y \to Y'$，其态射是交换图表
$$
\xymatrix{
& Y \ar[ld]_s \ar[rd]^t & \\
Y' \ar[rr]^a & & Y''
}
$$
其中 $a : Y' \to Y''$ 任意。我们断言范畴 $Y/S$ 滤过
（见定义 \ref{categories-definition-directed}）。事实上，LMS1 蕴含
$\text{id}_Y : Y \to Y$ 属于 $Y/S$，故 $Y/S$ 非空。
LMS2 蕴含：给定 $s_1 : Y \to Y_1$ 与 $s_2 : Y \to Y_2$，
可找到图表
$$
\xymatrix{
Y \ar[d]_{s_1} \ar[r]_{s_2} & Y_2 \ar[d]^t \\
Y_1 \ar[r]^a & Y_3
}
$$
且 $t \in S$。所以 $s_1 : Y \to Y_1$ 与 $s_2 : Y \to Y_2$
在 $Y/S$ 中都有到 $t \circ s_2 : Y \to Y_3$ 的映射。
最后，给定 $Y/S$ 中从 $s_1 : Y \to Y_1$ 到
$s_2 : Y \to Y_2$ 的两个态射 $a, b$，有
$a \circ s_1 = b \circ s_1$；故由 LMS3，存在 $S$ 中的
$t : Y_2 \to Y_3$，使 $t \circ a = t \circ b$。
现在，引理 \ref{categories-lemma-morphisms-left-localization} 与
\ref{categories-lemma-equality-morphisms-left-localization} 的合并结果告诉我们
\begin{equation}
\label{categories-equation-left-localization-morphisms-colimit}
\Mor_{S^{-1}\mathcal{C}}(X, Y) =
\colim_{(s : Y \to Y') \in Y/S} \Mor_\mathcal{C}(X, Y')
\end{equation}
这个把 $S^{-1}\mathcal{C}$ 中的态射集表示为 $\mathcal{C}$ 中态射集的
滤过余极限的公式，有时很有用。
\end{remark}

\begin{lemma}
\label{categories-lemma-properties-left-localization}
设 $\mathcal{C}$ 为范畴，$S$ 为 $\mathcal{C}$ 的态射构成的左乘法系。
\begin{enumerate}
\item 规则 $X \mapsto X$ 以及
$(f : X \to Y) \mapsto (f : X \to Y, \text{id}_Y : Y \to Y)$
定义函子 $Q : \mathcal{C} \to S^{-1}\mathcal{C}$。
\item 对任意 $s \in S$，态射 $Q(s)$ 是
$S^{-1}\mathcal{C}$ 中的同构。
\item 若 $G : \mathcal{C} \to \mathcal{D}$ 为任意函子，且对每个
$s \in S$，$G(s)$ 都可逆，则存在唯一函子
$H : S^{-1}\mathcal{C} \to \mathcal{D}$，使 $H \circ Q = G$。
\end{enumerate}
\end{lemma}

\begin{proof}
(1) 与 (2) 很清楚。（在 (2) 中，$Q(s)$ 的逆是偶对
$(\text{id}_Y, s)$ 的等价类。）为说明 (3)，只须令
$H(X) = G(X)$，并令
$H((f : X \to Y', s : Y \to Y')) = G(s)^{-1} \circ G(f)$。
细节略去。
\end{proof}

\begin{lemma}
\label{categories-lemma-left-localization-limits}
设 $\mathcal{C}$ 为范畴，$S$ 为 $\mathcal{C}$ 的态射构成的左乘法系。
局部化函子 $Q : \mathcal{C} \to S^{-1}\mathcal{C}$ 与有限余极限交换。
\end{lemma}

\begin{proof}
设 $\mathcal{I}$ 为有限范畴，并设
$\mathcal{I} \to \mathcal{C}$、$i \mapsto X_i$ 为余极限存在的函子。
利用 (\ref{categories-equation-left-localization-morphisms-colimit})、$Y/S$ 滤过这一事实
以及引理 \ref{categories-lemma-directed-commutes}，有
\begin{align*}
\Mor_{S^{-1}\mathcal{C}}(Q(\colim X_i), Q(Y))
& =
\colim_{(s : Y \to Y') \in Y/S} \Mor_\mathcal{C}(\colim X_i, Y') \\
& =
\colim_{(s : Y \to Y') \in Y/S} \lim_i \Mor_\mathcal{C}(X_i, Y') \\
& =
\lim_i \colim_{(s : Y \to Y') \in Y/S} \Mor_\mathcal{C}(X_i, Y') \\
& =
\lim_i \Mor_{S^{-1}\mathcal{C}}(Q(X_i), Q(Y))
\end{align*}
而且，对每个 $j \in \Ob(\mathcal{I})$，此同构与两边到集合
$\Mor_{S^{-1}\mathcal{C}}(Q(X_j), Q(Y))$ 的投影交换。因此，
$Q(\colim X_i)$ 满足函子 $i \mapsto Q(X_i)$ 之余极限的泛性质，
故如所需，它就是该余极限。
\end{proof}

\begin{lemma}
\label{categories-lemma-left-localization-diagram}
设 $\mathcal{C}$ 为范畴，$S$ 为左乘法系。若
$f : X \to Y$、$f' : X' \to Y'$ 是 $\mathcal{C}$ 的两个态射，并且
$$
\xymatrix{
Q(X) \ar[d]_{Q(f)} \ar[r]_a & Q(X') \ar[d]^{Q(f')} \\
Q(Y) \ar[r]^b & Q(Y')
}
$$
是 $S^{-1}\mathcal{C}$ 中的交换图表，则存在 $\mathcal{C}$ 中的态射
$f'' : X'' \to Y''$ 以及 $\mathcal{C}$ 中的交换图表
$$
\xymatrix{
X \ar[d]_f \ar[r]_g & X'' \ar[d]^{f''} & X' \ar[d]^{f'} \ar[l]^s \\
Y \ar[r]^h & Y'' & Y' \ar[l]_t
}
$$
其中 $s, t \in S$，且 $a = s^{-1}g$、$b = t^{-1}h$。
\end{lemma}

\begin{proof}
如下选择映射与对象。首先，将 $a = s^{-1}g$ 写成某个
$S$ 中的 $s : X' \to X''$ 与某个 $g : X \to X''$。
由 LMS2，可以找到 $S$ 中的 $t : Y' \to Y''$ 以及
$f'' : X'' \to Y''$，使图表
$$
\xymatrix{
X \ar[d]_f \ar[r]_g & X'' \ar[d]^{f''} & X' \ar[d]^{f'} \ar[l]^s \\
Y & Y'' & Y' \ar[l]_t
}
$$
可交换。接下来，在该图表中，我们将反复改变对
$$
X'' \xrightarrow{f''} Y'' \xleftarrow{t} Y'
$$
的选择，方法是用一个满足 $d \circ t \in S$ 的态射
$d : Y'' \to Y'''$ 与 $t$ 和 $f''$ 后复合。根据注
\ref{categories-remark-left-localization-morphisms-colimit}，经过这种替换后，
可假设存在态射 $h : Y \to Y''$，使 $b = t^{-1}h$\footnote{还可用
更直接的方式看出这一点：把 $b = q^{-1}i$ 写成 $S$ 中的某个
$q : Y' \to Z$ 及某个 $i : Y \to Z$。由 LMS2，可找到
$S$ 中的 $r : Y'' \to Y'''$ 以及 $j : Z \to Y'''$，使
$j \circ q = r \circ t$。现在令 $d = r$ 且 $h = j \circ i$。}。
此时我们已具备引理中的所有条件，唯独尚不知道图表的左方形是否交换。
但 $S^{-1} \mathcal{C}$ 中复合的定义表明，
$b \circ Q\left(f\right)$ 是偶对
$(h \circ f : X \to Y'', t : Y' \to Y'')$ 的等价类
（因为 $b$ 是偶对 $(h : Y \to Y'', t : Y' \to Y'')$ 的等价类，
而 $Q\left(f\right)$ 是偶对
$(f : X \to Y, \text{id} : Y \to Y)$ 的等价类）；另一方面，
$Q\left(f'\right) \circ a$ 是偶对
$(f'' \circ g : X \to Y'', t : Y' \to Y'')$ 的等价类
（因为 $a$ 是偶对 $(g : X \to X'', s : X' \to X'')$ 的等价类，
而 $Q\left(f'\right)$ 是偶对
$(f' : X' \to Y', \text{id} : Y' \to Y')$ 的等价类）。
已知 $b \circ Q\left(f\right) = Q\left(f'\right) \circ a$，
故偶对 $(h \circ f : X \to Y'', t : Y' \to Y'')$ 与
$(f'' \circ g : X \to Y'', t : Y' \to Y'')$ 的等价类相等。
因此，由引理 \ref{categories-lemma-equality-morphisms-left-localization}，
可找到态射 $d : Y'' \to Y'''$，使 $d \circ t \in S$ 且
$d \circ h \circ f = d \circ f'' \circ g$。最后再作一次上述类型的
替换，便得到所需结论。
\end{proof}

\noindent
{\bf 右分式演算。}
设 $\mathcal{C}$ 为范畴，$S$ 为右乘法系。如下定义一个新范畴
$S^{-1}\mathcal{C}$（我们将在引理 \ref{categories-lemma-right-localization}
的证明中验证此定义确实成立）：
\begin{enumerate}
\item 令 $\Ob(S^{-1}\mathcal{C}) = \Ob(\mathcal{C})$。
\item $S^{-1}\mathcal{C}$ 中的态射 $X \to Y$ 由偶对
$(f : X' \to Y, s : X' \to X)$ 给出，其中 $s \in S$，
并模去等价关系。（该等价关系定义如下。可以把偶对 $(f, s)$
的等价类看作 $fs^{-1} : X \to Y$。）
\item 称两个偶对 $(f_1 : X_1 \to Y, s_1 : X_1 \to X)$ 与
$(f_2 : X_2 \to Y, s_2 : X_2 \to X)$ 等价，如果存在第三个偶对
$(f_3 : X_3 \to Y, s_3 : X_3 \to X)$ 以及 $\mathcal{C}$ 的态射
$u : X_3 \to X_1$ 与 $v : X_3 \to X_2$，使它们嵌入交换图表
$$
\xymatrix{
 & X_1 \ar[ld]_{s_1} \ar[rd]^{f_1} & \\
X &
X_3 \ar[l]_{s_3} \ar[u]_u \ar[d]^v \ar[r]^{f_3} &
Y \\
& X_2 \ar[lu]^{s_2} \ar[ru]_{f_2} &
}
$$
\item 偶对 $(f : X' \to Y, s : X' \to X)$ 与
$(g : Y' \to Z, t : Y' \to Y)$ 的等价类之复合，定义为偶对
$(g \circ h : X'' \to Z, s \circ u : X'' \to X)$ 的等价类；
其中选择 $h$ 与 $u \in S$ 使其嵌入交换图表
$$
\xymatrix{
X'' \ar[d]_u \ar[r]^h & Y' \ar[d]^t \\
X' \ar[r]^f & Y
}
$$
这样的选择由假设存在。
\item $S^{-1} \mathcal{C}$ 中的恒等态射 $X \to X$ 是偶对
$(\text{id} : X \to X,
\text{id} : X \to X)$ 的等价类。
\end{enumerate}

\begin{lemma}
\label{categories-lemma-right-localization}
设 $\mathcal{C}$ 为范畴，$S$ 为右乘法系。
\begin{enumerate}
\item 上面在偶对上定义的关系是等价关系。
\item 上面给出的复合规则在等价类上良定义。
\item 复合满足结合律（且恒等态射满足恒等公理），
所以 $S^{-1}\mathcal{C}$ 是范畴。
\end{enumerate}
\end{lemma}

\begin{proof}
本引理与引理 \ref{categories-lemma-left-localization} 对偶。把
$\mathcal{C}$ 换成其对偶范畴（其中 $S$ 是左乘法系），
即可由该引理形式地推出本引理。
\end{proof}

\begin{definition}
\label{categories-definition-right-localization-as-fraction}
设 $\mathcal{C}$ 为范畴，$S$ 为 $\mathcal{C}$ 的态射构成的右乘法系。
给定 $\mathcal{C}$ 中任意态射 $f : X' \to Y$ 及 $S$ 中任意态射
$s : X' \to X$，以 {\it $f s^{-1}$} 表示偶对
$(f : X' \to Y, s : X' \to X)$ 的等价类。这是
$S^{-1} \mathcal{C}$ 中从 $X$ 到 $Y$ 的态射。
\end{definition}

\noindent
与定义 \ref{categories-definition-left-localization-as-fraction} 中恒等式类似
（事实上对偶）的恒等式成立。我们可以“把具有相同源的任意有限族态射
写成有公分母的分式”。

\begin{lemma}
\label{categories-lemma-morphisms-right-localization}
设 $\mathcal{C}$ 为范畴，$S$ 为 $\mathcal{C}$ 的态射构成的右乘法系。
给定 $S^{-1}\mathcal{C}$ 中任意有限族态射 $g_i : X \to Y_i$
（以 $i$ 为指标），可以找到 $S$ 的元素 $s : X' \to X$
以及 $\mathcal{C}$ 的一族态射 $f_i : X' \to Y_i$，使 $g_i$
是偶对 $(f_i : X' \to Y_i, s : X' \to X)$ 的等价类。
\end{lemma}

\begin{proof}
本引理是引理 \ref{categories-lemma-morphisms-left-localization} 的对偶；
把所涉及的所有范畴换成其对偶范畴，即可由该引理形式地推出。
\end{proof}

\noindent
具有相同分母的态射之相等性有一个简单刻画。

\begin{lemma}
\label{categories-lemma-equality-morphisms-right-localization}
设 $\mathcal{C}$ 为范畴，$S$ 为 $\mathcal{C}$ 的态射构成的右乘法系。
设 $A, B : X \to Y$ 是 $S^{-1}\mathcal{C}$ 的态射，分别为
$(f : X' \to Y, s : X' \to X)$ 与
$(g : X' \to Y, s : X' \to X)$ 的等价类。下列条件等价：
\begin{enumerate}
\item $A = B$；
\item 存在 $S$ 中的态射 $t : X'' \to X'$，使
$f \circ t = g \circ t$；
\item 存在态射 $a : X'' \to X'$，使
$f \circ a = g \circ a$ 且 $s \circ a \in S$。
\end{enumerate}
\end{lemma}

\begin{proof}
这与引理 \ref{categories-lemma-equality-morphisms-left-localization} 对偶。
\end{proof}

\begin{remark}
\label{categories-remark-right-localization-morphisms-colimit}
设 $\mathcal{C}$ 为范畴，$S$ 为右乘法系。给定 $\mathcal{C}$ 的对象
$X$，以 $S/X$ 表示如下范畴：其对象是满足 $s \in S$ 的
$s : X' \to X$，其态射是交换图表
$$
\xymatrix{
X' \ar[rd]_s \ar[rr]_a & & X'' \ar[ld]^t \\
& X
}
$$
其中 $a : X' \to X''$ 任意。范畴 $S/X$ 余滤过
（见定义 \ref{categories-definition-codirected}）。
（这与注 \ref{categories-remark-left-localization-morphisms-colimit} 中的相应断言对偶。）
现在，引理 \ref{categories-lemma-morphisms-right-localization} 与
\ref{categories-lemma-equality-morphisms-right-localization} 的合并结果告诉我们
\begin{equation}
\label{categories-equation-right-localization-morphisms-colimit}
\Mor_{S^{-1}\mathcal{C}}(X, Y) =
\colim_{(s : X' \to X) \in (S/X)^{opp}} \Mor_\mathcal{C}(X', Y)
\end{equation}
这个把 $S^{-1}\mathcal{C}$ 中的态射表示为 $\mathcal{C}$ 中态射的
滤过余极限的公式，有时很有用。
\end{remark}

\begin{lemma}
\label{categories-lemma-properties-right-localization}
设 $\mathcal{C}$ 为范畴，$S$ 为 $\mathcal{C}$ 的态射构成的右乘法系。
\begin{enumerate}
\item 规则 $X \mapsto X$ 以及
$(f : X \to Y) \mapsto (f : X \to Y, \text{id}_X : X \to X)$
定义函子 $Q : \mathcal{C} \to S^{-1}\mathcal{C}$。
\item 对任意 $s \in S$，态射 $Q(s)$ 是
$S^{-1}\mathcal{C}$ 中的同构。
\item 若 $G : \mathcal{C} \to \mathcal{D}$ 为任意函子，且对每个
$s \in S$，$G(s)$ 都可逆，则存在唯一函子
$H : S^{-1}\mathcal{C} \to \mathcal{D}$，使 $H \circ Q = G$。
\end{enumerate}
\end{lemma}

\begin{proof}
本引理是引理 \ref{categories-lemma-properties-left-localization} 的对偶；
把所涉及的所有范畴换成其对偶范畴，即可由该引理形式地推出。
\end{proof}

\begin{lemma}
\label{categories-lemma-right-localization-limits}
设 $\mathcal{C}$ 为范畴，$S$ 为 $\mathcal{C}$ 的态射构成的右乘法系。
局部化函子 $Q : \mathcal{C} \to S^{-1}\mathcal{C}$ 与有限极限交换。
\end{lemma}

\begin{proof}
这与引理 \ref{categories-lemma-left-localization-limits} 对偶。
\end{proof}

\begin{lemma}
\label{categories-lemma-right-localization-diagram}
设 $\mathcal{C}$ 为范畴，$S$ 为右乘法系。若
$f : X \to Y$、$f' : X' \to Y'$ 是 $\mathcal{C}$ 的两个态射，并且
$$
\xymatrix{
Q(X) \ar[d]_{Q(f)} \ar[r]_a & Q(X') \ar[d]^{Q(f')} \\
Q(Y) \ar[r]^b & Q(Y')
}
$$
是 $S^{-1}\mathcal{C}$ 中的交换图表，则存在 $\mathcal{C}$ 中的态射
$f'' : X'' \to Y''$ 以及 $\mathcal{C}$ 中的交换图表
$$
\xymatrix{
X \ar[d]_f & X'' \ar[l]^s \ar[d]^{f''} \ar[r]_g & X' \ar[d]^{f'} \\
Y & Y'' \ar[l]_t \ar[r]^h & Y'
}
$$
其中 $s, t \in S$，且 $a = gs^{-1}$、$b = ht^{-1}$。
\end{lemma}

\begin{proof}
本引理与引理 \ref{categories-lemma-left-localization-diagram} 对偶。
\end{proof}

\noindent
{\bf 乘法系与双侧分式演算。}
若 $S$ 是乘法系，则左、右分式演算给出典范同构的范畴。

\begin{lemma}
\label{categories-lemma-multiplicative-system}
设 $\mathcal{C}$ 为范畴，$S$ 为乘法系。左分式范畴与右分式范畴
$S^{-1}\mathcal{C}$ 典范同构。
\end{lemma}

\begin{proof}
以 $\mathcal{C}_{left}$、$\mathcal{C}_{right}$ 表示两个分式范畴。
由引理 \ref{categories-lemma-properties-left-localization} 与
\ref{categories-lemma-properties-right-localization} 的泛性质，得到函子
$\mathcal{C}_{left} \to \mathcal{C}_{right}$ 与
$\mathcal{C}_{right} \to \mathcal{C}_{left}$。
由泛性质中的唯一性断言，这两个函子互为逆。
\end{proof}

\begin{definition}
\label{categories-definition-saturated-multiplicative-system}
设 $\mathcal{C}$ 为范畴，$S$ 为乘法系。若除了 MS1、MS2、MS3 之外，
还有
\begin{enumerate}
\item[MS4] 给定三个可复合态射 $f, g, h$，若
$fg, gh \in S$，则 $g \in S$，
\end{enumerate}
则称 $S$ 是{\it 饱和}的。
\end{definition}

\noindent
注意，饱和乘法系包含所有同构。此外，若 $f, g, h$ 是范畴中
可复合的态射，且 $fg, gh$ 是同构，则 $g$ 是同构
（因为此时 $g$ 同时有左逆和右逆，故可逆）。

\begin{lemma}
\label{categories-lemma-what-gets-inverted}
设 $\mathcal{C}$ 为范畴，$S$ 为乘法系。以
$Q : \mathcal{C} \to S^{-1}\mathcal{C}$ 表示局部化函子。集合
$$
\hat S = \{f \in \text{Arrows}(\mathcal{C}) \mid
Q(f) \text{ 为同构}\}
$$
等于
$$
S' = \{f \in \text{Arrows}(\mathcal{C}) \mid
\text{存在 }g, h\text{ 使得 }gf, fh \in S\}
$$
并且是包含 $S$ 的最小饱和乘法系。特别地，若 $S$ 饱和，
则 $\hat S = S$。
\end{lemma}

\begin{proof}
显然 $S \subset S' \subset \hat S$，因为 $S'$ 的元素映到
$S^{-1}\mathcal{C}$ 中同时具有左逆和右逆的态射。注意 $S'$ 满足
MS4，而 $\hat S$ 满足 MS1。接下来证明 $S' = \hat S$。

\medskip\noindent
设 $f \in \hat S$。令 $s^{-1}g = ht^{-1}$ 为
$S^{-1}\mathcal{C}$ 中的逆态射。（由引理
\ref{categories-lemma-multiplicative-system}，可以同时使用左分式和右分式描述
$S^{-1}\mathcal{C}$ 中的态射。）关系
$\text{id}_X = s^{-1}gf$ 在 $S^{-1}\mathcal{C}$ 中成立，意味着存在
交换图表
$$
\xymatrix{
 & X' \ar[d]^u & \\
X \ar[ru]^{gf} \ar[r]^{f'} \ar[rd]_{\text{id}_X} &
X'' &
X \ar[lu]_s \ar[l]_{s'} \ar[ld]^{\text{id}_X} \\
& X \ar[u]_v &
}
$$
其中有某些态射 $f', u, v$ 以及 $s' \in S$。因此
$ugf = s' \in S$。类似地，利用 $\text{id}_Y = fht^{-1}$，
可证明对某个 $w$ 有 $fhw \in S$。由此 $f \in S'$，故
$S' = \hat S$。只要证明 $S' = \hat S$ 是乘法系，就显然可推出
$S' = \hat S$ 是包含 $S$ 的最小饱和乘法系。

\medskip\noindent
上面的讨论已经处理了 MS1 与 MS4，所以为完成引理的证明，
须证明 $\hat S$ 满足 LMS2、RMS2、LMS3、RMS3。
先验证 $\hat S$ 的 LMS2。假设有实线图表
$$
\xymatrix{
X \ar[d]_t \ar[r]_g & Y \ar@{..>}[d]^s \\
Z \ar@{..>}[r]^f & W
}
$$
其中 $t \in \hat S$。选择态射 $a : Z \to Z'$，使 $at \in S$。
于是可对 $S$ 使用 LMS2，找到交换图表
$$
\xymatrix{
X \ar[d]_t \ar[r]_g & Y \ar[dd]^s \\
Z \ar[d]_a \\
Z' \ar[r]^{f'} & W
}
$$
令 $f = f' \circ a$ 即得结论。RMS2 的证明与此对偶。
最后，假设给定一对态射 $f, g : X \to Y$ 以及靶为 $X$ 的
$t \in \hat S$，并且 $ft = gt$。选择态射 $b$ 使 $tb \in S$。
于是 $ftb = gtb$，由 $S$ 的 LMS3，存在源为 $Y$ 的
$s \in S$，使 $sf = sg$，正如所需。RMS3 的证明与此对偶。
\end{proof}








\section{形式性质}
\label{categories-section-formal-cat-cat}

\noindent
本节讨论范畴的 $2$-范畴的一些形式性质。这将引出稍后对（严格）
$2$-范畴的定义。

\medskip\noindent
以 $\Ob(\textit{Cat})$ 表示所有范畴构成的类。对每一对范畴
$\mathcal{A}, \mathcal{B} \in \Ob(\textit{Cat})$，都有函子构成的
“小”范畴 $\text{Fun}(\mathcal{A}, \mathcal{B})$。像
$$
\xymatrix{
\mathcal{A}
\rruppertwocell^{F''}{t'}
\ar[rr]_(.3){F'}
\rrlowertwocell_F{t}
& &
\mathcal{B}
}
\text{ 复合为 }
\xymatrix{
\mathcal{A}
\rrtwocell^{F''}_F{\ \ t \circ t'}
& &
\mathcal{B}
}
$$
这样的函子变换之复合称为{\it 竖直}复合。对此使用通常的符号
$\circ$。接下来定义{\it 水平}复合。为此，先进一步说明现有结构。

\medskip\noindent
具体地，对每三个范畴 $\mathcal{A}$、$\mathcal{B}$、$\mathcal{C}$，
由函子复合得到复合法则
$$
\circ : \Ob(\text{Fun}(\mathcal{B}, \mathcal{C}))
\times
\Ob(\text{Fun}(\mathcal{A}, \mathcal{B}))
\longrightarrow
\Ob(\text{Fun}(\mathcal{A}, \mathcal{C}))
$$
该复合法则满足结合律，恒等函子充当单位。换言之，若忘掉函子变换，
可见 $\textit{Cat}$ 构成一个范畴。这个结构如何与函子之间的态射相互作用？

\medskip\noindent
给定函子 $F, F' : \mathcal{A} \to \mathcal{B}$ 的变换
$t : F \to F'$ 以及函子 $G : \mathcal{B} \to \mathcal{C}$，
可以定义函子变换 $G\circ F \to G \circ F'$。将此变换记为
${}_Gt$。对所有 $x \in \mathcal{A}$，它由公式
$({}_Gt)_x = G(t_x) : G(F(x)) \to G(F'(x))$ 给出。
这样，与 $G$ 复合成为函子
$$
\text{Fun}(\mathcal{A}, \mathcal{B})
\longrightarrow
\text{Fun}(\mathcal{A}, \mathcal{C}).
$$
只须验证 ${}_G(\text{id}_F) = \text{id}_{G \circ F}$ 以及
${}_G(t_1 \circ t_2) = {}_Gt_1 \circ {}_Gt_2$ 即可看出这一点。
当然还有 ${}_{\text{id}_\mathcal{B}}t = t$。

\medskip\noindent
类似地，给定函子 $G, G' : \mathcal{B} \to \mathcal{C}$ 的变换
$s : G \to G'$ 以及函子 $F : \mathcal{A} \to \mathcal{B}$，
可定义 $s_F$ 为函子变换 $G\circ F \to G' \circ F$，其中对所有
$x \in \mathcal{A}$，有
$(s_F)_x = s_{F(x)} : G(F(x)) \to G'(F(x))$。
这样，与 $F$ 复合成为函子
$$
\text{Fun}(\mathcal{B}, \mathcal{C})
\longrightarrow
\text{Fun}(\mathcal{A}, \mathcal{C}).
$$
只须验证 $(\text{id}_G)_F = \text{id}_{G\circ F}$ 以及
$(s_1 \circ s_2)_F = s_{1, F} \circ s_{2, F}$ 即可看出这一点。
当然还有 $s_{\text{id}_\mathcal{B}} = s$。

\medskip\noindent
只要各式有意义，这些构造还满足
$$
{}_{G_1}({}_{G_2}t) = {}_{G_1\circ G_2}t,
\ (s_{F_1})_{F_2} = s_{F_1 \circ F_2},
\text{ 且 }{}_H(s_F) = ({}_Hs)_F
$$
最后，给定函子 $F, F' : \mathcal{A} \to \mathcal{B}$、
$G, G' : \mathcal{B} \to \mathcal{C}$ 以及变换
$t : F \to F'$、$s : G \to G'$，下图可交换：
$$
\xymatrix{
G \circ F \ar[r]^{{}_Gt} \ar[d]_{s_F}
&
G \circ F' \ar[d]^{s_{F'}} \\
G' \circ F \ar[r]_{{}_{G'}t}
&
G' \circ F'
}
$$
换言之，${}_{G'}t \circ s_F = s_{F'}\circ {}_Gt$。
为证明这一点，只须考察任意对象 $x \in \Ob(\mathcal{A})$ 上的情形：
$$
\xymatrix{
G(F(x)) \ar[r]^{G(t_x)} \ar[d]_{s_{F(x)}}
&
G(F'(x)) \ar[d]^{s_{F'(x)}} \\
G'(F(x)) \ar[r]_{G'(t_x)}
&
G'(F'(x))
}
$$
该图可交换，因为 $s$ 是函子变换。这个相容关系使我们可以定义水平复合。

\begin{definition}
\label{categories-definition-horizontal-composition}
给定下式左侧所示图表：
$$
\xymatrix{
\mathcal{A}
\rtwocell^F_{F'}{t}
&
\mathcal{B}
\rtwocell^G_{G'}{s}
&
\mathcal{C}
}
\text{ 给出 }
\xymatrix{
\mathcal{A}
\rrtwocell^{G \circ F} _{G' \circ F'}{\ \ s \star t}
& &
\mathcal{C}
}
$$
定义{\it 水平}复合 $s \star t$ 为函子变换
${}_{G'}t \circ s_F = s_{F'}\circ {}_Gt$。
\end{definition}

\noindent
现在可将先前构造的变换 ${}_Gt$ 与 $s_F$ 恢复为
$ {}_Gt = \text{id}_G \star t $ 与 $ s_F = s \star \text{id}_F $。
此外，上面得到的所有规则都是下一个引理所述性质的推论。

\begin{lemma}
\label{categories-lemma-properties-2-cat-cats}
水平复合与竖直复合具有下列性质：
\begin{enumerate}
\item $\circ$ 与 $\star$ 满足结合律；
\item 恒等变换 $\text{id}_F$ 是 $\circ$ 的单位；
\item 恒等函子的恒等变换
$\text{id}_{\text{id}_\mathcal{A}}$ 是 $\star$ 与 $\circ$ 的单位；
\item 给定图表
$$
\xymatrix{
\mathcal{A}
\rruppertwocell^F{t}
\ar[rr]_(.3){F'}
\rrlowertwocell_{F''}{t'}
& &
\mathcal{B}
\rruppertwocell^G{s}
\ar[rr]_(.3){G'}
\rrlowertwocell_{G''}{s'}
& &
\mathcal{C}
}
$$
有 $ (s' \circ s) \star (t' \circ t) = (s' \star t') \circ (s \star t)$。
\end{enumerate}
\end{lemma}

\begin{proof}
用先前的记号，最后一个断言化为等式
$$
s'_{F''}
\circ
{}_{G'}t'
\circ
s_{F'}
\circ
{}_Gt
=
(s' \circ s)_{F''}
\circ
{}_G(t' \circ t).
$$
把上面所得结果用于中间的复合，可将左边改写为
$
s'_{F''}
\circ
s_{F''}
\circ
{}_Gt'
\circ
{}_Gt
$，它显然等于右边。
\end{proof}

\noindent
引理条件 (4) 的另一种表述是：函子复合与函子变换的水平复合给出函子
$$
(\circ, \star) :
\text{Fun}(\mathcal{B}, \mathcal{C})
\times
\text{Fun}(\mathcal{A}, \mathcal{B})
\longrightarrow
\text{Fun}(\mathcal{A}, \mathcal{C})
$$
其源为乘积范畴，见定义 \ref{categories-definition-product-category}。

\section{2-范畴}
\label{categories-section-2-categories}

\noindent
我们将定义栈的语境中出现的（严格）$2$-范畴。阅读前，请先看
第 \ref{categories-section-formal-cat-cat} 节与例
\ref{categories-example-2-1-category-of-categories}。基本上，就是取该例，并把
其中对象、$1$-态射、$2$-态射所满足的全部规则写出来。

\begin{definition}
\label{categories-definition-2-category}
一个（严格）{\it $2$-范畴} $\mathcal{C}$ 由下列数据组成：
\begin{enumerate}
\item 对象集 $\Ob(\mathcal{C})$。
\item 对每一对 $x, y \in \Ob(\mathcal{C})$，给定范畴
$\Mor_\mathcal{C}(x, y)$。$\Mor_\mathcal{C}(x, y)$ 的对象称为
{\it $1$-态射}，记为 $F : x \to y$；这些 $1$-态射之间的态射
称为{\it $2$-态射}，记为 $t : F' \to F$。
$\Mor_\mathcal{C}(x, y)$ 中 $2$-态射的复合称为{\it 竖直}复合；
对 $t : F' \to F$ 与 $t' : F'' \to F'$，记为 $t \circ t'$。
\item 对每三个 $x, y, z\in \Ob(\mathcal{C})$，给定函子
$$
(\circ, \star) :
\Mor_\mathcal{C}(y, z) \times \Mor_\mathcal{C}(x, y)
\longrightarrow
\Mor_\mathcal{C}(x, z).
$$
左侧一对 $1$-态射 $(F, G)$ 的像称为 $F$ 与 $G$ 的
{\it 复合}，记为 $F\circ G$。一对 $2$-态射 $(t, s)$ 的像称为
{\it 水平}复合，记为 $t \star s$。
\end{enumerate}
这些数据须满足下列规则：
\begin{enumerate}
\item 对象集与 $1$-态射集连同 $1$-态射的复合构成一个范畴。
\item $2$-态射的水平复合满足结合律。
\item 恒等 $1$-态射 $\text{id}_x$ 的恒等 $2$-态射
$\text{id}_{\text{id}_x}$ 是水平复合的单位。
\end{enumerate}
\end{definition}

\noindent
显然，这并不是一种很便于使用的对象。不过，有许多例子中如何使用它
相当清楚。迄今唯一的例子是如下 $2$-范畴：其对象是给定的一族范畴，
$1$-态射是这些范畴之间的函子，$2$-态射是函子的自然变换，见
第 \ref{categories-section-formal-cat-cat} 节。就本文而言，所有 $2$-范畴都将是
该例的子 $2$-范畴。以下说明何谓子 $2$-范畴。

\begin{definition}
\label{categories-definition-sub-2-category}
设 $\mathcal{C}$ 为 $2$-范畴。$\mathcal{C}$ 的一个
{\it 子 $2$-范畴} $\mathcal{C}'$ 由 $\Ob(\mathcal{C})$ 的子集
$\Ob(\mathcal{C}')$ 以及对所有 $x, y \in \Ob(\mathcal{C}')$，
范畴 $\Mor_\mathcal{C}(x, y)$ 的子范畴
$\Mor_{\mathcal{C}'}(x, y)$ 给出；要求这些数据连同运算
$\circ$（$1$-态射的复合）、$\circ$（$2$-态射的竖直复合）以及
$\star$（水平复合）构成一个 $2$-范畴。
\end{definition}

\begin{remark}
\label{categories-remark-big-2-categories}
大 $2$-范畴。在许多文献中，允许 $2$-范畴具有对象类
（但愿不允许“类的类”）。我们也允许这些“大”$2$-范畴，
不过仅限于下列情形（以后随进展更新）：
\begin{enumerate}
\item 范畴的 $2$-范畴 $\textit{Cat}$。
\item 范畴的 $(2, 1)$-范畴 $\textit{Cat}$。
\item 群胚的 $2$-范畴 $\textit{Groupoids}$；这是一个
$(2, 1)$-范畴。
\item 固定范畴上的纤维化范畴的 $2$-范畴。
\item 固定范畴上的纤维化范畴的 $(2, 1)$-范畴。
\item 固定范畴上的群胚纤维化范畴的 $2$-范畴；这是一个
$(2, 1)$-范畴。
\item 固定站点上的栈的 $2$-范畴。
\item 固定站点上的栈的 $(2, 1)$-范畴。
\item 固定站点上的群胚值栈的 $2$-范畴；这是一个
$(2, 1)$-范畴。
\item 固定站点上的广群值栈的 $2$-范畴；这是一个
$(2, 1)$-范畴。
\item 固定概形上的代数栈的 $2$-范畴；这是一个
$(2, 1)$-范畴。
\end{enumerate}
见定义 \ref{categories-definition-2-1-category}。注意，在每种情形中，
$2$-范畴 $\mathcal{C}$ 的对象类都是真类，但对所有对象
$x, y \in \Ob(C)$，范畴 $\Mor_\mathcal{C}(x, y)$ 都是“小”的
（按照我们的约定）。
\end{remark}

\noindent
第 \ref{categories-section-definition-categories} 节中定义的范畴等价概念，
可如下推广到更一般的 $2$-范畴情形。

\begin{definition}
\label{categories-definition-equivalence}
若存在 $1$-态射 $F : x \to y$ 与 $G : y \to x$，使
$F \circ G$ 与 $\text{id}_y$ 之间 $2$-同构，且
$G \circ F$ 与 $\text{id}_x$ 之间 $2$-同构，则称 $2$-范畴的两个对象
$x, y$ {\it 等价}。
\end{definition}

\noindent
有时需要说明从一个范畴到一个 $2$-范畴的函子是什么意思。

\begin{definition}
\label{categories-definition-functor-into-2-category}
设 $\mathcal{A}$ 为范畴，$\mathcal{C}$ 为 $2$-范畴。
\begin{enumerate}
\item 除非另有说明，从通常范畴到 $2$-范畴的{\it 函子}忽略
$2$-态射。换言之，它是一个“通常”函子，其靶是从 2-范畴忘掉所有
2-态射后形成的范畴。
\item 从 $\mathcal{A}$ 到 2-范畴 $\mathcal{C}$ 的
{\it 弱函子}，或{\it 伪函子} $\varphi$，由下列数据给出：
\begin{enumerate}
\item 映射 $\varphi : \Ob(\mathcal{A}) \to \Ob(\mathcal{C})$；
\item 对每一对 $x, y\in \Ob(\mathcal{A})$ 以及每个态射
$f : x \to y$，给定 $1$-态射
$\varphi(f) : \varphi(x) \to \varphi(y)$；
\item 对每个 $x\in \Ob(A)$，给定 $2$-态射
$\alpha_x : \text{id}_{\varphi(x)} \to \varphi(\text{id}_x)$；
\item 对 $\mathcal{A}$ 中每一对可复合态射 $f : x \to y$、
$g : y \to z$，给定 $2$-态射
$\alpha_{g, f} : \varphi(g \circ f) \to \varphi(g) \circ \varphi(f)$。
\end{enumerate}
这些数据须满足下列条件：
\begin{enumerate}
\item $2$-态射 $\alpha_x$ 与 $\alpha_{g, f}$ 都是同构；
\item 对 $\mathcal{A}$ 中任意态射 $f : x \to y$，有
$\alpha_{\text{id}_y, f} = \alpha_y \star \text{id}_{\varphi(f)}$：
$$
\xymatrix{
\varphi(x)
\rrtwocell^{\varphi(f)}_{\varphi(f)}{\ \ \ \ \text{id}_{\varphi(f)}}
& &
\varphi(y)
\rrtwocell^{\text{id}_{\varphi(y)}}_{\varphi(\text{id}_y)}{\alpha_y}
& &
\varphi(y)
}
=
\xymatrix{
\varphi(x)
\rrtwocell^{\varphi(f)}_{\varphi(\text{id}_y) \circ \varphi(f)}{\ \ \ \ \alpha_{\text{id}_y, f}}
& &
\varphi(y)
}
$$
\item 对 $\mathcal{A}$ 中任意态射 $f : x \to y$，有
$\alpha_{f, \text{id}_x} = \text{id}_{\varphi(f)} \star \alpha_x$；
\item 对 $\mathcal{A}$ 中任意三个可复合态射
$f : w \to x$、$g : x \to y$、$h : y \to z$，有
$$
(\text{id}_{\varphi(h)} \star \alpha_{g, f})
\circ
\alpha_{h, g \circ f}
=
(\alpha_{h, g} \star \text{id}_{\varphi(f)})
\circ
\alpha_{h \circ g, f}
$$
换言之，下列以 $1$-态射为对象、以 $2$-态射为箭头的图表可交换：
$$
\xymatrix{
\varphi(h \circ g \circ f)
\ar[d]_{\alpha_{h, g \circ f}}
\ar[rr]_{\alpha_{h \circ g, f}}
& &
\varphi(h \circ g) \circ \varphi(f)
\ar[d]^{\alpha_{h, g} \star \text{id}_{\varphi(f)}} \\
\varphi(h) \circ \varphi(g \circ f)
\ar[rr]^{\text{id}_{\varphi(h)} \star \alpha_{g, f}}
& &
\varphi(h) \circ \varphi(g) \circ \varphi(f)
}
$$
\end{enumerate}
\end{enumerate}
\end{definition}

\noindent
这仍然不是很便于使用的概念，但有时确实会出现。有一个定理断言，
任意伪函子都同构于一个函子。最后，还有{\it $2$-范畴之间的函子}
与{\it $2$-范畴之间的伪函子}的概念。后一概念会把我们带入
$3$-范畴的领域。我们希望几乎不惜一切代价避免定义它！

\section{(2, 1)-范畴}
\label{categories-section-2-1-categories}

\noindent
某些 $2$-范畴具有所有 $2$-态射都是同构的性质。它们将在下文发挥
重要作用，而且更容易使用。

\begin{definition}
\label{categories-definition-2-1-category}
一个（严格）{\it $(2, 1)$-范畴}，是所有 $2$-态射都是同构的
$2$-范畴。
\end{definition}

\begin{example}
\label{categories-example-2-1-category-of-categories}
注 \ref{categories-remark-big-2-categories} 中的 $2$-范畴 $\textit{Cat}$，
若只允许函子同构作为 $2$-态射，就可成为一个 $(2, 1)$-范畴。
\end{example}

\noindent
事实上，更一般地，任意 $2$-范畴 $\mathcal{C}$ 都产生一个
$(2, 1)$-范畴：取具有相同对象与 $1$-态射的子 $2$-范畴
$\mathcal{C}'$，但其 $2$-态射是 $\mathcal{C}$ 的可逆 $2$-态射。
在这种情形下，我们会说“{\it 令 $\mathcal{C}'$ 为与
$\mathcal{C}$ 关联的 $(2, 1)$-范畴}”或类似的话。例如，
群胚的 $(2, 1)$-范畴是指如下 $2$-范畴：对象是群胚，
$1$-态射是函子，$2$-态射是函子同构。不过这是个不好的例子，
因为群胚之间的函子之间的变换自动就是同构！

\begin{remark}
\label{categories-remark-other-2-categories}
因此，例 \ref{categories-example-2-1-category-of-categories} 的上述构造还有一些
变体，其中考虑群胚的 $2$-范畴、固定范畴上的群胚纤维化范畴，
或栈，等等。
\end{remark}






\section{2-纤维积}
\label{categories-section-2-fibre-products}

\noindent
本节介绍 $2$-纤维积。设 $\mathcal{C}$ 是一个 2-范畴。若图表
$$
\xymatrix{
w \ar[r] \ar[d] & y \ar[d] \\
x \ar[r] & z }
$$
中的两个 1-态射 $w \to y \to z$ 与 $w \to x \to z$ 是
2-同构的，则称该图表是 2-交换的。在 2-范畴中，要求图表
2-交换，比要求其严格交换更为自然。（事实上，有人或许会主张
根本不应使用严格 2-范畴；而在“弱”2-范畴中，1-态射图表
严格交换这一概念甚至没有意义。）相应地，纤维积的概念也必须调整。

\medskip\noindent
令 $\mathcal{C}$ 为 $2$-范畴。令 $x, y, z\in \Ob(\mathcal{C})$，并令
$f\in \Mor_\mathcal{C}(x, z)$ 与 $g\in \Mor_{\mathcal C}(y, z)$。
为了定义 $f$ 与 $g$ 的 2-纤维积，我们将考察 2-交换图表
$$
\xymatrix{
& w \ar[r]_a \ar[d]_b & x \ar[d]^{f} \\
& y \ar[r]^{g} & z. }
$$
在普通范畴的情形，纤维积是这类图表所成范畴中的终对象。相应地，
2-纤维积是某个 2-范畴中的终对象（见下述定义）。{\it 位于
$f$ 与 $g$ 上方的 $2$-交换图表的 $2$-范畴}是定义如下的 $2$-范畴：
\begin{enumerate}
\item 对象是如上的四元组 $(w, a, b, \phi)$，其中 $\phi$
是可逆 2-态射 $\phi : f \circ a \to g \circ b$；
\item 从 $(w', a', b', \phi')$ 到 $(w, a, b, \phi)$ 的 1-态射由
$(k : w' \to w, \alpha : a' \to a \circ k, \beta : b' \to b \circ k)$
给出，并使
$$
\xymatrix{
f \circ a'
\ar[rr]_{\text{id}_f \star \alpha}
\ar[d]_{\phi'}
& &
f \circ a \circ k
\ar[d]^{\phi \star \text{id}_k}
\\
g \circ b'
\ar[rr]^{\text{id}_g \star \beta}
& &
g \circ b \circ k
}
$$
交换；
\item 给定第二个 $1$-态射
$(k', \alpha', \beta') : (w'', a'', b'', \phi'') \to
(w', \alpha', \beta', \phi')$，$1$-态射的复合由下式给出：
$$
(k, \alpha, \beta) \circ (k', \alpha', \beta') =
(k \circ k',
(\alpha \star \text{id}_{k'}) \circ \alpha',
(\beta \star \text{id}_{k'}) \circ \beta'),
$$
\item 具有相同源与靶的两个 $1$-态射
$(k_i, \alpha_i, \beta_i)$（$i = 1, 2$）之间的一个 2-态射，
由一个 2-态射 $\delta : k_1 \to k_2$ 给出，并使
$$
\xymatrix{
a'
\ar[rd]_{\alpha_2}
\ar[r]_{\alpha_1} &
a \circ k_1
\ar[d]^{\text{id}_a \star \delta} &
&
b \circ k_1
\ar[d]_{\text{id}_b \star \delta} &
b'
\ar[l]^{\beta_1}
\ar[ld]^{\beta_2}
\\
&
a \circ k_2
&
&
b \circ k_2
&
}
$$
交换；
\item $2$-态射的竖直复合，由 $\mathcal{C}$ 中态射 $\delta$ 的
竖直复合给出；并且
\item 图表
$$
\xymatrix{
(w'', a'', b'', \phi'')
\rrtwocell^{(k'_1, \alpha'_1, \beta'_1)}_{(k'_2, \alpha'_2, \beta'_2)}{\delta'}
& &
(w', a', b', \phi')
\rrtwocell^{(k_1, \alpha_1, \beta_1)}_{(k_2, \alpha_2, \beta_2)}{\delta}
& &
(w, a, b, \phi)
}
$$
的水平复合由图表
$$
\xymatrix@C=12pc{
(w'', a'', b'', \phi'')
\rtwocell^{(k_1 \circ k'_1, (\alpha_1 \star \text{id}_{k'_1}) \circ \alpha'_1, (\beta_1 \star \text{id}_{k'_1}) \circ \beta'_1)}_{(k_2 \circ k'_2, (\alpha_2 \star \text{id}_{k'_2}) \circ \alpha'_2, (\beta_2 \star \text{id}_{k'_2}) \circ \beta'_2)}{\ \ \ \delta \star \delta'}
&
(w, a, b, \phi)
}
$$
给出。
\end{enumerate}
注意，如果 $\mathcal{C}$ 实际上是一个 $(2, 1)$-范畴，则上述 (2) 中的
态射 $\alpha$ 与 $\beta$ 自动也是同构\footnote{事实上，在 $2$-范畴的
情形，似乎还可以定义另一个 2-交换图表的 2-范畴，其中箭头
$\alpha$、$\beta$ 的方向反转，甚至只反转其中一个箭头的方向。这就是
我们稍后限制到 $(2, 1)$-范畴的原因。}。此外，如果 $\mathcal{C}$ 是
$(2, 1)$-范畴，那么 $2$-交换图表的 $2$-范畴也是 $(2, 1)$-范畴。

\begin{definition}
\label{categories-definition-final-object-2-category}
$(2, 1)$-范畴 $\mathcal{C}$ 的一个{\it 终对象}，是满足下列条件的对象
$x$：
\begin{enumerate}
\item 对每个 $y \in \Ob(\mathcal{C})$，都存在一个态射 $y \to x$；并且
\item 任意两个态射 $y \to x$ 都由唯一的 2-态射彼此同构。
\end{enumerate}
\end{definition}

\noindent
在更一般的 $2$-范畴情形，终对象很可能有不同的变体。我们不想在此
深入讨论，因而只在 $(2, 1)$-情形定义 $2$-纤维积。





\begin{definition}
\label{categories-definition-2-fibre-products}
令 $\mathcal{C}$ 为一个 $(2, 1)$-范畴。令
$x, y, z\in \Ob(\mathcal{C})$，并令
$f\in \Mor_\mathcal{C}(x, z)$
与 $g\in \Mor_{\mathcal C}(y, z)$。所谓{\it $f$ 与 $g$ 的
2-纤维积}，是上述 2-交换图表所成范畴中的一个终对象。
若 2-纤维积存在，我们将其记为
$x \times_z y\in \Ob(\mathcal{C})$，并将使图表
$$
\xymatrix{
& x \times_z y \ar[r]^{p} \ar[d]_q & x \ar[d]^{f} \\
& y \ar[r]^{g} & z }
$$
2-交换的所需态射记为
$p\in \Mor_{\mathcal C}(x \times_z y, x)$ 与
$q\in \Mor_{\mathcal C}(x \times_z y, y)$，将给出这一
2-交换性的可逆 2-态射记为 $\psi : f \circ p \to g \circ q$。
\end{definition}

\noindent
于是下列泛性质成立：对任意
$w\in \Ob(\mathcal{C})$ 以及态射
$a \in \Mor_{\mathcal C}(w, x)$ 和
$b \in \Mor_\mathcal{C}(w, y)$，若给定一个 2-同构
$\phi : f \circ a \to g\circ b$，
则存在 $\gamma \in \Mor_{\mathcal C}(w, x \times_z y)$，使图表
$$
\xymatrix{
w\ar[rrrd]^a \ar@{-->}[rrd]_\gamma \ar[rrdd]_b & & \\
& & x \times_z y \ar[r]_p \ar[d]_q & x \ar[d]^{f} \\
& & y \ar[r]^{g} & z }
$$
2-交换，并且在适当选择
$a \to p \circ \gamma$ 与 $b \to q \circ \gamma$ 后，图表
$$
\xymatrix{
f \circ a \ar[r] \ar[d]_\phi &
f \circ p \circ \gamma
\ar[d]^{\psi \star \text{id}_\gamma}
\\
g\circ b
\ar[r]
&
g \circ q \circ \gamma
}
$$
交换。此外，$\gamma$ 在同构意义下唯一。当然，精确的泛性质比这更
细致。我们以后所需的全部 2-纤维积情形，都来自下例：范畴的
2-范畴中的 2-纤维积。

\begin{example}
\label{categories-example-2-fibre-product-categories}
令 $\mathcal{A}$、$\mathcal{B}$ 和 $\mathcal{C}$ 为范畴。
令 $F : \mathcal{A} \to \mathcal{C}$ 与
$G : \mathcal{B} \to \mathcal{C}$ 为函子。如下定义范畴
$\mathcal{A} \times_\mathcal{C} \mathcal{B}$：
\begin{enumerate}
\item $\mathcal{A} \times_\mathcal{C} \mathcal{B}$ 的一个对象是三元组
$(A, B, f)$，其中 $A\in \Ob(\mathcal{A})$、
$B\in \Ob(\mathcal{B})$，而
$f : F(A) \to G(B)$ 是 $\mathcal{C}$ 中的一个同构；
\item 一个态射 $(A, B, f) \to (A', B', f')$ 由一对 $(a, b)$ 给出，
其中 $a : A \to A'$ 是 $\mathcal{A}$ 中的态射，
$b : B \to B'$ 是 $\mathcal{B}$ 中的态射，并使图表
$$
\xymatrix{
F(A) \ar[r]^f \ar[d]^{F(a)} & G(B) \ar[d]^{G(b)} \\
F(A') \ar[r]^{f'} & G(B')
}
$$
交换。
\end{enumerate}
此外，我们通过令
$$
p(A, B, f) = A, \quad q(A, B, f) = B,
$$
定义函子
$p : \mathcal{A} \times_\mathcal{C}\mathcal{B} \to \mathcal{A}$
和
$q : \mathcal{A} \times_\mathcal{C}\mathcal{B} \to \mathcal{B}$；
换言之，它们是遗忘函子。定义函子变换
$\psi : F \circ p \to G \circ q$。在对象
$\xi = (A, B, f)$ 上，它由
$\psi_\xi = f : F(p(\xi)) = F(A) \to G(B) = G(q(\xi))$ 给出。
\end{example}

\begin{lemma}
\label{categories-lemma-2-fibre-product-categories}
在范畴的 $(2, 1)$-范畴中，$2$-纤维积存在，并由
例 \ref{categories-example-2-fibre-product-categories} 的构造给出。
\end{lemma}

\begin{proof}
我们来验证泛性质。令 $\mathcal{W}$ 为范畴，令
$a : \mathcal{W} \to \mathcal{A}$ 与
$b : \mathcal{W} \to \mathcal{B}$ 为函子，并令
$t : F \circ a \to G \circ b$ 为函子同构。

\medskip\noindent
考虑由 $W \mapsto (a(W), b(W), t_W)$ 给出的函子
$\gamma : \mathcal{W} \to \mathcal{A} \times_\mathcal{C}\mathcal{B}$。
（省略对它是函子的验证。）再考虑由恒等式
$p \circ \gamma = a$ 与 $q \circ \gamma = b$ 得到的
$\alpha : a \to p \circ \gamma$ 和
$\beta : b \to q \circ \gamma$。于是显然，
$(\gamma, \alpha, \beta)$ 是从 $(W, a, b, t)$ 到
$(\mathcal{A} \times_\mathcal{C} \mathcal{B}, p, q, \psi)$ 的态射。

\medskip\noindent
令
$(k, \alpha', \beta') :
(W, a, b, t) \to (\mathcal{A} \times_\mathcal{C} \mathcal{B}, p, q, \psi)$
为第二个这样的态射。对 $\mathcal{W}$ 的一个对象 $W$，记
$k(W) = (a_k(W), b_k(W), t_{k, W})$。于是
$p(k(W)) = a_k(W)$，其余类同。映射 $\alpha'$ 对应于函子性映射
$\alpha' : a(W) \to a_k(W)$。由于我们是在范畴的
$(2, 1)$-范畴中工作，实际上每个映射
$a(W) \to a_k(W)$ 都是同构。利用这些同构（以及对应的
$b(W) \to b_k(W)$），可得到同构
$$
\delta_W :
\gamma(W) = (a(W), b(W), t_W)
\longrightarrow
(a_k(W), b_k(W), t_{k, W}) = k(W).
$$
直接可证 $\delta$ 定义了 $\gamma$ 与 $k$ 之间的
$2$-同构，它位于所需的 $2$-交换图表的 $2$-范畴中。
\end{proof}


\begin{remark}
\label{categories-remark-other-description-2-fibre-product}
令 $\mathcal{A}$、$\mathcal{B}$ 和 $\mathcal{C}$ 为范畴。
令 $F : \mathcal{A} \to \mathcal{C}$ 与
$G : \mathcal{B} \to \mathcal{C}$ 为函子。如下还可给出
$2$-纤维积 $\mathcal{A} \times_\mathcal{C} \mathcal{B}$ 的另一种、
稍为对称的构造。一个对象是五元组 $(A, B, C, a, b)$，其中
$A, B, C$ 分别是 $\mathcal{A}, \mathcal{B}, \mathcal{C}$ 的对象，
而 $a : F(A) \to C$ 与 $b : G(B) \to C$ 是同构。一个态射
$(A, B, C, a, b) \to (A', B', C', a', b')$ 由与态射
$a, b, a', b'$ 相容的三个态射
$A \to A', B \to B', C \to C'$ 给出。可以直接证明这会给出一个
$2$-纤维积。不过，更容易的做法是注意到，函子
$(A, B, C, a, b) \mapsto (A, B, b^{-1} \circ a)$ 给出了从五元组的
范畴到例 \ref{categories-example-2-fibre-product-categories} 所构造范畴的等价。
\end{remark}

\begin{lemma}
\label{categories-lemma-functoriality-2-fibre-product}
设
$$
\xymatrix{
& \mathcal{Y} \ar[d]_I \ar[rd]^K & \\
\mathcal{X} \ar[r]^H \ar[rd]^L &
\mathcal{Z} \ar[rd]^M & \mathcal{B} \ar[d]^G \\
& \mathcal{A} \ar[r]^F & \mathcal{C}
}
$$
是范畴的一个 $2$-交换图表。选择同构
$\alpha : G \circ K \to M \circ I$ 和
$\beta : M \circ H \to F \circ L$，就确定了与此情形相关的
$2$-纤维积之间的一个态射
$$
\mathcal{X} \times_\mathcal{Z} \mathcal{Y}
\longrightarrow
\mathcal{A} \times_\mathcal{C} \mathcal{B}
$$
\end{lemma}

\begin{proof}
在对象上使用函子
$$
(X, Y, \phi) \longmapsto (L(X), K(Y),
\alpha^{-1}_Y \circ M(\phi) \circ \beta^{-1}_X)
$$
并在态射上使用
$$
(a, b) \longmapsto (L(a), K(b))
$$
即可。
\end{proof}

\begin{lemma}
\label{categories-lemma-equivalence-2-fibre-product}
假设同引理 \ref{categories-lemma-functoriality-2-fibre-product}。
\begin{enumerate}
\item 若 $K$ 与 $L$ 忠实，则态射
$\mathcal{X} \times_\mathcal{Z} \mathcal{Y} \to
\mathcal{A} \times_\mathcal{C} \mathcal{B}$
忠实。
\item 若 $K$ 与 $L$ 全忠实而 $M$ 忠实，则态射
$\mathcal{X} \times_\mathcal{Z} \mathcal{Y} \to
\mathcal{A} \times_\mathcal{C} \mathcal{B}$
全忠实。
\item 若 $K$ 与 $L$ 是等价而 $M$ 全忠实，则态射
$\mathcal{X} \times_\mathcal{Z} \mathcal{Y} \to
\mathcal{A} \times_\mathcal{C} \mathcal{B}$
是等价。
\end{enumerate}
\end{lemma}

\begin{proof}
令 $(X, Y, \phi)$ 与 $(X', Y', \phi')$ 为
$\mathcal{X} \times_\mathcal{Z} \mathcal{Y}$ 的对象。置
$Z = H(X)$，并通过 $\phi$ 将其与 $I(Y)$ 等同。还通过
$\alpha_X$ 将 $M(Z)$ 与 $F(L(X))$ 等同，通过
$\beta_Y$ 将 $M(Z)$ 与 $G(K(Y))$ 等同。对
$Z' = H(X')$ 与 $M(Z')$ 也作类似处理。态射上的映射是
$$
\xymatrix{
\Mor_\mathcal{X}(X, X')
\times_{\Mor_\mathcal{Z}(Z, Z')}
\Mor_\mathcal{Y}(Y, Y')
\ar[d] \\
\Mor_\mathcal{A}(L(X), L(X'))
\times_{\Mor_\mathcal{C}(M(Z), M(Z'))}
\Mor_\mathcal{B}(K(Y), K(Y'))
}
$$
这个映射。因此 (1) 与 (2) 成立。进一步，若 $K$ 与 $L$
是等价而 $M$ 全忠实，则任意对象 $(A, B, \phi)$ 都在本质像中，
理由如下：选取 $X$、$Y$，使
$L(X) \cong A$ 且 $K(Y) \cong B$。于是 $M$ 的全忠实性保证
我们可以找到一个同构 $H(X) \cong I(Y)$。略去若干细节。
\end{proof}

\begin{lemma}
\label{categories-lemma-associativity-2-fibre-product}
设
$$
\xymatrix{
\mathcal{A} \ar[rd] & & \mathcal{C} \ar[ld] \ar[rd] & & \mathcal{E} \ar[ld] \\
& \mathcal{B} & & \mathcal{D}
}
$$
是范畴和函子的一个图表。那么存在范畴的典范同构
$$
(\mathcal{A} \times_\mathcal{B} \mathcal{C}) \times_\mathcal{D} \mathcal{E}
\cong
\mathcal{A} \times_\mathcal{B} (\mathcal{C} \times_\mathcal{D} \mathcal{E})
$$
\end{lemma}

\begin{proof}
如果你明白这里的意思，只需使用函子
$$
((A, C, \phi), E, \psi)
\longmapsto
(A, (C, E, \psi), \phi)
$$
即可。
\end{proof}

\noindent
今后处理两个以上范畴的纤维积时，我们不再书写括号。


\begin{lemma}
\label{categories-lemma-triple-2-fibre-product-pr02}
设
$$
\xymatrix{
\mathcal{A} \ar[rd] & & \mathcal{C} \ar[ld] \ar[rd] & & \mathcal{E} \ar[ld] \\
& \mathcal{B} \ar[rd]_F & & \mathcal{D} \ar[ld]^G \\
& & \mathcal{F} &
}
$$
是范畴和函子的一个交换图表。那么存在范畴的典范函子
$$
\text{pr}_{02} :
\mathcal{A} \times_\mathcal{B} \mathcal{C} \times_\mathcal{D} \mathcal{E}
\longrightarrow
\mathcal{A} \times_\mathcal{F} \mathcal{E}
$$
\end{lemma}

\begin{proof}
若将
$\mathcal{A} \times_\mathcal{B} \mathcal{C}
\times_\mathcal{D} \mathcal{E}$
写成
$(\mathcal{A} \times_\mathcal{B} \mathcal{C})
\times_\mathcal{D} \mathcal{E}$，
那么，如果你明白我的意思，只需使用函子
$$
((A, C, \phi), E, \psi)
\longmapsto
(A, E, G(\psi) \circ F(\phi))
$$
即可。
\end{proof}

\begin{lemma}
\label{categories-lemma-2-fibre-product-erase-factor}
设
$$
\mathcal{A} \to
\mathcal{B} \leftarrow \mathcal{C} \leftarrow \mathcal{D}
$$
是范畴和函子的一个图表。那么存在范畴的典范同构
$$
\mathcal{A} \times_\mathcal{B} \mathcal{C} \times_\mathcal{C} \mathcal{D}
\cong
\mathcal{A} \times_\mathcal{B} \mathcal{D}
$$
\end{lemma}

\begin{proof}
略。
\end{proof}

\noindent
我们断言，这意味着可以像使用普通纤维积一样使用这些 $2$-纤维积。
下面是稍后确实会用到的一些进一步引理。

\begin{lemma}
\label{categories-lemma-diagonal-1}
设
$$
\xymatrix{
\mathcal{C}_3 \ar[r] \ar[d] & \mathcal{S} \ar[d]^\Delta \\
\mathcal{C}_1 \times \mathcal{C}_2 \ar[r]^{G_1 \times G_2} &
\mathcal{S} \times \mathcal{S}
}
$$
是范畴的一个 $2$-纤维积。那么存在典范同构
$\mathcal{C}_3 \cong
\mathcal{C}_1 \times_{G_1, \mathcal{S}, G_2} \mathcal{C}_2$。
\end{lemma}

\begin{proof}
可以假设 $\mathcal{C}_3$ 是例
\ref{categories-example-2-fibre-product-categories} 中构造的范畴
$(\mathcal{C}_1 \times \mathcal{C}_2)\times_{\mathcal{S} \times \mathcal{S}}
\mathcal{S}$。于是，一个对象是三元组
$((X_1, X_2), S, \phi)$，其中
$\phi = (\phi_1, \phi_2) : (G_1(X_1), G_2(X_2)) \to (S, S)$
是同构。因此可以将三元组
$(X_1, X_2, \phi_2^{-1} \circ \phi_1)$ 与之对应。反过来，若
$(X_1, X_2, \psi)$ 是
$\mathcal{C}_1 \times_{G_1, \mathcal{S}, G_2} \mathcal{C}_2$ 的对象，
则可以将三元组
$((X_1, X_2), G_2(X_2), (\psi, \text{id}_{G_2(X_2)}))$
与之对应。我们断言，这两个构造给出互逆函子。略去对态射的处理，
以及证明它们互逆的过程。
\end{proof}

\begin{lemma}
\label{categories-lemma-diagonal-2}
设
$$
\xymatrix{
\mathcal{C}' \ar[r] \ar[d] & \mathcal{S} \ar[d]^\Delta \\
\mathcal{C} \ar[r]^{(G_1, G_2)} &
\mathcal{S} \times \mathcal{S}
}
$$
是范畴的一个 $2$-纤维积。那么存在典范同构
$$
\mathcal{C}' \cong
(\mathcal{C} \times_{G_1, \mathcal{S}, G_2} \mathcal{C})
\times_{(p, q), \mathcal{C} \times \mathcal{C}, \Delta}
\mathcal{C}.
$$
\end{lemma}

\begin{proof}
右端的一个对象由
$((C_1, C_2, \phi), C_3, \psi)$ 给出，其中
$\phi : G_1(C_1) \to G_2(C_2)$ 是同构，而
$\psi = (\psi_1, \psi_2) : (C_1, C_2) \to (C_3, C_3)$
是同构。因此可以将三元组
$(C_3, G_1(C_1), (G_1(\psi_1^{-1}), \phi^{-1} \circ G_2(\psi_2^{-1})))$
与之对应；它是 $\mathcal{C}'$ 的对象。略去细节。
\end{proof}

\begin{lemma}
\label{categories-lemma-fibre-product-after-map}
令 $\mathcal{A} \to \mathcal{C}$、$\mathcal{B} \to \mathcal{C}$
以及 $\mathcal{C} \to \mathcal{D}$ 为范畴之间的函子。那么图表
$$
\xymatrix{
\mathcal{A} \times_\mathcal{C} \mathcal{B} \ar[d] \ar[r] &
\mathcal{A} \times_\mathcal{D} \mathcal{B} \ar[d] \\
\mathcal{C} \ar[r]^-{\Delta_{\mathcal{C}/\mathcal{D}}} \ar[r] &
\mathcal{C} \times_\mathcal{D} \mathcal{C}
}
$$
是一个 $2$-纤维积图表。
\end{lemma}

\begin{proof}
略。
\end{proof}

\begin{lemma}
\label{categories-lemma-base-change-diagonal}
设
$$
\xymatrix{
\mathcal{U} \ar[d] \ar[r] & \mathcal{V} \ar[d] \\
\mathcal{X} \ar[r] & \mathcal{Y}
}
$$
是范畴的一个 $2$-纤维积。那么图表
$$
\xymatrix{
\mathcal{U} \ar[d] \ar[r] &
\mathcal{U} \times_\mathcal{V} \mathcal{U} \ar[d] \\
\mathcal{X} \ar[r] &
\mathcal{X} \times_\mathcal{Y} \mathcal{X}
}
$$
是 $2$-笛卡尔的。
\end{lemma}

\begin{proof}
这是一个纯粹的 $2$-范畴论陈述，在任何具有 $2$-纤维积的
$(2, 1)$-范畴中都成立。具体地说，它由下列等价链推出：
\begin{align*}
\mathcal{X} \times_{(\mathcal{X} \times_\mathcal{Y} \mathcal{X})}
(\mathcal{U} \times_\mathcal{V} \mathcal{U})
& =
\mathcal{X} \times_{(\mathcal{X} \times_\mathcal{Y} \mathcal{X})}
((\mathcal{X} \times_\mathcal{Y} \mathcal{V})
\times_\mathcal{V} (\mathcal{X} \times_\mathcal{Y} \mathcal{V})) \\
& =
\mathcal{X} \times_{(\mathcal{X} \times_\mathcal{Y} \mathcal{X})}
(\mathcal{X} \times_\mathcal{Y} \mathcal{X}
\times_\mathcal{Y} \mathcal{V}) \\
& =
\mathcal{X} \times_\mathcal{Y} \mathcal{V} = \mathcal{U}
\end{align*}
参见引理 \ref{categories-lemma-associativity-2-fibre-product} 和
\ref{categories-lemma-2-fibre-product-erase-factor}。
\end{proof}


\section{范畴上的范畴}
\label{categories-section-categories-over-categories}

\noindent
本节给定一个函子 $p : \mathcal{S} \to \mathcal{C}$。我们把
$\mathcal{S}$ 看作位于上方，把 $\mathcal{C}$ 看作位于下方。
为了明确所谈的对象，我们定义 $\mathcal{C}$ 上的范畴所成的
$2$-范畴。

\begin{definition}
\label{categories-definition-categories-over-C}
令 $\mathcal{C}$ 为范畴。所谓{\it $\mathcal{C}$ 上的范畴所成的
$2$-范畴}，是按如下方式定义的 $2$-范畴：
\begin{enumerate}
\item 它的对象是函子 $p : \mathcal{S} \to \mathcal{C}$。
\item 它的 $1$-态射
$(\mathcal{S}, p) \to (\mathcal{S}', p')$ 是满足
$p' \circ G = p$ 的函子 $G : \mathcal{S} \to \mathcal{S}'$。
\item 对
$G, H : (\mathcal{S}, p) \to (\mathcal{S}', p')$，它的 $2$-态射
$t : G \to H$ 是满足下述条件的函子态射：
对所有 $x \in \Ob(\mathcal{S})$，均有
$p'(t_x) = \text{id}_{p(x)}$。
\end{enumerate}
在这种情形，我们将 $(\mathcal{S}, p)$ 与
$(\mathcal{S}', p')$ 之间 $1$-态射所成的范畴记为
$$
\Mor_{\textit{Cat}/\mathcal{C}}(\mathcal{S}, \mathcal{S}')
$$
\end{definition}

\noindent
在这个 $2$-范畴中，我们完全按照
第 \ref{categories-section-formal-cat-cat} 节中对 $\textit{Cat}$ 的做法，
定义水平复合与竖直复合。$2$-范畴的公理成立，理由与它们在
$\textit{Cat}$ 中成立的理由相同。也可以利用这些公理在
$\textit{Cat}$ 中成立这一事实，再验证诸如“$\mathcal{C}$ 上的
$2$-态射的竖直复合仍给出 $\mathcal{C}$ 上的 $2$-态射”之类的
结论。此事显然。

\medskip\noindent
类似于空间之间映射的纤维，这里有纤维范畴的概念，以及与此情形
相关的一些提升概念。

\begin{definition}
\label{categories-definition-fibre-category}
令 $\mathcal{C}$ 为范畴。令
$p : \mathcal{S} \to \mathcal{C}$ 为 $\mathcal{C}$ 上的范畴。
\begin{enumerate}
\item 对象 $U\in \Ob(\mathcal{C})$ 上的{\it 纤维范畴}，是具有对象
$$
\Ob(\mathcal{S}_U) = \{x\in \Ob(\mathcal{S}) :
p(x) = U\}
$$
和态射
$$
\Mor_{\mathcal{S}_U}(x, y) = \{ \phi \in \Mor_\mathcal{S}(x, y) :
p(\phi) = \text{id}_U\}.
$$
的范畴 $\mathcal{S}_U$。
\item 对象 $U \in \Ob(\mathcal{C})$ 的一个{\it 提升}，是满足
$p(x) = U$ 的对象 $x\in \Ob(\mathcal{S})$；也就是说，
$x\in \Ob(\mathcal{S}_U)$。有时也说{\it $x$ 位于 $U$ 之上}。
\item 类似地，$\mathcal{C}$ 中态射 $f : V \to U$ 的一个
{\it 提升}，是 $\mathcal{S}$ 中满足 $p(\phi) = f$ 的态射
$\phi : y \to x$。有时也说{\it $\phi$ 位于 $f$ 之上}。
\end{enumerate}
\end{definition}

\noindent
这里可以作出一些观察。例如，若
$F : (\mathcal{S}, p) \to (\mathcal{S}', p')$ 是
$\mathcal{C}$ 上范畴的一个 $1$-态射，则 $F$ 诱导纤维范畴之间的
函子 $F : \mathcal{S}_U \to \mathcal{S}'_U$。对 $2$-态射亦同。

\medskip\noindent
下面给出必不可少的引理，它描述 $\mathcal{C}$ 上范畴所成的
$(2, 1)$-范畴中的 $2$-纤维积。

\begin{lemma}
\label{categories-lemma-2-product-categories-over-C}
令 $\mathcal{C}$ 为范畴。$\mathcal{C}$ 上的范畴所成的
$(2, 1)$-范畴具有 2-纤维积。设
$F : \mathcal{X} \to \mathcal{S}$ 与
$G : \mathcal{Y} \to \mathcal{S}$ 是 $\mathcal{C}$ 上范畴的态射。
一个显式的 2-纤维积
$\mathcal{X} \times_\mathcal{S}\mathcal{Y}$ 由下述描述给出：
\begin{enumerate}
\item $\mathcal{X} \times_\mathcal{S} \mathcal{Y}$ 的一个对象是四元组
$(U, x, y, f)$，其中 $U \in \Ob(\mathcal{C})$、
$x\in \Ob(\mathcal{X}_U)$、$y\in \Ob(\mathcal{Y}_U)$，且
$f : F(x) \to G(y)$ 是 $\mathcal{S}_U$ 中的同构；
\item 一个态射
$(U, x, y, f) \to (U', x', y', f')$ 由一对 $(a, b)$ 给出，
其中 $a : x \to x'$ 是 $\mathcal{X}$ 中的态射，
$b : y \to y'$ 是 $\mathcal{Y}$ 中的态射，并满足：
\begin{enumerate}
\item $a$ 与 $b$ 诱导同一个态射 $U \to U'$；并且
\item 图表
$$
\xymatrix{
F(x) \ar[r]^f \ar[d]^{F(a)} & G(y) \ar[d]^{G(b)} \\
F(x') \ar[r]^{f'} & G(y')
}
$$
交换。
\end{enumerate}
\end{enumerate}
在这种情形，函子
$p : \mathcal{X} \times_\mathcal{S}\mathcal{Y} \to \mathcal{X}$
和 $q : \mathcal{X} \times_\mathcal{S}\mathcal{Y} \to \mathcal{Y}$
是遗忘函子。变换 $\psi : F \circ p \to G \circ q$ 在对象
$\xi = (U, x, y, f)$ 上由
$\psi_\xi = f : F(p(\xi)) = F(x) \to G(y) = G(q(\xi))$ 给出。
\end{lemma}

\begin{proof}
我们来验证泛性质。令
$p_\mathcal{W} : \mathcal{W}\to \mathcal{C}$ 为
$\mathcal{C}$ 上的范畴，令
$X : \mathcal{W} \to \mathcal{X}$ 和
$Y : \mathcal{W} \to \mathcal{Y}$ 为 $\mathcal{C}$ 上的函子，
并令 $t : F \circ X \to G \circ Y$ 为 $\mathcal{C}$ 上的函子同构。
所需函子
$\gamma : \mathcal{W} \to \mathcal{X} \times_\mathcal{S} \mathcal{Y}$
由 $W \mapsto (p_\mathcal{W}(W), X(W), Y(W), t_W)$ 给出。
略去细节；比较引理 \ref{categories-lemma-2-fibre-product-categories}。
\end{proof}

\begin{example}
\label{categories-example-different-2-fibre-products}
引理 \ref{categories-lemma-2-product-categories-over-C} 中给出的范畴上的范畴之
$2$-纤维积构造，与引理 \ref{categories-lemma-2-fibre-product-categories}
中给出的范畴之 2-纤维积构造（如例
\ref{categories-example-2-fibre-product-categories}），一般会产生不等价的
结果。具体地，令 $\mathcal{S}$ 为仅有一个对象和两个箭头的群胚，
令 $\mathcal{X}$ 为仅有一个对象的离散范畴。将
$2$-纤维积 $\mathcal{X} \times_\mathcal{S} \mathcal{X}$
取作范畴，会得到含两个对象的离散范畴。然而，如果把它们全都视为
$\mathcal{S}$ 上的范畴，那么作为 $\mathcal{S}$ 上范畴的
$2$-纤维积 $\mathcal{X} \times_\mathcal{S} \mathcal{X}$
是仅有一个对象的离散范畴。区别在于（沿用引理
\ref{categories-lemma-2-product-categories-over-C} 的记号），在第一种情形
可以选择任意比较同构 $f$，而在第二种情形只能选择恒等箭头。
\end{example}

\begin{lemma}
\label{categories-lemma-fibre-2-fibre-product-categories-over-C}
令 $\mathcal{C}$ 为范畴。令
$f : \mathcal{X} \to \mathcal{S}$ 和
$g : \mathcal{Y} \to \mathcal{S}$ 为 $\mathcal{C}$ 上范畴的态射。
对 $\mathcal{C}$ 的任意对象 $U$，纤维范畴有下列恒等式：
$$
\left(\mathcal{X} \times_\mathcal{S}\mathcal{Y}\right)_U
=
\mathcal{X}_U \times_{\mathcal{S}_U} \mathcal{Y}_U
$$
\end{lemma}

\begin{proof}
略。
\end{proof}


\section{纤维化范畴}
\label{categories-section-fibred-categories}

\noindent
这里有必要非常简略地讨论纤维化范畴。

\medskip\noindent
令 $p : \mathcal{S} \to \mathcal{C}$ 为 $\mathcal{C}$ 上的范畴。
给定满足 $p(x) = U$ 的对象 $x \in \mathcal{S}$，以及态射
$f : V \to U$，可以尝试取某种“纤维积 $V \times_U x$”
（或者说，通过 $V \to U$ 对 $x$ 作一次{\it 基变换}）。具体地，
从对象 $z \in \mathcal{S}$ 到“$V \times_U x$”的一个态射，
应由一对 $(\varphi, g)$ 给出，其中
$\varphi : z \to x$、$g : p(z) \to V$，且
$p(\varphi) = f \circ g$。图示如下：
$$
\xymatrix{
z \ar@{~>}[d]^p \ar@{-}[r] &
? \ar[r] \ar@{~>}[d]^p &
x \ar@{~>}[d]^p \\
p(z) \ar[r] & V \ar[r]^f & U
}
$$
若这样的态射 $V \times_U x \to x$ 存在，则称其为强笛卡尔态射。

\begin{definition}
\label{categories-definition-cartesian-over-C}
令 $\mathcal{C}$ 为范畴。令
$p : \mathcal{S} \to \mathcal{C}$ 为 $\mathcal{C}$ 上的范畴。
一个{\it 强笛卡尔态射}，更精确地说，一个
{\it 强 $\mathcal{C}$-笛卡尔态射}，是 $\mathcal{S}$ 的态射
$\varphi : y \to x$，使得对每个 $z \in \Ob(\mathcal{S})$，
由 $\psi \longmapsto (\varphi \circ \psi, p(\psi))$ 给出的映射
$$
\Mor_\mathcal{S}(z, y)
\longrightarrow
\Mor_\mathcal{S}(z, x)
\times_{\Mor_\mathcal{C}(p(z), p(x))}
\Mor_\mathcal{C}(p(z), p(y)),
$$
是双射。
\end{definition}

\noindent
注意，由米田引理 \ref{categories-lemma-yoneda}，给定一个位于
$U \in \Ob(\mathcal{C})$ 之上的对象
$x \in \Ob(\mathcal{S})$，以及 $\mathcal{C}$ 的态射
$f : V \to U$，如果存在满足 $p(\varphi) = f$ 的强笛卡尔态射
$\varphi : y \to x$，那么 $(y, \varphi)$ 在唯一同构意义下唯一。
这从上述定义即显然，因为函子
$$
z
\longmapsto
\Mor_\mathcal{S}(z, x)
\times_{\Mor_\mathcal{C}(p(z), U)}
\Mor_\mathcal{C}(p(z), V)
$$
只依赖于数据 $(x, U, f : V \to U)$。因此，我们有时会用
$V \times_U x \to x$ 或 $f^*x \to x$ 表示作为 $f$ 的提升的
强笛卡尔态射。

\begin{lemma}
\label{categories-lemma-composition-cartesian}
令 $\mathcal{C}$ 为范畴。令
$p : \mathcal{S} \to \mathcal{C}$ 为 $\mathcal{C}$ 上的范畴。
\begin{enumerate}
\item 两个强笛卡尔态射的复合是强笛卡尔的。
\item $\mathcal{S}$ 的任意同构都是强笛卡尔的。
\item 若强笛卡尔态射 $\varphi$ 的 $p(\varphi)$ 是同构，则
该态射是同构。
\end{enumerate}
\end{lemma}

\begin{proof}
证明 (1)。令 $\varphi : y \to x$ 与 $\psi : z \to y$ 为
强笛卡尔态射。令 $t$ 为 $\mathcal{S}$ 的任意对象。于是
\begin{align*}
& \Mor_\mathcal{S}(t, z) \\
& =
\Mor_\mathcal{S}(t, y)
\times_{\Mor_\mathcal{C}(p(t), p(y))}
\Mor_\mathcal{C}(p(t), p(z)) \\
& =
\Mor_\mathcal{S}(t, x)
\times_{\Mor_\mathcal{C}(p(t), p(x))}
\Mor_\mathcal{C}(p(t), p(y))
\times_{\Mor_\mathcal{C}(p(t), p(y))}
\Mor_\mathcal{C}(p(t), p(z)) \\
& =
\Mor_\mathcal{S}(t, x)
\times_{\Mor_\mathcal{C}(p(t), p(x))}
\Mor_\mathcal{C}(p(t), p(z))
\end{align*}
故 $z \to x$ 是强笛卡尔的。

\medskip\noindent
证明 (2)。令 $y \to x$ 为同构。那么 $p(y) \to p(x)$ 也是同构。
因此
$\Mor_\mathcal{C}(p(z), p(y)) \to
\Mor_\mathcal{C}(p(z), p(x))$
是双射。于是
$\Mor_\mathcal{S}(z, x)
\times_{\Mor_\mathcal{C}(p(z), p(x))}
\Mor_\mathcal{C}(p(z), p(y))$
与 $\Mor_\mathcal{S}(z, x)$ 双射。因此，定义
\ref{categories-definition-cartesian-over-C} 中展示的映射是双射，因为
$y \to x$ 是同构；由此可知 $y \to x$ 是强笛卡尔的。

\medskip\noindent
证明 (3)。设 $\varphi : y \to x$ 是强笛卡尔的，且
$p(\varphi) : p(y) \to p(x)$ 是同构。将定义用于
$z = x$，可知 $(\text{id}_x, p(\varphi)^{-1})$ 来自唯一态射
$\chi : x \to y$。略去对 $\chi$ 是 $\varphi$ 之逆的验证。
\end{proof}

\begin{lemma}
\label{categories-lemma-cartesian-over-cartesian}
令 $F : \mathcal{A} \to \mathcal{B}$ 与
$G : \mathcal{B} \to \mathcal{C}$ 为范畴之间可复合的函子。
令 $x \to y$ 为 $\mathcal{A}$ 的态射。若 $x \to y$ 是强
$\mathcal{B}$-笛卡尔的，且 $F(x) \to F(y)$ 是强
$\mathcal{C}$-笛卡尔的，则 $x \to y$ 是强
$\mathcal{C}$-笛卡尔的。
\end{lemma}

\begin{proof}
这直接由定义推出。
\end{proof}

\begin{lemma}
\label{categories-lemma-strongly-cartesian-fibre-product}
令 $\mathcal{C}$ 为范畴。令
$p : \mathcal{S} \to \mathcal{C}$ 为 $\mathcal{C}$ 上的范畴。
令 $x \to y$ 与 $z \to y$ 为 $\mathcal{S}$ 的态射。假设
\begin{enumerate}
\item $x \to y$ 是强笛卡尔的；
\item $p(x) \times_{p(y)} p(z)$ 存在；并且
\item 在 $\mathcal{S}$ 中存在强笛卡尔态射 $a : w \to z$，满足
$p(w) = p(x) \times_{p(y)} p(z)$ 且
$p(a) = \text{pr}_2 : p(x) \times_{p(y)} p(z) \to p(z)$。
\end{enumerate}
那么纤维积 $x \times_y z$ 存在，并同构于 $w$。
\end{lemma}

\begin{proof}
由于 $x \to y$ 是强笛卡尔的，存在唯一态射
$b : w \to x$，使 $p(b) = \text{pr}_1$。为了说明 $w$ 是纤维积，
我们计算
\begin{align*}
& \Mor_\mathcal{S}(t, w) \\
& = \Mor_\mathcal{S}(t, z)
\times_{\Mor_\mathcal{C}(p(t), p(z))}
\Mor_\mathcal{C}(p(t), p(w)) \\
& = \Mor_\mathcal{S}(t, z)
\times_{\Mor_\mathcal{C}(p(t), p(z))}
(\Mor_\mathcal{C}(p(t), p(x))
\times_{\Mor_\mathcal{C}(p(t), p(y))}
\Mor_\mathcal{C}(p(t), p(z))) \\
& = \Mor_\mathcal{S}(t, z)
\times_{\Mor_\mathcal{C}(p(t), p(y))}
\Mor_\mathcal{C}(p(t), p(x)) \\
& = \Mor_\mathcal{S}(t, z)
\times_{\Mor_\mathcal{S}(t, y)}
\Mor_\mathcal{S}(t, y)
\times_{\Mor_\mathcal{C}(p(t), p(y))}
\Mor_\mathcal{C}(p(t), p(x)) \\
& = \Mor_\mathcal{S}(t, z)
\times_{\Mor_\mathcal{S}(t, y)}
\Mor_\mathcal{S}(t, x)
\end{align*}
正如所需。第一个等式成立，是因为 $a : w \to z$ 是强笛卡尔的；
最后一个等式成立，是因为 $x \to y$ 是强笛卡尔的。
\end{proof}

\begin{definition}
\label{categories-definition-fibred-category}
令 $\mathcal{C}$ 为范畴。令
$p : \mathcal{S} \to \mathcal{C}$ 为 $\mathcal{C}$ 上的范畴。
若对任意位于 $U \in \Ob(\mathcal{C})$ 之上的
$x \in \Ob(\mathcal{S})$，以及 $\mathcal{C}$ 的任意态射
$f : V \to U$，都存在位于 $f$ 之上的强笛卡尔态射
$f^*x \to x$，则称 $\mathcal{S}$ 是一个
{\it $\mathcal{C}$ 上的纤维化范畴}。
\end{definition}


\noindent
假设 $p : \mathcal{S} \to \mathcal{C}$ 是纤维化范畴。对定义中那样的
每个 $f : V \to U$ 和 $x\in \Ob(\mathcal{S}_U)$，可以选择一个
位于 $f$ 之上的强笛卡尔态射 $f^\ast x \to x$。由选择公理，
可以同时为所有 $f: V \to U = p(x)$ 选择
$f^*x \to x$。我们断言，对 $\mathcal{S}_U$ 中的每个态射
$\phi : x \to x'$ 以及每个 $f : V \to U$，都存在唯一的
$\mathcal{S}_V$ 中的态射
$f^\ast \phi : f^\ast x \to f^\ast x'$，使图表
$$
\xymatrix{
f^\ast x \ar[r]_{f^\ast \phi} \ar[d] & f^\ast x' \ar[d] \\
x \ar[r]^{\phi} & x' }
$$
交换。事实上，该箭头存在且唯一，是因为 $f^*x' \to x'$
是强笛卡尔的。这个箭头的唯一性保证 $f^\ast$（现在也已在态射上
定义）是函子 $ f^\ast : \mathcal{S}_U \to \mathcal{S}_V$。

\begin{definition}
\label{categories-definition-pullback-functor-fibred-category}
假设 $p : \mathcal{S} \to \mathcal{C}$ 是纤维化范畴。
\begin{enumerate}
\item 对 $p : \mathcal{S} \to \mathcal{C}$ 的一个
{\it 拉回选择}\footnote{这很可能不是标准术语。在一些文献中，这被
称为“cleavage”，但它会唤起错误的形象。也许“cleaving”会是更好的
词。一个相关概念是“splitting”，但在许多文献中，“splitting”是指
满足如下条件的拉回选择：对任意一对可复合态射，都有
$g^*f^* = (f \circ g)^*$。另见定义
\ref{categories-definition-split-fibred-category}。}
是指：对 $\mathcal{C}$ 的任意态射 $f: V \to U$ 和任意
$x \in \Ob(\mathcal{S}_U)$，选择一个位于 $f$ 之上的强笛卡尔态射
$f^\ast x \to x$。
\item 给定一个拉回选择，对 $\mathcal{C}$ 的任意态射
$f : V \to U$，上述函子
$f^* : \mathcal{S}_U \to \mathcal{S}_V$ 称为一个
{\it 拉回函子}（它与上面作出的选择 $f^*x \to x$ 相关）。
\end{enumerate}
\end{definition}

\noindent
当然，总可以假设我们的拉回选择具有
$\text{id}_U^*x = x$ 这一性质，不过在实践中，如果不对拉回施加
进一步假设，这个性质没有用处。

\begin{lemma}
\label{categories-lemma-fibred}
假设 $p : \mathcal{S} \to \mathcal{C}$ 是纤维化范畴，并且已经
给定 $p : \mathcal{S} \to \mathcal{C}$ 的一个拉回选择。
\begin{enumerate}
\item 对任意一对可复合态射 $f : V \to U$、
$g : W \to V$，存在唯一的函子同构
$$
\alpha_{g, f} :
(f \circ g)^\ast
\longrightarrow
g^\ast \circ f^\ast
$$
（作为函子 $\mathcal{S}_U \to \mathcal{S}_W$），使得对每个
$y\in \Ob(\mathcal{S}_U)$，下图交换：
$$
\xymatrix{
g^\ast f^\ast y \ar[r]
&
f^\ast y \ar[d] \\
(f \circ g)^\ast y \ar[r]
\ar[u]^{(\alpha_{g, f})_y}
&
y
}
$$
\item 若 $f = \text{id}_U$，则存在函子
$\mathcal{S}_U \to \mathcal{S}_U$ 的典范同构
$\alpha_U : \text{id} \to (\text{id}_U)^*$。
\item 四元组
$(U \mapsto \mathcal{S}_U, f \mapsto f^*, \alpha_{g, f}, \alpha_U)$
定义了从 $\mathcal{C}^{opp}$ 到范畴的 $(2, 1)$-范畴的一个
伪函子；见定义 \ref{categories-definition-functor-into-2-category}。
\end{enumerate}
\end{lemma}

\begin{proof}
事实上，显然 (1) 中的交换图表唯一确定纤维范畴
$\mathcal{S}_W$ 中的态射 $(\alpha_{g, f})_y$。它是同构，因为态射
$(f \circ g)^*y \to y$ 与复合
$g^*f^*y \to f^*y \to y$ 都是提升 $f \circ g$ 的强笛卡尔态射
（见定义 \ref{categories-definition-cartesian-over-C} 后的讨论和引理
\ref{categories-lemma-composition-cartesian}）。同理，因为
$\text{id}_x : x \to x$ 显然在 $\text{id}_U$ 上方是强笛卡尔的
（其中 $U = p(x)$），所以存在同构
$(\alpha_U)_x : x \to (\text{id}_U)^*x$。（当然，我们本来可以
预先假设，只要 $f$ 是恒等态射就有 $f^*x = x$；但为了保持一般性，
最好不作此假设。）略去对这样得到的 $\alpha_{g, f}$ 和
$\alpha_U$ 是函子变换的验证，也略去对 (3) 的验证。
\end{proof}

\begin{lemma}
\label{categories-lemma-fibred-equivalent}
令 $\mathcal{C}$ 为范畴。令 $\mathcal{S}_1$、$\mathcal{S}_2$
为 $\mathcal{C}$ 上的范畴。假设 $\mathcal{S}_1$ 与
$\mathcal{S}_2$ 作为 $\mathcal{C}$ 上的范畴等价。那么
$\mathcal{S}_1$ 在 $\mathcal{C}$ 上纤维化，当且仅当
$\mathcal{S}_2$ 在 $\mathcal{C}$ 上纤维化。
\end{lemma}

\begin{proof}
将给定函子记为 $p_i : \mathcal{S}_i \to \mathcal{C}$。令
$F : \mathcal{S}_1 \to \mathcal{S}_2$、
$G : \mathcal{S}_2 \to \mathcal{S}_1$ 为 $\mathcal{C}$ 上的函子，
并令
$i : F \circ G \to \text{id}_{\mathcal{S}_2}$、
$j : G \circ F \to \text{id}_{\mathcal{S}_1}$ 为
$\mathcal{C}$ 上的函子同构。我们断言，在这种情形，$F$ 将强
笛卡尔态射映为强笛卡尔态射。事实上，设
$\varphi : y \to x$ 是 $\mathcal{S}_1$ 中的强笛卡尔态射。
令 $f : V \to U$ 等于 $p_1(\varphi)$。假设
$z' \in \Ob(\mathcal{S}_2)$，其中 $W = p_2(z')$，并且给定
$g : W \to V$ 和 $\psi' : z' \to F(x)$，使
$p_2(\psi') = f \circ g$。那么
$$
\psi = j \circ G(\psi') : G(z') \to G(F(x)) \to x
$$
是 $\mathcal{S}_1$ 中满足 $p_1(\psi) = f \circ g$ 的态射。
因此，由假设，存在唯一的位于 $g$ 之上的态射
$\xi : G(z') \to y$，使 $\psi = \varphi \circ \xi$。这又给出态射
$$
\xi' = F(\xi) \circ i^{-1} : z' \to F(G(z')) \to F(y)
$$
它位于 $g$ 之上，并满足 $\psi' = F(\varphi) \circ \xi'$。
略去对 $\xi'$ 唯一性的验证。
\end{proof}

\noindent
由引理 \ref{categories-lemma-fibred-equivalent} 可知，等价将强笛卡尔态射映为
强笛卡尔态射。但对 $\mathcal{C}$ 上纤维化范畴之间的任意函子，
这未必成立。因此，我们如下定义纤维化范畴所成的 $2$-范畴。

\begin{definition}
\label{categories-definition-fibred-categories-over-C}
令 $\mathcal{C}$ 为范畴。所谓{\it $\mathcal{C}$ 上纤维化范畴
所成的 $2$-范畴}，是 $\mathcal{C}$ 上范畴所成的 $2$-范畴
（见定义 \ref{categories-definition-categories-over-C}）的子 $2$-范畴，
定义如下：
\begin{enumerate}
\item 它的对象是纤维化范畴
$p : \mathcal{S} \to \mathcal{C}$。
\item 它的 $1$-态射
$(\mathcal{S}, p) \to (\mathcal{S}', p')$ 是满足
$p' \circ G = p$ 的函子 $G : \mathcal{S} \to \mathcal{S}'$，
并且 $G$ 将强笛卡尔态射映为强笛卡尔态射。
\item 对
$G, H : (\mathcal{S}, p) \to (\mathcal{S}', p')$，它的 $2$-态射
$t : G \to H$ 是满足如下条件的函子态射：对所有
$x \in \Ob(\mathcal{S})$，均有
$p'(t_x) = \text{id}_{p(x)}$。
\end{enumerate}
在这种情形，我们将 $(\mathcal{S}, p)$ 与
$(\mathcal{S}', p')$ 之间 $1$-态射所成的范畴记为
$$
\Mor_{\textit{Fib}/\mathcal{C}}(\mathcal{S}, \mathcal{S}')
$$
\end{definition}

\noindent
注意 $1$-态射所满足的条件。还要注意，这是真正的 $2$-范畴，
而不是 $(2, 1)$-范畴。因此，在取 $2$-纤维积时，我们先转到
与之关联的 $(2, 1)$-范畴。


\begin{lemma}
\label{categories-lemma-2-product-fibred-categories-over-C}
令 $\mathcal{C}$ 为范畴。$\mathcal{C}$ 上纤维化范畴所成的
$(2, 1)$-范畴具有 2-纤维积，它们如引理
\ref{categories-lemma-2-product-categories-over-C} 所述。
\end{lemma}

\begin{proof}
这里基本上需要证明的是：给定 $\mathcal{C}$ 上纤维化范畴的态射
$F : \mathcal{X} \to \mathcal{S}$ 和
$G : \mathcal{Y} \to \mathcal{S}$，引理
\ref{categories-lemma-2-product-categories-over-C} 中描述的范畴
$\mathcal{X} \times_\mathcal{S} \mathcal{Y}$ 是纤维化的。
我们来证明 $\mathcal{X} \times_\mathcal{S} \mathcal{Y}$
有足够多的强笛卡尔态射。设 $(U, x, y, \phi)$ 是
$\mathcal{X} \times_\mathcal{S} \mathcal{Y}$ 的对象，并设
$f : V \to U$ 是 $\mathcal{C}$ 中的态射。选择 $\mathcal{X}$
中位于 $f$ 之上的强笛卡尔态射 $a : f^*x \to x$，以及
$\mathcal{Y}$ 中位于 $f$ 之上的强笛卡尔态射
$b : f^*y \to y$。由假设，$F(a)$ 与 $G(b)$ 是强笛卡尔的。
由于 $\phi : F(x) \to G(y)$ 是同构，强笛卡尔态射的唯一性给出
唯一同构 $f^*\phi : F(f^*x) \to G(f^*y)$，使
$G(b) \circ f^*\phi = \phi \circ F(a)$。换言之，
$(a, b) : (V, f^*x, f^*y, f^*\phi) \to (U, x, y, \phi)$
是 $\mathcal{X} \times_\mathcal{S} \mathcal{Y}$ 中的态射。
略去对它是强笛卡尔态射的验证（也略去对这些事实上是仅有的强
笛卡尔态射的验证）。
\end{proof}

\begin{lemma}
\label{categories-lemma-cute}
令 $\mathcal{C}$ 为范畴，令 $U \in \Ob(\mathcal{C})$。
若 $p : \mathcal{S} \to \mathcal{C}$ 是纤维化范畴，并且
$p$ 经由 $p' : \mathcal{S} \to \mathcal{C}/U$ 分解，那么
$p' : \mathcal{S} \to \mathcal{C}/U$ 是纤维化范畴。
\end{lemma}

\begin{proof}
设 $\varphi : x' \to x$ 关于 $p$ 是强笛卡尔的。我们断言，
$\varphi$ 关于 $p'$ 也是强笛卡尔的。置
$g = p'(\varphi)$，则对某些态射 $f : V \to U$ 和
$f' : V' \to U$，有 $g : V'/U \to V/U$。令
$z \in \Ob(\mathcal{S})$，并置 $p'(z) = (W \to U)$。
为了证明 $\varphi$ 对 $p'$ 是强笛卡尔的，需要证明由
$\psi' \longmapsto (\varphi \circ \psi', p'(\psi'))$ 给出的映射
$$
\Mor_\mathcal{S}(z, x')
\longrightarrow
\Mor_\mathcal{S}(z, x)
\times_{\Mor_{\mathcal{C}/U}(W/U, V/U)}
\Mor_{\mathcal{C}/U}(W/U, V'/U),
$$
是双射。设给定右端的元素 $(\psi, h)$。特别地，
$g \circ h = p(\psi)$；由 $\varphi$ 是强笛卡尔的这一条件，
得到唯一态射 $\psi' : z \to x'$，使
$\psi = \varphi \circ \psi'$ 且 $p(\psi') = h$。好，现在
$p'(\psi') : W/U \to V/U$ 是一个态射，其对应映射
$W \to V$ 是 $h$，因而它作为 $\mathcal{C}/U$ 中的态射等于
$h$。所以 $\psi'$ 是映到给定有序对 $(\psi, h)$ 的唯一态射
$z \to x'$。这证明了断言。

\medskip\noindent
最后，设给定 $g : V'/U \to V/U$ 以及满足
$p'(x) = V/U$ 的 $x$。由于
$p : \mathcal{S} \to \mathcal{C}$ 是纤维化范畴，存在满足
$p(\varphi) = g$ 的强笛卡尔态射 $\varphi : x' \to x$。
同样由上述论证可得
$p'(\varphi) = g : V'/U \to V/U$。而且如上所见，态射
$\varphi$ 是强笛卡尔的。因此定义
\ref{categories-definition-fibred-category} 的条件均得到满足，结论即得。
\end{proof}

\begin{lemma}
\label{categories-lemma-fibred-over-fibred}
令 $\mathcal{A} \to \mathcal{B} \to \mathcal{C}$ 为范畴之间的函子。
若 $\mathcal{A}$ 在 $\mathcal{B}$ 上纤维化，而
$\mathcal{B}$ 在 $\mathcal{C}$ 上纤维化，则
$\mathcal{A}$ 在 $\mathcal{C}$ 上纤维化。
\end{lemma}

\begin{proof}
这由定义和引理 \ref{categories-lemma-cartesian-over-cartesian} 推出。
\end{proof}

\begin{lemma}
\label{categories-lemma-fibred-category-representable-goes-up}
令 $p : \mathcal{S} \to \mathcal{C}$ 为纤维化范畴。
令 $x \to y$ 与 $z \to y$ 为 $\mathcal{S}$ 的态射，其中
$x \to y$ 是强笛卡尔的。若
$p(x) \times_{p(y)} p(z)$ 存在，则 $x \times_y z$ 存在，
$p(x \times_y z) = p(x) \times_{p(y)} p(z)$，并且
$x \times_y z \to z$ 是强笛卡尔的。
\end{lemma}

\begin{proof}
选取位于
$\text{pr}_2 : p(x) \times_{p(y)} p(z) \to p(z)$ 之上的强
笛卡尔态射 $\text{pr}_2^*z \to z$。于是由引理
\ref{categories-lemma-strongly-cartesian-fibre-product}，
$\text{pr}_2^*z = x \times_y z$。
\end{proof}

\begin{lemma}
\label{categories-lemma-ameliorate-morphism-fibred-categories}
令 $\mathcal{C}$ 为范畴。令 $F : \mathcal{X} \to \mathcal{Y}$
为 $\mathcal{C}$ 上纤维化范畴的一个 $1$-态射。存在
$\mathcal{C}$ 上纤维化范畴的 $1$-态射
$$
\xymatrix{
\mathcal{X} \ar@<1ex>[r]^u &
\mathcal{X}' \ar[r]^v \ar@<1ex>[l]^w & \mathcal{Y}
}
$$
使 $F = v \circ u$，并且
\begin{enumerate}
\item $u : \mathcal{X} \to \mathcal{X}'$ 全忠实；
\item $w$ 是 $u$ 的左伴随；并且
\item $v : \mathcal{X}' \to \mathcal{Y}$ 是纤维化范畴。
\end{enumerate}
\end{lemma}

\begin{proof}
将结构函子记为 $p : \mathcal{X} \to \mathcal{C}$ 和
$q : \mathcal{Y} \to \mathcal{C}$。我们如下显式构造
$\mathcal{X}'$。$\mathcal{X}'$ 的一个对象是四元组
$(U, x, y, f)$，其中 $x \in \Ob(\mathcal{X}_U)$、
$y \in \Ob(\mathcal{Y}_U)$，而
$f : y \to F(x)$ 是 $\mathcal{Y}_U$ 中的态射。一个态射
$(a, b) : (U, x, y, f) \to (U', x', y', f')$ 由
$a : x \to x'$ 和 $b : y \to y'$ 给出，其中
$p(a) = q(b) : U \to U'$ 且 $f' \circ b = F(a) \circ f$。

\medskip\noindent
同时为 $p$ 和 $q$ 选定拉回，并使用相同记号表示它们。令
$(U, x, y, f)$ 为对象，令 $h : V \to U$ 为态射。考虑由给定的
强笛卡尔映射 $h^*x \to x$ 与 $h^*y \to y$ 所产生的态射
$c : (V, h^*x, h^*y, h^*f) \to (U, x, y, f)$。
我们断言，$c$ 在 $\mathcal{X}'$ 中关于 $\mathcal{C}$ 是强
笛卡尔的。事实上，设给定 $\mathcal{X}'$ 的一个对象
$(W, x', y', f')$，一个位于 $W \to U$ 之上的态射
$(a, b) : (W, x', y', f') \to (U, x, y, f)$，以及
$W \to U$ 经由 $h$ 的分解 $W \to V \to U$。由于
$h^*x \to x$ 和 $h^*y \to y$ 是强笛卡尔的，得到位于给定态射
$W \to V$ 之上的态射
$a' : x' \to h^*x$ 与 $b' : y' \to h^*y$。考虑图表
$$
\xymatrix{
y' \ar[d]_{f'} \ar[r] & h^*y \ar[r] \ar[d]_{h^*f} & y \ar[d]_f \\
F(x') \ar[r] & F(h^*x) \ar[r] & F(x)
}
$$
外部矩形和右方块交换。由于 $F$ 是纤维化范畴的一个 $1$-态射，
态射 $F(h^*x) \to F(x)$ 是强笛卡尔的。于是由强笛卡尔态射的
泛性质，左方块交换。这证明 $\mathcal{X}'$ 在 $\mathcal{C}$
上纤维化。

\medskip\noindent
函子 $u : \mathcal{X} \to \mathcal{X}'$ 由
$x \mapsto (p(x), x, F(x), \text{id})$ 给出。它是全忠实的。
函子 $\mathcal{X}' \to \mathcal{Y}$ 由
$(U, x, y, f) \mapsto y$ 给出。函子
$w : \mathcal{X}' \to \mathcal{X}$ 由
$(U, x, y, f) \mapsto x$ 给出。根据我们对
$\mathcal{X}'$ 中关于 $\mathcal{C}$ 的强笛卡尔态射的描述，
这些函子都是 $\mathcal{C}$ 上纤维化范畴的 $1$-态射。
$w$ 与 $u$ 的伴随性意味着
$$
\Mor_\mathcal{X}(x, x') =
\Mor_{\mathcal{X}'}((U, x, y, f), (p(x'), x', F(x'), \text{id})),
$$
这直接由定义推出。

\medskip\noindent
最后，需要证明 $\mathcal{X}' \to \mathcal{Y}$ 是纤维化范畴。
令 $c : y' \to y$ 为 $\mathcal{Y}$ 中的态射，并令
$(U, x, y, f)$ 为 $\mathcal{X}'$ 中位于 $y$ 之上的对象。
置 $V = q(y')$，并令 $h = q(c) : V \to U$。令
$a : h^*x \to x$ 与 $b : h^*y \to y$ 为覆盖 $h$ 的强笛卡尔态射。
由于 $F$ 是纤维化范畴的一个 $1$-态射，我们可以将
$h^*F(x) = F(h^*x)$ 等同，并取强笛卡尔态射
$F(a) : F(h^*x) \to F(x)$。由 $b : h^*y \to y$ 的泛性质，
存在 $\mathcal{Y}_V$ 中的态射 $c' : y' \to h^*y$，使
$c = b \circ c'$。我们断言
$$
(a, c) : (V, h^*x, y', h^*f \circ c') \longrightarrow (U, x, y, f)
$$
在 $\mathcal{X}'$ 中关于 $\mathcal{Y}$ 是强笛卡尔的。为此，
令 $(W, x_1, y_1, f_1)$ 为 $\mathcal{X}'$ 的对象，令
$(a_1, b_1) : (W, x_1, y_1, f_1) \to (U, x, y, f)$ 为态射，
并设对某个态射 $b_1' : y_1 \to y'$ 有
$b_1 = c \circ b_1'$。那么
$$
(a_1', b_1') : (W, x_1, y_1, f_1) \longrightarrow (V, h^*x, y', h^*f \circ c')
$$
（其中 $a_1' : x_1 \to h^*x$ 是位于给定态射
$q(b_1') : W \to V$ 之上的唯一态射，使
$a_1 = a \circ a_1'$）就是所需态射。
\end{proof}


\section{惯性}
\label{categories-section-inertia}

\noindent
给定范畴 $\mathcal{C}$ 上的纤维化范畴
$p : \mathcal{S} \to \mathcal{C}$ 与
$p' : \mathcal{S}' \to \mathcal{C}$，以及一个 $1$-态射
$F : \mathcal{S} \to \mathcal{S}'$，在 $\mathcal{C}$ 上纤维化范畴
所成的 $(2, 1)$-范畴中有对角态射
$$
\Delta = \Delta_{\mathcal{S}/\mathcal{S}'} :
\mathcal{S} \longrightarrow \mathcal{S} \times_{\mathcal{S}'} \mathcal{S}
$$

\begin{lemma}
\label{categories-lemma-inertia-fibred-category}
令 $\mathcal{C}$ 为范畴。令
$p : \mathcal{S} \to \mathcal{C}$ 与
$p' : \mathcal{S}' \to \mathcal{C}$ 为纤维化范畴。令
$F : \mathcal{S} \to \mathcal{S}'$ 为 $\mathcal{C}$ 上纤维化范畴的
一个 $1$-态射。考虑 $\mathcal{C}$ 上的范畴
$\mathcal{I}_{\mathcal{S}/\mathcal{S}'}$，其
\begin{enumerate}
\item 对象是有序对 $(x, \alpha)$，其中
$x \in \Ob(\mathcal{S})$，且 $\alpha : x \to x$ 是满足
$F(\alpha) = \text{id}$ 的自同构；
\item 态射 $(x, \alpha) \to (y, \beta)$ 由使图表
$$
\xymatrix{
x\ar[r]_\phi\ar[d]_\alpha &
y\ar[d]^{\beta} \\
x\ar[r]^\phi &
y \\
}
$$
交换的态射 $\phi : x \to y$ 给出；并且
\item 函子 $\mathcal{I}_{\mathcal{S}/\mathcal{S}'} \to \mathcal{C}$
由 $(x, \alpha) \mapsto p(x)$ 给出。
\end{enumerate}
那么
\begin{enumerate}
\item 在 $\mathcal{C}$ 上范畴所成的 $(2, 1)$-范畴中存在等价
$$
\mathcal{I}_{\mathcal{S}/\mathcal{S}'} \longrightarrow
\mathcal{S}
\times_{\Delta, (\mathcal{S} \times_{\mathcal{S}'} \mathcal{S}), \Delta}
\mathcal{S}
$$
并且
\item $\mathcal{I}_{\mathcal{S}/\mathcal{S}'}$ 是
$\mathcal{C}$ 上的纤维化范畴。
\end{enumerate}
\end{lemma}

\begin{proof}
由引理 \ref{categories-lemma-2-product-fibred-categories-over-C} 和
\ref{categories-lemma-fibred-equivalent}，(2) 从 (1) 推出。因此只需证明 (1)。
以下不再特别说明地使用引理
\ref{categories-lemma-2-product-fibred-categories-over-C} 中的
$2$-纤维积构造。特别地，
$\mathcal{S}
\times_{\Delta, (\mathcal{S} \times_{\mathcal{S}'} \mathcal{S}), \Delta}
\mathcal{S}$
的一个对象是三元组 $(x, y, (\iota, \kappa))$，其中 $x$ 和 $y$
是 $\mathcal{S}$ 的对象，而
$(\iota, \kappa) : (x, x, \text{id}_{F(x)}) \to
(y, y, \text{id}_{F(y)})$
是 $\mathcal{S} \times_{\mathcal{S}'} \mathcal{S}$ 中的同构。
这恰好意味着 $\iota, \kappa : x \to y$ 是同构，且
$F(\iota) = F(\kappa)$。考虑函子
$$
I_{\mathcal{S}/\mathcal{S}'}
\longrightarrow
\mathcal{S}
\times_{\Delta, (\mathcal{S} \times_{\mathcal{S}'} \mathcal{S}), \Delta}
\mathcal{S}
$$
它将左端的对象 $(x, \alpha)$ 送到右端的对象
$(x, x, (\alpha, \text{id}_x))$，并将左端的态射 $\phi$
送到右端的态射 $(\phi, \phi)$。我们断言，该态射的一个拟逆由函子
$$
\mathcal{S}
\times_{\Delta, (\mathcal{S} \times_{\mathcal{S}'} \mathcal{S}), \Delta}
\mathcal{S}
\longrightarrow
I_{\mathcal{S}/\mathcal{S}'}
$$
给出；它将左端的对象 $(x, y, (\iota, \kappa))$ 送到右端的对象
$(x, \kappa^{-1} \circ \iota)$，并将左端的态射
$(\phi, \phi') : (x, y, (\iota, \kappa)) \to
(z, w, (\lambda, \mu))$
送到态射 $\phi$。事实上，复合上述两个函子所诱导的
$I_{\mathcal{S}/\mathcal{S}'}$ 的自函子严格等于恒等函子；而在
$\mathcal{S}
\times_{\Delta, (\mathcal{S} \times_{\mathcal{S}'} \mathcal{S}), \Delta}
\mathcal{S}$
上诱导的自函子，通过自然同构
$$
(\text{id}_x, \kappa) :
(x, x, (\kappa^{-1} \circ \iota, \text{id}_x))
\longrightarrow
(x, y, (\iota, \kappa)).
$$
同构于恒等函子。略去若干细节。
\end{proof}

\begin{definition}
\label{categories-definition-inertia-fibred-category}
令 $\mathcal{C}$ 为范畴。
\begin{enumerate}
\item 令 $F : \mathcal{S} \to \mathcal{S}'$ 为 $\mathcal{C}$
上纤维化范畴的一个 $1$-态射。$\mathcal{S}$ 在
$\mathcal{S}'$ 上的{\it 相对惯性}，是引理
\ref{categories-lemma-inertia-fibred-category} 的纤维化范畴
$\mathcal{I}_{\mathcal{S}/\mathcal{S}'} \to \mathcal{C}$。
\item 所谓 $\mathcal{S}$ 的{\it 惯性纤维化范畴
$\mathcal{I}_\mathcal{S}$}，是指
$\mathcal{I}_\mathcal{S} = \mathcal{I}_{\mathcal{S}/\mathcal{C}}$。
\end{enumerate}
\end{definition}

\noindent
注意，$\mathcal{C}$ 上的纤维化范畴之间有典范 $1$-态射
\begin{equation}
\label{categories-equation-inertia-structure-map}
\mathcal{I}_{\mathcal{S}/\mathcal{S}'} \longrightarrow \mathcal{S}
\quad\text{且}\quad
\mathcal{I}_\mathcal{S} \longrightarrow \mathcal{S}
\end{equation}
按照引理 \ref{categories-lemma-inertia-fibred-category} 的描述，它们只是将对象
$(x, \alpha)$ 映到对象 $x$，并将态射
$\phi : (x, \alpha) \to (y, \beta)$ 映到态射
$\phi : x \to y$。还有一个{\it 中性截面}
\begin{equation}
\label{categories-equation-neutral-section}
e : \mathcal{S} \to \mathcal{I}_{\mathcal{S}/\mathcal{S}'}
\quad\text{且}\quad
e : \mathcal{S} \to \mathcal{I}_\mathcal{S}
\end{equation}
由规则 $x \mapsto (x, \text{id}_x)$ 和
$(\phi : x \to y) \mapsto \phi$ 定义。它是
(\ref{categories-equation-inertia-structure-map}) 的右逆。给定一个
$2$-交换方块
$$
\xymatrix{
\mathcal{S}_1 \ar[d]_{F_1} \ar[r]_G & \mathcal{S}_2 \ar[d]^{F_2} \\
\mathcal{S}'_1 \ar[r]^{G'} & \mathcal{S}'_2
}
$$
有{\it 函子性映射}
\begin{equation}
\label{categories-equation-functorial}
\mathcal{I}_{\mathcal{S}_1/\mathcal{S}'_1}
\longrightarrow
\mathcal{I}_{\mathcal{S}_2/\mathcal{S}'_2}
\quad\text{且}\quad
\mathcal{I}_{\mathcal{S}_1}
\longrightarrow
\mathcal{I}_{\mathcal{S}_2}
\end{equation}
由规则 $(x, \alpha) \mapsto (G(x), G(\alpha))$ 和
$\phi \mapsto G(\phi)$ 定义。特别地，总有比较映射
\begin{equation}
\label{categories-equation-comparison}
\mathcal{I}_{\mathcal{S}/\mathcal{S}'}
\longrightarrow
\mathcal{I}_\mathcal{S}
\end{equation}
而且上述所有映射均与之相容。

\begin{lemma}
\label{categories-lemma-relative-inertia-as-fibre-product}
令 $F : \mathcal{S} \to \mathcal{S}'$ 为在范畴
$\mathcal{C}$ 上纤维化的范畴之间的一个 $1$-态射。那么图表
$$
\xymatrix{
\mathcal{I}_{\mathcal{S}/\mathcal{S}'}
\ar[d]_{F \circ (\ref{categories-equation-inertia-structure-map})}
\ar[rr]_{(\ref{categories-equation-comparison})} & &
\mathcal{I}_\mathcal{S} \ar[d]^{(\ref{categories-equation-functorial})} \\
\mathcal{S}' \ar[rr]^e & &
\mathcal{I}_{\mathcal{S}'}
}
$$
是一个 $2$-纤维积。
\end{lemma}

\begin{proof}
略。
\end{proof}


\section{群胚纤维化范畴}
\label{categories-section-fibred-groupoids}

\noindent
本节说明应如何理解群胚纤维化范畴，并看到它们基本上等同于取值于
群胚的 $(2, 1)$-范畴的函子。

\begin{definition}
\label{categories-definition-fibred-groupoids}
令 $p : \mathcal{S} \to \mathcal{C}$ 为函子。如果下列两个条件成立，
则称 $\mathcal{S}$ 在 $\mathcal{C}$ 上{\it 以群胚为纤维}：
\begin{enumerate}
\item 对 $\mathcal{C}$ 中的每个态射 $f : V \to U$ 以及
$U$ 的每个提升 $x$，都存在靶为 $x$ 的 $f$ 的提升
$\phi : y \to x$。
\item 对每一对态射 $\phi : y \to x$ 与 $ \psi : z \to x$，
以及任意满足 $p(\phi) \circ f = p(\psi)$ 的态射
$f : p(z) \to p(y)$，都存在唯一的 $f$ 的提升
$\chi : z \to y$，使 $\phi \circ \chi = \psi$。
\end{enumerate}
\end{definition}

\noindent
条件 (2) 换一种说法是：应用函子 $p$，会在下列交换图表中虚线箭头
所成的集合之间给出双射：
$$
\xymatrix{
y \ar[r] & x & p(y) \ar[r] & p(x) \\
z \ar@{-->}[u] \ar[ru] & & p(z) \ar@{-->}[u]\ar[ru] & \\
}
$$
还可以用下列方式理解第二个条件。设 $g : W \to V$ 和
$f : V \to U$ 是 $\mathcal{C}$ 中的态射。令
$x \in \Ob(\mathcal{S}_U)$。由第一个条件，可以将 $f$ 提升为
$ \phi : y \to x$，然后将 $g$ 提升为 $\psi : z \to y$。
也可以不作这个两步过程，而将 $g \circ f$ 直接提升为
$\gamma : z' \to x$。这给出图表
\begin{equation}
\label{categories-equation-fibred-groupoids}
\vcenter{
\xymatrix{
z' \ar@{-->}[d]\ar[rrd]^\gamma & & \\
z \ar@{-->}[u] \ar[r]^\psi \ar@{~>}[d]^p &
y \ar[r]^\phi \ar@{~>}[d]^p &
x \ar@{~>}[d]^p
\\
W \ar[r]^g & V \ar[r]^f & U \\
}
}
\end{equation}
中的实线箭头，其中波浪箭头表示的不是态射，而是函子 $p$。
把第二个条件应用于箭头 $\phi \circ \psi$、$\gamma$ 和
$\text{id}_W$，得到 $\mathcal{S}_W$ 中唯一的态射
$\chi : z \to z'$，使
$\gamma \circ \chi = \phi \circ \psi$。类似地，存在唯一态射
$z' \to z$。唯一性蕴含态射 $z' \to z$ 与 $z\to z'$ 互为逆，
换言之，它们是同构。

\medskip\noindent
从上述讨论应当已经清楚，群胚纤维化范畴与纤维化范畴关系非常
密切。下面给出相应结论。

\begin{lemma}
\label{categories-lemma-fibred-groupoids}
令 $p : \mathcal{S} \to \mathcal{C}$ 为函子。下列条件等价：
\begin{enumerate}
\item $p : \mathcal{S} \to \mathcal{C}$ 是群胚纤维化范畴；并且
\item 所有纤维范畴都是群胚，而且 $\mathcal{S}$ 是
$\mathcal{C}$ 上的纤维化范畴。
\end{enumerate}
进一步，在这种情形，$\mathcal{S}$ 的每个态射都是强笛卡尔的。
此外，如果对所有 $f: V \to U = p(x)$ 都给定一个位于 $f$
之上的 $f^\ast x \to x$，那么引理 \ref{categories-lemma-fibred} 中构造的数据
$(U \mapsto \mathcal{S}_U, f \mapsto f^*, \alpha_{f, g}, \alpha_U)$
定义了从 $\mathcal{C}^{opp}$ 到群胚的 $(2, 1)$-范畴的一个伪函子。
\end{lemma}

\begin{proof}
假设 $p : \mathcal{S} \to \mathcal{C}$ 是群胚纤维化范畴。
为了证明对 $U \in \Ob(\mathcal{C})$ 的所有纤维范畴
$\mathcal{S}_U$ 都是群胚，必须为 $\mathcal{S}_U$ 中每个
$f : y \to x$ 展示一个逆态射。左侧图表（在 $\mathcal{S}_U$ 中）
经 $p$ 映到右侧图表：
$$
\xymatrix{
y \ar[r]^f & x & U \ar[r]^{\text{id}_U} & U \\
x \ar@{-->}[u] \ar[ru]_{\text{id}_x} & &
U \ar@{-->}[u]\ar[ru]_{\text{id}_U} & \\
}
$$
因为只有 $\text{id}_U$ 能使右侧图表交换，所以存在唯一的
$g : x \to y$ 使左侧图表交换，故
$fg = \text{id}_x$。类似地，存在唯一的 $h : y \to x$，使
$gh = \text{id}_y$。于是 $fgh = f : y \to x$。由于
$fg = \text{id}_x$，得到 $h = f$。定义
\ref{categories-definition-fibred-groupoids} 的条件 (2) 恰好说明
$\mathcal{S}$ 的每个态射都是强笛卡尔的。因此，定义
\ref{categories-definition-fibred-groupoids} 的条件 (1) 蕴含
$\mathcal{S}$ 是 $\mathcal{C}$ 上的纤维化范畴。

\medskip\noindent
反过来，假设所有纤维范畴都是群胚，并且 $\mathcal{S}$ 是
$\mathcal{C}$ 上的纤维化范畴。需要验证定义
\ref{categories-definition-fibred-groupoids} 的条件 (1) 与 (2)。
第一个条件显然成立。令 $\phi : y \to x$、
$\psi : z \to x$ 和 $f : p(z) \to p(y)$ 满足
$p(\phi) \circ f = p(\psi)$，如定义
\ref{categories-definition-fibred-groupoids} 的条件 (2) 所述。记
$U = p(x)$、$V = p(y)$、$W = p(z)$、
$p(\phi) = g : V \to U$、$p(\psi) = h : W \to U$。
选择一个位于 $g$ 之上的强笛卡尔态射 $g^*x \to x$。
于是得到 $\mathcal{S}_V$ 中的态射 $i : y \to g^*x$，
因而它是同构。由于 $g^*x \to x$ 是强笛卡尔的，还得到对应于
有序对 $(\psi, f)$ 的态射 $j : z \to g^*x$。
容易验证 $\chi = i^{-1} \circ j$ 是一个解。

\medskip\noindent
我们已经在 (1) $\Rightarrow$ (2) 的证明中看到，
$\mathcal{S}$ 的每个态射都是强笛卡尔的。最后一个陈述直接由
引理 \ref{categories-lemma-fibred} 推出。
\end{proof}


\begin{lemma}
\label{categories-lemma-fibred-gives-fibred-groupoids}
令 $\mathcal{C}$ 为范畴。令
$p : \mathcal{S} \to \mathcal{C}$ 为纤维化范畴。如下定义
$\mathcal{S}$ 的子范畴 $\mathcal{S}'$：
\begin{enumerate}
\item $\Ob(\mathcal{S}') = \Ob(\mathcal{S})$；并且
\item 对 $x, y \in \Ob(\mathcal{S}')$，$x$ 与 $y$ 之间
$\mathcal{S}'$ 中的态射集合，是 $x$ 与 $y$ 之间
$\mathcal{S}$ 中的强笛卡尔态射集合。
\end{enumerate}
令 $p' : \mathcal{S}' \to \mathcal{C}$ 为 $p$ 在
$\mathcal{S}'$ 上的限制。那么
$p' : \mathcal{S}' \to \mathcal{C}$ 是群胚纤维化范畴。
\end{lemma}

\begin{proof}
注意，该构造确实有意义，因为由引理
\ref{categories-lemma-composition-cartesian}，$\mathcal{S}$ 任意对象的恒等
态射是强笛卡尔的，而且强笛卡尔态射的复合仍是强笛卡尔的。
定义 \ref{categories-definition-fibred-groupoids} 的第一个提升性质，来自
纤维化范畴的下述条件：给定任意态射 $f : V \to U$ 以及位于
$U$ 之上的 $x$，都存在位于 $f$ 之上的强笛卡尔态射
$\varphi : y \to x$。我们来验证 $\mathcal{C}$ 上的范畴
$p' : \mathcal{S}' \to \mathcal{C}$ 满足定义
\ref{categories-definition-fibred-groupoids} 的第二个提升性质。为此，仿照
定义 \ref{categories-definition-fibred-groupoids} 后的讨论论证。因此，在图表
\ref{categories-equation-fibred-groupoids} 中，态射 $\phi$、$\psi$ 和
$\gamma$ 是 $\mathcal{S}$ 的强笛卡尔态射。于是
$\gamma$ 与 $\phi \circ \psi$ 是 $\mathcal{S}$ 中位于
$\mathcal{C}$ 同一箭头之上的强笛卡尔态射，并且在
$\mathcal{S}$ 中有相同的靶。由定义
\ref{categories-definition-cartesian-over-C} 后的讨论，这意味着这两个箭头
按所需方式同构（这里还再次使用引理
\ref{categories-lemma-composition-cartesian}：$\mathcal{S}$ 中任意同构都是
强笛卡尔的）。
\end{proof}

\begin{example}
\label{categories-example-group-homomorphism-fibreedingroupoids}
群同态 $p : G \to H$ 产生例
\ref{categories-example-group-homomorphism-functor} 中那样的函子
$p : \mathcal{S}\to\mathcal{C}$。这个函子
$p : \mathcal{S}\to\mathcal{C}$ 是群胚纤维化范畴，当且仅当
$p$ 是满射。位于（唯一的）对象
$U\in \Ob(\mathcal{C})$ 之上的纤维范畴 $\mathcal{S}_U$，
是例 \ref{categories-example-group-groupoid} 中与 $p$ 的核关联的范畴。
\end{example}

\noindent
给定 $p : \mathcal{S} \to \mathcal{C}$，可以问：如果对所有
$U \in \Ob(\mathcal{C})$，纤维范畴 $\mathcal{S}_U$ 都是群胚，
那么 $\mathcal{S}$ 必然是 $\mathcal{C}$ 上的群胚纤维化范畴吗？
如下可见，答案是否定的。从一个群胚纤维化范畴
$p : \mathcal{S} \to \mathcal{C}$ 开始。改变 $\mathcal{S}$
中不映到某个对象的恒等态射的那些态射，并不会改变范畴
$\mathcal{S}_U$。所以可以破坏提升的存在性和唯一性条件。
一个例子是例 \ref{categories-example-group-homomorphism-fibreedingroupoids}
中的函子，其中 $G \to H$ 不是满射。下面是另一个例子。

\begin{example}
\label{categories-example-not-fibred-in-groupoids-but-fibre-cats-are}
令 $\Ob(\mathcal{C}) = \{A, B, T\}$，
$\Mor_\mathcal{C}(A, B) = \{f\}$、
$\Mor_\mathcal{C}(B, T) = \{g\}$、
$\Mor_\mathcal{C}(A, T) = \{h\} = \{gf\}, $，再加上每个对象的
恒等态射。下图画出了这个范畴。现在令
$\Ob(\mathcal{S}) = \{A', B', T'\}$，
$\Mor_\mathcal{S}(A', B') = \emptyset$、
$\Mor_\mathcal{S}(B', T') = \{g'\}$、
$\Mor_\mathcal{S}(A', T') = \{h'\}, $，再加上恒等态射。函子
$p : \mathcal{S} \to \mathcal{C}$ 是显然的。那么对每个
$U \in \Ob(\mathcal{C})$，$\mathcal{S}_U$ 都是只含一个对象及其
恒等态射的范畴，因而是群胚；但态射 $f: A \to B$ 无法提升。
类似地，若令
$\Mor_\mathcal{S}(A', B') = \{f'_1, f'_2\}$ 和
$ \Mor_\mathcal{S}(A', T') = \{h'\} = \{g'f'_1 \} = \{g'f'_2\}$，
那么纤维范畴不变，而下图中的 $f: A \to B$ 有两个提升。
$$
\xymatrix{
B' \ar[r]^{g'} & T' & & B \ar[r]^g & T & \\
A' \ar@{-->}[u]^{f'_1,\,f'_2} \ar[ru]_{h'} & & \ar@{}[u]^{\text{上图}} &
A \ar[u]^f \ar[ru]_{gf = h} & \\
}
$$
\end{example}

\noindent
稍后，我们想作出诸如“$\mathcal{C}$ 上的任意群胚纤维化范畴都等价
于一个分裂的范畴”，或“纤维范畴均为集合状的任意群胚纤维化范畴
都等价于一个以集合为纤维的范畴”这样的断言。等价的概念取决于
我们所使用的 $2$-范畴。

\begin{definition}
\label{categories-definition-categories-fibred-in-groupoids-over-C}
令 $\mathcal{C}$ 为范畴。所谓{\it $\mathcal{C}$ 上群胚纤维化范畴
所成的 $2$-范畴}，是 $\mathcal{C}$ 上纤维化范畴所成的
$2$-范畴（见定义 \ref{categories-definition-fibred-categories-over-C}）的
子 $2$-范畴，定义如下：
\begin{enumerate}
\item 它的对象是群胚纤维化范畴
$p : \mathcal{S} \to \mathcal{C}$。
\item 它的 $1$-态射
$(\mathcal{S}, p) \to (\mathcal{S}', p')$ 是满足
$p' \circ G = p$ 的函子 $G : \mathcal{S} \to \mathcal{S}'$
（由于每个态射都是强笛卡尔的，$G$ 自动保持它们）。
\item 对
$G, H : (\mathcal{S}, p) \to (\mathcal{S}', p')$，它的 $2$-态射
$t : G \to H$ 是满足下述条件的函子态射：对所有
$x \in \Ob(\mathcal{S})$，均有
$p'(t_x) = \text{id}_{p(x)}$。
\end{enumerate}
\end{definition}

\noindent
注意，每个 $2$-态射自动是同构！所以这实际上是
$(2, 1)$-范畴，而不仅是 $2$-范畴。下面给出关于 $2$-纤维积的
必备引理。

\begin{lemma}
\label{categories-lemma-2-product-fibred-categories}
令 $\mathcal{C}$ 为范畴。$\mathcal{C}$ 上群胚纤维化范畴所成的
$2$-范畴具有 2-纤维积，它们如引理
\ref{categories-lemma-2-product-categories-over-C} 所述。
\end{lemma}

\begin{proof}
由引理 \ref{categories-lemma-2-product-fibred-categories-over-C}，引理
\ref{categories-lemma-2-product-categories-over-C} 所描述的纤维积是
纤维化范畴。因此，只需证明纤维范畴都是群胚；见引理
\ref{categories-lemma-fibred-groupoids}。由引理
\ref{categories-lemma-fibre-2-fibre-product-categories-over-C}，只需证明
群胚的 $2$-纤维积是群胚；这显然成立（例如由引理
\ref{categories-lemma-2-fibre-product-categories} 中的构造可见）。
\end{proof}


\begin{remark}
\label{categories-remark-2-product-fibred-categories-same}
令 $\mathcal{C}$ 为范畴。令 $f : \mathcal{X} \to \mathcal{S}$
与 $g : \mathcal{Y} \to \mathcal{S}$ 为 $\mathcal{C}$ 上群胚纤维化
范畴的 $1$-态射。令 $p : \mathcal{S} \to \mathcal{C}$ 为给定函子。
我们断言，引理 \ref{categories-lemma-2-product-fibred-categories} 的
$2$-纤维积作为范畴典范等价于例
\ref{categories-example-2-fibre-product-categories} 中的纤维积。前者的对象是
满足 $p(\alpha) = \text{id}_U$ 的四元组
$(U, x, y, \alpha)$（见引理
\ref{categories-lemma-2-product-categories-over-C}），后者的对象是三元组
$(x, y, \alpha)$（见例
\ref{categories-example-2-fibre-product-categories}）。两个范畴之间的等价
由规则 $(U, x, y, \alpha) \mapsto (x, y, \alpha)$ 和
$(x, y, \alpha) \mapsto (p(f(x)), x, y', \alpha')$ 给出，其中
$\alpha' = g(\gamma)^{-1} \circ \alpha$，而
$\gamma : y' \to y$ 是箭头
$p(\alpha) : p(f(x)) \to p(g(y))$ 的一个提升。略去细节。
\end{remark}

\begin{lemma}
\label{categories-lemma-equivalence-fibred-categories}
令 $p : \mathcal{S}\to \mathcal{C}$ 与
$p' : \mathcal{S'}\to \mathcal{C}$ 为群胚纤维化范畴，并设
$G : \mathcal{S}\to \mathcal {S}'$ 是 $\mathcal{C}$ 上的函子。
\begin{enumerate}
\item $G$ 忠实（相应地，全忠实；相应地，是等价），当且仅当对每个
$U\in\Ob(\mathcal{C})$，诱导函子
$G_U : \mathcal{S}_U\to \mathcal{S}'_U$ 忠实
（相应地，全忠实；相应地，是等价）。
\item 若 $G$ 是等价，则 $G$ 是 $\mathcal{C}$ 上群胚纤维化范畴所成
$2$-范畴中的等价。
\end{enumerate}
\end{lemma}

\begin{proof}
令 $x, y$ 为 $\mathcal{S}$ 中位于同一对象 $U$ 之上的对象。
考虑交换图表
$$
\xymatrix{
\Mor_\mathcal{S}(x, y) \ar[rd]_p \ar[rr]_G & &
\Mor_{\mathcal{S}'}(G(x), G(y)) \ar[ld]^{p'} \\
& \Mor_\mathcal{C}(U, U) &
}
$$
从该图表可清楚看出，如果 $G$ 忠实（相应地，全忠实），那么每个
$G_U$ 也如此。

\medskip\noindent
假设 $G$ 是等价。对 $\mathcal{S}'$ 的每个对象 $x'$，都存在
$\mathcal{S}$ 的对象 $x$，使 $G(x)$ 同构于 $x'$。设 $x'$
位于 $U'$ 之上，而 $x$ 位于 $U$ 之上。那么 $\mathcal{C}$
中存在同构 $f : U' \to U$，即把 $p'$ 应用于同构
$x' \to G(x)$ 所得的同构。由群胚纤维化范畴的公理，
$\mathcal{S}$ 中存在位于 $f$ 之上的箭头 $f^*x \to x$。
因此，存在满足 $p'(\alpha) = \text{id}_{U'}$ 的同构
$\alpha : x' \to G(f^*x)$（这次使用 $\mathcal{S}'$ 的公理）。
综上，对 $\mathcal{S}'$ 的每个对象 $x'$，可以选择一个有序对
$(o_{x'}, \alpha_{x'})$，由 $\mathcal{S}$ 的对象 $o_{x'}$
以及同构 $\alpha_{x'} : x' \to G(o_{x'})$ 构成，并满足
$p'(\alpha_{x'}) = \text{id}_{p'(x')}$。从这里开始，按通常方式
进行（见引理 \ref{categories-lemma-equivalence-categories} 的证明），构造
逆函子 $F : \mathcal{S}' \to \mathcal{S}$：按
$x' \mapsto o_{x'}$，并将 $\varphi' : x' \to y'$ 映到唯一箭头
$\varphi_{\varphi'} : o_{x'} \to o_{y'}$，它满足
$\alpha_{y'}^{-1} \circ G(\varphi_{\varphi'}) \circ \alpha_{x'} = \varphi'$。
在这些选择下，$F$ 是 $\mathcal{C}$ 上的函子。略去对
$G \circ F$ 与 $F \circ G$ 分别和相应恒等函子 $2$-同构的验证
（在 $\mathcal{C}$ 上群胚纤维化范畴所成的 $2$-范畴中）。

\medskip\noindent
假设对所有 $U\in\Ob(\mathcal C)$，$G_U$ 都忠实
（相应地，全忠实）。为了证明 $G$ 忠实（相应地，全忠实），
需要证明，对任意对象 $x, y\in\Ob(\mathcal{S})$，$G$ 在
$\Mor_\mathcal{S}(x, y)$ 与
$\Mor_{\mathcal{S}'}(G(x), G(y))$ 之间诱导单射（相应地，双射）。
置 $U = p(x)$、$V = p(y)$。只需证明，对任意态射
$f : U \to V$，$G$ 在位于 $f$ 之上的态射 $x \to y$ 与位于
$f$ 之上的态射 $G(x) \to G(y)$ 之间诱导单射（相应地，双射）。
现在固定 $f : U \to V$，将一个拉回记为 $f^*y \to y$。
那么 $G(f^*y) \to G(y)$ 也是拉回。位于 $f$ 之上的从 $x$ 到
$y$ 的态射集合，与位于 $\text{id}_U$ 之上的从 $x$ 到
$f^*y$ 的态射集合双射。（由群胚纤维化范畴的第二个公理。）
类似地，位于 $f$ 之上的从 $G(x)$ 到 $G(y)$ 的态射集合，
与位于 $\text{id}_U$ 之上的从 $G(x)$ 到 $G(f^*y)$ 的态射集合
双射。因此，$G_U$ 忠实（相应地，全忠实）这一事实给出所需结论。

\medskip\noindent
最后，假设对所有 $U$，$G_U$ 都是等价，因而全忠实且本质满。
我们已经看到这蕴含 $G$ 全忠实，所以要证明它是等价，只需证明它
本质满。此事显然：若 $z'\in \Ob(\mathcal{S}')$，则
$z'\in \Ob(\mathcal{S}'_U)$，其中 $U = p'(z')$。由于 $G_U$
本质满，$z'$ 在 $\mathcal{S}'_U$ 中同构于形如 $G_U(z)$ 的对象，
其中 $z\in \Ob(\mathcal{S}_U)$。但 $\mathcal{S}'_U$ 中的态射
也是 $\mathcal{S}'$ 中的态射，所以 $z'$ 在 $\mathcal{S}'$ 中
同构于 $G(z)$。
\end{proof}


\begin{lemma}
\label{categories-lemma-fully-faithful-diagonal-equivalence}
令 $\mathcal{C}$ 为范畴。令 $p : \mathcal{S}\to \mathcal{C}$ 与
$p' : \mathcal{S'}\to \mathcal{C}$ 为群胚纤维化范畴。令
$G : \mathcal{S}\to \mathcal {S}'$ 为 $\mathcal{C}$ 上的函子。
那么 $G$ 全忠实，当且仅当对角函子
$$
\Delta_G :
\mathcal{S}
\longrightarrow
\mathcal{S} \times_{G, \mathcal{S}', G} \mathcal{S}
$$
是等价。
\end{lemma}

\begin{proof}
由引理 \ref{categories-lemma-equivalence-fibred-categories}，只需考察位于
$\mathcal{C}$ 的一个对象 $U$ 之上的纤维范畴。右端的一个对象
是三元组 $(x, x', \alpha)$，其中
$\alpha : G(x) \to G(x')$ 是 $\mathcal{S}'_U$ 中的态射。
函子 $\Delta_G$ 将 $\mathcal{S}_U$ 的对象 $x$ 映到三元组
$(x, x, \text{id}_{G(x)})$。注意，$(x, x', \alpha)$ 在
$\Delta_G$ 的本质像中，当且仅当对 $\mathcal{S}_U$ 中某个态射
$\beta : x \to x'$ 有 $\alpha = G(\beta)$（略去细节）。
因此，要使 $\Delta_G$ 成为等价，每个 $\alpha$ 都必须是某个态射
$\beta : x \to x'$ 的像，而且任意两个不同的态射
$\beta, \beta' : x \to x'$ 必须给出不同的态射
$G(\beta), G(\beta')$。这证明了引理。
\end{proof}

\begin{lemma}
\label{categories-lemma-morphisms-equivalent-fibred-groupoids}
令 $\mathcal{C}$ 为范畴。令 $\mathcal{S}_i$（$i = 1, 2, 3, 4$）
为 $\mathcal{C}$ 上的群胚纤维化范畴。设
$\varphi : \mathcal{S}_1 \to \mathcal{S}_2$ 与
$\psi : \mathcal{S}_3 \to \mathcal{S}_4$ 是
$\mathcal{C}$ 上的等价。那么
$$
\Mor_{\textit{Cat}/\mathcal{C}}(\mathcal{S}_2, \mathcal{S}_3)
\longrightarrow
\Mor_{\textit{Cat}/\mathcal{C}}(\mathcal{S}_1, \mathcal{S}_4),
\quad \alpha \longmapsto \psi \circ \alpha \circ \varphi
$$
是范畴等价。
\end{lemma}

\begin{proof}
这是一个一般事实，在任意 $2$-范畴中都成立。
\end{proof}

\begin{lemma}
\label{categories-lemma-inertia-fibred-groupoids}
令 $\mathcal{C}$ 为范畴。若
$p : \mathcal{S} \to \mathcal{C}$ 是群胚纤维化范畴，则惯性
纤维化范畴 $\mathcal{I}_\mathcal{S} \to \mathcal{C}$ 也是如此。
\end{lemma}

\begin{proof}
这从引理 \ref{categories-lemma-inertia-fibred-category} 中的构造即显然；
也可使用同一引理给出的等价
$I_\mathcal{S} \to \mathcal{S}
\times_{\Delta, \mathcal{S} \times_\mathcal{C} \mathcal{S}, \Delta}\mathcal{S}$，
再应用引理 \ref{categories-lemma-2-product-fibred-categories}。
\end{proof}

\begin{lemma}
\label{categories-lemma-cute-groupoids}
令 $\mathcal{C}$ 为范畴，令 $U \in \Ob(\mathcal{C})$。
若 $p : \mathcal{S} \to \mathcal{C}$ 是群胚纤维化范畴，并且
$p$ 经由 $p' : \mathcal{S} \to \mathcal{C}/U$ 分解，那么
$p' : \mathcal{S} \to \mathcal{C}/U$ 是群胚纤维化范畴。
\end{lemma}

\begin{proof}
我们已经在引理 \ref{categories-lemma-cute} 中看到，$p'$ 是纤维化范畴。
因此，只需证明纤维范畴都是群胚；见引理
\ref{categories-lemma-fibred-groupoids}。对 $V \in \Ob(\mathcal{C})$，有
$$
\mathcal{S}_V = \coprod\nolimits_{f : V \to U} \mathcal{S}_{(f : V \to U)}
$$
其中左端是 $p$ 的纤维范畴，右端是 $p'$ 的纤维范畴的不交并。
结论随即成立。
\end{proof}

\begin{lemma}
\label{categories-lemma-fibred-in-groupoids-over-fibred-in-groupoids}
令 $\mathcal{A} \to \mathcal{B} \to \mathcal{C}$ 为范畴之间的函子。
若 $\mathcal{A}$ 在 $\mathcal{B}$ 上以群胚为纤维，且
$\mathcal{B}$ 在 $\mathcal{C}$ 上以群胚为纤维，则
$\mathcal{A}$ 在 $\mathcal{C}$ 上以群胚为纤维。
\end{lemma}

\begin{proof}
可以直接由定义证明这一点。不过，我们将使用引理
\ref{categories-lemma-fibred-groupoids} 的判据。由引理
\ref{categories-lemma-fibred-over-fibred}，$\mathcal{A}$ 在
$\mathcal{C}$ 上纤维化。为了完成证明，只需说明对
$\mathcal{C}$ 中的 $U$，纤维范畴 $\mathcal{A}_U$ 是群胚。
事实上，若 $x \to y$ 是 $\mathcal{A}_U$ 的态射，则它在
$\mathcal{B}$ 中的像是同构，因为 $\mathcal{B}_U$ 是群胚。
但于是 $x \to y$ 是同构；例如，这可由引理
\ref{categories-lemma-composition-cartesian} 以及 $\mathcal{A}$ 的每个态射
都是强 $\mathcal{B}$-笛卡尔的这一事实推出（见引理
\ref{categories-lemma-fibred-groupoids}）。
\end{proof}

\begin{lemma}
\label{categories-lemma-fibred-groupoids-fibre-product-goes-up}
令 $p : \mathcal{S} \to \mathcal{C}$ 为群胚纤维化范畴。
令 $x \to y$ 与 $z \to y$ 为 $\mathcal{S}$ 的态射。若
$p(x) \times_{p(y)} p(z)$ 存在，则 $x \times_y z$ 存在，且
$p(x \times_y z) = p(x) \times_{p(y)} p(z)$。
\end{lemma}

\begin{proof}
由引理 \ref{categories-lemma-fibred-category-representable-goes-up} 推出。
\end{proof}


\begin{lemma}
\label{categories-lemma-ameliorate-morphism-categories-fibred-groupoids}
令 $\mathcal{C}$ 为范畴。令 $F : \mathcal{X} \to \mathcal{Y}$
为 $\mathcal{C}$ 上群胚纤维化范畴的一个 $1$-态射。存在由
$\mathcal{C}$ 上群胚纤维化范畴的 $1$-态射构成的分解
$\mathcal{X} \to \mathcal{X}' \to \mathcal{Y}$，使得
$\mathcal{X} \to \mathcal{X}'$ 是 $\mathcal{C}$ 上的等价，
并且 $\mathcal{X}'$ 是 $\mathcal{Y}$ 上的群胚纤维化范畴。
\end{lemma}

\begin{proof}
将结构函子记为 $p : \mathcal{X} \to \mathcal{C}$ 和
$q : \mathcal{Y} \to \mathcal{C}$。我们如下显式构造
$\mathcal{X}'$。$\mathcal{X}'$ 的一个对象是四元组
$(U, x, y, f)$，其中 $x \in \Ob(\mathcal{X}_U)$、
$y \in \Ob(\mathcal{Y}_U)$，而
$f : F(x) \to y$ 是 $\mathcal{Y}_U$ 中的同构。一个态射
$(a, b) : (U, x, y, f) \to (U', x', y', f')$ 由
$a : x \to x'$ 和 $b : y \to y'$ 给出，其中
$p(a) = q(b)$ 且 $f' \circ F(a) = b \circ f$。换言之，
$\mathcal{X}' = \mathcal{X} \times_{F, \mathcal{Y}, \text{id}} \mathcal{Y}$，
这里采用引理 \ref{categories-lemma-2-product-categories-over-C} 中的
$2$-纤维积构造。由引理 \ref{categories-lemma-2-product-fibred-categories}，
$\mathcal{X}'$ 是 $\mathcal{C}$ 上的群胚纤维化范畴，而且
$\mathcal{X}' \to \mathcal{Y}$ 是 $\mathcal{C}$ 上范畴的态射。
取函子 $\mathcal{X} \to \mathcal{X}'$，它在对象上由
$x \mapsto (p(x), x, F(x), \text{id}_{F(x)})$ 给出，在态射上由
$(a : x \to x') \mapsto (a, F(a))$ 给出。显然，复合
$\mathcal{X} \to \mathcal{X}' \to \mathcal{Y}$ 等于 $F$。
略去对 $\mathcal{X} \to \mathcal{X}'$ 是 $\mathcal{C}$ 上
纤维化范畴等价的验证。

\medskip\noindent
最后，需要证明 $\mathcal{X}' \to \mathcal{Y}$ 是群胚纤维化范畴。
令 $b : y' \to y$ 为 $\mathcal{Y}$ 中的态射，并令
$(U, x, y, f)$ 为 $\mathcal{X}'$ 中位于 $y$ 之上的对象。
因为 $\mathcal{X}$ 在 $\mathcal{C}$ 上以群胚为纤维，可以找到
位于 $U' = q(y') \to q(y) = U$ 之上的态射 $a : x' \to x$。
因为 $\mathcal{Y}$ 在 $\mathcal{C}$ 上以群胚为纤维，而且
$F(x') \to F(x)$ 与 $y' \to y$ 都位于同一个态射
$U' \to U$ 之上，可以找到位于 $\text{id}_{U'}$ 之上的
$f' : F(x') \to y'$，使 $f \circ F(a) = b \circ f'$。
由此得到
$(a, b) : (U', x', y', f') \to (U, x, y, f)$。
这验证了定义 \ref{categories-definition-fibred-groupoids} 的第一个条件 (1)。
为了验证 (2)，令
$(a, b) : (U', x', y', f') \to (U, x, y, f)$ 和
$(a', b') : (U'', x'', y'', f'') \to (U, x, y, f)$ 为
$\mathcal{X}'$ 的态射，并令 $b'' : y' \to y''$ 为
$\mathcal{Y}$ 中满足 $b' \circ b'' = b$ 的态射。需要证明存在唯一态射
$a'' : x' \to x''$，使
$f'' \circ F(a'') = b'' \circ f'$，并使
$(a', b') \circ (a'', b'') = (a, b)$。因为 $\mathcal{X}$
以群胚为纤维，存在唯一态射 $a'' : x' \to x''$，使
$a' \circ a'' = a$ 且 $p(a'') = q(b'')$。因为 $\mathcal{Y}$
以群胚为纤维，$F(a'')$ 是满足
$F(a') \circ F(a'') = F(a)$ 与 $q(F(a'')) = q(b'')$ 的唯一态射
$F(x') \to F(x'')$。关系
$f'' \circ F(a'') = b'' \circ f'$ 由此以及给定关系
$f \circ F(a) = b \circ f'$ 和
$f \circ F(a') = b' \circ f''$ 推出。
\end{proof}

\begin{lemma}
\label{categories-lemma-amelioration-unique}
令 $\mathcal{C}$ 为范畴。令 $F : \mathcal{X} \to \mathcal{Y}$
为 $\mathcal{C}$ 上群胚纤维化范畴的一个 $1$-态射。假设有
$2$-交换图表
$$
\xymatrix{
\mathcal{X}' \ar[rd]_f &
\mathcal{X} \ar[l]^a \ar[d]^F \ar[r]_b &
\mathcal{X}'' \ar[ld]^g \\
& \mathcal{Y}
}
$$
其中 $a$ 和 $b$ 是 $\mathcal{C}$ 上范畴的等价，而 $f$ 和 $g$
是群胚纤维化范畴。那么存在 $\mathcal{Y}$ 上范畴的等价
$h : \mathcal{X}'' \to \mathcal{X}'$，使 $h \circ b$ 与 $a$
作为 $\mathcal{C}$ 上范畴的 $1$-态射是 $2$-同构的。若上述图表
实际上交换，则可作出安排，使 $h \circ b$ 与 $a$ 作为
$\mathcal{Y}$ 上范畴的 $1$-态射是 $2$-同构的。
\end{lemma}

\begin{proof}
我们将证明，$\mathcal{Y}$ 上的 $\mathcal{X}'$ 和
$\mathcal{X}''$ 都等价于 $\mathcal{Y}$ 上的群胚纤维化范畴
$\mathcal{X} \times_{F, \mathcal{Y}, \text{id}} \mathcal{Y}$；
见引理
\ref{categories-lemma-ameliorate-morphism-categories-fibred-groupoids} 的证明。
在 $\mathcal{C}$ 上范畴所成的 $2$-范畴中，选择一个拟逆
$b^{-1} : \mathcal{X}'' \to \mathcal{X}$。由于图表的右三角形
$2$-交换，可见
$$
\xymatrix{
\mathcal{X} \ar[d]_F & \mathcal{X}'' \ar[l]^{b^{-1}} \ar[d]^g \\
\mathcal{Y} & \mathcal{Y} \ar[l]
}
$$
是 $2$-交换的。因此，由 $2$-纤维积的泛性质，得到一个 $1$-态射
$c : \mathcal{X}'' \to
\mathcal{X} \times_{F, \mathcal{Y}, \text{id}} \mathcal{Y}$。
此外，$c$ 是 $\mathcal{Y}$ 上范畴的态射（！），而且是等价
（由 $b$ 是等价这一假设；见引理
\ref{categories-lemma-equivalence-2-fibre-product}）。所以由引理
\ref{categories-lemma-equivalence-fibred-categories}，$c$ 是
$\mathcal{Y}$ 上群胚纤维化范畴所成的 $2$-范畴中的等价。

\medskip\noindent
还需构造 $c \circ b$ 与函子
$d : \mathcal{X} \to
\mathcal{X} \times_{F, \mathcal{Y}, \text{id}} \mathcal{Y}$ 之间的
$2$-同构；后者是在引理
\ref{categories-lemma-ameliorate-morphism-categories-fibred-groupoids} 的证明中
由 $x \mapsto (p(x), x, F(x), \text{id}_{F(x)})$ 构造的。
令 $\alpha : F \to g \circ b$ 与
$\beta : b^{-1} \circ b \to \text{id}$ 为 $\mathcal{C}$ 上范畴的
$1$-态射之间的 $2$-同构。注意，$c \circ b$ 在对象上由规则
$$
x \mapsto (p(x), b^{-1}(b(x)), g(b(x)), \alpha_x \circ F(\beta_x))
$$
给出。于是可见
$$
(\beta_x, \alpha_x) :
(p(x), x, F(x), \text{id}_{F(x)})
\longrightarrow
(p(x), b^{-1}(b(x)), g(b(x)), \alpha_x \circ F(\beta_x))
$$
是一个函子性同构，给出我们的 $2$-态射
$d \to b \circ c$。最后，如果图表交换，那么对所有 $x$，
$\alpha_x$ 都是恒等态射，从而这个 $2$-态射是在
$\mathcal{Y}$ 上范畴所成的 $2$-范畴中的 $2$-态射。
\end{proof}


\section{范畴预层}
\label{categories-section-presheaves-categories}

\noindent
本节将纤维化范畴的概念与关系密切的“范畴预层”概念作比较。
下面的例子说明基本构造。

\begin{example}
\label{categories-example-functor-categories}
令 $\mathcal{C}$ 为范畴。假设
$F : \mathcal{C}^{opp} \to \textit{Cat}$ 是到范畴的 $2$-范畴的函子；
见定义 \ref{categories-definition-functor-into-2-category}。对
$\mathcal{C}$ 中的 $f : V \to U$，我们启发性地将从 $F(U)$
到 $F(V)$ 的函子记为 $F(f) = f^\ast$。由此可以如下构造
$\mathcal{C}$ 上的纤维化范畴 $\mathcal{S}_F$。定义
$$
\Ob(\mathcal{S}_F) =
\{(U, x) \mid U\in \Ob(\mathcal{C}), x\in \Ob(F(U))\}.
$$
对 $(U, x), (V, y) \in \Ob(\mathcal{S}_F)$，定义
\begin{align*}
\Mor_{\mathcal{S}_F}((V, y), (U, x)) & =
\{ (f, \phi) \mid f \in \Mor_\mathcal{C}(V, U),
\phi \in \Mor_{F(V)}(y, f^\ast x)\} \\
& =
\coprod\nolimits_{f \in \Mor_\mathcal{C}(V, U)}
\Mor_{F(V)}(y, f^\ast x)
\end{align*}
为了定义复合，使用可复合的一对 $\mathcal{C}$ 中态射所满足的
$g^\ast \circ f^\ast = (f \circ g)^\ast$
（由取值于 $2$-范畴的函子的定义）。具体地，把
$\psi : z \to g^\ast y$ 与 $ \phi : y \to f^\ast x$ 的复合定义为
$ g^\ast(\phi) \circ \psi$。函子
$p_F : \mathcal{S}_F \to \mathcal{C}$ 由规则
$(U, x) \mapsto U$ 给出。我们来验证，这确实是纤维化范畴。
给定 $\mathcal{C}$ 中的 $f: V \to U$ 以及 $U$ 的提升
$(U, x)$，我们断言
$(f, \text{id}_{f^\ast x}): (V, {f^\ast x}) \to (U, x)$
是 $f$ 的强笛卡尔提升。需要证明左图中的 $h$ 确定右图中的
$(h, \nu)$：
$$
\xymatrix{
V \ar[r]^f &
U &
(V, f^*x) \ar[r]^{(f, \text{id}_{f^*x})} &
(U, x) \\
W \ar@{-->}[u]^h \ar[ru]_g & &
(W, z) \ar@{-->}[u]^{(h, \nu)} \ar[ru]_{(g, \psi)} &
}
$$
只需取 $\nu = \psi$；这是可行的，因为
$f \circ h = g$，从而 $g^*x = h^*f^*x$。此外，这是使右侧图表
交换的唯一提升。
\end{example}

\begin{definition}
\label{categories-definition-split-fibred-category}
令 $\mathcal{C}$ 为范畴。假设
$F : \mathcal{C}^{opp} \to \textit{Cat}$ 是到范畴的 $2$-范畴的
函子。把例 \ref{categories-example-functor-categories} 中构造的纤维化范畴
记为 $p_F : \mathcal{S}_F \to \mathcal{C}$。所谓
{\it 分裂纤维化范畴}，是一个在 $\mathcal{C}$ 上同构（！）于某个
这种范畴 {\it $\mathcal{S}_F$} 的纤维化范畴。
\end{definition}

\begin{lemma}
\label{categories-lemma-when-split}
令 $\mathcal{C}$ 为范畴。令 $\mathcal{S}$ 为
$\mathcal{C}$ 上的纤维化范畴。那么 $\mathcal{S}$ 是分裂的，
当且仅当对某个拉回选择（见定义
\ref{categories-definition-pullback-functor-fibred-category}），拉回函子
$(f \circ g)^*$ 与 $g^* \circ f^*$ 相等。
\end{lemma}

\begin{proof}
这直接由定义推出。
\end{proof}

\begin{lemma}
\label{categories-lemma-fibred-strict}
令 $ p : \mathcal{S} \to \mathcal{C}$ 为纤维化范畴。
存在一个反变函子 $F : \mathcal{C} \to \textit{Cat}$，使
$\mathcal{S}$ 与 $\mathcal{S}_F$ 在 $\mathcal{C}$ 上纤维化范畴所成
的 $2$-范畴中等价。换言之，每个纤维化范畴都等价于一个分裂的
纤维化范畴。
\end{lemma}

\begin{proof}
作一个拉回选择（见定义
\ref{categories-definition-pullback-functor-fibred-category}）。由引理
\ref{categories-lemma-fibred}，对 $\mathcal{C}$ 的每个态射 $f$，
得到拉回函子 $f^*$。

\medskip\noindent
如下构造一个新范畴 $\mathcal{S}'$。$\mathcal{S}'$ 的对象是有序对
$(x, f)$，它由 $\mathcal{C}$ 的态射 $f : V \to U$ 和位于 $U$
之上的 $\mathcal{S}$ 的对象 $x$ 组成；也就是说，
$x\in \Ob(\mathcal{S}_U)$。函子
$p' : \mathcal{S}' \to \mathcal{C}$ 将有序对 $(x, f)$ 映到态射
$f$ 的源；换言之，$p'(x, f : V\to U) = V$。一个态射
$\varphi : (x_1, f_1: V_1 \to U_1) \to
(x_2, f_2 : V_2 \to U_2)$
由有序对 $(\varphi, g)$ 给出，它由态射 $g : V_1 \to V_2$
以及满足 $p(\varphi) = g$ 的态射
$\varphi : f_1^\ast x_1 \to f_2^\ast x_2$ 组成。
定义复合法则毫无困难：对任意一对可复合态射，置
$(\varphi, g) \circ (\psi, h) =
(\varphi \circ \psi, g\circ h)$。有一个自然函子
$\mathcal{S} \to \mathcal{S}'$，它只是把位于 $U$ 之上的 $x$
映到有序对 $(x, \text{id}_U)$。

\medskip\noindent
此时需要验证 $p'$ 使 $\mathcal{S}'$ 成为 $\mathcal{C}$ 上的
纤维化范畴，并且需要验证 $\mathcal{S} \to \mathcal{S}'$
是 $\mathcal{C}$ 上范畴的一个等价，它将强笛卡尔态射映为
强笛卡尔态射。略去这些验证。

\medskip\noindent
最后，可以如下定义 $\mathcal{S}'$ 上的拉回函子：若
$g : V' \to V$ 且 $f : V \to U$，则在对象上置
$g^\ast(x, f) = (x, f \circ g)$。对于 $\mathcal{S}'_V$ 中态射之间
的态射
$(\varphi, \text{id}_V) : (x_1, f_1) \to (x_2, f_2)$，在态射上置
$g^\ast(\varphi, \text{id}_V) =
(g^\ast\varphi, \text{id}_{V'})$；这里使用引理
\ref{categories-lemma-fibred} 中的唯一等同
$g^\ast f_i^\ast x_i = (f_i \circ g)^\ast x_i$，将
$g^\ast\varphi$ 看作从 $(f_1 \circ g)^\ast x_1$ 到
$(f_2 \circ g)^\ast x_2$ 的态射。显然，这些拉回函子 $g^\ast$
满足 $g_1^\ast \circ g_2^\ast = (g_2\circ g_1)^\ast$；
换言之，$\mathcal{S}'$ 按所需方式分裂。
\end{proof}


\section{群胚预层}
\label{categories-section-presheaves-groupoids}

\noindent
本节将群胚纤维化范畴的概念与关系密切的“群胚预层”概念作比较。
下面的例子说明基本构造。

\begin{example}
\label{categories-example-functor-groupoids}
这是例 \ref{categories-example-functor-categories} 对“群胚预层”而非“范畴预层”
的类似物。输出将是群胚纤维化范畴，而非纤维化范畴。假设
$F : \mathcal{C}^{opp} \to \textit{Groupoids}$ 是到群胚范畴的函子；
见定义 \ref{categories-definition-functor-into-2-category}。对
$\mathcal{C}$ 中的 $f : V \to U$，我们启发性地将从 $F(U)$
到 $F(V)$ 的函子记为 $F(f) = f^\ast$。如下构造
$\mathcal{C}$ 上的群胚纤维化范畴 $\mathcal{S}_F$。定义
$$
\Ob(\mathcal{S}_F) =
\{(U, x) \mid U\in \Ob(\mathcal{C}), x\in \Ob(F(U))\}.
$$
对 $(U, x), (V, y) \in \Ob(\mathcal{S}_F)$，定义
\begin{align*}
\Mor_{\mathcal{S}_F}((V, y), (U, x))
& =
\{ (f, \phi) \mid f \in \Mor_\mathcal{C}(V, U),
\phi \in \Mor_{F(V)}(y, f^\ast x)\} \\
& =
\coprod\nolimits_{f \in \Mor_\mathcal{C}(V, U)}
\Mor_{F(V)}(y, f^\ast x)
\end{align*}
为了定义复合，使用可复合的一对 $\mathcal{C}$ 中态射所满足的
$g^\ast \circ f^\ast = (f \circ g)^\ast$
（由取值于 $2$-范畴的函子的定义）。具体地，把
$\psi : z \to g^\ast y$ 与 $ \phi : y \to f^\ast x$ 的复合定义为
$ g^\ast(\phi) \circ \psi$。函子
$p_F : \mathcal{S}_F \to \mathcal{C}$ 由规则
$(U, x) \mapsto U$ 给出。对每个 $U$，$F(U)$ 是群胚这一条件保证
$\mathcal{S}_F$ 在 $\mathcal{C}$ 上以群胚为纤维，因为我们已在
例 \ref{categories-example-functor-categories} 中看到 $\mathcal{S}_F$
是纤维化范畴；见引理 \ref{categories-lemma-fibred-groupoids}。
也可以如下直接证明定义 \ref{categories-definition-fibred-groupoids} 的条件
(1)、(2)：(1) 态射的提升存在，因为给定
$\mathcal{C}$ 中的 $f: V \to U$ 和 $\mathcal{S}_F$ 中位于 $U$
之上的对象 $(U, x)$，则
$(f, \text{id}_{f^\ast x}): (V, {f^\ast x}) \to (U, x)$
是 $f$ 的一个提升。(2) 假设给定下列实线图表：
$$
\xymatrix{
V \ar[r]^f & U & (V, y) \ar[r]^{(f, \phi)} & (U, x) \\
W \ar@{-->}[u]^h \ar[ru]_g & &
(W, z) \ar@{-->}[u]^{(h, \nu)} \ar[ru]_{(g, \psi)} & \\
}
$$
那么对虚线箭头有
$\nu = (h^\ast \phi)^{-1} \circ \psi$，所以给定 $h$ 后，
存在 $\nu$，并且由逆的唯一性，它是唯一的。
\end{example}

\begin{definition}
\label{categories-definition-split-category-fibred-in-groupoids}
令 $\mathcal{C}$ 为范畴。假设
$F : \mathcal{C}^{opp} \to \textit{Groupoids}$ 是到群胚的
$2$-范畴的函子。把例 \ref{categories-example-functor-groupoids} 中构造的
群胚纤维化范畴记为
$p_F : \mathcal{S}_F \to \mathcal{C}$。所谓
{\it 分裂群胚纤维化范畴}，是一个在 $\mathcal{C}$ 上同构（！）
于某个这种范畴 {\it $\mathcal{S}_F$} 的群胚纤维化范畴。
\end{definition}

\begin{lemma}
\label{categories-lemma-fibred-groupoids-strict}
令 $ p : \mathcal{S} \to \mathcal{C}$ 为群胚纤维化范畴。
存在一个反变函子 $F : \mathcal{C} \to \textit{Groupoids}$，
使 $\mathcal{S}$ 与 $\mathcal{S}_F$ 在 $\mathcal{C}$ 上等价。
换言之，每个群胚纤维化范畴都等价于一个分裂的群胚纤维化范畴。
\end{lemma}

\begin{proof}
作一个拉回选择（见定义
\ref{categories-definition-pullback-functor-fibred-category}）。由引理
\ref{categories-lemma-fibred} 和 \ref{categories-lemma-fibred-groupoids}，对
$\mathcal{C}$ 的每个态射 $f$，得到拉回函子 $f^*$。

\medskip\noindent
如下构造一个新范畴 $\mathcal{S}'$。$\mathcal{S}'$ 的对象是有序对
$(x, f)$，它由 $\mathcal{C}$ 的态射 $f : V \to U$ 和位于
$U$ 之上的 $\mathcal{S}$ 的对象 $x$ 组成；也就是说，
$x\in \Ob(\mathcal{S}_U)$。函子
$p' : \mathcal{S}' \to \mathcal{C}$ 将有序对 $(x, f)$ 映到态射
$f$ 的源；换言之，$p'(x, f : V\to U) = V$。一个态射
$\varphi : (x_1, f_1: V_1 \to U_1) \to
(x_2, f_2 : V_2 \to U_2)$
由有序对 $(\varphi, g)$ 给出，它由态射 $g : V_1 \to V_2$
以及满足 $p(\varphi) = g$ 的态射
$\varphi : f_1^\ast x_1 \to f_2^\ast x_2$ 组成。
定义复合法则毫无困难：对任意一对可复合态射，置
$(\varphi, g) \circ (\psi, h) =
(\varphi \circ \psi, g\circ h)$。有一个自然函子
$\mathcal{S} \to \mathcal{S}'$，它只是把位于 $U$ 之上的 $x$
映到有序对 $(x, \text{id}_U)$。

\medskip\noindent
此时需要验证 $p'$ 使 $\mathcal{S}'$ 成为 $\mathcal{C}$ 上的
群胚纤维化范畴，并且需要验证 $\mathcal{S} \to \mathcal{S}'$
是 $\mathcal{C}$ 上范畴的等价。略去这些验证。

\medskip\noindent
最后，可以如下定义 $\mathcal{S}'$ 上的拉回函子：若
$g : V' \to V$ 且 $f : V \to U$，则在对象上置
$g^\ast(x, f) = (x, f \circ g)$。对于 $\mathcal{S}'_V$ 中态射之间
的态射
$(\varphi, \text{id}_V) : (x_1, f_1) \to (x_2, f_2)$，在态射上置
$g^\ast(\varphi, \text{id}_V) =
(g^\ast\varphi, \text{id}_{V'})$；这里使用引理
\ref{categories-lemma-fibred-groupoids} 中的唯一等同
$g^\ast f_i^\ast x_i = (f_i \circ g)^\ast x_i$，将
$g^\ast\varphi$ 看作从 $(f_1 \circ g)^\ast x_1$ 到
$(f_2 \circ g)^\ast x_2$ 的态射。显然，这些拉回函子 $g^\ast$
满足 $g_1^\ast \circ g_2^\ast = (g_2\circ g_1)^\ast$；
换言之，$\mathcal{S}'$ 按所需方式分裂。
\end{proof}

\noindent
我们将在第 \ref{categories-section-representable-1-morphisms} 节看到这个引理的
另一种证明。


\section{集合纤维化范畴}
\label{categories-section-fibred-in-sets}

\begin{definition}
\label{categories-definition-discrete}
如果一个范畴仅有恒等态射，则称其为{\it 离散的}。
\end{definition}

\noindent
离散范畴只有一项有意义的信息：它的对象集合。因此，我们有时
混同离散范畴与集合。

\begin{definition}
\label{categories-definition-category-fibred-sets}
令 $\mathcal{C}$ 为范畴。所谓{\it 集合纤维化范畴}，或者
{\it 离散范畴纤维化范畴}，是所有纤维范畴均为离散范畴的
群胚纤维化范畴。
\end{definition}

\noindent
我们想阐明集合纤维化范畴与预层（见定义
\ref{categories-definition-presheaf}）之间的关系。为此，先作如下定义是
合理的。

\begin{definition}
\label{categories-definition-categories-fibred-in-sets-over-C}
令 $\mathcal{C}$ 为范畴。所谓{\it $\mathcal{C}$ 上集合纤维化范畴
所成的 $2$-范畴}，是 $\mathcal{C}$ 上群胚纤维化范畴所成范畴
（见定义 \ref{categories-definition-categories-fibred-in-groupoids-over-C}）
的子 $2$-范畴，定义如下：
\begin{enumerate}
\item 它的对象是集合纤维化范畴
$p : \mathcal{S} \to \mathcal{C}$。
\item 它的 $1$-态射
$(\mathcal{S}, p) \to (\mathcal{S}', p')$ 是满足
$p' \circ G = p$ 的函子 $G : \mathcal{S} \to \mathcal{S}'$
（由于每个态射都是强笛卡尔的，$G$ 自动保持它们）。
\item 对
$G, H : (\mathcal{S}, p) \to (\mathcal{S}', p')$，它的 $2$-态射
$t : G \to H$ 是满足下述条件的函子态射：对所有
$x \in \Ob(\mathcal{S})$，均有
$p'(t_x) = \text{id}_{p(x)}$。
\end{enumerate}
\end{definition}

\noindent
注意，每个 $2$-态射自动是同构。因此，这个 $2$-范畴实际上是
$(2, 1)$-范畴。下面给出关于 $2$-纤维积存在性的必备引理。

\begin{lemma}
\label{categories-lemma-2-product-categories-fibred-sets}
令 $\mathcal{C}$ 为范畴。$\mathcal{C}$ 上集合纤维化范畴所成的
2-范畴具有 2-纤维积。更精确地说，如果从集合纤维化范畴出发，
则引理 \ref{categories-lemma-2-product-categories-over-C} 中描述的
2-纤维积仍给出集合纤维化范畴。
\end{lemma}

\begin{proof}
略。
\end{proof}

\begin{example}
\label{categories-example-presheaf}
这是例 \ref{categories-example-functor-categories} 和
\ref{categories-example-functor-groupoids} 对预层而非“范畴预层”的类似物。
输出将是集合纤维化范畴，而非纤维化范畴。假设
$F : \mathcal{C}^{opp} \to \textit{Sets}$ 是预层。对
$\mathcal{C}$ 中的 $f : V \to U$，我们启发性地记
$F(f) = f^\ast : F(U) \to F(V)$。如下构造 $\mathcal{C}$ 上的
集合纤维化范畴 $\mathcal{S}_F$。定义
$$
\Ob(\mathcal{S}_F) =
\{(U, x) \mid U \in \Ob(\mathcal{C}), x \in \Ob(F(U))\}.
$$
对 $(U, x), (V, y) \in \Ob(\mathcal{S}_F)$，定义
\begin{align*}
\Mor_{\mathcal{S}_F}((V, y), (U, x))
& =
\{f \in \Mor_\mathcal{C}(V, U) \mid f^*x = y\}
\end{align*}
复合继承自 $\mathcal{C}$ 中的复合；这是可行的，因为对
$\mathcal{C}$ 中可复合的一对态射，有
$g^\ast \circ f^\ast = (f \circ g)^\ast$。函子
$p_F : \mathcal{S}_F \to \mathcal{C}$ 由规则
$(U, x) \mapsto U$ 给出。由于每个纤维范畴
$\mathcal{S}_{F, U}$ 都是以 $F(U)$ 为底层集合的离散范畴，
并且我们已在例 \ref{categories-example-functor-groupoids} 中看到
$\mathcal{S}_F$ 是群胚纤维化范畴，所以可知
$\mathcal{S}_F$ 以集合为纤维。
\end{example}

\begin{lemma}
\label{categories-lemma-2-category-fibred-sets}
\begin{slogan}
集合纤维化范畴恰好就是预层。
\end{slogan}
令 $\mathcal{C}$ 为范畴。集合纤维化范畴之间仅有的
$2$-态射是恒等态射。换言之，集合纤维化范畴所成的 $2$-范畴是
一个范畴。此外，存在范畴等价
$$
\left\{
\begin{matrix}
\text{预层范畴}\\
\text{集合范畴，位于 }\mathcal{C}
\end{matrix}
\right\}
\leftrightarrow
\left\{
\begin{matrix}
\text{范畴的范畴}\\
\text{在下列对象之上以集合为纤维： }\mathcal{C}
\end{matrix}
\right\}
$$
从左到右的函子是例 \ref{categories-example-presheaf} 中讨论的构造
$F \to \mathcal{S}_F$。从右到左的函子，把
$p : \mathcal{S} \to \mathcal{C}$ 映到对象预层
$U \mapsto \Ob(\mathcal{S}_U)$。
\end{lemma}

\begin{proof}
第一个断言显然，因为纤维范畴中仅有的态射是恒等态射。

\medskip\noindent
假设 $p :
\mathcal{S} \to \mathcal{C}$ 以集合为纤维。令
$f : V \to U$ 为 $\mathcal{C}$ 中的态射，令
$x \in \Ob(\mathcal{S}_U)$。那么对象 $f^\ast x$ 恰有一种选择。
因此，对引理 \ref{categories-lemma-fibred-groupoids} 中的 $f, g$，有
$(f \circ g)^\ast x = g^\ast(f^\ast x)$。所以可以把赋值
$U \mapsto \Ob(\mathcal{S}_U)$ 和 $f \mapsto f^\ast$
看作 $\mathcal{C}$ 上的预层。
\end{proof}

\noindent
下面给出集合纤维化范畴的一个重要例子。

\begin{example}
\label{categories-example-fibred-category-from-functor-of-points}
令 $\mathcal{C}$ 为范畴，令 $X \in \Ob(\mathcal{C})$。
考虑可表预层 $h_X = \Mor_\mathcal{C}(-, X)$
（见例 \ref{categories-example-hom-functor}）。另一方面，考虑例
\ref{categories-example-category-over-X} 中的范畴
$p : \mathcal{C}/X \to \mathcal{C}$。纤维范畴
$(\mathcal{C}/X)_U$ 的对象是态射 $h : U \to X$，而且只有
恒等态射。因此，在引理 \ref{categories-lemma-2-category-fibred-sets} 的
对应下，有
$$
h_X \longleftrightarrow \mathcal{C}/X.
$$
换言之，范畴 $\mathcal{C}/X$ 典范等价于例
\ref{categories-example-presheaf} 中与 $h_X$ 关联的范畴
$\mathcal{S}_{h_X}$。
\end{example}

\noindent
因此，很容易想把群胚纤维化范畴的 2-范畴中的“可表”对象定义为：
其关联预层可表的集合纤维化范畴。然而，这并不是适用的好定义，
因为我们更希望有一个在等价下不变的概念。为使这一点精确化，
我们将研究究竟哪些群胚纤维化范畴等价于集合纤维化范畴。


\section{集合胚纤维化范畴}
\label{categories-section-fibred-in-setoids}

\begin{definition}
\label{categories-definition-setoid}
如果一个范畴是群胚，且每个对象恰有一个自同构——恒等态射，
则称其为{\it 集合胚（setoid）}\footnote{吃了类固醇的集合！？}。
\end{definition}

\noindent
若 $C$ 是带等价关系 $\sim$ 的集合，则可以如下构造集合胚
$\mathcal{C}$：$\Ob(\mathcal{C}) = C$；除非 $x \sim y$，
否则 $\Mor_\mathcal{C}(x, y) = \emptyset$；而在该关系成立时，
置 $\Mor_\mathcal{C}(x, y) = \{1\}$。
$\sim$ 的传递性意味着可以复合态射。反过来，任意集合胚范畴在其
对象上定义一个等价关系（同构关系），而通过刚才描述的过程，
可以从这个关系恢复原范畴（在唯一同构意义下，而不只是在等价
意义下）。

\medskip\noindent
离散范畴是集合胚。对任意集合胚 $\mathcal{C}$，有一个典范过程
产生与之等价的离散范畴，即用同构类集合替代
$\Ob(\mathcal{C})$（并添加恒等态射）。用带等价关系的集合来说，
这对应于对等价关系取商。

\begin{definition}
\label{categories-definition-category-fibred-setoids}
令 $\mathcal{C}$ 为范畴。所谓{\it 集合胚纤维化范畴}，是所有纤维
范畴均为集合胚的群胚纤维化范畴。
\end{definition}

\noindent
下面将阐明集合胚纤维化范畴与集合纤维化范畴之间的关系。

\begin{definition}
\label{categories-definition-categories-fibred-in-setoids-over-C}
令 $\mathcal{C}$ 为范畴。所谓{\it $\mathcal{C}$ 上集合胚纤维化范畴
所成的 $2$-范畴}，是 $\mathcal{C}$ 上群胚纤维化范畴所成范畴
（见定义 \ref{categories-definition-categories-fibred-in-groupoids-over-C}）
的子 $2$-范畴，定义如下：
\begin{enumerate}
\item 它的对象是集合胚纤维化范畴
$p : \mathcal{S} \to \mathcal{C}$。
\item 它的 $1$-态射
$(\mathcal{S}, p) \to (\mathcal{S}', p')$ 是满足
$p' \circ G = p$ 的函子 $G : \mathcal{S} \to \mathcal{S}'$
（由于每个态射都是强笛卡尔的，$G$ 自动保持它们）。
\item 对
$G, H : (\mathcal{S}, p) \to (\mathcal{S}', p')$，它的 $2$-态射
$t : G \to H$ 是满足下述条件的函子态射：对所有
$x \in \Ob(\mathcal{S})$，均有
$p'(t_x) = \text{id}_{p(x)}$。
\end{enumerate}
\end{definition}

\noindent
注意，每个 $2$-态射自动是同构。因此，这个 $2$-范畴实际上是
$(2, 1)$-范畴。

\noindent
下面给出关于 $2$-纤维积存在性的必备引理。

\begin{lemma}
\label{categories-lemma-2-product-categories-fibred-setoids}
令 $\mathcal{C}$ 为范畴。$\mathcal{C}$ 上集合胚纤维化范畴所成的
2-范畴具有 2-纤维积。更精确地说，如果从集合胚纤维化范畴出发，
则引理 \ref{categories-lemma-2-product-categories-over-C} 中描述的
2-纤维积仍给出集合胚纤维化范畴。
\end{lemma}

\begin{proof}
略。
\end{proof}

\begin{lemma}
\label{categories-lemma-setoid-fibres}
令 $\mathcal{C}$ 为范畴。令 $\mathcal{S}$ 为
$\mathcal{C}$ 上的范畴。
\begin{enumerate}
\item 若 $\mathcal{S} \to \mathcal{S}'$ 是 $\mathcal{C}$ 上的等价，
而 $\mathcal{S}'$ 是 $\mathcal{C}$ 上的集合纤维化范畴，则
\begin{enumerate}
\item $\mathcal{S}$ 是 $\mathcal{C}$ 上的集合胚纤维化范畴；并且
\item 对每个 $U \in \Ob(\mathcal{C})$，映射
$\Ob(\mathcal{S}_U) \to \Ob(\mathcal{S}'_U)$ 将靶标识为源的
同构类集合。
\end{enumerate}
\item 若 $p : \mathcal{S} \to \mathcal{C}$ 是集合胚纤维化范畴，
则存在集合纤维化范畴
$p' : \mathcal{S}' \to \mathcal{C}$，以及 $\mathcal{C}$ 上的等价
$\text{can} : \mathcal{S} \to \mathcal{S}'$。
\end{enumerate}
\end{lemma}

\begin{proof}
我们证明 (2)。范畴 $\mathcal{S}'$ 的一个对象是有序对
$(U, \xi)$，其中 $U \in \Ob(\mathcal{C})$，而 $\xi$ 是
$\mathcal{S}_U$ 的对象的一个同构类。一个态射
$(U, \xi) \to (V , \psi)$ 由态射 $x \to y$ 给出，其中
$x \in \xi$ 且 $y \in \psi$。这里，如果两个态射
$x \to y$ 与 $x' \to y'$ 诱导相同的态射 $U \to V$，并且对
$\mathcal{S}_U$ 中同构 $x \to x'$ 和 $\mathcal{S}_V$ 中同构
$y \to y'$ 的某些选择，复合 $x \to x' \to y'$ 与
$x \to y \to y'$ 相同，就把这两个态射等同。由构造，
$\mathcal{S} \to \mathcal{S}'$ 在对象和态射上都有满映射。
把 $\mathcal{S}'$ 中的态射复合定义为使
$\mathcal{S} \to \mathcal{S}'$ 成为函子的唯一法则。
略去若干细节。
\end{proof}


\noindent
所以，集合胚纤维化范畴恰好就是等价于集合纤维化范畴的群胚
纤维化范畴。此外，由引理 \ref{categories-lemma-2-category-fibred-sets}，
集合纤维化范畴之间的等价是同构。

\begin{lemma}
\label{categories-lemma-2-category-fibred-setoids}
令 $\mathcal{C}$ 为范畴。引理 \ref{categories-lemma-setoid-fibres} 第 (2) 部分
的构造给出函子
$$
F :
\left\{
\begin{matrix}
\text{范畴的 2-范畴}\\
\text{在下列对象之上以集合胚为纤维： }\mathcal{C}
\end{matrix}
\right\}
\longrightarrow
\left\{
\begin{matrix}
\text{范畴的范畴}\\
\text{在下列对象之上以集合为纤维： }\mathcal{C}
\end{matrix}
\right\}
$$
（见定义 \ref{categories-definition-functor-into-2-category}）。这个函子在下述
意义下是等价：
\begin{enumerate}
\item 对满足 $F(f) = F(g)$ 的任意两个 1-态射
$f, g : \mathcal{S}_1 \to \mathcal{S}_2$，存在唯一的 2-同构
$f \to g$；
\item 对任意态射 $h : F(\mathcal{S}_1) \to F(\mathcal{S}_2)$，
存在满足 $F(f) = h$ 的 1-态射
$f : \mathcal{S}_1 \to \mathcal{S}_2$；并且
\item 任意集合纤维化范畴 $\mathcal{S}$ 都等于
$F(\mathcal{S})$。
\end{enumerate}
特别地，按规则
$F_i(U) = \Ob(\mathcal{S}_{i, U})/\cong$
定义 $F_i \in \textit{PSh}(\mathcal{C})$，则有
$$
\Mor_{\textit{Cat}/\mathcal{C}}(\mathcal{S}_1, \mathcal{S}_2)
\Big/
2\text{-同构}
=
\Mor_{\textit{PSh}(\mathcal{C})}(F_1, F_2)
$$
更精确地说，给定任意映射 $\phi : F_1 \to F_2$，存在一个
$1$-态射 $f : \mathcal{S}_1 \to \mathcal{S}_2$，它在对象的
同构类上诱导 $\phi$，并且在唯一 $2$-同构意义下唯一。
\end{lemma}

\begin{proof}
由引理 \ref{categories-lemma-2-category-fibred-sets}，$F$ 的靶是一个范畴，
所以这个断言有意义。引理 \ref{categories-lemma-setoid-fibres} 第 (2) 部分
的构造，把 $\mathcal{S}$ 映到一个集合纤维化范畴，其在 $U$
上方的值是 $\mathcal{S}_U$ 中同构类的集合。因此，显然它定义了
所述函子。令 $f, g : \mathcal{S}_1 \to \mathcal{S}_2$ 满足
$F(f) = F(g)$，如 (1) 所述。对 $\mathcal{C}$ 的每个对象 $U$
以及 $\mathcal{S}_{1, U}$ 的每个对象 $x$，由假设有
$f(x) \cong g(x)$。由于 $\mathcal{S}_2$ 以集合胚为纤维，
在 $\mathcal{S}_{2, U}$ 中存在唯一同构
$t_x : f(x) \to g(x)$。显然，规则 $x \mapsto t_x$ 给出所需的
$2$-同构 $f \to g$。略去 (2) 与 (3) 的证明。为了得到最后一个
断言，使用引理 \ref{categories-lemma-2-category-fibred-sets} 看出右端等于
$\Mor_{\textit{Cat}/\mathcal{C}}(F(\mathcal{S}_1), F(\mathcal{S}_2))$，
再应用上述 (1) 与 (2)。
\end{proof}

\noindent
下面给出集合胚纤维化范畴在所有群胚纤维化范畴中的另一种刻画。

\begin{lemma}
\label{categories-lemma-characterize-fibred-setoids-inertia}
令 $\mathcal{C}$ 为范畴。令
$p : \mathcal{S} \to \mathcal{C}$ 为群胚纤维化范畴。下列条件等价：
\begin{enumerate}
\item $p : \mathcal{S} \to \mathcal{C}$ 是集合胚纤维化范畴；并且
\item 典范 $1$-态射 $\mathcal{I}_\mathcal{S} \to \mathcal{S}$
（见 (\ref{categories-equation-inertia-structure-map})）是等价
（$\mathcal{C}$ 上范畴的等价）。
\end{enumerate}
\end{lemma}

\begin{proof}
假设 (2)。范畴 $\mathcal{I}_\mathcal{S}$ 的对象是
$(x, \alpha)$，其中 $x \in \mathcal{S}$，比如
$p(x) = U$，而 $\alpha : x \to x$ 是 $\mathcal{S}_U$ 中的态射。
因此，若
$\mathcal{I}_\mathcal{S} \to \mathcal{S}$ 是 $\mathcal{C}$
上的等价，则每一对对象 $(x, \alpha)$、$(x, \alpha')$ 都在
$\mathcal{I}_\mathcal{S}$ 位于 $U$ 上方的纤维范畴中同构。
考察 $\mathcal{I}_\mathcal{S}$ 中态射的定义，可知
$\alpha$、$\alpha'$ 在 $x$ 的自同构群中共轭。因此，取
$\alpha' = \text{id}_x$，便知 $x$ 的每个自同构都等于恒等态射。
由于 $\mathcal{S} \to \mathcal{C}$ 以群胚为纤维，这蕴含
$\mathcal{S} \to \mathcal{C}$ 以集合胚为纤维。略去
(1) $\Rightarrow$ (2) 的证明。
\end{proof}

\begin{lemma}
\label{categories-lemma-category-fibred-setoids-presheaves-products}
令 $\mathcal{C}$ 为范畴。引理
\ref{categories-lemma-2-category-fibred-setoids} 中将预层关联于集合胚
纤维化范畴的构造与积相容；其含义是，与 $2$-纤维积
$\mathcal{X} \times_\mathcal{Y} \mathcal{Z}$ 关联的预层，
是分别与 $\mathcal{X}, \mathcal{Y}, \mathcal{Z}$ 关联的预层的
纤维积。
\end{lemma}

\begin{proof}
令 $U \in \Ob(\mathcal{C})$。引理所说的恰好是
$$
\Ob((\mathcal{X} \times_\mathcal{Y} \mathcal{Z})_U)/\!\cong
\quad \text{等于} \quad
\Ob(\mathcal{X}_U)/\!\cong
\ \times_{\Ob(\mathcal{Y}_U)/\!\cong}
\ \Ob(\mathcal{Z}_U)/\!\cong
$$
略去其证明。（但注意，如果范畴 $\mathcal{Y}_U$ 不是集合胚，
这在一般情形并不成立。）
\end{proof}


\section{可表群胚纤维化范畴}
\label{categories-section-representable-fibred-groupoids}

\noindent
下面给出可表群胚纤维化范畴的定义。正如先前承诺的，这个定义在
等价下不变。

\begin{definition}
\label{categories-definition-representable-fibred-category}
令 $\mathcal{C}$ 为范畴。如果存在 $\mathcal{C}$ 的对象 $X$
以及一个等价 $j : \mathcal{S} \to \mathcal{C}/X$
（在 $\mathcal{C}$ 上群胚纤维化范畴所成的 $2$-范畴中），
则称群胚纤维化范畴
$p : \mathcal{S} \to \mathcal{C}$ 是{\it 可表的}。
\end{definition}

\noindent
通常会滥用记号说{\it $X$ 表示 $\mathcal{S}$}，而不提等价 $j$。
下面详述这意味着什么。

\begin{lemma}
\label{categories-lemma-characterize-representable-fibred-category}
令 $\mathcal{C}$ 为范畴。令
$p : \mathcal{S} \to \mathcal{C}$ 为群胚纤维化范畴。
\begin{enumerate}
\item $\mathcal{S}$ 可表，当且仅当下列条件成立：
\begin{enumerate}
\item $\mathcal{S}$ 以集合胚为纤维；并且
\item 预层 $U \mapsto \Ob(\mathcal{S}_U)/\cong$ 可表。
\end{enumerate}
\item 若 $\mathcal{S}$ 可表，则有序对 $(X, j)$ 在同构意义下唯一
确定，其中 $j$ 是等价 $j : \mathcal{S} \to \mathcal{C}/X$。
\end{enumerate}
\end{lemma}

\begin{proof}
第一个断言直接由引理 \ref{categories-lemma-setoid-fibres} 推出。
对第二个断言，设 $j' : \mathcal{S} \to \mathcal{C}/X'$
是第二个这样的有序对。选择一个 $1$-态射
$t' : \mathcal{C}/X' \to \mathcal{S}$，使
$j' \circ t' \cong \text{id}_{\mathcal{C}/X'}$ 且
$t' \circ j' \cong \text{id}_\mathcal{S}$。那么
$j \circ t' : \mathcal{C}/X' \to \mathcal{C}/X$ 是等价。
由引理 \ref{categories-lemma-2-category-fibred-sets}，它因而是同构。
所以由米田引理 \ref{categories-lemma-yoneda}（例如经由例
\ref{categories-example-fibred-category-from-functor-of-points}），
它由一个同构 $X' \to X$ 给出。
\end{proof}

\begin{lemma}
\label{categories-lemma-morphisms-representable-fibred-categories}
令 $\mathcal{C}$ 为范畴。令 $\mathcal{X}$、$\mathcal{Y}$
为 $\mathcal{C}$ 上的群胚纤维化范畴。假设
$\mathcal{X}$、$\mathcal{Y}$ 分别由 $\mathcal{C}$ 的对象
$X$、$Y$ 表示。那么
$$
\Mor_{\textit{Cat}/\mathcal{C}}(\mathcal{X}, \mathcal{Y})
\Big/
2\text{-同构}
=
\Mor_\mathcal{C}(X, Y)
$$
更精确地说，给定 $\phi : X \to Y$，存在一个 $1$-态射
$f : \mathcal{X} \to \mathcal{Y}$，它在对象的同构类上诱导
$\phi$，并且在唯一 $2$-同构意义下唯一。
\end{lemma}

\begin{proof}
由例 \ref{categories-example-fibred-category-from-functor-of-points}，有
$\mathcal{C}/X = \mathcal{S}_{h_X}$ 和
$\mathcal{C}/Y = \mathcal{S}_{h_Y}$。由引理
\ref{categories-lemma-2-category-fibred-setoids}，有
$$
\Mor_{\textit{Cat}/\mathcal{C}}(\mathcal{X}, \mathcal{Y})
\Big/
2\text{-同构}
=
\Mor_{\textit{PSh}(\mathcal{C})}(h_X, h_Y)
$$
由米田引理 \ref{categories-lemma-yoneda}，有
$\Mor_{\textit{PSh}(\mathcal{C})}(h_X, h_Y)
= \Mor_\mathcal{C}(X, Y)$。
\end{proof}


\section{2-米田引理}
\label{categories-section-2-yoneda}

\noindent
令 $\mathcal{C}$ 为范畴。$\mathcal{C}$ 上纤维范畴所成的 $2$-范畴
已在定义 \ref{categories-definition-fibred-categories-over-C} 中构造并定义。
若 $\mathcal{S}$、$\mathcal{S}'$ 是 $\mathcal{C}$ 上的纤维范畴，则
$$
\Mor_{\textit{Fib}/\mathcal{C}}(\mathcal{S}, \mathcal{S}')
$$
表示这个 $2$-范畴中 $1$-态射所成的范畴。下面给出纤维范畴情形下
米田引理的 $2$-范畴类似物。

\begin{lemma}[纤维范畴的 2-米田引理]
\label{categories-lemma-yoneda-2category-fibred}
令 $\mathcal{C}$ 为范畴。令 $\mathcal{S} \to \mathcal{C}$
为 $\mathcal{C}$ 上的纤维范畴。令 $U \in \Ob(\mathcal{C})$。
由 $G \mapsto G(\text{id}_U)$ 给出的函子
$$
\Mor_{\textit{Fib}/\mathcal{C}}(\mathcal{C}/U, \mathcal{S})
\longrightarrow
\mathcal{S}_U
$$
是等价。
\end{lemma}

\begin{proof}
为 $\mathcal{S}$ 选择拉回（见定义
\ref{categories-definition-pullback-functor-fibred-category}）。我们如下定义函子
$$
\mathcal{S}_U
\longrightarrow
\Mor_{\textit{Fib}/\mathcal{C}}(\mathcal{C}/U, \mathcal{S})
$$
给定 $x \in \Ob(\mathcal{S}_U)$，相应的函子满足：
\begin{enumerate}
\item 在对象上：$(f : V \to U) \mapsto f^*x$；
\item 在态射上：箭头 $(g : V'/U \to V/U)$ 映到复合
$$
(f \circ g)^*x \xrightarrow{(\alpha_{g, f})_x} g^*f^*x \rightarrow f^*x
$$
其中 $\alpha_{g, f}$ 如引理 \ref{categories-lemma-fibred} 中所定义。
\end{enumerate}
略去验证它是引理中函子的逆。
\end{proof}

\noindent
令 $\mathcal{C}$ 为范畴。$\mathcal{C}$ 上群胚纤维化范畴所成的
$2$-范畴，是 $\mathcal{C}$ 上范畴所成 $2$-范畴的一个“满”子
$2$-范畴（见定义
\ref{categories-definition-categories-fibred-in-groupoids-over-C}）。因此，若
$\mathcal{S}$、$\mathcal{S}'$ 是 $\mathcal{C}$ 上的群胚纤维化范畴，则
$$
\Mor_{\textit{Cat}/\mathcal{C}}(\mathcal{S}, \mathcal{S}')
$$
表示这个 $2$-范畴中 $1$-态射所成的范畴（见定义
\ref{categories-definition-categories-over-C}）。这些范畴全都是群胚，见定义
\ref{categories-definition-categories-fibred-in-groupoids-over-C} 后的说明。
下面给出米田引理的 $2$-范畴类似物。

\begin{lemma}[2-米田引理]
\label{categories-lemma-yoneda-2category}
令 $\mathcal{S}\to \mathcal{C}$ 为群胚纤维化范畴。
令 $U \in \Ob(\mathcal{C})$。由 $G \mapsto G(\text{id}_U)$ 给出的函子
$$
\Mor_{\textit{Cat}/\mathcal{C}}(\mathcal{C}/U, \mathcal{S})
\longrightarrow
\mathcal{S}_U
$$
是等价。
\end{lemma}

\begin{proof}
为 $\mathcal{S}$ 选择拉回（见定义
\ref{categories-definition-pullback-functor-fibred-category}）。我们如下定义函子
$$
\mathcal{S}_U
\longrightarrow
\Mor_{\textit{Cat}/\mathcal{C}}(\mathcal{C}/U, \mathcal{S})
$$
给定 $x \in \Ob(\mathcal{S}_U)$，相应的函子满足：
\begin{enumerate}
\item 在对象上：$(f : V \to U) \mapsto f^*x$；
\item 在态射上：箭头 $(g : V'/U \to V/U)$ 映到复合
$$
(f \circ g)^*x \xrightarrow{(\alpha_{g, f})_x} g^*f^*x \rightarrow f^*x
$$
其中 $\alpha_{g, f}$ 如引理 \ref{categories-lemma-fibred-groupoids} 中所定义。
\end{enumerate}
略去验证它是引理中函子的逆。
\end{proof}

\begin{remark}
\label{categories-remark-alternative-fibred-groupoids-strict}
我们可以用 $2$-米田引理给出引理
\ref{categories-lemma-fibred-groupoids-strict} 的另一证明。令
$p : \mathcal{S} \to \mathcal{C}$ 为群胚纤维化范畴。如下定义从
$\mathcal{C}$ 到群胚范畴的反变函子 $F$：对
$U\in \Ob(\mathcal{C})$，令
$$
F(U) = \Mor_{\textit{Cat}/\mathcal{C}}(\mathcal{C}/U, \mathcal{S}).
$$
若 $f : U \to V$，则诱导函子 $\mathcal{C}/U \to \mathcal{C}/V$
诱导态射 $F(f) : F(V) \to F(U)$。显然 $F$ 是函子。令
$\mathcal{S}'$ 为例 \ref{categories-example-functor-groupoids} 中与之相伴的
群胚纤维化范畴。有一个显然的 $\mathcal{C}$ 上函子
$G : \mathcal{S}' \to \mathcal{S}$，它把有序对 $(U, x)$——其中
$U \in \Ob(\mathcal{C})$ 且 $x \in F(U)$——映到
$x(\text{id}_U) \in \mathcal{S}$。现在引理
\ref{categories-lemma-yoneda-2category} 蕴含对每个 $U$，
$$
G_U : \mathcal{S}'_U = F(U)=
\Mor_{\textit{Cat}/\mathcal{C}}(\mathcal{C}/U, \mathcal{S})
\to
\mathcal{S}_U
$$
是等价，故由引理 \ref{categories-lemma-equivalence-fibred-categories}，
$G$ 是 $\mathcal{S}$ 与 $\mathcal{S}'$ 之间的等价。
\end{remark}





\section{可表 1-态射}
\label{categories-section-representable-1-morphisms}

\noindent
令 $\mathcal{C}$ 为范畴。本节说明 $\mathcal{C}$ 上两个群胚纤维化范畴之间
的 $1$-态射可表是什么意思。

\medskip\noindent
令 $\mathcal{C}$ 为范畴。令 $\mathcal{X}$、$\mathcal{Y}$ 为
$\mathcal{C}$ 上的群胚纤维化范畴。令 $U \in \Ob(\mathcal{C})$。
令 $F : \mathcal{X} \to \mathcal{Y}$ 和
$G : \mathcal{C}/U \to \mathcal{Y}$ 为 $\mathcal{C}$ 上群胚纤维化范畴
的 $1$-态射。我们想描述 $2$-纤维积
$$
\xymatrix{
(\mathcal{C}/U) \times_\mathcal{Y} \mathcal{X} \ar[r] \ar[d] &
\mathcal{X} \ar[d]^F \\
\mathcal{C}/U \ar[r]^G &
\mathcal{Y}
}
$$
令 $y = G(\text{id}_U) \in \mathcal{Y}_U$。为 $\mathcal{Y}$
选择拉回（见定义 \ref{categories-definition-pullback-functor-fibred-category}）。
那么 $G$ 同构于函子 $(f : V \to U) \mapsto f^*y$，见引理
\ref{categories-lemma-yoneda-2category} 及其证明。我们可以把
$(\mathcal{C}/U) \times_\mathcal{Y} \mathcal{X}$ 的对象看作四元组
$(V, f : V \to U, x, \phi)$，见引理
\ref{categories-lemma-2-product-categories-over-C}。利用上面对 $G$ 的描述，
可以把 $\phi$ 看作 $\mathcal{Y}_V$ 中的同构
$\phi : f^*y \to F(x)$。

\begin{lemma}
\label{categories-lemma-identify-fibre-product}
在上述情形中，
$(\mathcal{C}/U) \times_\mathcal{Y} \mathcal{X}$ 在
$\mathcal{C}/U$ 的对象 $f : V \to U$ 上的纤维范畴描述如下：
\begin{enumerate}
\item 对象是有序对 $(x, \phi)$，其中
$x \in \Ob(\mathcal{X}_V)$，且
$\phi : f^*y \to F(x)$ 是 $\mathcal{Y}_V$ 中的态射；
\item $(x, \phi)$ 与 $(x', \phi')$ 之间的态射集，是
$\mathcal{X}_V$ 中满足 $F(\psi) = \phi' \circ \phi^{-1}$ 的态射
$\psi : x \to x'$ 所成的集合。
\end{enumerate}
\end{lemma}

\begin{proof}
见上面的讨论。
\end{proof}

\begin{lemma}
\label{categories-lemma-prepare-representable-map-stack-in-groupoids}
令 $\mathcal{C}$ 为范畴。令 $\mathcal{X}$、$\mathcal{Y}$ 为
$\mathcal{C}$ 上的群胚纤维化范畴。令
$F : \mathcal{X} \to \mathcal{Y}$ 为 $1$-态射。令
$G : \mathcal{C}/U \to \mathcal{Y}$ 为 $1$-态射。那么
$$
(\mathcal{C}/U) \times_\mathcal{Y} \mathcal{X}
\longrightarrow
\mathcal{C}/U
$$
是群胚纤维化范畴。
\end{lemma}

\begin{proof}
我们已在引理 \ref{categories-lemma-2-product-fibred-categories} 中看到，复合
$$
(\mathcal{C}/U) \times_\mathcal{Y} \mathcal{X}
\longrightarrow
\mathcal{C}/U
\longrightarrow
\mathcal{C}
$$
是群胚纤维化范畴。于是该引理由引理
\ref{categories-lemma-cute-groupoids} 推出。
\end{proof}

\begin{definition}
\label{categories-definition-representable-map-categories-fibred-in-groupoids}
令 $\mathcal{C}$ 为范畴。令 $\mathcal{X}$、$\mathcal{Y}$ 为
$\mathcal{C}$ 上的群胚纤维化范畴。令
$F : \mathcal{X} \to \mathcal{Y}$ 为 $1$-态射。若对每个
$U \in \Ob(\mathcal{C})$ 及任意
$G : \mathcal{C}/U \to \mathcal{Y}$，群胚纤维化范畴
$$
(\mathcal{C}/U) \times_\mathcal{Y} \mathcal{X}
\longrightarrow
\mathcal{C}/U
$$
都是可表的，则称 $F$ 是{\it 可表的}，或称
{\it $\mathcal{X}$ 在 $\mathcal{Y}$ 上相对可表}。
\end{definition}

\begin{lemma}
\label{categories-lemma-spell-out-representable-map-stack-in-groupoids}
令 $\mathcal{C}$ 为范畴。令 $\mathcal{X}$、$\mathcal{Y}$ 为
$\mathcal{C}$ 上的群胚纤维化范畴。令
$F : \mathcal{X} \to \mathcal{Y}$ 为 $1$-态射。若 $F$ 可表，
则纤维范畴之间的每个函子
$$
F_U : \mathcal{X}_U \longrightarrow \mathcal{Y}_U
$$
都是忠实的。
\end{lemma}

\begin{proof}
由引理 \ref{categories-lemma-identify-fibre-product} 中对纤维范畴的描述，
以及引理 \ref{categories-lemma-characterize-representable-fibred-category}
中对可表纤维化范畴的刻画，这是显然的。
\end{proof}

\begin{lemma}
\label{categories-lemma-criterion-representable-map-stack-in-groupoids}
令 $\mathcal{C}$ 为范畴。令 $\mathcal{X}$、$\mathcal{Y}$ 为
$\mathcal{C}$ 上的群胚纤维化范畴。令
$F : \mathcal{X} \to \mathcal{Y}$ 为 $1$-态射。为
$\mathcal{Y}$ 选择拉回。假设
\begin{enumerate}
\item 纤维范畴之间的每个函子
$F_U : \mathcal{X}_U \longrightarrow \mathcal{Y}_U$ 都是忠实的；
\item 对每个 $U$ 及每个 $y \in \mathcal{Y}_U$，$\mathcal{C}/U$
上的预层
$$
(f : V \to U)
\longmapsto
\{(x, \phi) \mid x \in \mathcal{X}_V, \phi : f^*y \to F(x)\}/\cong
$$
是可表预层。
\end{enumerate}
那么 $F$ 可表。
\end{lemma}

\begin{proof}
由引理 \ref{categories-lemma-identify-fibre-product} 中对纤维范畴的描述，
以及引理 \ref{categories-lemma-characterize-representable-fibred-category}
中对可表纤维化范畴的刻画，这是显然的。
\end{proof}

\noindent
在陈述下一个引理之前，先指出群胚纤维化范畴所成的 $2$-范畴
是 $(2, 1)$-范畴，因而我们知道说它有终对象是什么意思（见定义
\ref{categories-definition-final-object-2-category}）。它确有终对象，即
$\text{id} : \mathcal{C} \to \mathcal{C}$。因此，我们把
$\mathcal{C}$ 上群胚纤维化范畴的 {\it $2$-积}定义为 $2$-纤维积
$$
\mathcal{X} \times \mathcal{Y} :=
\mathcal{X} \times_\mathcal{C} \mathcal{Y}.
$$
有了这个定义，下面的引理便有意义。

\begin{lemma}
\label{categories-lemma-representable-diagonal-groupoids}
令 $\mathcal{C}$ 为范畴。令 $\mathcal{S} \to \mathcal{C}$
为群胚纤维化范畴。假设 $\mathcal{C}$ 有对象对的积及纤维积。
下列条件等价：
\begin{enumerate}
\item 对角态射 $\mathcal{S} \to \mathcal{S} \times \mathcal{S}$ 可表；
\item 对 $\mathcal{C}$ 中每个 $U$，任意
$G : \mathcal{C}/U \to \mathcal{S}$ 都可表。
\end{enumerate}
\end{lemma}

\begin{proof}
假设对角态射可表，并给定 $U, G$。任取
$V \in \Ob(\mathcal{C})$ 和
$G' : \mathcal{C}/V \to \mathcal{S}$。注意
$\mathcal{C}/U \times \mathcal{C}/V = \mathcal{C}/U \times V$
可表。因此按假设，纤维积
$$
\xymatrix{
(\mathcal{C}/U \times V)
\times_{(\mathcal{S} \times \mathcal{S})}
\mathcal{S}
\ar[r] \ar[d] &
\mathcal{S} \ar[d] \\
\mathcal{C}/U \times V \ar[r]^{(G, G')} &
\mathcal{S} \times \mathcal{S}
}
$$
可表。这意味着存在 $\mathcal{C}$ 中的 $W \to U \times V$，使得
$$
\xymatrix{
\mathcal{C}/W \ar[d] \ar[r] & \mathcal{S} \ar[d] \\
\mathcal{C}/U \times \mathcal{C}/V \ar[r] & \mathcal{S} \times \mathcal{S}
}
$$
为笛卡尔图。由此得到所需的
$\mathcal{C}/W \cong \mathcal{C}/U \times_\mathcal{S} \mathcal{C}/V$
（见引理 \ref{categories-lemma-diagonal-1} 和注
\ref{categories-remark-2-product-fibred-categories-same}）。

\medskip\noindent
假设 (2) 成立。任取 $V \in \Ob(\mathcal{C})$ 及
$(G, G') : \mathcal{C}/V \to \mathcal{S} \times \mathcal{S}$。
我们必须证明
$\mathcal{C}/V \times_{\mathcal{S} \times \mathcal{S}} \mathcal{S}$
可表。已知
$\mathcal{C}/V \times_{G, \mathcal{S}, G'} \mathcal{C}/V$
可表，设它由 $\mathcal{C}/V$ 中的 $a : W \to V$ 表示。
等价
$$
\mathcal{C}/W \to \mathcal{C}/V \times_{G, \mathcal{S}, G'} \mathcal{C}/V
$$
再接第二投影到 $\mathcal{C}/V$，给出第二个态射
$a' : W \to V$。考虑
$W' = W \times_{(a, a'), V \times V} V$。存在等价
$$
\mathcal{C}/W' \cong
\mathcal{C}/V \times_{\mathcal{S} \times \mathcal{S}} \mathcal{S}
$$
即
\begin{eqnarray*}
\mathcal{C}/W' & \cong &
\mathcal{C}/W \times_{(\mathcal{C}/V \times \mathcal{C}/V)} \mathcal{C}/V \\
& \cong &
\left(\mathcal{C}/V \times_{(G, \mathcal{S}, G')} \mathcal{C}/V\right)
\times_{(\mathcal{C}/V \times \mathcal{C}/V)} \mathcal{C}/V \\
& \cong &
\mathcal{C}/V \times_{(\mathcal{S} \times \mathcal{S})} \mathcal{S}
\end{eqnarray*}
（最后一个同构见引理 \ref{categories-lemma-diagonal-2} 和注
\ref{categories-remark-2-product-fibred-categories-same}），这就证明了引理。
\end{proof}

\noindent
{\bf 文献说明：}
本节部分内容取自 Vistoli 的讲义 \cite{Vis2}。





\section{幺半范畴}
\label{categories-section-monoidal}

\noindent
令 $\mathcal{C}$ 为范畴。假设给定函子
$$
\otimes : \mathcal{C} \times \mathcal{C} \longrightarrow \mathcal{C}
$$
我们常常想知道 $\otimes$ 是否满足结合律，以及 $\otimes$ 是否有单位元。

\medskip\noindent
$(\mathcal{C}, \otimes)$ 的一个{\it 结合约束}，是函子性的同构
$$
\phi_{X, Y, Z} : X \otimes (Y \otimes Z) \to (X \otimes Y) \otimes Z
$$
使得对所有对象 $X, Y, Z, W$，图
\begin{equation}
\label{categories-equation-pentagon}
\vcenter{
\xymatrix{
X \otimes (Y \otimes ( Z \otimes W)) \ar[r] \ar[d] &
X \otimes ((Y \otimes Z) \otimes W) \ar[d] \\
(X \otimes Y) \otimes (Z \otimes W) \ar[d] &
(X \otimes (Y \otimes Z)) \otimes W \ar[ld] \\
((X \otimes Y) \otimes Z) \otimes W
}
}
\end{equation}
交换，其中每个箭头都由对 $\phi$ 的适当应用和 $\otimes$ 的函子性确定。

\medskip\noindent
文献 \cite[Theorem 3.1]{associativity} 说明了上述三元组
$(\mathcal{C}, \otimes, \phi)$ 如何具有协调性。换一种表述，这样的三元组
对所有 $n \geq 1$ 给出良定义的函子系
$$
\mathcal{C} \times \ldots \times \mathcal{C} \longrightarrow \mathcal{C},
\quad
(X_1, \ldots, X_n) \longmapsto X_1 \otimes \ldots \otimes X_n
$$
并对 $n, m \geq 1$ 给出函子性同构
$$
(X_1 \otimes \ldots \otimes X_n) \otimes (Y_1 \otimes \ldots \otimes Y_m)
\to
X_1 \otimes \ldots \otimes X_n \otimes Y_1 \otimes \ldots \otimes Y_m
$$
使得由这些同构组成的一切可能的图都交换（上面的五角星图是一个例子）。
当 $n = 5$ 时，这意味着各表达式
$$
\begin{matrix}
(((X_1 \otimes X_2) \otimes X_3) \otimes X_4) \otimes X_5 &
((X_1 \otimes (X_2 \otimes X_3)) \otimes X_4) \otimes X_5 \\
((X_1 \otimes X_2) \otimes (X_3 \otimes X_4)) \otimes X_5 &
(X_1 \otimes ((X_2 \otimes X_3) \otimes X_4)) \otimes X_5 \\
(X_1 \otimes (X_2 \otimes (X_3 \otimes X_4))) \otimes X_5 &
((X_1 \otimes X_2) \otimes X_3) \otimes (X_4 \otimes X_5) \\
(X_1 \otimes (X_2 \otimes X_3)) \otimes (X_4 \otimes X_5) &
(X_1 \otimes X_2) \otimes ((X_3 \otimes X_4) \otimes X_5) \\
(X_1 \otimes X_2) \otimes (X_3 \otimes (X_4 \otimes X_5)) &
X_1 \otimes (((X_2 \otimes X_3) \otimes X_4) \otimes X_5) \\
X_1 \otimes ((X_2 \otimes (X_3 \otimes X_4)) \otimes X_5) &
X_1 \otimes ((X_2 \otimes X_3) \otimes (X_4 \otimes X_5)) \\
X_1 \otimes (X_2 \otimes ((X_3 \otimes X_4) \otimes X_5)) &
X_1 \otimes (X_2 \otimes (X_3 \otimes (X_4 \otimes X_5)))
\end{matrix}
$$
都函子性同构于 $X_1 \otimes \ldots \otimes X_5$，而且这些同构与
利用 $\phi$ 在三个连续位置（若可能）得到的这些函子之间的一切可能同构
全都相容。以下我们将不加说明地使用这一点，并且在书写 $\otimes$ 的迭代时
不再写括号。

\medskip\noindent
上述三元组 $(\mathcal{C}, \otimes, \phi)$ 的一个{\it 单位元}，是
$\mathcal{C}$ 的对象 $\mathbf{1}$，连同函子性同构
$$
l : \mathbf{1} \otimes X \to X
\quad\text{且}\quad
r : X \otimes \mathbf{1} \to X
$$
使得对所有对象 $X, Y$，图
\begin{equation}
\label{categories-equation-triangle}
\vcenter{
\xymatrix{
X \otimes (\mathbf{1} \otimes Y) \ar[rr]_\phi \ar[rd]_{\text{id} \otimes l} & &
(X \otimes \mathbf{1}) \otimes Y \ar[ld]^{r \otimes \text{id}} \\
& X \otimes Y
}
}
\end{equation}
交换。我们常把单位元看作下面引理中的有序对 $(\mathbf{1}, 1)$。

\begin{lemma}
\label{categories-lemma-monoidal-unit}
令 $(\mathcal{C}, \otimes, \phi)$ 如上。$\mathcal{C}$ 中的单位元
$(\mathbf{1}, l, r)$ 与有序对 $(\mathbf{1}, 1)$ 一一对应，其中
$\mathbf{1}$ 是 $\mathcal{C}$ 的对象，且
$1 : \mathbf{1} \otimes \mathbf{1} \to \mathbf{1}$ 是同构，并且函子
$L : X \mapsto \mathbf{1} \otimes X$ 与
$R : X \mapsto X \otimes \mathbf{1}$ 都是等价。
\end{lemma}

\begin{proof}
给定单位元 $(\mathbf{1}, l, r)$，我们得到同构
$r : \mathbf{1} \otimes \mathbf{1} \to \mathbf{1}$；又因为 $L$ 与 $R$
都同构于恒等函子，所以它们是等价。反过来，假设给定
$(\mathbf{1}, 1)$，并且 $L$ 与 $R$ 是等价。我们得到函子性同构
$l_X : \mathbf{1} \otimes X \to X$ 和
$r_X : X \otimes \mathbf{1} \to X$，它们分别由
$L(l_X) = 1 \otimes \text{id}_X$ 和
$R(r_X) = \text{id}_X \otimes 1$ 所刻画。于是必须证明，对所有
$X$ 和 $Y$，从 $X \otimes \mathbf{1} \otimes Y$ 到
$X \otimes Y$ 的两个箭头 $\text{id}_X \otimes l_Y$ 与
$r_X \otimes \text{id}_Y$ 相同。这个性质只依赖于 $X$ 与 $Y$ 的
同构类。由于 $R$ 与 $L$ 是等价，只需对
$X = Z \otimes \mathbf{1}$ 和 $Y = \mathbf{1} \otimes W$ 证明它，
其中 $Z$ 和 $W$ 是某些对象。换言之，我们必须证明
$$
\text{id}_Z \otimes \text{id}_\mathbf{1} \otimes l_{\mathbf{1} \otimes W}
=
r_{Z \otimes \mathbf{1}} \otimes \text{id}_\mathbf{1} \otimes \text{id}_Y
$$
按构造，这两个映射分别等于
$\text{id}_Z \otimes 1 \otimes \text{id}_\mathbf{1} \otimes \text{id}_W$
与
$\text{id}_Z \otimes \text{id}_\mathbf{1} \otimes 1 \otimes \text{id}_W$。
因此只需证明
$1 \otimes \text{id}_\mathbf{1} = \text{id}_\mathbf{1} \otimes 1$。

\medskip\noindent
我们有
$\text{id}_\mathbf{1} \otimes 1 = r_\mathbf{1} \otimes \text{id}_\mathbf{1}$
以及
$l_\mathbf{1} \otimes \text{id}_\mathbf{1} = \text{id}_\mathbf{1} \otimes 1$。
可对 $\mathbf{1}$ 的某些自同构 $a, b$ 写成
$r_\mathbf{1} = a \circ 1$ 和 $l_\mathbf{1} = b \circ 1$。于是
$\text{id}_\mathbf{1} \otimes 1 =
(a \otimes \text{id}_\mathbf{1}) \circ (1 \otimes \text{id}_\mathbf{1})$
且
$1 \otimes \text{id}_\mathbf{1} =
(\text{id}_\mathbf{1} \otimes b) \circ (\text{id}_\mathbf{1} \otimes 1)$。
因此可写成
\begin{align*}
(1 \otimes \text{id}_\mathbf{1}) \circ
(\text{id}_\mathbf{1} \otimes 1 \otimes \text{id}_\mathbf{1})
& =
(1 \otimes \text{id}_\mathbf{1}) \circ
(\text{id}_\mathbf{1} \otimes \text{id}_\mathbf{1} \otimes b)
\circ
(\text{id}_\mathbf{1} \otimes \text{id}_\mathbf{1} \otimes 1) \\
& =
(\text{id}_\mathbf{1} \otimes b) \circ
(1 \otimes \text{id}_\mathbf{1}) \circ
(\text{id}_\mathbf{1} \otimes \text{id}_\mathbf{1} \otimes 1) \\
& =
(\text{id}_\mathbf{1} \otimes b) \circ
(1 \otimes 1)
\end{align*}
并且还有
\begin{align*}
(1 \otimes \text{id}_\mathbf{1}) \circ
(\text{id}_\mathbf{1} \otimes 1 \otimes \text{id}_\mathbf{1})
& =
(\text{id}_\mathbf{1} \otimes b) \circ
(\text{id}_\mathbf{1} \otimes 1) \circ
(a \otimes \text{id}_\mathbf{1} \otimes \text{id}_\mathbf{1}) \circ
(1 \otimes \text{id}_\mathbf{1} \otimes \text{id}_\mathbf{1}) \\
& =
(\text{id}_\mathbf{1} \otimes b) \circ
(a \otimes \text{id}_\mathbf{1}) \circ
(1 \otimes 1)
\end{align*}
这证明 $a \otimes \text{id}_\mathbf{1}$ 是恒等态射，从而如所需，
$a$ 是恒等态射。

\medskip\noindent
为完成证明，注意上述规则在单位元所成的范畴（适当定义）与有序对
$(\mathbf{1}, 1)$ 所成的范畴之间给出互逆等价。
\end{proof}

\begin{lemma}
\label{categories-lemma-monoidal-unit-unique}
令 $(\mathcal{C}, \otimes, \phi)$ 如上。令 $(\mathbf{1}, 1)$ 为单位元
（见引理 \ref{categories-lemma-monoidal-unit}）。那么
\begin{enumerate}
\item $1 \otimes \text{id}_\mathbf{1} = \text{id}_\mathbf{1} \otimes 1$
\item $\Gamma = \text{Mor}(\mathbf{1}, \mathbf{1})$ 是交换幺半群；
\item  $a = 1 \circ (a \otimes \text{id}_\mathbf{1}) \circ 1^{-1} =
1 \circ (\text{id}_\mathbf{1} \otimes a) \circ 1^{-1}$ 对所有
$a \in \Gamma$ 成立；
\item 任何其他单位元都经唯一同构同构于 $(\mathbf{1}, 1)$。
\end{enumerate}
\end{lemma}

\begin{proof}
第 (1) 部分已在引理 \ref{categories-lemma-monoidal-unit} 的证明中证得。
对 $a \in \Gamma$，有
\begin{align*}
(1 \circ (a \otimes \text{id}_\mathbf{1})) \otimes \text{id}_\mathbf{1}
& =
(1 \otimes \text{id}_\mathbf{1}) \circ
(a \otimes \text{id}_\mathbf{1} \otimes \text{id}_\mathbf{1}) \\
& =
(\text{id}_\mathbf{1} \otimes 1) \circ
(a \otimes \text{id}_\mathbf{1} \otimes \text{id}_\mathbf{1}) \\
& =
(a \otimes \text{id}_\mathbf{1}) \circ
(\text{id}_\mathbf{1} \otimes 1) \\
& =
(a \otimes \text{id}_\mathbf{1}) \circ
(1 \otimes \text{id}_\mathbf{1}) \\
& =
(a \circ 1) \otimes \text{id}_\mathbf{1}
\end{align*}
所以 $1 \circ (a \otimes \text{id}_\mathbf{1}) = a \circ 1$，
这给出 (3) 的一半。另一半同理。为证 (2)，注意对
$a, b \in \Gamma$，有
$$
a \circ b =
1 \circ (a \otimes \text{id}_\mathbf{1}) \circ 1^{-1}
\circ
1 \circ (\text{id}_\mathbf{1} \otimes b) \circ 1^{-1}
=
1 \circ (a \otimes b) \circ \circ 1^{-1}
$$
而 $b \circ a$ 也计算为同一表达式。因此 (2) 成立。为证 (4)，
设 $(\mathbf{1}', 1')$ 是第二个单位元。利用 $r$ 与 $l'$，有同构
$\mathbf{1}' \otimes \mathbf{1} \to \mathbf{1}'$ 和
$\mathbf{1}' \otimes \mathbf{1} \to \mathbf{1}$。因此存在同构
$t : \mathbf{1}' \to \mathbf{1}$。图
$$
\xymatrix{
\mathbf{1}' \otimes \mathbf{1}' \ar[r]_{t \otimes t} \ar[d]_{1'} & 
\mathbf{1} \otimes \mathbf{1} \ar[d]^1 \\
\mathbf{1}' \ar[r]^t & \mathbf{1}
}
$$
至多相差 $\mathbf{1}$ 的一个自同构 $a$ 而交换。把 $t$ 替换为
$a \circ t$ 后，该图将交换（提示：利用 (3) 看出
$1 \circ (a \otimes a) = a^2 \circ 1$）。
\end{proof}

\begin{definition}
\label{categories-definition-monoidal-category}
三元组 $(\mathcal{C}, \otimes, \phi)$ 中，$\mathcal{C}$ 是范畴，
$\otimes : \mathcal{C} \times \mathcal{C} \to \mathcal{C}$ 是函子，
$\phi$ 是结合约束；若存在单位元 $\mathbf{1}$，则称该三元组为
{\it 幺半范畴}。
\end{definition}

\noindent
我们总以 $\mathbf{1}$ 表示幺半范畴的一个单位元，并以
$1 : \mathbf{1} \otimes \mathbf{1} \to \mathbf{1}$ 表示一个选定同构；
由于有序对 $(\mathbf{1}, 1)$ 在唯一同构意义下确定（引理
\ref{categories-lemma-monoidal-unit-unique}），这样选择并无妨碍。

\begin{lemma}
\label{categories-lemma-monoidal-coherence}
在幺半范畴 $\mathcal{C}, \otimes, \phi, \mathbf{1}, 1$ 中，沿用引理
\ref{categories-lemma-monoidal-unit} 证明中的记号，我们有
\begin{enumerate}
\item 箭头 $1, r_\mathbf{1}, l_\mathbf{1} :
\mathbf{1} \otimes \mathbf{1} \to \mathbf{1}$ 相同；
\item 箭头
$l_X \otimes \text{id}_Y, l_{X \otimes Y} :
\mathbf{1} \otimes X \otimes Y \to X \otimes Y$ 相同；
\item 箭头
$\text{id}_X \otimes r_Y , r_{X \otimes Y} :
X \otimes Y \otimes \mathbf{1} \to X \otimes Y$ 相同。
\end{enumerate}
幺半范畴满足 \cite[Theorem 5.2]{associativity} 的假设。
\end{lemma}

\begin{proof}
引理 \ref{categories-lemma-monoidal-unit} 的证明已经给出 (1)。在引理
\ref{categories-lemma-monoidal-unit} 的证明中还看到，
$l_X$ 与 $l_{X \otimes Y}$ 分别是满足
$\text{id}_\mathbf{1} \otimes l_{X \otimes Y} =
1 \otimes \text{id}_{X \otimes Y}$ 和
$\text{id}_\mathbf{1} \otimes l_X = 1 \otimes \text{id}_X$ 的唯一态射。
第 (2) 部分立即推出。第 (3) 部分同理可证。结合
(\ref{categories-equation-pentagon}) 与 (\ref{categories-equation-triangle}) 的交换性，
这说明引理的最后一个断言成立。
\end{proof}

\noindent
文献 \cite[Theorem 5.2]{associativity} 说明了上述引理中的五元组
$(\mathcal{C}, \otimes, \phi, \mathbf{1}, 1)$ 如何具有协调性。
换一种表述，在我们所定义的幺半范畴中，对每个
$n \geq i \geq 0$ 都有函子性同构
$$
X_1 \otimes \ldots \otimes X_i \otimes \mathbf{1}
\otimes X_{i + 1} \otimes \ldots \otimes X_n
\to
X_1 \otimes \ldots \otimes X_n
$$
使得由这些同构组成的一切可能的图都交换。例如，从插入两个
$\mathbf{1}$ 开始，以不同次序使用这些同构，所得态射相同。除了书写
$\otimes$ 的迭代时省略括号这一约定外，从现在起还把
$X \otimes \mathbf{1}$ 和 $\mathbf{1} \otimes X$ 与 $X$ 等同，
不再另作说明。此外，我们会说“令 $\mathcal{C}$ 为幺半范畴”，
并默认给定 $\otimes, \phi, \mathbf{1}$。

\begin{definition}
\label{categories-definition-functor-monoidal-categories}
令 $\mathcal{C}$ 与 $\mathcal{C}'$ 为幺半范畴。幺半范畴的一个
{\it 函子} $F : \mathcal{C} \to \mathcal{C}'$，由所示函子 $F$
以及关于 $X$、$Y$ 函子性的同构
$$
F(X) \otimes F(Y) \to F(X \otimes Y)
$$
给出，使得对所有对象 $X$、$Y$、$Z$，图
$$
\xymatrix{
F(X) \otimes (F(Y) \otimes F(Z)) \ar[r] \ar[d] &
F(X) \otimes F(Y \otimes Z) \ar[r] &
F(X \otimes (Y \otimes Z)) \ar[d] \\
(F(X) \otimes F(Y)) \otimes F(Z) \ar[r] &
F(X \otimes Y) \otimes F(Z) \ar[r] &
F((X \otimes Y) \otimes Z)
}
$$
交换，并且 $F(\mathbf{1})$ 是 $\mathcal{C}'$ 中的单位元。
\end{definition}

\noindent
依照我们关于单位元的约定，若 $F$ 是幺半范畴的函子，总可假定
$F(\mathbf{1}) = \mathbf{1}$。例如，若 $A \to B$ 是环同态，则函子
$M \mapsto M \otimes_A B$ 是从 $\text{Mod}_A$ 到 $\text{Mod}_B$
的幺半范畴函子。

\begin{lemma}
\label{categories-lemma-invertible}
令 $\mathcal{C}$ 为幺半范畴。令 $X$ 为 $\mathcal{C}$ 的对象。
下列条件等价：
\begin{enumerate}
\item 函子 $L : Y \mapsto X \otimes Y$ 是等价；
\item 函子 $R : Y \mapsto Y \otimes X$ 是等价；
\item 存在对象 $X'$，使得
$X \otimes X' \cong X' \otimes X \cong \mathbf{1}$。
\end{enumerate}
\end{lemma}

\begin{proof}
假设 (1)。选择 $X'$ 使 $L(X') = \mathbf{1}$，即
$X \otimes X' \cong \mathbf{1}$。以 $L'$ 和 $R'$ 表示与 $X'$
对应的函子。等式 $X \otimes X' \cong \mathbf{1}$ 蕴含
$L \circ L' \cong \text{id}$。因此 $L'$ 必为 $L$ 的拟逆
（它依假设存在）。所以 $L' \circ L \cong \text{id}$，进而
$X' \otimes X \cong \mathbf{1}$。故 (3) 成立。

\medskip\noindent
(2) $\Rightarrow$ (3) 的证明与刚才的论证对偶。

\medskip\noindent
假设 (3)。显然 $L'$ 与 $L$ 互为拟逆，且 $R'$ 与 $R$ 互为拟逆。
故 (1) 与 (2) 成立。
\end{proof}

\begin{definition}
\label{categories-definition-invertible}
令 $\mathcal{C}$ 为幺半范畴。若引理 \ref{categories-lemma-invertible}
中的任一（从而全部）等价条件成立，则称 $\mathcal{C}$ 的对象 $X$
是{\it 可逆的}。
\end{definition}

\noindent
注意，若 $F : \mathcal{C} \to \mathcal{C}'$ 是幺半范畴的函子，
则 $F$ 把可逆对象映到可逆对象。

\begin{definition}
\label{categories-definition-dual}
给定幺半范畴 $(\mathcal{C}, \otimes, \phi)$ 及对象 $X$，一个
{\it 左对偶}是对象 $Y$，连同态射
$\eta : \mathbf{1} \to X \otimes Y$ 和
$\epsilon : Y \otimes X \to \mathbf{1}$，使得图
$$
\vcenter{
\xymatrix{
X \ar[rd]_1 \ar[r]_-{\eta \otimes 1} &
X \otimes Y \otimes X \ar[d]^{1 \otimes \epsilon} \\
& X
}
}
\quad\text{且}\quad
\vcenter{
\xymatrix{
Y \ar[rd]_1 \ar[r]_-{1 \otimes \eta} &
Y \otimes X \otimes Y \ar[d]^{\epsilon \otimes 1} \\
& Y
}
}
$$
交换。在这种情形下，我们称 $X$ 是 $Y$ 的一个{\it 右对偶}。
\end{definition}

\noindent
注意，若 $F : \mathcal{C} \to \mathcal{C}'$ 是幺半范畴的函子，
且 $Y$ 是 $X$ 的左对偶，则 $F(Y)$ 是 $F(X)$ 的左对偶。

\begin{lemma}
\label{categories-lemma-left-dual}
令 $\mathcal{C}$ 为幺半范畴。若 $Y$ 是 $X$ 的左对偶，则
$$
\Mor(Z' \otimes X, Z) = \Mor(Z', Z \otimes Y)
\quad\text{且}\quad
\Mor(Y \otimes Z', Z) = \Mor(Z', X \otimes Z)
$$
关于 $Z$ 与 $Z'$ 函子性地成立。
\end{lemma}

\begin{proof}
考虑映射
$$
\Mor(Z' \otimes X, Z) \to
\Mor(Z' \otimes X \otimes Y, Z \otimes Y) \to
\Mor(Z', Z \otimes Y)
$$
其中第二个箭头使用 $\eta$；以及映射序列
$$
\Mor(Z', Z \otimes Y) \to
\Mor(Z' \otimes X, Z \otimes Y \otimes X) \to
\Mor(Z' \otimes X, Z)
$$
其中第二个箭头使用 $\epsilon$。为证明这些箭头互逆，考虑映射
$a : Z' \to Z \otimes Y$。我们必须证明
$$
Z' \xrightarrow{\text{id}_{Z'} \otimes \eta}
Z' \otimes X \otimes Y \xrightarrow{a \otimes \text{id}_{X \otimes Y}}
Z \otimes Y \otimes X \otimes Y
\xrightarrow{\text{id}_Z \otimes \epsilon \otimes \text{id}_Y}
Z \otimes Y
$$
等于 $a$。前两个箭头的复合等于
$(\text{id}_{Z \otimes Y} \otimes \eta) \circ a :
Z' \to Z \otimes Y \to Z \otimes Y \otimes X \otimes Y$。
而 $\text{id}_{Z \otimes Y} \otimes \eta$ 与
$\text{id}_Z \otimes \epsilon \otimes \text{id}_Y$ 的复合，
依对偶的定义等于恒等态射。另一个复合同理。略去第二个等式的证明。
\end{proof}

\begin{remark}
\label{categories-remark-left-dual-adjoint}
引理 \ref{categories-lemma-left-dual} 特别说明，$Z \mapsto Z \otimes Y$ 是
$Z' \mapsto Z' \otimes X$ 的右伴随。特别地，伴随函子的唯一性保证，
若 $X$ 的左对偶存在，则它在唯一同构意义下唯一。反过来，假设函子
$Z \mapsto Z \otimes Y$ 是函子 $Z' \mapsto Z' \otimes X$ 的右伴随，
即给定关于 $Z$ 与 $Z'$ 均函子性的双射
$$
\Mor(Z' \otimes X, Z) \longrightarrow \Mor(Z', Z \otimes Y)
$$
伴随的单位给出关于 $Z$ 函子性的映射
$$
\eta_Z : Z \to Z \otimes X \otimes Y
$$
而伴随的余单位给出关于 $Z'$ 函子性的映射
$$
\epsilon_{Z'} : Z' \otimes Y \otimes X \to Z'
$$
特别地，我们得到
$\eta = \eta_\mathbf{1} : \mathbf{1} \to X \otimes Y$ 和
$\epsilon = \epsilon_\mathbf{1} : Y \otimes X \to \mathbf{1}$。
作为关于单位、余单位与伴随同构之关系的一个练习，读者可证明
$$
(\epsilon \otimes \text{id}_Y) \circ \eta_Y = \text{id}_Y
\quad\text{且}\quad
\epsilon_X \circ (\eta \otimes \text{id}_X) = \text{id}_X
$$
然而，这不足以证明
$(\epsilon \otimes \text{id}_Y) \circ (\text{id}_Y \otimes \eta) =
\text{id}_Y$ 以及
$(\text{id}_X \otimes \epsilon) \circ (\eta \otimes \text{id}_X) =
\text{id}_X$，因为一般而言，我们不知道
$\eta_Y = \text{id}_Y \otimes \eta$，也不知道
$\epsilon_X = \epsilon \otimes \text{id}_X$。为此，只要知道我们的伴随同构
具有下述性质便足够：对每个 $W, Z, Z'$，图
$$
\xymatrix{
\Mor(Z' \otimes X, Z) \ar[r] \ar[d]_{\text{id}_W \otimes -} &
\Mor(Z', Z \otimes Y) \ar[d]^{\text{id}_W \otimes -} \\
\Mor(W \otimes Z' \otimes X, W \otimes Z) \ar[r] &
\Mor(W \otimes Z', W \otimes Z \otimes Y)
}
$$
交换。若此条件成立，我们就说{\it 该伴随与给定张量结构相容}。
因此，要求 $Z \mapsto Z \otimes Y$ 是
$Z' \mapsto Z' \otimes X$ 的、与给定张量结构相容的右伴随，
等价于左对偶这一性质。
\end{remark}

\begin{lemma}
\label{categories-lemma-tensor-dual}
令 $\mathcal{C}$ 为幺半范畴。若 $Y_i$（$i = 1, 2$）分别是
$X_i$（$i = 1, 2$）的左对偶，则 $Y_2 \otimes Y_1$ 是
$X_1 \otimes X_2$ 的左对偶。
\end{lemma}

\begin{proof}
由伴随的唯一性和注 \ref{categories-remark-left-dual-adjoint} 得证。
\end{proof}

\noindent
$(\mathcal{C}, \otimes)$ 的一个{\it 交换约束}，是函子性同构
$$
\psi : X \otimes Y \longrightarrow Y \otimes X
$$
使得复合
$$
X \otimes Y \xrightarrow{\psi} Y \otimes X \xrightarrow{\psi} X \otimes Y
$$
是恒等态射。若对所有对象 $X, Y, Z$，图
\begin{equation}
\label{categories-equation-hexagon}
\vcenter{
\xymatrix{
X \otimes (Y \otimes Z) \ar[r]_\phi \ar[d]^\psi &
(X \otimes Y) \otimes Z \ar[r]_\psi &
Z \otimes (X \otimes Y) \ar[d]^\phi \\
X \otimes (Z \otimes Y) \ar[r]^\phi &
(X \otimes Z) \otimes Y \ar[r]^\psi &
(Z \otimes X) \otimes Y
}
}
\end{equation}
交换，则称 $\psi$ 与给定结合约束 $\phi$ {\it 相容}。

\begin{definition}
\label{categories-definition-symmetric-monoidal-category}
四元组 $(\mathcal{C}, \otimes, \phi, \psi)$ 中，$\mathcal{C}$ 是范畴，
$\otimes : \mathcal{C} \otimes \mathcal{C} \to \mathcal{C}$ 是函子，
$\phi$ 是结合约束，$\psi$ 是与 $\phi$ 相容的交换约束；若存在单位元，
则称该四元组为{\it 对称幺半范畴}。
\end{definition}

\noindent
确切地说，若 $(\mathcal{C}, \otimes, \phi, \psi)$ 是对称幺半范畴，
则 $(\mathcal{C}, \otimes, \phi)$ 是幺半范畴，因而可以使用上面讨论的
语言和记号。

\begin{lemma}
\label{categories-lemma-symmetric-monoidal-coherence}
在对称幺半范畴
$\mathcal{C}, \otimes, \phi, \psi, \mathbf{1}, 1$ 中，我们有
\begin{enumerate}
\item 箭头 $1 \circ \psi, 1 :
\mathbf{1} \otimes \mathbf{1} \to \mathbf{1}$ 相同；
\item 箭头 $\text{id}_X \otimes l_Y,
(l_X \otimes \text{id}) \circ (\psi \otimes \text{id}_Y):
X \otimes \mathbf{1} \otimes Y \to X \otimes Y$ 相同；
\end{enumerate}
对称幺半范畴满足 \cite[Theorem 5.1]{associativity} 的假设。
\end{lemma}

\begin{proof}
可写成 $\psi = a \otimes \text{id}_\mathbf{1}$，其中
$a : \mathbf{1} \to \mathbf{1}$ 是唯一同构。引理
\ref{categories-lemma-monoidal-unit-unique} 蕴含
$a \otimes \text{id}_\mathbf{1} = \text{id}_\mathbf{1} \otimes a$。
$\psi$ 的函子性说明图
$$
\xymatrix{
(\mathbf{1} \otimes \mathbf{1}) \otimes \mathbf{1}
\ar[r]_\psi
\ar[d]_{1 \otimes \text{id}_\mathbf{1}} &
\mathbf{1} \otimes (\mathbf{1} \otimes \mathbf{1})
\ar[d]^{\text{id}_\mathbf{1} \otimes 1} \\
\mathbf{1} \otimes \mathbf{1}  \ar[r]^\psi &
\mathbf{1} \otimes \mathbf{1}
}
$$
交换。因此上方箭头等于
$a \otimes \text{id}_\mathbf{1} \otimes \text{id}_\mathbf{1}$。
于是对 $X = Y = Z = \mathbf{1}$，(\ref{categories-equation-hexagon}) 表明
$a \otimes \text{id}_\mathbf{1} \otimes \text{id}_\mathbf{1}$ 等于自身的
平方。所以 $a = \text{id}_\mathbf{1}$。这就证明了 (1)。

\medskip\noindent
第 (2) 部分断言，如果在我们的幺半范畴中把源和靶都与 $X$ 等同，
则 $\psi : X \otimes \mathbf{1} \to \mathbf{1} \otimes X$ 是恒等态射。
这由 (\ref{categories-equation-hexagon}) 对 $\mathbf{1}, X, \mathbf{1}$ 的交换性推出，
即
$$
\xymatrix{
\mathbf{1} \otimes (X \otimes \mathbf{1}) \ar[r]_\phi \ar[d]^\psi &
(\mathbf{1} \otimes X) \otimes \mathbf{1} \ar[r]_\psi &
\mathbf{1} \otimes (\mathbf{1} \otimes X) \ar[d]^\phi \\
\mathbf{1} \otimes (\mathbf{1} \otimes X) \ar[r]^\phi &
(\mathbf{1} \otimes \mathbf{1}) \otimes X \ar[r]^\psi &
(\mathbf{1} \otimes \mathbf{1}) \otimes X
}
$$
交换，而我们知道此图中除一个以外的所有态射 $\psi$ 都等于恒等态射。

\medskip\noindent
除 (1) 与 (2) 外，图 (\ref{categories-equation-pentagon})、
(\ref{categories-equation-triangle})、(\ref{categories-equation-hexagon}) 的交换性以及引理
\ref{categories-lemma-monoidal-coherence} 的结果，蕴含引理的最后一个断言。
\end{proof}

\noindent
文献 \cite[Theorem 5.1]{associativity} 说明了上述引理中的六元组
$(\mathcal{C}, \otimes, \phi, \psi, \mathbf{1}, 1)$ 如何具有协调性。
换一种表述，在我们所定义的对称幺半范畴中，对每个 $n \geq 1$
以及 $\{1, \ldots, n\}$ 的每个置换 $\sigma$，都有函子性同构
$$
X_1 \otimes \ldots \otimes X_n
\to
X_{\sigma(1)} \otimes \ldots \otimes X_{\sigma(n)}
$$
使得由这些同构组成的一切可能的图都交换，并且这些同构与幺半范畴的
结构相容（例如，与在某一位置插入 $\mathbf{1}$ 时的同构相容）。

\begin{lemma}
\label{categories-lemma-dual-symmetric}
令 $(\mathcal{C}, \otimes, \phi, \psi)$ 为对称幺半范畴。
令 $X$ 为 $\mathcal{C}$ 的对象，并令 $Y$、
$\eta : \mathbf{1} \to X \otimes Y$、
$\epsilon : Y \otimes X \to \mathbf{1}$ 为定义
\ref{categories-definition-dual} 中 $X$ 的一个左对偶。那么
$\eta' = \psi \circ \eta : \mathbf{1} \to Y \otimes X$ 和
$\epsilon' = \epsilon \circ \psi : X \otimes Y \to \mathbf{1}$
使 $X$ 成为 $Y$ 的一个左对偶。
\end{lemma}

\begin{proof}
略。提示：这是一个愉快的定义练习。
\end{proof}

\begin{definition}
\label{categories-definition-functor-symmetric-monoidal-categories}
令 $\mathcal{C}$ 与 $\mathcal{C}'$ 为对称幺半范畴。对称幺半范畴的一个
{\it 函子} $F : \mathcal{C} \to \mathcal{C}'$，由所示函子 $F$ 以及
关于 $X$、$Y$ 函子性的同构
$$
F(X) \otimes F(Y) \to F(X \otimes Y)
$$
给出，使得 $F$ 是幺半范畴的函子，并且对所有对象 $X$ 与 $Y$，图
$$
\xymatrix{
F(X) \otimes F(Y) \ar[r] \ar[d] &
F(X \otimes Y) \ar[d] \\
F(Y) \otimes F(X) \ar[r] &
F(Y \otimes X)
}
$$
交换。
\end{definition}

\begin{remark}
\label{categories-remark-internal-hom-monoidal}
令 $\mathcal{C}$ 为幺半范畴。若对 $\mathcal{C}$ 的每对对象 $X, Y$，
都有 $\mathcal{C}$ 的对象 $hom(X, Y)$，使得
$$
\Mor(X, hom(Y, Z)) = \Mor(X \otimes Y, Z)
$$
关于 $X, Y, Z$ 函子性地成立，则称 $\mathcal{C}$ 有
{\it 内部 hom}。由米田引理，如果双函子
$(X, Y) \mapsto hom(X, Y)$ 存在，则它在唯一同构意义下确定。
给定内部 hom，我们得到典范映射
\begin{enumerate}
\item $hom(X, Y) \otimes X \to Y$；
\item $hom(Y, Z) \otimes hom(X, Y) \to hom(X, Z)$；
\item $Z \otimes hom(X, Y) \to hom(X, Z \otimes Y)$；
\item $Y \to hom(X, Y \otimes X)$；
\item 若 $\mathcal{C}$ 是对称幺半范畴，则有
$hom(Y, Z) \otimes X \to hom(hom(X, Y), Z)$。
\end{enumerate}
具体而言，(1) 中的映射是 $\text{id}_{hom(X, Y)}$ 在映射
$\Mor(hom(X, Y), hom(X, Y)) \to \Mor(hom(X, Y) \otimes X, Y)$
下的像。为用 $hom(X, Z)$ 的定义性质构造 (2) 中的映射，需要构造映射
$$
hom(Y, Z) \otimes hom(X, Y) \otimes X \longrightarrow Z
$$
而这样的映射存在，因为由 (1) 有映射
$hom(X, Y) \otimes X \to Y$ 与 $hom(Y, Z) \otimes Y \to Z$。
为用 $hom(X, Z \otimes Y)$ 的定义性质构造 (3) 中的映射，需要构造映射
$$
Z \otimes hom(X, Y) \otimes X \to Z \otimes Y
$$
为此使用 $\text{id}_Z \otimes a$，其中 $a$ 是 (1) 中的映射。
为构造 (4) 中的映射，注意由 (3) 已有映射
$Y \otimes hom(X, X) \to hom(X, Y \otimes X)$。所以只需构造映射
$\mathbf{1} \to hom(X, X)$；为此取
$\Mor(\mathbf{1}, hom(X, X))$ 中与
$\Mor(\mathbf{1} \otimes X, X)$ 中典范同构
$\mathbf{1} \otimes X \to X$ 对应的元素。最后来看 (5)。
由 $hom(hom(X, Y), Z)$ 的泛性质，只需构造映射
$$
hom(Y, Z) \otimes X \otimes hom(X, Y) \longrightarrow Z
$$
我们用交换约束交换最后两个张量积，然后使用映射
$hom(X, Y) \otimes X \to Y$ 与 $hom(Y, Z) \otimes Y \to Z$。
\end{remark}





\section{点线箭头范畴}
\label{categories-section-dotted-arrows}

\noindent 
我们讨论 $(2,1)$-范畴中的某些“点线箭头范畴”。在表述代数栈的各种
提升判据时会用到它们，例如见栈的态射，第
\ref{stacks-morphisms-section-valuative} 节，以及栈的态射进阶，第
\ref{stacks-more-morphisms-section-formally-smooth} 节。

\begin{definition}
\label{categories-definition-dotted-arrows}
令 $\mathcal{C}$ 为 $(2,1)$-范畴。考虑 $\mathcal{C}$ 中的
$2$-交换实线图
\begin{equation}
\label{categories-equation-dotted-arrows}
\vcenter{
\xymatrix{
S \ar[r]_-x \ar[d]_j & X \ar[d]^f \\
T \ar[r]^-y \ar@{..>}[ru] & Y
}
}
\end{equation}
固定一个见证此图 $2$-交换的 $2$-同构
$$
\gamma : y \circ j \rightarrow f \circ x
$$
给定 (\ref{categories-equation-dotted-arrows}) 与 $\gamma$，一个\emph{点线箭头}
是三元组 $(a, \alpha, \beta)$，其中包括态射 $a \colon T \to X$
及 $2$-同构
$\alpha : a \circ j \to x$、$\beta : y \to f \circ a$，并且
$\gamma = (\text{id}_f \star \alpha) \circ (\beta \star \text{id}_j)$；
换言之，图
$$
\xymatrix{
& f \circ a \circ j \ar[rd]^{\text{id}_f \star \alpha} \\
y \circ j \ar[ru]^{\beta \star \text{id}_j} \ar[rr]^\gamma & &
f \circ x
}
$$
交换。点线箭头的一个{\it 态射}
$(a, \alpha, \beta) \to (a', \alpha', \beta')$ 是满足
$\alpha = \alpha' \circ (\theta \star \text{id}_j)$ 且
$\beta' = (\text{id}_f \star \theta) \circ \beta$ 的
$2$-箭头 $\theta : a \to a'$。
\end{definition}

\noindent
在定义 \ref{categories-definition-dotted-arrows} 的情形中，有一个相伴的
\emph{点线箭头范畴}。这个范畴是群胚。一般而言，它可能依赖于
$\gamma$。下面两个引理说明，点线箭头范畴对基变换以及 $f$ 的复合
表现良好。

\begin{lemma}
\label{categories-lemma-cat-dotted-arrows-base-change}
令 $\mathcal{C}$ 为 $(2,1)$-范畴。假设给定 $\mathcal{C}$ 中的
$2$-交换图
$$
\xymatrix{
S \ar[r]_-{x'} \ar[d]_j &
X' \ar[d]^p \ar[r]_q &
X \ar[d]^f \\
T \ar[r]^-{y'} &
Y' \ar[r]^g &
Y
}
$$
其中右方块关于 $2$-同构
$\phi \colon g \circ p \to f \circ q$ 是 $2$-笛卡尔的。
选择 $2$-箭头
$\gamma' : y' \circ j \to p \circ x'$。令
$x = q \circ x'$、$y = g \circ y'$，并令
$\gamma : y \circ j \to f \circ x$ 为 $2$-同构
$\gamma = (\phi \star \text{id}_{x'}) \circ (\text{id}_g \star \gamma')$。
那么，左方块及 $\gamma'$ 的点线箭头范畴 $\mathcal{D}'$，等价于
外侧矩形及 $\gamma$ 的点线箭头范畴 $\mathcal{D}$。
\end{lemma}

\begin{proof}
有一个函子 $\mathcal{D}' \to \mathcal{D}$，它在对象上为
$(a, \alpha, \beta) \mapsto (q \circ a, \text{id}_q \star \alpha,
(\phi \star \text{id}_a) \circ (\text{id}_g \star \beta))$，
在箭头上为 $\theta \mapsto \text{id}_q \star \theta$。
由第 \ref{categories-section-2-fibre-products} 节所述 $2$-纤维积的泛性质，
形式地即可验证此函子 $\mathcal{D}' \to \mathcal{D}$ 是等价。
细节略去。
\end{proof}

\begin{lemma}
\label{categories-lemma-cat-dotted-arrows-composition}
令 $\mathcal{C}$ 为 $(2,1)$-范畴。假设给定 $\mathcal{C}$ 中的实线
$2$-交换图
$$
\xymatrix{
S \ar[r]_-x \ar[dd]_j & X \ar[d]^f \\
& Y \ar[d]^g \\
T \ar[r]^-z \ar@{..>}[ruu] & Z
}
$$
选择 $2$-同构 $\gamma \colon z \circ j \to g \circ f \circ x$。
令 $\mathcal{D}$ 为外侧矩形及 $\gamma$ 的点线箭头范畴。令
$\mathcal{D}'$ 为实线方块
$$
\xymatrix{
S \ar[r]_-{f \circ x} \ar[d]_j & Y \ar[d]^g \\
T \ar[r]^-z \ar@{..>}[ru] & Z
}
$$
及 $\gamma$ 的点线箭头范畴。那么 $\mathcal{D}$ 等价于一个范畴
$\mathcal{D}''$，后者具有下述性质：有函子
$\mathcal{D}'' \to \mathcal{D}'$，使 $\mathcal{D}''$ 成为
$\mathcal{D}'$ 上的群胚纤维化范畴；它的纤维范畴同构于形如
$$
\xymatrix{
S \ar[r]_-x \ar[d]_j & X \ar[d]^f \\
T \ar[r]^-y \ar@{..>}[ru] & Y
}
$$
的某些实线方块在选择了某个 $2$-同构
$y \circ j \to f \circ x$ 后所得的点线箭头范畴。
\end{lemma}

\begin{proof}
构造范畴 $\mathcal{D}''$，其对象是元组
$(a,\alpha,\beta,b,\eta)$，其中 $(a,\alpha,\beta)$ 是
$\mathcal{D}$ 的对象，$b \colon T \rightarrow Y$ 是 $1$-态射，
且 $\eta \colon b \rightarrow f \circ a$ 是 $2$-同构。
$\mathcal{D}''$ 中的态射
$(a,\alpha,\beta,b,\eta) \rightarrow (a',\alpha',\beta',b',\eta')$
是有序对 $(\theta_1,\theta_2)$，其中
$\theta_1 \colon a \rightarrow a'$ 在 $\mathcal{D}$ 中定义箭头
$(a, \alpha, \beta) \rightarrow (a', \alpha', \beta')$，且
$\theta_2 \colon b \rightarrow b'$ 是满足相容条件
$\eta' \circ \theta_2 = (\text{id}_f \star \theta_1) \circ \eta$
的 $2$-同构。

\medskip\noindent
有函子 $\mathcal{D}'' \rightarrow \mathcal{D}'$，它在对象上为
$(a, \alpha, \beta, b, \eta) \mapsto
(b, (\text{id}_f \star \alpha) \circ (\eta \star \text{id}_j),
(\text{id}_g \star \eta^{-1}) \circ \beta)$，在箭头上为
$(\theta_1,\theta_2) \mapsto \theta_2$。那么
$\mathcal{D}'' \rightarrow \mathcal{D}'$ 是群胚纤维化的。

\medskip\noindent
若 $(y, \delta, \epsilon)$ 是 $\mathcal{D}'$ 的对象，以
$\mathcal{D}_{y,\delta}$ 表示最后一个所示图在
$y \circ j \rightarrow f \circ x$ 取为 $\delta$ 时的点线箭头范畴。
有函子 $\mathcal{D}_{y,\delta} \rightarrow \mathcal{D}''$，它在对象上为
$(a, \alpha, \eta) \mapsto
(a, \alpha, (\text{id}_g \star \eta) \circ \epsilon, y, \eta)$，
在箭头上为 $\theta \mapsto (\theta, \text{id}_y)$。这给出从
$\mathcal{D}_{y,\delta}$ 到 $\mathcal{D}'' \rightarrow \mathcal{D}'$
在 $(y,\delta,\epsilon)$ 上的纤维范畴的同构。

\medskip\noindent
还有一个函子 $\mathcal{D} \rightarrow \mathcal{D}''$，它在对象上为
$(a,\alpha,\beta) \mapsto (a,\alpha,\beta,f \circ a, \text{id}_{f \circ a})$，
在箭头上为 $\theta \mapsto (\theta, \text{id}_f \star \theta)$。
此函子全忠实且本质满，因而是等价。细节略去。
\end{proof}












% P01 cursor: frozen categories.tex translated through source line 9855 (chapter complete).
